%%%=====================================================================%%%
%%%                                                                      %%%
%%%   V-JEPA 2.1 PAS À PAS                                               %%%
%%%   Comprendre, coder et déployer un modèle vidéo auto-supervisé       %%%
%%%                                                                      %%%
%%%   Fichier maître. Compilation :                                      %%%
%%%       latexmk        (recommandé, tout automatique)                  %%%
%%%     ou :  make                                                       %%%
%%%     ou à la main :                                                   %%%
%%%       pdflatex -shell-escape main                                    %%%
%%%       bibtex   main                                                  %%%
%%%       pdflatex -shell-escape main                                    %%%
%%%       pdflatex -shell-escape main                                    %%%
%%%                                                                      %%%
%%%   IMPORTANT : l'option -shell-escape est indispensable (minted).     %%%
%%%                                                                      %%%
%%%   La classe et les styles viennent du template ML4E, repris tels     %%%
%%%   quels. Seule l'identité du livre est redéfinie ci-dessous, dans    %%%
%%%   le préambule : le template d'origine n'est pas modifié.            %%%
%%%=====================================================================%%%
\documentclass{ml4ebook}     % ajoutez l'option [brouillon] pour un filigrane

%----------------------------------------------------------------------
% Identité de CE livre (le fichier de classe reste celui du template)
%----------------------------------------------------------------------
\renewcommand{\mbTitre}{V-JEPA 2.1 pas à pas}
\renewcommand{\mbTitreCourt}{V-JEPA 2.1 pas à pas}
\renewcommand{\mbSousTitre}{Comprendre, coder et déployer un modèle vidéo
  auto-supervisé, en partant vraiment de zéro}
\renewcommand{\mbAuteur}{Doctor Mokira}

\begin{document}

%----------------------------------------------------------------------
% Pages liminaires
%----------------------------------------------------------------------
\frontmatter
%%%=====================================================================%%%
%%%  frontmatter/cover.tex — Page de titre                              %%%
%%%                                                                      %%%
%%%  Le titre, le sous-titre et l'auteur viennent des macros \mbTitre,   %%%
%%%  \mbSousTitre et \mbAuteur redéfinies en tête de main.tex.           %%%
%%%                                                                      %%%
%%%  PARTI PRIS : la couverture raconte, sans un mot, la boucle          %%%
%%%  complète du livre :                                                 %%%
%%%                                                                      %%%
%%%     une vidéo  ·  on en cache des morceaux  ·  on devine l'idée      %%%
%%%              un seul encodeur gelé  ·  un petit modèle              %%%
%%%                                                                      %%%
%%%  Composition en registres horizontaux :                              %%%
%%%     1. la rangée haute : les trois temps de l'apprentissage          %%%
%%%     2. le bandeau bleu qui porte le titre                            %%%
%%%     3. le filet ambre, le sous-titre, la promesse                    %%%
%%%     4. la rangée basse : ce qu'on en fait ensuite                    %%%
%%%     5. l'auteur                                                      %%%
%%%                                                                      %%%
%%%  Les teintes sont plus soutenues que dans le corps du livre : une    %%%
%%%  couverture est vue de loin, souvent imprimée, parfois en noir et    %%%
%%%  blanc.                                                              %%%
%%%=====================================================================%%%
\begin{titlepage}
\thispagestyle{empty}

\tikzset{
  cov formule/.style={font=\small, text=mbInk},
  cov domaine/.style={font=\sffamily\scriptsize\bfseries, text=mbSlate},
}

\begin{tikzpicture}[remember picture,overlay,mbfig]

  %--------------------------------------------------------------------
  % Fond
  %--------------------------------------------------------------------
  \fill[mbPaper] (current page.south west) rectangle (current page.north east);

  %====================================================================
  % 1. RANGÉE HAUTE — les trois temps de l'apprentissage
  %====================================================================
  \coordinate (rhaut) at ($(current page.north)+(0cm,-4.6cm)$);

  %--------------------------------------------------------------------
  % 1a. UNE VIDÉO — une pile d'images qui avance dans le temps.
  %--------------------------------------------------------------------
  \begin{scope}[shift={($(rhaut)+(-4.5cm,0cm)$)}]
    \foreach \k [count=\n from 0] in {0,1,2,3}{
      \pgfmathsetmacro{\dx}{-0.95+0.55*\k}
      \pgfmathsetmacro{\dy}{-0.72+0.40*\k}
      \fill[mbBlue!\the\numexpr18+14*\n\relax]
            (\dx,\dy) rectangle ++(1.5,1.05);
      \draw[mbBlue!70, line width=0.7pt] (\dx,\dy) rectangle ++(1.5,1.05);
    }
    \draw[-{Stealth[round,length=2.2mm,width=1.7mm]},
          draw=mbAmber!85!black, line width=1pt]
          (-1.15,-1.05) -- (1.35,-1.05);
    \node[font=\scriptsize, text=mbAmber!80!black, below] at (0.1,-1.08)
         {temps};
    \node[cov formule] at (0,-1.85) {$(T,H,W,C)$};
    \node[cov domaine] at (0,-2.35) {UNE VIDÉO};
  \end{scope}

  %--------------------------------------------------------------------
  % 1b. ON EN CACHE DES MORCEAUX — la grille des jetons, dont une partie
  %     est masquée. C'est le geste central de tout le livre.
  %--------------------------------------------------------------------
  \begin{scope}[shift={($(rhaut)+(0cm,0cm)$)}]
    \foreach \ri in {0,...,5}{
      \foreach \ci in {0,...,5}{
        \pgfmathsetmacro{\dd}{(\ci-1.6)*(\ci-1.6)+(\ri-3.2)*(\ri-3.2)}
        \pgfmathsetmacro{\ee}{(\ci-4.3)*(\ci-4.3)+(\ri-1.4)*(\ri-1.4)}
        \pgfmathsetmacro{\cache}{(\dd<2.2 || \ee<1.6) ? 1 : 0}
        \ifdim\cache pt>0.5pt
          \fill[mbSlate!35] (-1.32+\ci*0.44, 1.32-\ri*0.44) rectangle ++(0.44,-0.44);
          \draw[mbSlate!55, line width=0.4pt, dashed]
                (-1.32+\ci*0.44, 1.32-\ri*0.44) rectangle ++(0.44,-0.44);
        \else
          \fill[mbAzure!\the\numexpr20+7*\ri+5*\ci\relax]
                (-1.32+\ci*0.44, 1.32-\ri*0.44) rectangle ++(0.44,-0.44);
          \draw[mbAzure!60, line width=0.4pt]
                (-1.32+\ci*0.44, 1.32-\ri*0.44) rectangle ++(0.44,-0.44);
        \fi
      }
    }
    \node[cov formule] at (0,-1.85) {$x \longrightarrow E_{\theta}(x)$};
    \node[cov domaine] at (0,-2.35) {ON EN CACHE};
  \end{scope}

  %--------------------------------------------------------------------
  % 1c. ON DEVINE L'IDÉE — pas les pixels : la carte de traits, celle
  %     qu'on lira toute l'année sous forme de carte en couleurs.
  %--------------------------------------------------------------------
  \begin{scope}[shift={($(rhaut)+(4.5cm,0cm)$)}]
    \foreach \ri in {0,...,5}{
      \foreach \ci in {0,...,5}{
        \pgfmathsetmacro{\da}{(\ci-1.3)*(\ci-1.3)+(\ri-3.4)*(\ri-3.4)}
        \pgfmathsetmacro{\db}{(\ci-4.2)*(\ci-4.2)+(\ri-1.5)*(\ri-1.5)}
        \pgfmathsetmacro{\za}{\da<2.6 ? 1 : 0}
        \pgfmathsetmacro{\zb}{\db<2.2 ? 1 : 0}
        \ifdim\za pt>0.5pt
          \fill[mbAmber!70] (-1.32+\ci*0.44, 1.32-\ri*0.44) rectangle ++(0.44,-0.44);
        \else
          \ifdim\zb pt>0.5pt
            \fill[mbTeal!65] (-1.32+\ci*0.44, 1.32-\ri*0.44) rectangle ++(0.44,-0.44);
          \else
            \fill[mbViolet!35] (-1.32+\ci*0.44, 1.32-\ri*0.44) rectangle ++(0.44,-0.44);
          \fi
        \fi
        \draw[white, line width=0.5pt]
              (-1.32+\ci*0.44, 1.32-\ri*0.44) rectangle ++(0.44,-0.44);
      }
    }
    \node[cov formule] at (0,-1.85) {$\lVert \hat{z} - z \rVert_{1}$};
    \node[cov domaine] at (0,-2.35) {ON DEVINE L'IDÉE};
  \end{scope}

  %====================================================================
  % 2. Bandeau du titre
  %====================================================================
  \fill[mbBlue] (current page.west|-current page.center)
        rectangle ([yshift=3.7cm]current page.east|-current page.center);
  \fill[mbAmber] (current page.west|-current page.center)
        rectangle ([yshift=-0.28cm]current page.east|-current page.center);

  \node[anchor=west,text width=11.5cm,align=left]
        at ([yshift=1.85cm,xshift=1.9cm]current page.west) {%
    \sffamily\bfseries\color{white}\fontsize{34}{39}\selectfont
    V-JEPA 2.1\\ pas à pas};

  %====================================================================
  % 3. Sous-titre et promesse, sous le filet — donc sur le fond crème
  %====================================================================
  \node[anchor=north west,text width=12.4cm,align=left]
        at ([yshift=-1.05cm,xshift=1.9cm]current page.west) {%
    \sffamily\color{mbInk}\large \mbSousTitre};
  \node[anchor=north west,text width=12.4cm,align=left]
        at ([yshift=-2.35cm,xshift=1.9cm]current page.west) {%
    \sffamily\bfseries\color{mbAmber!80!black}\normalsize
    Aucun prérequis : on part vraiment de zéro.};

  %====================================================================
  % 4. RANGÉE BASSE — ce qu'on fait du modèle une fois qu'il est appris
  %====================================================================
  \coordinate (rbas) at ($(current page.south)+(0cm,6.5cm)$);

  %--------------------------------------------------------------------
  % 4a. UN ENCODEUR GELÉ — un seul gros modèle, qu'on ne touche plus, et
  %     derrière lui de petites têtes, une par tâche.
  %--------------------------------------------------------------------
  \begin{scope}[shift={($(rbas)+(-3.4cm,0cm)$)}]
    \fill[mbBlue!25] (-1.7,-0.95) rectangle (-0.5,0.95);
    \draw[mbBlue, line width=1pt] (-1.7,-0.95) rectangle (-0.5,0.95);
    \node[font=\sffamily\scriptsize\bfseries, text=mbBlue!80!black, rotate=90]
          at (-1.1,0) {ENCODEUR};
    \node[font=\normalsize, text=mbAmber!85!black] at (-1.1,1.32)
          {\faIcon{lock}};
    \foreach \k [count=\n from 0] in {1,0,-1}{
      \draw[-{Stealth[round,length=1.8mm,width=1.4mm]},
            draw=mbSlate!70, line width=0.7pt]
            (-0.45,\k*0.6) -- (0.45,\k*0.6);
      \fill[mbAmber!35] (0.5,{\k*0.6-0.22}) rectangle ++(1.1,0.44);
      \draw[mbAmber!85!black, line width=0.7pt]
            (0.5,{\k*0.6-0.22}) rectangle ++(1.1,0.44);
    }
    \node[cov formule] at (0,-1.85) {une sonde par tâche};
    \node[cov domaine] at (0,-2.35) {ENCODEUR GELÉ};
  \end{scope}

  %--------------------------------------------------------------------
  % 4b. UN PETIT MODÈLE — la distillation : le gros apprend, le petit
  %     l'imite, et c'est le petit qui part en production.
  %--------------------------------------------------------------------
  \begin{scope}[shift={($(rbas)+(3.4cm,0cm)$)}]
    \fill[mbBlue!25] (-1.75,-1.1) rectangle (-0.35,1.1);
    \draw[mbBlue, line width=1pt] (-1.75,-1.1) rectangle (-0.35,1.1);
    \node[font=\sffamily\tiny\bfseries, text=mbBlue!80!black] at (-1.05,0)
          {2\,000\\ millions};
    \fill[mbTeal!30] (0.75,-0.42) rectangle (1.6,0.42);
    \draw[mbTeal!90!black, line width=1pt] (0.75,-0.42) rectangle (1.6,0.42);
    \node[font=\sffamily\tiny\bfseries, text=mbTeal!60!black] at (1.175,0)
          {300\\ millions};
    \draw[-{Stealth[round,length=2.2mm,width=1.7mm]},
          draw=mbAmber!85!black, line width=1.1pt] (-0.25,0) -- (0.65,0);
    \node[font=\sffamily\scriptsize, text=mbSlate, above] at (0.2,0.1)
          {imite};
    \node[cov formule] at (0,-1.85) {presque la même qualité};
    \node[cov domaine] at (0,-2.35) {DISTILLATION};
  \end{scope}

  %====================================================================
  % 5. Auteur
  %====================================================================
  \node[anchor=south west,text width=12.4cm,align=left]
        at ([yshift=1.5cm,xshift=1.9cm]current page.south west) {%
    \sffamily
    {\bfseries\color{mbInk}\large\mbAuteur}\\[6pt]
    {\color{mbBlue}\footnotesize\faIcon{calendar-alt}\ \the\year}};
\end{tikzpicture}
\end{titlepage}

% Verso de la page de titre : colophon
\thispagestyle{empty}
\null\vfill
{\footnotesize\sffamily\color{mbSlate}
\noindent Édition \the\year. Document composé avec \LaTeX{} (pdf\TeX), polices
Palatino (texte et mathématiques), TeX~Gyre~Heros (titres) et DejaVu~Sans~Mono
(code).
Figures et schémas réalisés avec \textsf{TikZ} et \textsf{pgfplots}, coloration
syntaxique assurée par \textsf{minted}/Pygments, encadrés par
\textsf{tcolorbox}, bibliographie par \textsf{biblatex}.\par
\medskip
\noindent Ce livre explique de bout en bout l'article \emph{V-JEPA 2.1:
Unlocking Dense Features in Video Self-Supervised Learning}, publié par
l'équipe FAIR de Meta et l'Université de Saragosse. Il est écrit pour être lu
sans aucun prérequis : on part de zéro, on explique avec des mots avant
d'écrire des symboles, et on n'avance jamais sans avoir d'abord dit pourquoi.
Il ne s'arrête pas à la compréhension : il conduit le lecteur jusqu'au code
qui tourne, puis jusqu'au modèle mis en service.\par}
\clearpage

\chapter*{Avant-propos}
\addcontentsline{toc}{chapter}{Avant-propos}
\markboth{Avant-propos}{}

\begin{quote}
\itshape « Un article de recherche ne devrait effrayer personne. »
\end{quote}

Ce livre raconte un seul article scientifique, du premier mot jusqu'à la
dernière ligne de code : \emph{V-JEPA~2.1: Unlocking Dense Features in Video
Self-Supervised Learning}~\parencite{murlabadia2026vjepa21}, écrit par
l'équipe FAIR de Meta et l'Université de Saragosse. Cet article présente un
programme qui apprend à comprendre des vidéos tout seul, sans qu'un humain
lui explique jamais ce qu'il regarde. Il fait suite à un premier modèle,
V-JEPA~2~\parencite{assran2025vjepa2}, dont nous parlerons beaucoup : c'est
en corrigeant ses défauts que V-JEPA~2.1 est né.

Un article de recherche est écrit par des spécialistes, pour des
spécialistes. Il va vite. Il emploie des mots que personne ne vous a jamais
définis. Il suppose que vous connaissez déjà dix autres articles. Résultat :
beaucoup de gens curieux, capables, motivés, s'arrêtent à la première page.

\textbf{Ici, on ne suppose rien du tout.} Chaque idée est d'abord racontée
avec des mots de tous les jours, puis dessinée, puis seulement écrite en
langage mathématique, puis enfin mise en pratique avec du code que vous
pouvez faire tourner vous-même.

\begin{enclair}
Notre promesse : si vous savez compter et si vous savez lire, vous pouvez
comprendre comment fonctionne V-JEPA~2.1. Le reste, nous le construisons
ensemble, dans l'ordre, sans jamais sauter une marche.
\end{enclair}

\section*{Pour qui est ce livre}

Pour toute personne qui veut trois choses à la fois :

\begin{enumerate}
  \item \textbf{comprendre} ce que fait V-JEPA~2.1, et pourquoi il fait
        mieux que le modèle qui existait avant lui ;
  \item \textbf{écrire elle-même} le modèle et le programme qui l'entraîne ;
  \item \textbf{mettre le modèle au travail} pour de vrai, tous les jours,
        au service d'une application.
\end{enumerate}

Aucun diplôme n'est demandé. Vous n'avez pas besoin de savoir ce qu'est un
réseau de neurones, ni un transformeur, ni même une image numérique : tout
cela est expliqué à partir de rien, dans les quatre premiers chapitres.

\section*{Comment ce livre est fabriqué}

Huit règles ont guidé chaque page, et elles ne sont pas des intentions : ce
sont des engagements que vous pouvez vérifier.

\begin{enumerate}
  \item Aucun mot technique n'apparaît sans avoir été défini juste avant.
  \item On commence toujours par les prérequis, jamais par la nouveauté.
  \item Un dessin pour chaque idée qui bouge, qui circule ou qui se compare.
  \item Un exemple après chaque notion, avec des chiffres quand c'est
        possible.
  \item Des mots simples, ceux de tous les jours.
  \item Le texte vous parle et vous accompagne.
  \item Une phrase de transition à la fin de chaque section, qui annonce la
        suite.
  \item Des phrases entières, comme dans un vrai livre.
\end{enumerate}

Le texte utilise quelques encadrés qui reviennent souvent, faciles à
reconnaître à leur pictogramme :

\begin{itemize}
  \item \textbf{\faIcon{comments}~En clair} — l'idée expliquée avec des
        mots de tous les jours.
  \item \textbf{\faIcon{brain}~L'intuition} — l'image mentale à garder en
        tête.
  \item \textbf{\faIcon{robot}~À quoi ça sert en IA ?} — où cette notion
        sert vraiment dans V-JEPA~2.1.
  \item \textbf{\faIcon{chart-line}~Sur des données réelles} — la même idée,
        appliquée aux chiffres réels de l'article.
  \item \textbf{\faIcon{seedling}~Ce qu'il faut savoir avant} — les seuls
        prérequis de la section.
  \item \textbf{\faIcon{exclamation-triangle}~Piège à éviter} — l'erreur
        classique que tout le monde fait.
  \item \textbf{\faIcon{key}~À retenir} — l'essentiel, si vous ne deviez
        garder qu'une chose.
  \item \textbf{\faIcon{rocket}~Pour aller plus loin} — un pont vers un
        sujet plus avancé, que vous pouvez sauter en première lecture.
\end{itemize}

Le code, lui, se présente de deux façons, et il ne faut pas les confondre.
Ce que \emph{vous} écrivez apparaît sur fond clair, avec ses numéros de
ligne. Ce que la \emph{machine} répond apparaît dans une petite fenêtre
noire, comme celle d'un vrai terminal, avec ses trois pastilles de couleur
en haut à gauche. Vous saurez ainsi, d'un coup d'œil et sans lire le titre,
si vous avez sous les yeux une instruction à taper ou une réponse à
comparer avec la vôtre.

Chaque chapitre se termine par des exercices, tous corrigés. Ils ne servent
pas à vous noter : ils servent à vous montrer, à vous-même, que vous avez
compris.

\section*{Un mot d'honnêteté}

L'article explique comment \emph{entraîner} le modèle et comment
l'\emph{évaluer}. Il ne parle presque pas de sa \emph{mise en service} : pas
de serveur, pas de temps de réponse, pas de surveillance. La dernière partie
de ce livre sera donc de deux couleurs, et nous vous le dirons franchement à
chaque fois :

\begin{itemize}
  \item \textbf{Ce que l'article dit} — la taille des modèles, la résolution
        des images, la vitesse de calcul. Ce sont des résultats mesurés.
  \item \textbf{Ce que nous ajoutons} — le serveur, la mesure du temps de
        réponse, la surveillance. Ce sont des conseils de métier, utiles et
        éprouvés, mais qui ne viennent pas de l'article.
\end{itemize}

Nous dirons aussi, aussi clairement que ses réussites, ce que ce modèle ne
sait pas faire, et les tâches sur lesquelles d'autres modèles restent
meilleurs que lui. Un livre qui ne raconte que les victoires ne vous aide
pas à choisir.

\vspace{1em}
\hfill\textit{Bonne lecture.}
\clearpage


\cleardoublepage
\pdfbookmark[0]{Table des matières}{toc}
\tableofcontents

%----------------------------------------------------------------------
% Corps du livre
%----------------------------------------------------------------------
\mainmatter

% Progression voulue : on ne suppose rien au départ, et chaque partie ne
% s'appuie que sur les précédentes. Les chapitres sont décommentés au fur
% et à mesure de leur rédaction.

\part{Les fondations : partir vraiment de zéro}
\chapter{Une image, une vidéo, vues par un ordinateur}
\label{chap:image-video}

\begin{quote}\itshape
Avant de parler d'intelligence artificielle, de vidéos cachées et de
modèles à deux milliards de réglages, mettons-nous d'accord sur une chose
toute simple : qu'est-ce qu'une image, pour une machine ? Ce chapitre pose
le décor, tranquillement, sans rien supposer.
\end{quote}

\begin{prerequis}
Rien du tout. Vraiment rien. Il faut savoir compter, savoir qu'une
multiplication sert à répéter une addition, et c'est tout. Chaque mot
nouveau sera expliqué au moment où il arrive.
\end{prerequis}

Vous regardez une photo de chat. Vous voyez un chat. C'est immédiat, c'est
gratuit, vous n'avez rien eu à faire. Un ordinateur, lui, ne voit pas de
chat. Il ne voit même pas d'image. Il voit une très longue liste de
nombres, et rien d'autre.

Tout ce livre part de là. Si nous voulons qu'une machine finisse par
comprendre une vidéo, il faut d'abord savoir sous quelle forme la vidéo
arrive chez elle. C'est exactement ce que nous allons faire dans ce
chapitre : ouvrir le capot et regarder les nombres.

\section{Un ordinateur ne voit pas des objets, il voit des nombres}

\begin{enclair}
Une image, pour un ordinateur, c'est une grille de petits carrés. Chaque
carré porte un nombre, et ce nombre dit une seule chose : à quel point ce
petit carré est clair ou sombre.
\end{enclair}

Pensez à une mosaïque, comme celles qu'on voit sur les murs des vieilles
maisons. De loin, vous voyez un dessin. De près, vous voyez des centaines
de petits carreaux collés les uns aux autres, chacun d'une seule couleur.
Une image numérique, c'est exactement cela : une mosaïque de petits
carrés, tous de la même taille, rangés en lignes et en colonnes.

\begin{definition}[Le pixel][pixel]
Un \emph{pixel} est un de ces petits carrés. C'est le plus petit morceau
d'une image : on ne peut pas le couper en deux. Le mot vient de l'anglais
\emph{picture element}, c'est-à-dire « élément d'image ». Chaque pixel
porte un nombre qui décrit sa couleur.
\end{definition}

Commençons par le cas le plus simple : une image en \emph{niveaux de
gris}, c'est-à-dire une image sans couleur, comme une vieille photo en
noir et blanc. Dans ce cas, un seul nombre suffit par pixel. On se met
d'accord une fois pour toutes sur ce que ce nombre veut dire :

\begin{itemize}
  \item le nombre \textbf{0} veut dire « ce carré est complètement noir » ;
  \item le nombre \textbf{255} veut dire « ce carré est complètement
        blanc » ;
  \item tous les nombres entre les deux sont des gris, de plus en plus
        clairs à mesure qu'on monte.
\end{itemize}

La figure~\ref{fig:image-grille} montre les deux faces de la même image : à
gauche ce que la machine a réellement en mémoire, à droite ce que nous
voyons apparaître à l'écran.

\begin{schema}{Les deux visages d'une même image. À gauche, les seize
nombres que l'ordinateur garde en mémoire. À droite, le dessin obtenu en
peignant chaque carré : petit nombre, carré sombre ; grand nombre, carré
clair. Il n'y a rien d'autre dans une image que ces nombres.}[image-grille]
  \begin{scope}[x=1cm, y=1cm]
    \node[font=\sffamily\small\bfseries, text=mbSlate] at (1.8,0.62)
         {Ce que la machine garde};
    \foreach \mbrow [count=\ri from 0] in
      {{240,235,210,245},{235,120,60,210},{215,60,20,120},{245,205,115,235}}{
      \foreach \mbval [count=\ci from 0] in \mbrow {
        \draw[draw=mbSlate!45, line width=0.5pt]
              (\ci*0.9,-\ri*0.9) rectangle ++(0.9,-0.9);
        \node[font=\small, text=mbInk]
              at (\ci*0.9+0.45,-\ri*0.9-0.45) {\mbval};
      }
    }
    \draw[fleche accent] (3.8,-1.8) -- (5.2,-1.8);
    \node[etiquette, above] at (4.5,-1.74) {on affiche};
    \begin{scope}[xshift=5.4cm]
      \foreach \mbrow [count=\ri from 0] in
        {{240,235,210,245},{235,120,60,210},{215,60,20,120},{245,205,115,235}}{
        \foreach \mbval [count=\ci from 0] in \mbrow {
          \pgfmathsetmacro{\mblum}{100-\mbval/2.55}
          \fill[black!\mblum] (\ci*0.9,-\ri*0.9) rectangle ++(0.9,-0.9);
        }
      }
      \draw[draw=mbSlate!45, line width=0.5pt] (0,0) rectangle (3.6,-3.6);
      \node[font=\sffamily\small\bfseries, text=mbSlate] at (1.8,0.62)
           {Ce que nous voyons};
    \end{scope}
  \end{scope}
\end{schema}

\begin{exemple}[Lire trois pixels]
Dans la grille de gauche de la figure~\ref{fig:image-grille}, le pixel tout en haut
à gauche vaut $240$ : il est presque blanc. Le pixel du milieu vaut $20$ :
il est presque noir. Un pixel qui vaudrait $128$, soit à peu près la
moitié de $255$, serait un gris moyen, ni clair ni foncé.
\end{exemple}

\begin{intuition}
Voyez chaque pixel comme un petit interrupteur à gradateur, du genre de
ceux qui règlent la lumière d'un salon. Tourné à fond dans un sens, la
lampe est éteinte : c'est $0$. Tourné à fond dans l'autre, la lampe éclaire
au maximum : c'est $255$. Une image, c'est un mur couvert de milliers de
ces petites lampes, chacune réglée à sa valeur.
\end{intuition}

\begin{piege}
Pourquoi $255$ et pas $100$, ou $1000$ ? C'est une simple convention, née
de la façon dont les ordinateurs rangent les nombres : avec huit cases de
mémoire, on peut écrire exactement $256$ valeurs différentes, de $0$ à
$255$. Retenez surtout le sens de lecture : \textbf{plus le nombre est
grand, plus le pixel est clair}. Beaucoup de débutants inversent, dessinent
des images en négatif, et cherchent l'erreur pendant une heure.
\end{piege}

Nous savons maintenant ce que contient une case de la mosaïque. La question
suivante vient toute seule : combien y a-t-il de cases ?

\section{La résolution : la finesse de la mosaïque}

\begin{definition}[La résolution][resolution]
La \emph{résolution} d'une image, c'est le nombre de pixels qu'elle
contient, donné sous la forme « largeur $\times$ hauteur ». Une image de
résolution $256 \times 256$ possède $256$ colonnes de pixels et $256$
lignes de pixels.
\end{definition}

Plus la mosaïque a de carreaux, plus le dessin est fin. Avec très peu de
carreaux, un rond ressemble à un escalier ; avec beaucoup, il redevient un
rond. La figure~\ref{fig:resolution} montre le même disque dessiné deux fois, la
première fois grossièrement, la seconde finement.

\begin{schema}{Le même disque, dessiné avec deux finesses différentes. À
gauche, la mosaïque est grossière et les bords deviennent des marches
d'escalier. À droite, les carreaux sont seize fois plus nombreux et l'\oe{}il
retrouve un vrai rond. Rien n'a changé dans le dessin : seule la finesse de
la grille a changé.}[resolution]
  \begin{scope}[x=1cm, y=1cm]
    % --- grille grossière : 6 x 6 ---
    \foreach \ri in {0,...,5}{
      \foreach \ci in {0,...,5}{
        \pgfmathsetmacro{\mbcx}{(\ci+0.5)/6-0.5}
        \pgfmathsetmacro{\mbcy}{(\ri+0.5)/6-0.5}
        \pgfmathsetmacro{\mbdd}{\mbcx*\mbcx+\mbcy*\mbcy}
        \pgfmathsetmacro{\mbin}{\mbdd<0.2100 ? 1 : 0}
        \ifdim\mbin pt>0.5pt
          \fill[mbBlue!65] (\ci*0.5,-\ri*0.5) rectangle ++(0.5,-0.5);
        \fi
        \draw[draw=mbSlate!30, line width=0.4pt]
              (\ci*0.5,-\ri*0.5) rectangle ++(0.5,-0.5);
      }
    }
    \node[font=\sffamily\small, text=mbSlate] at (1.5,-3.35)
         {$6 \times 6 = 36$ pixels};
    % --- grille fine : 24 x 24 ---
    \begin{scope}[xshift=5cm]
      \foreach \ri in {0,...,23}{
        \foreach \ci in {0,...,23}{
          \pgfmathsetmacro{\mbcx}{(\ci+0.5)/24-0.5}
          \pgfmathsetmacro{\mbcy}{(\ri+0.5)/24-0.5}
          \pgfmathsetmacro{\mbdd}{\mbcx*\mbcx+\mbcy*\mbcy}
          \pgfmathsetmacro{\mbin}{\mbdd<0.2100 ? 1 : 0}
          \ifdim\mbin pt>0.5pt
            \fill[mbBlue!65] (\ci*0.125,-\ri*0.125) rectangle ++(0.125,-0.125);
          \fi
        }
      }
      \draw[draw=mbSlate!30, line width=0.4pt] (0,0) rectangle (3,-3);
      \node[font=\sffamily\small, text=mbSlate] at (1.5,-3.35)
           {$24 \times 24 = 576$ pixels};
    \end{scope}
  \end{scope}
\end{schema}

Le nombre total de pixels s'obtient par une simple multiplication : on
multiplie la largeur par la hauteur, comme on calculerait le nombre de
carreaux d'un sol.

\begin{exemple}[Compter les pixels d'une vraie image]
Une image de résolution $256 \times 256$, qui est la taille utilisée
pendant la plus grande partie de l'entraînement de V-JEPA~2.1, contient
\[
  256 \times 256 = 65\,536 \text{ pixels.}
\]
Soit soixante-cinq mille petits carrés. Et ce n'est encore qu'une image en
gris : nous n'avons même pas parlé de couleur.
\end{exemple}

\begin{aretenir}
Deux images peuvent montrer exactement la même chose et ne pas avoir la
même résolution. Quand la résolution monte, le dessin gagne en finesse,
mais le nombre de nombres à traiter grimpe très vite : il grimpe comme le
\emph{carré} de la taille. Doubler la largeur et la hauteur, c'est
multiplier le travail par quatre. Cette remarque, qui paraît anodine ici,
expliquera plus tard pourquoi la dernière phase d'entraînement de
V-JEPA~2.1 est si courte.
\end{aretenir}

Nos images sont encore grises. Ajoutons la couleur : vous allez voir que
cela ne change presque rien à l'histoire.

\section{La couleur, c'est trois grilles au lieu d'une}

Un écran ne sait fabriquer que trois couleurs de lumière : du rouge, du
vert et du bleu. Toutes les autres couleurs sont obtenues en mélangeant ces
trois-là, en dosant chacune plus ou moins fort. Du rouge à fond plus du
vert à fond donne du jaune. Les trois à fond donnent du blanc. Les trois
éteintes donnent du noir.

Un pixel en couleur a donc besoin non pas d'un nombre, mais de trois : un
pour le rouge, un pour le vert, un pour le bleu. Chacun se lit comme avant,
de $0$ (éteint) à $255$ (à fond).

\begin{definition}[Le canal][canal]
Un \emph{canal} est une grille complète de nombres qui ne s'occupe que
d'une seule couleur de base. Une image en couleur possède trois canaux
superposés : le canal rouge, le canal vert et le canal bleu. On les appelle
souvent « les canaux R, V, B », ou « RGB » à l'anglaise.
\end{definition}

\begin{schema}{Une image en couleur, c'est trois mosaïques posées les unes
sur les autres. La première dit combien il y a de rouge dans chaque carré,
la deuxième combien de vert, la troisième combien de bleu. En regardant les
trois ensemble, l'\oe{}il reconstitue la couleur.}[canaux-rvb]
  \begin{scope}[x=1cm, y=1cm]
    \foreach \mbcol [count=\ri from 0] in {red, green, blue}{
      \begin{scope}[shift={(\ri*0.55,-\ri*0.5)}]
        \fill[\mbcol!18] (0,0) rectangle (2.4,2.4);
        \draw[\mbcol!70!black, line width=0.8pt] (0,0) rectangle (2.4,2.4);
        \foreach \mbk in {0.6,1.2,1.8}{
          \draw[\mbcol!35!white, line width=0.4pt] (\mbk,0) -- (\mbk,2.4);
          \draw[\mbcol!35!white, line width=0.4pt] (0,\mbk) -- (2.4,\mbk);
        }
      \end{scope}
    }
    \node[font=\sffamily\scriptsize\bfseries, text=red!60!black]
         at (2.75,2.15) {canal rouge};
    \node[font=\sffamily\scriptsize\bfseries, text=green!45!black]
         at (3.30,1.65) {canal vert};
    \node[font=\sffamily\scriptsize\bfseries, text=blue!60!black]
         at (3.85,1.15) {canal bleu};
    \draw[fleche accent] (4.9,0.7) -- (6.1,0.7);
    \node[etiquette, above] at (5.5,0.78) {on mélange};
    \begin{scope}[xshift=6.4cm, yshift=-0.5cm]
      \foreach \mbrow [count=\ri from 0] in
        {{blue!25,blue!25,blue!30,blue!25},
         {blue!25,orange!70,orange!80,blue!30},
         {blue!30,orange!80,orange!60,orange!70},
         {teal!35,teal!40,teal!35,teal!45}}{
        \foreach \mbval [count=\ci from 0] in \mbrow {
          \fill[\mbval] (\ci*0.6,{2.4-\ri*0.6}) rectangle ++(0.6,-0.6);
        }
      }
      \draw[draw=mbSlate!45, line width=0.6pt] (0,0) rectangle (2.4,2.4);
      \node[font=\sffamily\scriptsize\bfseries, text=mbSlate] at (1.2,-0.32)
           {l'image en couleur};
    \end{scope}
  \end{scope}
\end{schema}

\begin{exemple}[Quelques couleurs, en trois nombres]
Voici comment se lisent quelques couleurs courantes, dans l'ordre
rouge, vert, bleu :

\begin{center}
\begin{tabular}{lccc}
\toprule
\textbf{Couleur} & \textbf{Rouge} & \textbf{Vert} & \textbf{Bleu} \\
\midrule
Noir           & 0   & 0   & 0   \\
Blanc          & 255 & 255 & 255 \\
Rouge vif      & 255 & 0   & 0   \\
Jaune          & 255 & 255 & 0   \\
Gris moyen     & 128 & 128 & 128 \\
\bottomrule
\end{tabular}
\end{center}

Vous remarquez que le gris s'obtient en donnant la même valeur aux trois
canaux. C'est logique : aucune couleur ne tire la couverture à elle.
\end{exemple}

Comptons maintenant, pour de bon, le nombre de cases d'une image en
couleur.

\begin{exemple}[Le poids d'une image de $256 \times 256$]
Notre image possède $65\,536$ pixels, et chaque pixel réclame trois
nombres. Au total :
\[
  256 \times 256 \times 3 = 196\,608 \text{ nombres.}
\]
Presque deux cent mille nombres pour une seule petite image carrée, plus
petite que la moitié d'un écran de téléphone. Gardez ce chiffre en tête :
il va grossir très vite.
\end{exemple}

Nous avons beaucoup de nombres sur les bras. Il est temps de les ranger
proprement, et de leur donner un nom.

\section{Ranger les nombres : le tenseur et sa forme}

Pour parler de tous ces nombres sans se perdre, les informaticiens
utilisent un mot unique, qui a l'air savant mais qui ne cache rien de
compliqué : le \emph{tenseur}. Approchons-le par étapes, en montant une
marche à la fois.

\begin{schema}{La même idée, à quatre étages. Un nombre tout seul ; une
liste de nombres ; un tableau de nombres ; puis une pile de tableaux. Le
mot « tenseur » désigne tout cela d'un coup, quel que soit le nombre
d'étages.}[echelle-tenseur]
  \begin{scope}[x=1cm, y=1cm]
    % 1. un nombre
    \draw[draw=mbSlate!50, fill=mbBlue!10, line width=0.6pt]
          (0,-0.75) rectangle (0.5,-1.25);
    \node[font=\scriptsize, text=mbInk] at (0.25,-1) {7};
    \node[font=\sffamily\scriptsize, text=mbSlate, align=center]
          at (0.25,-2.9) {un nombre\\ (0 étage)};
    % 2. une liste
    \begin{scope}[xshift=2cm]
      \foreach \ci in {0,...,3}{
        \draw[draw=mbSlate!50, fill=mbBlue!10, line width=0.6pt]
              (\ci*0.5,-0.75) rectangle ++(0.5,-0.5);
      }
      \node[font=\sffamily\scriptsize, text=mbSlate, align=center]
            at (1,-2.9) {une liste\\ (1 étage)};
    \end{scope}
    % 3. un tableau
    \begin{scope}[xshift=5.2cm]
      \foreach \ri in {0,...,3}{
        \foreach \ci in {0,...,3}{
          \draw[draw=mbSlate!50, fill=mbBlue!10, line width=0.6pt]
                (\ci*0.5,-\ri*0.5) rectangle ++(0.5,-0.5);
        }
      }
      \node[font=\sffamily\scriptsize, text=mbSlate, align=center]
            at (1,-2.9) {un tableau\\ (2 étages)};
    \end{scope}
    % 4. une pile de tableaux
    \begin{scope}[xshift=8.6cm]
      \foreach \mbk in {2,1,0}{
        \begin{scope}[shift={(\mbk*0.28,-\mbk*0.28)}]
          \fill[mbAmber!\the\numexpr14+10*\mbk\relax] (0,0) rectangle (2,-2);
          \draw[draw=mbAmber!80!black, line width=0.6pt] (0,0) rectangle (2,-2);
          \foreach \mbj in {0.5,1,1.5}{
            \draw[draw=mbAmber!45!white, line width=0.4pt] (\mbj,0) -- (\mbj,-2);
            \draw[draw=mbAmber!45!white, line width=0.4pt] (0,-\mbj) -- (2,-\mbj);
          }
        \end{scope}
      }
      \node[font=\sffamily\scriptsize, text=mbSlate, align=center]
            at (1.3,-2.9) {une pile\\ (3 étages)};
    \end{scope}
  \end{scope}
\end{schema}

\begin{definition}[Le tenseur][tenseur]
Un \emph{tenseur} est un tableau de nombres qui peut avoir autant de
directions que l'on veut. Un nombre seul est un tenseur sans direction ;
une liste est un tenseur à une direction ; un tableau à lignes et colonnes
est un tenseur à deux directions ; une pile de tableaux est un tenseur à
trois directions, et ainsi de suite. Chaque direction s'appelle un
\emph{axe}.
\end{definition}

Pour décrire un tenseur, il suffit de dire combien de cases il compte le
long de chaque axe. Cette liste de tailles porte un nom, et c'est le mot le
plus utile de tout le chapitre.

\begin{notation}[La forme d'un tenseur][forme]
La \emph{forme} d'un tenseur est la liste du nombre de cases sur chaque
axe. On l'écrit entre parenthèses. Pour une image en couleur, on écrit
\[
  (\text{hauteur},\ \text{largeur},\ \text{canaux}),
\]
et notre image de tout à l'heure a donc la forme $(256, 256, 3)$.
\end{notation}

\begin{exemple}[Lire une forme]
Un tenseur de forme $(64, 384, 384, 3)$ se lit de gauche à droite comme une
phrase : « soixante-quatre choses, chacune haute de $384$ cases, large de
$384$ cases, et épaisse de $3$ canaux ». Nous verrons dans un instant que
ces soixante-quatre choses sont les images d'une vidéo.
\end{exemple}

\begin{piege}
L'ordre des axes n'est pas le même partout, et c'est la source d'erreur
numéro un des débutants. Deux conventions se partagent le monde :

\begin{itemize}
  \item $(\text{hauteur},\ \text{largeur},\ \text{canaux})$ — l'ordre des
        outils d'affichage et de photo ;
  \item $(\text{canaux},\ \text{hauteur},\ \text{largeur})$ — l'ordre des
        bibliothèques d'apprentissage comme PyTorch.
\end{itemize}

Les nombres sont exactement les mêmes, seule leur façon d'être rangés
change. Mais si vous vous trompez, votre programme croira qu'une image de
$256 \times 256$ en couleur est une image de $3$ pixels de haut sur $256$ de
large avec $256$ canaux, et vous obtiendrez une bouillie. Prenez
l'habitude, dès aujourd'hui, d'écrire la forme attendue en commentaire
au-dessus de chaque ligne de code qui manipule une image.
\end{piege}

Nous savons ranger une image. Passons à ce qui nous intéresse vraiment dans
ce livre : la vidéo.

\section{Une vidéo, c'est une pile d'images}

Une vidéo n'est rien d'autre qu'une suite d'images fixes, montrées les unes
après les autres, très vite. Si elles défilent assez vite, notre \oe{}il ne
voit plus des images séparées : il voit du mouvement. C'est le principe du
petit carnet à feuilleter, celui où l'on dessine un bonhomme dans le coin
de chaque page et où l'on fait défiler les pages avec le pouce.

\begin{definition}[L'image d'une vidéo][frame]
Chacune des images fixes qui composent une vidéo s'appelle une
\emph{image}, ou une \emph{trame} ; en anglais on dit \emph{frame}, et ce
mot revient tellement souvent dans les articles et dans le code qu'il faut
le connaître. Une vidéo est donc une pile d'images empilées non pas dans
l'espace, mais dans le temps.
\end{definition}

\begin{definition}[Les images par seconde][fps]
Le nombre d'\emph{images par seconde} dit combien d'images sont montrées
pendant une seconde. On l'abrège souvent en \emph{fps}, de l'anglais
\emph{frames per second}. Un film de cinéma tourne à $24$ images par
seconde ; les vidéos de téléphone, souvent à $30$ ou $60$.
\end{definition}

\begin{schema}{Une vidéo vue comme une pile de feuilles. Chaque feuille est
une image complète, avec ses trois canaux de couleur. La flèche du bas est
le temps : elle avance d'un cran à chaque nouvelle feuille. Ici, quatre
images par seconde, donc un quart de seconde entre deux
feuilles.}[pile-video]
  \begin{scope}[x=1cm, y=1cm]
    \foreach \mbk in {4,3,2,1,0}{
      \begin{scope}[shift={(\mbk*1.35,\mbk*0.42)}]
        \fill[mbBlue!\the\numexpr12+8*\mbk\relax] (0,0) rectangle (1.9,1.35);
        \draw[draw=mbBlue!75, line width=0.7pt] (0,0) rectangle (1.9,1.35);
        \foreach \mbj in {0.475,0.95,1.425}{
          \draw[draw=mbBlue!30!white, line width=0.35pt] (\mbj,0) -- (\mbj,1.35);
        }
        \foreach \mbj in {0.45,0.9}{
          \draw[draw=mbBlue!30!white, line width=0.35pt] (0,\mbj) -- (1.9,\mbj);
        }
      \end{scope}
    }
    \draw[-{Stealth[round,length=2.6mm,width=2mm]}, draw=mbAmber!85!black,
          line width=1pt] (0.2,-0.55) -- (7.6,-0.55);
    \foreach \mbk in {0,1,2,3,4}{
      \draw[draw=mbAmber!85!black, line width=0.7pt]
            (\mbk*1.35+0.95,-0.45) -- (\mbk*1.35+0.95,-0.65);
    }
    \node[font=\sffamily\scriptsize, text=mbSlate] at (0.95,-0.95) {$0$~s};
    \node[font=\sffamily\scriptsize, text=mbSlate] at (2.30,-0.95) {$0{,}25$~s};
    \node[font=\sffamily\scriptsize, text=mbSlate] at (3.65,-0.95) {$0{,}50$~s};
    \node[font=\sffamily\scriptsize, text=mbSlate] at (5.00,-0.95) {$0{,}75$~s};
    \node[font=\sffamily\scriptsize, text=mbSlate] at (6.35,-0.95) {$1$~s};
    \node[font=\sffamily\small\bfseries, text=mbAmber!80!black]
          at (7.9,-0.95) {temps};
    \node[annotation bulle] (mbune) at (-0.9,2.6)
         {une image\\ $(H, W, 3)$};
    \draw[trait annotation] (mbune) to[bend right=15] (0.55,1.0);
  \end{scope}
\end{schema}

Puisque la vidéo est une pile d'images, sa forme s'obtient en ajoutant un
axe devant celle d'une image : l'axe du temps.

\begin{notation}[La forme d'une vidéo][forme-video]
Un extrait de vidéo s'écrit sous la forme
\[
  (\text{temps},\ \text{hauteur},\ \text{largeur},\ \text{canaux}),
\]
où le premier nombre est le nombre d'images de l'extrait. On appelle
souvent un tel extrait un \emph{clip} : c'est un petit morceau découpé dans
une vidéo plus longue, exactement comme on découpe un morceau dans une
bobine de film.
\end{notation}

La durée d'un clip se calcule sans effort. Si vous avez un certain nombre
d'images, et que vous savez combien d'images passent en une seconde, alors
la division vous donne les secondes :
\[
  \text{durée en secondes} \;=\; \frac{\text{nombre d'images}}
                                      {\text{images par seconde}}.
\]

\begin{exemple}[La durée d'un clip d'entraînement]
Pendant la plus grande partie de son entraînement, V-JEPA~2.1 regarde des
clips de $16$ images prises à raison de $4$ images par seconde. Ces clips
durent donc
\[
  \frac{16}{4} = 4 \text{ secondes.}
\]
Quatre secondes, c'est court, mais c'est bien assez pour voir un geste
commencer et se terminer : une main qui saisit une tasse, une porte qui
s'ouvre. Et c'est justement ce que le modèle doit apprendre à comprendre.
\end{exemple}

\begin{surlesdonnees}
L'article de V-JEPA~2.1 utilise deux réglages, l'un après l'autre. Pendant
la longue phase d'entraînement principale, les clips font $16$ images en
$256 \times 256$. Pendant la courte phase finale, appelée phase de
refroidissement, on monte à $64$ images en $384 \times 384$. Toujours à
quatre images par seconde, ces clips-là durent
\[
  \frac{64}{4} = 16 \text{ secondes.}
\]
Le modèle regarde donc, à la fin, des extraits quatre fois plus longs et
beaucoup plus fins. Nous verrons au chapitre consacré au refroidissement
pourquoi cet ordre — d'abord grossier et court, ensuite fin et long — n'est
pas un caprice.
\end{surlesdonnees}

Comptons enfin les nombres contenus dans un clip, en multipliant tout
simplement les quatre tailles entre elles.

\begin{exemple}[Le nombre de cases d'un clip]
Pour un clip de $16$ images en $256 \times 256$ couleur :
\[
  16 \times 256 \times 256 \times 3 = 3\,145\,728 \text{ nombres.}
\]
Plus de trois millions. Et pour un clip de la phase de refroidissement,
$64$ images en $384 \times 384$ :
\[
  64 \times 384 \times 384 \times 3 = 28\,311\,552 \text{ nombres,}
\]
soit environ vingt-huit millions, c'est-à-dire neuf fois plus que le
précédent.
\end{exemple}

\begin{aretenir}
Une image est un tenseur à trois axes, une vidéo un tenseur à quatre axes.
Rien de plus mystérieux : on a seulement ajouté le temps devant. Toutes les
formes que vous croiserez dans ce livre se lisent avec cette seule règle.
\end{aretenir}

Trois millions de nombres, cela reste abstrait. Regardons ce que cela
représente vraiment dans la mémoire d'une machine.

\section{Ce que tout cela pèse vraiment}

Un ordinateur range ses nombres dans des cases de mémoire. L'unité de
mesure de cette mémoire est l'\emph{octet}.

\begin{definition}[L'octet][octet]
Un \emph{octet} est un petit paquet de mémoire, celui qui suffit tout juste
à écrire un nombre entier compris entre $0$ et $255$. C'est aussi l'unité
dans laquelle on mesure la taille des fichiers. Mille octets font environ un
kilo-octet, un million d'octets font environ un méga-octet (noté Mo), un
milliard d'octets font environ un giga-octet (noté Go).
\end{definition}

Nos pixels vont de $0$ à $255$ : ils tiennent donc chacun dans un seul
octet. Mais dès qu'un modèle se met à calculer, il a besoin de nombres à
virgule — parce qu'une moyenne, une multiplication ou une petite correction
tombent rarement sur un compte rond. Ces nombres à virgule s'appellent des
\emph{nombres flottants}, et le format le plus courant en occupe quatre
octets chacun, soit quatre fois plus.

\begin{exemple}[Du fichier à la mémoire de calcul]
Notre clip de $3\,145\,728$ nombres pèse :

\begin{itemize}
  \item $3\,145\,728$ octets, soit environ \textbf{3 Mo}, tant que les
        pixels restent des entiers de $0$ à $255$ ;
  \item $3\,145\,728 \times 4 = 12\,582\,912$ octets, soit environ
        \textbf{12,6 Mo}, une fois convertis en nombres à virgule pour le
        calcul.
\end{itemize}
\end{exemple}

Or un modèle ne regarde jamais un seul clip à la fois : il en regarde tout
un paquet d'un coup, pour aller plus vite. Ce paquet s'appelle un
\emph{lot}, et sa taille est un réglage du programme d'entraînement.

\begin{exemple}[Le poids d'un lot entier]
V-JEPA~2.1 traite $128$ clips à la fois. Pour la phase principale, cela
donne
\[
  128 \times 12\,582\,912 \approx 1\,610\,000\,000 \text{ octets,}
\]
soit environ \textbf{1,6 Go} — rien que pour les vidéos d'entrée, avant
même le moindre calcul. Pour la phase de refroidissement, où chaque clip
pèse neuf fois plus, on arrive à environ \textbf{14,5 Go}. Voilà, en une
ligne de calcul, pourquoi cette phase-là ne dure que très peu de temps :
elle coûte cher.
\end{exemple}

\begin{enia}
Ce chapitre a l'air de ne parler que de pixels, mais il pose déjà les trois
problèmes que tout le reste du livre s'emploie à résoudre.

\begin{itemize}
  \item \textbf{Il y a beaucoup trop de nombres.} Aucun modèle ne peut
        traiter trois millions de valeurs une par une. Il faudra donc
        d'abord les regrouper en petits carrés de $16 \times 16$ pixels :
        c'est ce que nous ferons au chapitre~4, consacré au transformeur de
        vision.
  \item \textbf{La finesse coûte cher.} Chaque fois que la résolution
        monte, la facture explose. C'est pour cela que l'entraînement
        commence en basse résolution et ne monte qu'à la toute fin.
  \item \textbf{Lire une vidéo prend du temps.} Une vidéo est stockée
        compressée ; il faut la décompresser image par image avant de
        pouvoir en faire quoi que ce soit. En production, ce travail-là
        devient souvent plus lent que le modèle lui-même. Nous y
        reviendrons longuement dans la partie consacrée à la mise en
        service.
\end{itemize}
\end{enia}

\begin{plusloin}
Dans le même temps où il regarde $128$ clips, V-JEPA~2.1 regarde aussi
$2\,304$ images fixes. En reprenant nos calculs, ces images pèsent
$2\,304 \times 196\,608 \times 4 \approx 1,8$~Go, c'est-à-dire à peu près
autant que les $1,6$~Go de vidéo. Ces deux nombres, $128$ et $2\,304$,
n'ont donc pas été tirés au hasard : ils ont été choisis ensemble pour que
les deux moitiés du travail coûtent à peu près pareil. Nous y reviendrons
au chapitre sur les données.
\end{plusloin}

Assez de calculs sur le papier. Vérifions tout cela sur une vraie machine.

\section{Passons au code}

Nous allons manipuler nos images et nos clips avec Python, le langage
utilisé par presque tous les travaux d'intelligence artificielle, et avec
\co{python}{numpy}, la bibliothèque qui sait manipuler les tenseurs. Vous
n'avez pas besoin de comprendre chaque ligne dès maintenant : lisez d'abord
les commentaires, et comparez les nombres affichés avec ceux que nous
venons de calculer à la main.

\begin{code}{python}[Fabriquer une image et lire sa forme]
import numpy as np

# Une image en couleur de 256 x 256 pixels, entierement noire.
# La forme se lit : (hauteur, largeur, canaux).
image = np.zeros((256, 256, 3), dtype=np.uint8)

print(image.shape)   # -> (256, 256, 3)
print(image.size)    # -> 196608   le nombre total de cases
print(image.dtype)   # -> uint8    un entier de 0 a 255
print(image.nbytes)  # -> 196608   un octet par case

# Peignons le pixel de la ligne 10, colonne 20, en rouge vif.
image[10, 20] = [255, 0, 0]
print(image[10, 20])  # -> [255   0   0]
\end{code}

Le mot \co{python}{uint8} se lit « entier non signé sur huit bits » : c'est
exactement le type de nombre qui va de $0$ à $255$, celui dont nous avons
parlé. Passons maintenant à la vidéo.

\begin{code}{python}[Un clip, sa forme, sa durée et son poids]
import numpy as np

# Un clip de 16 images de 256 x 256 en couleur.
# La forme se lit : (temps, hauteur, largeur, canaux).
clip = np.zeros((16, 256, 256, 3), dtype=np.uint8)

print(clip.shape)   # -> (16, 256, 256, 3)
print(clip.size)    # -> 3145728   les trois millions annonces

images_par_seconde = 4.0
duree = clip.shape[0] / images_par_seconde
print(duree, "secondes")            # -> 4.0 secondes

# Le meme clip converti en nombres a virgule, ramenes entre 0 et 1.
clip_flottant = clip.astype(np.float32) / 255.0
print(clip_flottant.nbytes / 1e6, "Mo")   # -> 12.582912 Mo
\end{code}

Enfin, voici le changement d'ordre des axes dont nous parlions plus haut.
C'est une opération que vous ferez des centaines de fois, alors autant la
voir tout de suite.

\begin{code}{python}[Passer de l'ordre « affichage » à l'ordre « PyTorch »]
import numpy as np
import torch

clip = np.zeros((16, 256, 256, 3), dtype=np.uint8)   # (T, H, W, C)

x = torch.from_numpy(clip).float() / 255.0
x = x.permute(3, 0, 1, 2)                            # (C, T, H, W)

print(x.shape)   # -> torch.Size([3, 16, 256, 256])

# Aucun nombre n'a ete perdu : seul l'ordre des axes a change.
print(x.numel())  # -> 3145728
\end{code}

\begin{piege}
Remarquez la dernière ligne : le nombre de cases n'a pas bougé. Changer
l'ordre des axes ne crée et ne détruit aucune donnée, cela réorganise
seulement la façon de les parcourir. Si jamais ce nombre change après une
telle opération, c'est que vous avez commis une erreur ailleurs — et cette
vérification en une ligne vous fera gagner des heures.
\end{piege}

\begin{aretenir}
Une image est une grille de nombres ; la couleur ajoute une troisième
dimension ; la vidéo en ajoute une quatrième, le temps. Tout le reste de ce
livre travaille sur ces tenseurs-là, et sur rien d'autre. Quand vous serez
perdu dans un chapitre plus difficile, revenez à cette question toute
simple : « quelle est la forme de ce que je manipule ? » Elle remet presque
toujours les idées en place.
\end{aretenir}

Nous savons désormais sous quelle forme une image et une vidéo arrivent
dans une machine. Mais nous n'avons encore rien appris à personne : nos
tenseurs restent inertes. La question suivante est donc la bonne :
qu'est-ce que cela veut dire, au juste, qu'une machine « apprend » ? C'est
tout le sujet du chapitre suivant.

\section*{Exercices}
\addcontentsline{toc}{section}{Exercices}

\begin{exercice}[Compter les cases]
Combien de nombres contient une image en couleur de $512 \times 512$
pixels ? C'est la taille des images utilisées pendant la phase de
refroidissement de V-JEPA~2.1.
\end{exercice}

\begin{solution}
On multiplie la hauteur, la largeur et le nombre de canaux :
$512 \times 512 \times 3 = 262\,144 \times 3 = 786\,432$ nombres. C'est
quatre fois plus qu'une image de $256 \times 256$, ce qui est normal
puisque nous avons doublé les deux côtés.
\end{solution}

\begin{exercice}[La durée d'un clip]
Un clip contient $32$ images, prises à raison de $8$ images par seconde.
Combien de temps dure-t-il ? Et combien de temps durerait-il si les mêmes
$32$ images étaient prises à $4$ images par seconde ?
\end{exercice}

\begin{solution}
Dans le premier cas, $32 / 8 = 4$ secondes. Dans le second,
$32 / 4 = 8$ secondes. Le nombre d'images est le même, mais elles sont
prélevées plus loin les unes des autres : le clip couvre donc une tranche
de temps deux fois plus longue. C'est un réglage important, car il décide
de ce que le modèle a le droit de voir.
\end{solution}

\begin{exercice}[Le poids d'un lot]
Un lot contient $64$ clips de $16$ images en $256 \times 256$ couleur,
convertis en nombres à virgule de quatre octets. Combien de méga-octets
occupe ce lot ?
\end{exercice}

\begin{solution}
Un clip contient $16 \times 256 \times 256 \times 3 = 3\,145\,728$ nombres,
donc $3\,145\,728 \times 4 = 12\,582\,912$ octets, soit environ $12{,}6$~Mo.
Pour $64$ clips : $64 \times 12{,}6 \approx 806$~Mo, c'est-à-dire un peu
plus de huit cents méga-octets. Ce lot fait exactement la moitié de celui
de l'article, qui en contient $128$, et l'on retrouve bien les $1{,}6$~Go
calculés plus haut en doublant ce résultat.
\end{solution}

\begin{exercice}[Lire une forme]
Un collègue vous tend un tenseur de forme $(3, 64, 384, 384)$ et vous dit
que c'est une vidéo. Que représente chaque nombre ? Quelle est la
convention utilisée, et quelle serait la forme du même tenseur dans l'autre
convention ?
\end{exercice}

\begin{solution}
Le $3$ ne peut être que le nombre de canaux de couleur : aucune vidéo ne
fait trois pixels de haut. Il s'agit donc de l'ordre
$(\text{canaux},\ \text{temps},\ \text{hauteur},\ \text{largeur})$, celui de
PyTorch : une vidéo en couleur de $64$ images en $384 \times 384$. Dans la
convention d'affichage, le même tenseur s'écrirait
$(64, 384, 384, 3)$. Les nombres qu'il contient sont exactement les mêmes ;
seul leur rangement change.
\end{solution}

\chapter{Ce que veut dire « apprendre » pour une machine}
\label{chap:apprendre}

\begin{quote}\itshape
Une image, nous venons de le voir, n'est rien d'autre qu'un grand tas de
nombres. Mais un tas de nombres ne fait rien tout seul. Dans ce chapitre,
nous allons voir comment une machine part d'un réglage pris au hasard, se
trompe, remarque qu'elle s'est trompée, et se corrige toute seule. C'est
cela, et rien d'autre, qu'on appelle « apprendre ».
\end{quote}

\begin{prerequis}
Le chapitre~\ref{chap:image-video}, et seulement lui. Il faut savoir qu'une
image est une grille de nombres, et savoir ce qu'est la forme d'un tenseur.
Pour le reste, il suffit de savoir additionner, multiplier, et diviser. Nous
n'utiliserons aucune formule que nous n'aurons pas d'abord dessinée.
\end{prerequis}

Quand on dit qu'une machine « apprend », on imagine souvent quelque chose de
mystérieux, presque magique. La réalité est beaucoup plus simple, et
beaucoup plus rassurante : une machine qui apprend, c'est une machine qui
essaie, qui se trompe, qui regarde de combien elle s'est trompée, et qui
recommence un tout petit peu mieux. Des milliers de fois. Des millions de
fois.

Tout ce chapitre raconte cette boucle-là. Nous allons la suivre pas à pas
sur un exemple minuscule, que nous calculerons entièrement à la main, avant
de vérifier sur une vraie machine que nos calculs de papier tombent juste.

\section{Une machine qui apprend, c'est une machine à boutons}

\begin{enclair}
Imaginez un appareil avec quelques boutons à tourner. Vous lui posez une
question, il donne une réponse. Si la réponse est mauvaise, vous tournez un
peu les boutons et vous redemandez. Apprendre, pour une machine, c'est
exactement cela : trouver la bonne position des boutons.
\end{enclair}

Vous avez sûrement déjà réglé la température d'une douche. Il y a un
robinet chaud et un robinet froid. Vous les tournez un peu, vous mettez la
main sous l'eau, c'est trop froid, alors vous tournez encore. Au bout de
quelques essais, l'eau est parfaite. Personne ne vous a donné la position
exacte des robinets : vous l'avez trouvée en essayant et en corrigeant.

Une machine qui apprend fait la même chose, sauf qu'elle a beaucoup plus de
robinets, et qu'elle ne trempe pas la main dans l'eau : elle calcule.

\begin{definition}[Le modèle][modele]
Un \emph{modèle} est une machine à calculer qui prend quelque chose en
entrée et rend une réponse en sortie. Ce qu'elle rend ne dépend pas
seulement de l'entrée : cela dépend aussi de la position de ses boutons. Le
même modèle, avec des boutons réglés autrement, donnera une autre réponse à
la même question.
\end{definition}

\begin{definition}[Le paramètre, ou poids][parametre]
Un \emph{paramètre} est un de ces boutons. C'est simplement un nombre, rangé
dans la mémoire de la machine, que l'on a le droit de changer. On l'appelle
aussi un \emph{poids}, parce qu'il décide de l'importance que la machine
donne à telle ou telle information. Régler un modèle, c'est choisir la
valeur de tous ses paramètres.
\end{definition}

La figure~\ref{fig:boutons} montre cette idée dans son entier : une entrée
qui arrive, des boutons qu'on peut tourner, une réponse qui sort.

\begin{schema}{Un modèle, c'est une boîte à calculer munie de boutons. On
lui donne une entrée par la gauche, elle rend une réponse par la droite.
Entre les deux, ce sont les boutons — les paramètres — qui décident du
résultat. Tourner les boutons change la réponse sans rien changer à la
question.}[boutons]
  \begin{scope}[x=1cm, y=1cm]
    % --- la boîte du modèle ---
    \draw[draw=mbBlue, fill=mbBlue!8, line width=1pt, rounded corners=4pt]
          (0,0) rectangle (5,2.5);
    \node[font=\sffamily\small\bfseries, text=mbBlue] at (2.5,0.45)
         {le modèle};
    % --- les trois boutons ---
    % Les boutons sont numérotés, et non nommés a, b, c : le modèle de la
    % page suivante n'en a que deux, et il ne faut pas que le lecteur croie
    % qu'il en manque un.
    \foreach \mbk/\mbang/\mbnom in {0/55/1, 1/-30/2, 2/120/3}{
      \begin{scope}[shift={(1.3+\mbk*1.2, 1.6)}]
        \draw[draw=mbSlate!65, fill=white, line width=0.9pt] (0,0) circle (0.34);
        \draw[draw=mbAmber!85!black, line width=1.3pt] (0,0) -- (\mbang:0.27);
        \fill[mbInk] (0,0) circle (0.045);
        \node[font=\sffamily\scriptsize, text=mbSlate] at (0,-0.62) {\mbnom};
      \end{scope}
    }
    \node[font=\sffamily\scriptsize\itshape, text=mbAmber!75!black]
         at (2.5,2.85) {les boutons, c'est-à-dire les paramètres};
    % --- entrée ---
    \draw[fleche] (-2.1,1.25) -- (-0.12,1.25);
    \node[font=\sffamily\small, text=mbInk, align=center] at (-2.9,1.25)
         {la question\\ $x$};
    % --- sortie ---
    \draw[fleche accent] (5.12,1.25) -- (7.1,1.25);
    \node[font=\sffamily\small, text=mbInk, align=center] at (7.95,1.25)
         {la réponse\\ $\pred{y}$};
  \end{scope}
\end{schema}

Prenons tout de suite un exemple, et gardons-le jusqu'à la fin du chapitre.

\begin{exemple}[La machine à bonbons]
Dans le hall d'une école, il y a une machine à bonbons. On y glisse des
pièces, et elle rend des bonbons. Trois enfants ont essayé et ont noté ce
qu'ils ont obtenu :

\begin{center}
\begin{tabular}{ccc}
\toprule
\textbf{Enfant} & \textbf{Pièces glissées} & \textbf{Bonbons reçus} \\
\midrule
Léa   & 1 & 2 \\
Sami  & 2 & 5 \\
Nina  & 3 & 8 \\
\bottomrule
\end{tabular}
\end{center}

Nous aimerions une machine à calculer qui devine le nombre de bonbons à
partir du nombre de pièces. Appelons $x$ le nombre de pièces et $\pred{y}$
le nombre de bonbons deviné. Le modèle le plus simple du monde s'écrit
\[
  \pred{y} \;=\; a \times x \;+\; b .
\]
Il a exactement deux boutons : $a$ et $b$. Le bouton $a$ dit combien de
bonbons on gagne par pièce ; le bouton $b$ ajoute ou enlève un nombre fixe
de bonbons, quel que soit le nombre de pièces.
\end{exemple}

\begin{intuition}
Le bouton $a$, c'est l'inclinaison de la droite : plus il est grand, plus la
machine est généreuse à chaque pièce. Le bouton $b$, c'est le cadeau de départ : il
décale toute la réponse vers le haut ou vers le bas, même quand on n'a mis
aucune pièce.
\end{intuition}

Vous avez peut-être déjà trouvé la règle de tête : trois bonbons par pièce,
moins un bonbon que la machine garde toujours pour elle. Autrement dit
$a = 3$ et $b = -1$. Vérifions sur les trois enfants : $3 \times 1 - 1 = 2$
pour Léa, $3 \times 2 - 1 = 5$ pour Sami, $3 \times 3 - 1 = 8$ pour Nina.
Tout tombe juste.

\begin{piege}
Nous connaissons la bonne réponse, mais \textbf{la machine, elle, ne la
connaît pas}. C'est tout l'intérêt de l'exercice : nous allons la laisser
partir de boutons réglés à zéro, complètement à côté de la plaque, et la
regarder retrouver toute seule le $3$ et le $-1$. Connaître la réponse à
l'avance nous sert seulement à vérifier qu'elle ne se trompe pas de chemin.
\end{piege}

\begin{definition}[La prédiction][prediction]
Une \emph{prédiction} est la réponse que le modèle donne avec les boutons
tels qu'ils sont réglés en ce moment. On la note $\pred{y}$, avec un petit
chapeau, pour bien la distinguer de la vraie valeur $y$ que l'on a observée
dans la réalité.
\end{definition}

\begin{exemple}[Une première prédiction, très mauvaise]
Réglons les deux boutons à zéro : $a = 0$ et $b = 0$. Pour Léa, qui a mis
une pièce :
\[
  \pred{y} \;=\; 0 \times 1 + 0 \;=\; 0 .
\]
La machine annonce zéro bonbon, alors que Léa en a reçu deux. Elle se trompe
donc de deux bonbons. Pour Sami elle annonce encore zéro, alors qu'il en a
reçu cinq. Et pour Nina, zéro contre huit. Notre modèle est pour l'instant
parfaitement nul, et c'est normal : nous n'avons rien réglé du tout.
\end{exemple}

Nous avons une machine, des boutons, et des réponses fausses. Il nous
manque quelque chose d'essentiel : un moyen de dire \emph{à quel point}
elles sont fausses.

\section{Se tromper, et mesurer de combien}

\begin{enclair}
Pour se corriger, il ne suffit pas de savoir qu'on s'est trompé. Il faut un
nombre qui dise si l'on s'est trompé beaucoup ou un tout petit peu. Ce
nombre s'appelle l'erreur.
\end{enclair}

Reprenons Léa. La machine a dit $0$, la réalité était $2$. L'écart entre les
deux vaut $0 - 2 = -2$. Nous pourrions nous arrêter là et additionner les
écarts des trois enfants. Mais un problème apparaît aussitôt.

\begin{piege}
Additionner les écarts tels quels ne marche pas. Imaginez une machine qui se
trompe de $+5$ bonbons pour Léa et de $-5$ bonbons pour Sami. En
additionnant, on obtient $5 + (-5) = 0$ : on croirait que la machine est
parfaite, alors qu'elle s'est lourdement trompée deux fois ! Les erreurs se
sont annulées entre elles.
\end{piege}

Il faut donc empêcher les écarts de se compenser. La solution retenue
presque partout est de \textbf{mettre chaque écart au carré}, c'est-à-dire de
le multiplier par lui-même. Un carré est toujours positif : $(+5)^2 = 25$ et
$(-5)^2 = 25$ aussi. Plus rien ne s'annule.

\begin{intuition}
Mettre au carré a un deuxième avantage, qui est même le principal : cela
punit très fort les grosses erreurs. Se tromper de $1$ coûte $1$. Se tromper
de $2$ ne coûte pas $2$ mais $4$. Se tromper de $10$ coûte $100$. La machine
va donc s'occuper en priorité de ses plus grosses bêtises, ce qui est
exactement ce qu'on veut.
\end{intuition}

\begin{definition}[L'erreur, ou perte][perte]
L'\emph{erreur} est un seul nombre qui résume à quel point le modèle se
trompe sur l'ensemble des exemples. On l'appelle aussi la \emph{perte}, et
on la note souvent $\mathcal{L}$, de l'anglais \emph{loss}. On la calcule en
mettant chaque écart au carré, puis en prenant la moyenne :
\[
  \mathcal{L} \;=\; \frac{1}{N}\sum_{i=1}^{N}
                    \bigl(\pred{y}_i - y_i\bigr)^{2},
\]
où $N$ est le nombre d'exemples, $\pred{y}_i$ la réponse du modèle sur
l'exemple numéro $i$, et $y_i$ la vraie valeur. Une erreur de $0$ veut dire
que le modèle tombe juste partout. Plus le nombre est grand, plus le modèle
est mauvais.
\end{definition}

\begin{notation}[Lire le signe somme][somme]
Le grand symbole $\sum$ se lit « somme de ». L'écriture
$\sum_{i=1}^{3} z_i$ veut simplement dire $z_1 + z_2 + z_3$ : on fait
prendre à $i$ toutes les valeurs de $1$ à $3$, et on additionne tout. C'est
une façon courte d'écrire une addition longue, rien de plus. Si l'addition
ne comporte que trois termes, vous avez parfaitement le droit de les écrire
à la main : c'est ce que nous ferons.
\end{notation}

\begin{exemple}[L'erreur de notre machine réglée à zéro]
Avec $a = 0$ et $b = 0$, les trois prédictions valent $0$. Calculons les
écarts, puis leurs carrés :

\begin{center}
\begin{tabular}{lccccc}
\toprule
\textbf{Enfant} & $x$ & $y$ & $\pred{y}$ & \textbf{écart } $\pred{y}-y$
& \textbf{carré} \\
\midrule
Léa  & 1 & 2 & 0 & $-2$ & $4$  \\
Sami & 2 & 5 & 0 & $-5$ & $25$ \\
Nina & 3 & 8 & 0 & $-8$ & $64$ \\
\bottomrule
\end{tabular}
\end{center}

La somme des carrés vaut $4 + 25 + 64 = 93$. Comme il y a trois exemples, la
moyenne vaut
\[
  \mathcal{L} \;=\; \frac{93}{3} \;=\; 31 .
\]
Retenez ce nombre : \textbf{l'erreur de départ vaut exactement $31$}. Nous
allons passer le reste du chapitre à la faire descendre.
\end{exemple}

La figure~\ref{fig:erreur-droite} donne à voir ce que ces écarts
représentent : ce sont des distances verticales entre les vrais points et la
droite que dessine le modèle.

\begin{schema}{Les trois enfants sont les points bleus. La droite en trait
plein est la réponse du modèle. Ici, ses deux boutons sont encore à zéro :
elle est couchée sur l'axe et annonce zéro bonbon à tout le monde. Les
traits rouges verticaux sont les écarts, c'est-à-dire de combien la machine
s'est trompée sur chaque enfant : $2$, $5$ et $8$ bonbons. Régler les
boutons, ce sera faire pivoter et monter cette droite jusqu'à ce que les
traits rouges soient les plus courts possible.}[erreur-droite]
  \begin{scope}[x=1.25cm, y=0.36cm]
    % --- axes ---
    \draw[->, draw=mbSlate, line width=0.7pt] (0,0) -- (3.9,0)
          node[right, font=\sffamily\scriptsize, text=mbSlate] {pièces $x$};
    \draw[->, draw=mbSlate, line width=0.7pt] (0,0) -- (0,10.2)
          node[above, font=\sffamily\scriptsize, text=mbSlate] {bonbons};
    \foreach \mbk in {1,2,3}{
      \draw[draw=mbSlate!60, line width=0.6pt] (\mbk,-0.35) -- (\mbk,0.35);
      \node[font=\sffamily\scriptsize, text=mbSlate] at (\mbk,-1.1) {\mbk};
    }
    \foreach \mbk in {2,4,6,8}{
      \draw[draw=mbSlate!60, line width=0.6pt] (-0.07,\mbk) -- (0.07,\mbk);
      \node[font=\sffamily\scriptsize, text=mbSlate, left] at (-0.09,\mbk) {\mbk};
    }
    % --- la droite du modèle, boutons à zéro : y = 0 ---
    \draw[draw=mbAmber!85!black, line width=1.4pt] (0,0) -- (3.6,0);
    \draw[draw=mbAmber!65, line width=0.5pt] (0.45,-1.55) -- (0.45,-0.15);
    \node[font=\sffamily\scriptsize, text=mbAmber!75!black, align=center]
          at (0.45,-2.45) {la droite du modèle,\\ boutons à zéro};
    % --- écarts verticaux ---
    \foreach \mbx/\mby/\mbp/\mbe in {1/2/0/2, 2/5/0/5, 3/8/0/8}{
      \draw[draw=mbRose, line width=1.1pt] (\mbx,\mby) -- (\mbx,\mbp);
      \node[font=\sffamily\scriptsize, text=mbRose, anchor=west]
            at (\mbx+0.08,{(\mby+\mbp)/2}) {\mbe};
    }
    % --- points réels ---
    \foreach \mbx/\mby/\mbnom in {1/2/Léa, 2/5/Sami, 3/8/Nina}{
      \fill[mbBlue] (\mbx,\mby) circle (0.09cm);
      \node[font=\sffamily\scriptsize, text=mbBlue, anchor=north east]
            at (\mbx-0.05,\mby-0.2) {\mbnom};
    }
  \end{scope}
\end{schema}

\begin{aretenir}
Toute la suite du livre repose sur ce couple d'idées : un modèle donne une
réponse, et une erreur mesure de combien cette réponse est fausse.
Apprendre, ce sera toujours et uniquement : \textbf{faire descendre
l'erreur en changeant les paramètres}.
\end{aretenir}

Nous savons mesurer notre bêtise. Reste la question qui compte vraiment :
dans quel sens faut-il tourner les boutons pour que ce nombre baisse ?

\section{De quel côté tourner le bouton ? La pente}

\begin{enclair}
Pour savoir dans quel sens tourner un bouton, il suffit de le pousser un tout
petit peu et de regarder ce que devient l'erreur. Si elle monte, on tournait
du mauvais côté. Si elle descend, on continue.
\end{enclair}

C'est le jeu du « tu chauffes, tu refroidis ». On cache un objet dans une
pièce, l'enfant fait un pas, et on lui dit s'il se rapproche ou s'il
s'éloigne. Avec cette seule information, répétée assez souvent, il finit
toujours par trouver.

Faisons-le pour de vrai sur notre machine à bonbons. Nos boutons sont à
$a = 0$ et $b = 0$, et l'erreur vaut $31$. Poussons le bouton $a$ d'un tout
petit peu, disons de $0{,}01$, sans toucher à $b$.

\begin{exemple}[Pousser un bouton et regarder]
Avec $a = 0{,}01$ et $b = 0$, les trois prédictions deviennent $0{,}01$,
$0{,}02$ et $0{,}03$. Les écarts valent donc $-1{,}99$, $-4{,}98$ et
$-7{,}97$. Leurs carrés valent $3{,}9601$, $24{,}8004$ et $63{,}5209$, dont
la somme fait $92{,}2814$. La nouvelle erreur vaut
\[
  \mathcal{L} \;=\; \frac{92{,}2814}{3} \;\approx\; 30{,}7605 .
\]
Elle est passée de $31$ à $30{,}7605$ : elle a \textbf{baissé} de
$0{,}2395$. Nous poussions donc dans le bon sens. Pour savoir à quelle
vitesse elle baisse, on divise la baisse par la taille de la poussée :
\[
  \frac{-0{,}2395}{0{,}01} \;=\; -23{,}95 .
\]
Ce nombre, $-23{,}95$, dit deux choses à la fois. Son \emph{signe} négatif
dit : « en augmentant $a$, l'erreur diminue ». Sa \emph{taille}, presque
$24$, dit : « et elle diminue vite ».
\end{exemple}

\begin{definition}[La pente][pente]
La \emph{pente} par rapport à un bouton est le nombre qui dit de combien
l'erreur change quand on pousse ce bouton d'une unité. Une pente négative
veut dire que l'erreur descend quand on augmente le bouton. Une pente
positive veut dire qu'elle monte. Une pente proche de zéro veut dire que ce
bouton ne change presque rien.
\end{definition}

\begin{definition}[Le gradient][gradient]
Le \emph{gradient} est la liste de toutes les pentes, une par bouton. S'il y
a deux boutons, le gradient contient deux nombres ; s'il y en a deux
milliards, il en contient deux milliards. Le gradient est donc une sorte de
boussole : il indique, pour chaque bouton à la fois, dans quel sens tourner.
\end{definition}

Pousser chaque bouton à la main pour mesurer sa pente serait épouvantable :
un modèle comme V-JEPA~2.1 a deux milliards de boutons, il faudrait deux
milliards de petites poussées à chaque étape. Heureusement, les
mathématiques savent calculer ces pentes directement, d'un coup, par une
formule. Pour notre droite, les deux formules sont les suivantes :

\begin{notation}[Les deux pentes de notre modèle][pentes-droite]
En notant $e_i = \pred{y}_i - y_i$ l'écart sur l'exemple $i$ :
\[
  \text{pente selon } a \;=\; \frac{2}{N}\sum_i e_i \times x_i,
  \qquad
  \text{pente selon } b \;=\; \frac{2}{N}\sum_i e_i .
\]
\end{notation}

Ces deux formules ont l'air de tomber du ciel. Elles ne sont pourtant pas
mystérieuses, et il vaut la peine de comprendre ce qu'elles racontent.

\begin{intuition}
Regardez la pente selon $b$ : c'est simplement la moyenne des écarts,
multipliée par deux. C'est logique, car le bouton $b$ déplace toute la
droite vers le haut ou vers le bas, de la même quantité pour tout le monde.
Si en moyenne la machine annonce trop de bonbons, il faut baisser $b$.

La pente selon $a$ est la même chose, mais chaque écart y est d'abord
multiplié par $x_i$, le nombre de pièces. C'est l'effet du levier. Le
bouton $a$ fait pivoter la droite autour du point de départ : un enfant
placé loin, comme Nina avec ses trois pièces, monte ou descend trois fois
plus qu'un enfant placé près, comme Léa. Son écart pèse donc trois fois plus
lourd dans la décision.
\end{intuition}

\begin{exemple}[La formule donne-t-elle le même résultat que la poussée ?]
Vérifions. Avec $a = 0$ et $b = 0$, les écarts valent $-2$, $-5$ et $-8$, et
les pièces valent $1$, $2$ et $3$. La pente selon $a$ vaut donc
\begin{align*}
  \frac{2}{3}\bigl[(-2)\times 1 + (-5)\times 2 + (-8)\times 3\bigr]
    &= \frac{2}{3}\bigl[-2 - 10 - 24\bigr] \\[2pt]
    &= \frac{2}{3}\times(-36) = -24 .
\end{align*}
La formule annonce $-24$, et notre petite poussée avait mesuré $-23{,}95$.
Les deux méthodes s'accordent. Le minuscule écart vient de ce que notre
poussée de $0{,}01$ n'était pas infiniment petite ; en poussant encore
moins fort, on collerait encore mieux au $-24$.

Calculons aussi la pente selon $b$ :
\[
  \frac{2}{3}\bigl[(-2) + (-5) + (-8)\bigr]
  = \frac{2}{3}\times(-15) = -10 .
\]
\end{exemple}

Les deux pentes sont négatives. La boussole nous dit donc : augmente $a$,
et augmente $b$ aussi. Voyons maintenant comment transformer ce conseil en
mouvement.

\section{La descente de gradient, calculée à la main}

\begin{enclair}
La règle tient en une phrase : on tourne chaque bouton dans le sens opposé
à sa pente, d'une petite quantité. Puis on recommence.
\end{enclair}

Pourquoi le sens \emph{opposé} ? Parce qu'une pente négative signifie « en
augmentant, l'erreur descend ». Prendre l'opposé du négatif donne du
positif : on augmente. Et une pente positive signifie « en augmentant,
l'erreur monte » ; son opposé nous fait diminuer. Dans les deux cas, on va
vers le bas. C'est la même règle qu'une bille lâchée sur une colline : elle
part toujours dans le sens contraire de la montée.

\begin{definition}[La descente de gradient][descente]
La \emph{descente de gradient} est la méthode qui consiste à répéter
l'opération suivante : calculer l'erreur, calculer le gradient, puis
déplacer chaque paramètre d'un petit pas dans le sens opposé à sa pente.
Pour un paramètre $p$ quelconque, la mise à jour s'écrit
\[
  p_{\text{nouveau}} \;=\; p_{\text{ancien}}
       \;-\; \pas \times \text{pente selon } p ,
\]
où $\pas$ est un petit nombre positif choisi par nous, qui règle la
longueur du pas.
\end{definition}

La figure~\ref{fig:vallee} montre cette idée sous sa forme la plus parlante :
l'erreur dessine une vallée, et le modèle est une bille qui la descend.

\begin{schema}{L'erreur vue comme une vallée. En allant vers la droite ou
vers la gauche, on change la valeur d'un bouton ; la hauteur de la courbe
donne l'erreur correspondante. La bille est notre modèle. La droite en
pointillés est la pente à l'endroit où se trouve la bille : elle monte vers
la droite, donc la pente est positive, donc il faut aller à gauche. Le
meilleur réglage est tout en bas de la vallée, là où la courbe est
plate.}[vallee]
  \begin{scope}[x=1.15cm, y=0.72cm]
    % --- la vallée ---
    \draw[draw=mbSlate!55, line width=1.2pt, smooth, samples=80,
          domain=-2.9:2.9] plot (\x, {0.5*\x*\x});
    % --- axes discrets ---
    \draw[->, draw=mbSlate!70, line width=0.6pt] (-3.3,0) -- (3.4,0)
          node[right, font=\sffamily\scriptsize, text=mbSlate]
          {valeur du bouton};
    \draw[->, draw=mbSlate!70, line width=0.6pt] (-3.3,-0.3) -- (-3.3,4.6)
          node[above, font=\sffamily\scriptsize, text=mbSlate] {erreur};
    % --- tangente au point (2 ; 2), de pente 2 ---
    \draw[draw=mbAmber!85!black, line width=1pt, dashed]
          (1.25,0.5) -- (2.75,3.5);
    \node[font=\sffamily\scriptsize, text=mbAmber!75!black, anchor=west]
          at (2.8,3.4) {la pente ici};
    % --- la bille ---
    \fill[mbBlue] (2,2) circle (0.13cm);
    \node[font=\sffamily\scriptsize, text=mbBlue, anchor=south west]
          at (2.12,2.12) {le modèle};
    % --- sens du déplacement (en vert, pour ne pas le confondre avec
    %     la tangente, qui est ambre et en pointillés) ---
    \draw[->, draw=mbTeal, line width=1.1pt] (1.85,1.55) -- (1.05,1.05);
    \node[font=\sffamily\scriptsize\itshape, text=mbTeal,
          anchor=east] at (1.0,1.35) {on va par là};
    % --- le fond ---
    \fill[mbGreen] (0,0) circle (0.1cm);
    \node[font=\sffamily\scriptsize, text=mbGreen!60!black, align=center]
          at (0,-0.85) {le meilleur\\ réglage};
  \end{scope}
\end{schema}

Nous avons tout ce qu'il faut. Choisissons une longueur de pas de $0{,}1$ et
faisons deux étapes complètes, en écrivant chaque calcul.

\begin{exemple}[Première étape, tout en détail]
Nous partons de $a = 0$ et $b = 0$. Nous avons déjà tout calculé :
l'erreur vaut $31$, la pente selon $a$ vaut $-24$, la pente selon $b$ vaut
$-10$. Appliquons la règle avec un pas de $0{,}1$ :
\[
  a_{\text{nouveau}} = 0 - 0{,}1 \times (-24) = 0 + 2{,}4 = 2{,}4 ,
\]
\[
  b_{\text{nouveau}} = 0 - 0{,}1 \times (-10) = 0 + 1 = 1 .
\]
Nos deux boutons sont maintenant réglés à $a = 2{,}4$ et $b = 1$.
\end{exemple}

\begin{exemple}[A-t-on vraiment progressé ?]
Vérifions en recalculant l'erreur avec les nouveaux boutons. Les trois
prédictions valent :
\[
  2{,}4 \times 1 + 1 = 3{,}4 ; \quad
  2{,}4 \times 2 + 1 = 5{,}8 ; \quad
  2{,}4 \times 3 + 1 = 8{,}2 .
\]
Les vraies valeurs étaient $2$, $5$ et $8$. Les écarts valent donc $1{,}4$,
$0{,}8$ et $0{,}2$, et leurs carrés $1{,}96$, $0{,}64$ et $0{,}04$. La
nouvelle erreur vaut
\[
  \mathcal{L} = \frac{1{,}96 + 0{,}64 + 0{,}04}{3}
              = \frac{2{,}64}{3} = 0{,}88 .
\]
Nous sommes passés de $31$ à $0{,}88$ en une seule étape. L'erreur a été
divisée par plus de trente-cinq. La machine, qui annonçait zéro bonbon à
tout le monde, en annonce maintenant trois et demi à Léa et huit à Nina :
elle a déjà presque compris.
\end{exemple}

\begin{exemple}[Deuxième étape]
Recommençons exactement le même travail depuis $a = 2{,}4$ et $b = 1$. Les
écarts, que nous venons de calculer, valent $1{,}4$, $0{,}8$ et $0{,}2$. Ils
sont maintenant positifs : la machine annonce trop de bonbons.

Pente selon $a$ :
\[
  \frac{2}{3}\bigl[1{,}4 \times 1 + 0{,}8 \times 2 + 0{,}2 \times 3\bigr]
  = \frac{2}{3}\bigl[1{,}4 + 1{,}6 + 0{,}6\bigr]
  = \frac{2}{3} \times 3{,}6 = 2{,}4 .
\]
Pente selon $b$ :
\[
  \frac{2}{3}\bigl[1{,}4 + 0{,}8 + 0{,}2\bigr]
  = \frac{2}{3} \times 2{,}4 = 1{,}6 .
\]
Les deux pentes sont cette fois positives : il faut donc \emph{baisser} les
deux boutons.
\[
  a_{\text{nouveau}} = 2{,}4 - 0{,}1 \times 2{,}4 = 2{,}4 - 0{,}24 = 2{,}16 ,
\]
\[
  b_{\text{nouveau}} = 1 - 0{,}1 \times 1{,}6 = 1 - 0{,}16 = 0{,}84 .
\]
Et la nouvelle erreur vaut $0{,}496$, contre $0{,}88$ juste avant : elle a
encore baissé.
\end{exemple}

Récapitulons les deux étapes dans un seul tableau, que vous pourrez
comparer tout à l'heure à la sortie du programme.

\begin{center}
\begin{tabular}{cccccc}
\toprule
\textbf{Étape} & $a$ & $b$ & \textbf{erreur} &
\textbf{pente selon } $a$ & \textbf{pente selon } $b$ \\
\midrule
0 & $0$      & $0$      & $31$      & $-24$   & $-10$  \\
1 & $2{,}4$  & $1$      & $0{,}88$  & $2{,}4$ & $1{,}6$ \\
2 & $2{,}16$ & $0{,}84$ & $0{,}496$ & $-0{,}48$ & $0{,}32$ \\
\bottomrule
\end{tabular}
\end{center}

\begin{piege}
Regardez la colonne du bouton $a$ : il monte à $2{,}4$, puis \emph{redescend}
à $2{,}16$, alors que la bonne réponse est $3$. On dirait qu'il s'éloigne !
Ne vous inquiétez pas, et surtout ne jugez jamais un entraînement en
surveillant un paramètre isolé. Les boutons bougent ensemble, ils se
gênent et se rattrapent l'un l'autre. Le seul chiffre qui compte est
l'erreur, et elle, elle descend sans arrêt : $31$, puis $0{,}88$, puis
$0{,}496$. Au bout de deux cents étapes, $a$ vaut $2{,}993$ et $b$ vaut
$-0{,}985$ : la machine a bien retrouvé le $3$ et le $-1$ que nous
cherchions.
\end{piege}

\begin{enia}
Ce que nous venons de faire à la main sur deux boutons est \emph{exactement}
ce que fait l'entraînement de V-JEPA~2.1, sans un seul changement d'idée.
Les seules différences sont des différences de taille : deux milliards de
boutons au lieu de deux, une erreur calculée sur des morceaux de vidéo au
lieu de trois enfants, et cent trente-cinq mille étapes au lieu de deux.
La boucle, elle, est identique : je prédis, je mesure mon erreur, je calcule
mes pentes, je fais un petit pas, je recommence.
\end{enia}

Il nous reste à parler du seul réglage que nous avons choisi arbitrairement,
et qui est pourtant le plus délicat de tous : la longueur du pas.

\section{La taille des pas : le taux d'apprentissage}

\begin{definition}[Le taux d'apprentissage][taux]
Le \emph{taux d'apprentissage}, noté $\pas$, est le petit nombre par lequel
on multiplie la pente avant de bouger un paramètre. Il décide de la longueur
des pas. C'est nous qui le choisissons : la machine ne l'apprend pas.
\end{definition}

\begin{intuition}
Vous descendez une colline dans le brouillard. Vous ne voyez que le sol sous
vos pieds, donc vous ne connaissez que la pente à l'endroit où vous êtes.
Avec des pas minuscules, vous descendrez, mais la nuit tombera avant que
vous arriviez. Avec des pas de sept lieues, vous sauterez par-dessus la
vallée et vous vous retrouverez plus haut qu'avant, sur l'autre versant. Il
existe une bonne taille de pas entre les deux, et la trouver fait partie du
métier.
\end{intuition}

La figure~\ref{fig:taille-des-pas} met les trois situations côte à côte.

\begin{schema}{Trois façons de descendre la même vallée. À gauche, les pas
sont trop courts : la bille avance à peine et il faudra une éternité. Au
milieu, les pas sont bien choisis : la bille rejoint le fond en quelques
sauts. À droite, les pas sont trop longs : la bille saute par-dessus le fond,
remonte plus haut de l'autre côté, et s'échappe de la vallée au lieu d'y
descendre.}[taille-des-pas]
  \begin{scope}[x=0.42cm, y=0.30cm]
    % ================= panneau 1 : pas trop petit =================
    % depart en x = 3, chaque pas ne retire que 12 % du chemin restant
    \begin{scope}
      \draw[draw=mbSlate!45, line width=1pt, smooth, samples=60,
            domain=-3.7:3.7] plot (\x, {0.5*\x*\x});
      \foreach \mbx/\mby in {3.0/4.5, 2.64/3.485, 2.323/2.698,
                             2.044/2.089, 1.799/1.618}{
        \fill[mbBlue] (\mbx,\mby) circle (0.07cm);
      }
      \draw[fleche accent] (3.05,5.35) -- (2.0,5.35);
      \node[font=\sffamily\scriptsize\bfseries, text=mbSlate]
            at (0,-1.6) {pas trop petit};
      \node[font=\sffamily\scriptsize, text=mbSlate!85, align=center]
            at (0,-3.4) {on descend,\\ mais si lentement};
    \end{scope}
    % ================= panneau 2 : pas bien choisi =================
    % meme depart, mais chaque pas retire la moitie du chemin restant
    \begin{scope}[xshift=4.0cm]
      \draw[draw=mbSlate!45, line width=1pt, smooth, samples=60,
            domain=-3.7:3.7] plot (\x, {0.5*\x*\x});
      \foreach \mbx/\mby in {3.0/4.5, 1.5/1.125, 0.75/0.281,
                             0.375/0.070, 0.188/0.018}{
        \fill[mbGreen] (\mbx,\mby) circle (0.07cm);
      }
      \draw[fleche accent] (3.05,5.35) -- (0.55,5.35);
      \node[font=\sffamily\scriptsize\bfseries, text=mbGreen!55!black]
            at (0,-1.6) {pas bien choisi};
      \node[font=\sffamily\scriptsize, text=mbSlate!85, align=center]
            at (0,-3.4) {on atteint le fond\\ en quelques pas};
    \end{scope}
    % ================= panneau 3 : pas trop grand =================
    % chaque pas depasse le fond et repart plus haut de l'autre cote
    \begin{scope}[xshift=8.0cm]
      \draw[draw=mbSlate!45, line width=1pt, smooth, samples=60,
            domain=-3.7:3.7] plot (\x, {0.5*\x*\x});
      \foreach \mbx/\mby in {1.5/1.125, -1.95/1.901, 2.535/3.213,
                             -3.296/5.434}{
        \fill[mbRose] (\mbx,\mby) circle (0.07cm);
      }
      \draw[fleche accent, draw=mbRose] (1.28,1.30) -- (-1.68,1.76);
      \draw[fleche accent, draw=mbRose] (-1.70,2.16) -- (2.28,3.09);
      \draw[fleche accent, draw=mbRose] (2.30,3.52) -- (-3.02,5.29);
      \node[font=\sffamily\scriptsize\bfseries, text=mbRose]
            at (0,-1.6) {pas trop grand};
      \node[font=\sffamily\scriptsize, text=mbSlate!85, align=center]
            at (0,-3.4) {on saute par-dessus\\ et l'on remonte};
    \end{scope}
  \end{scope}
\end{schema}

Ce ne sont pas des dessins pour faire joli : refaisons une seule étape de
notre machine à bonbons avec trois pas différents, et regardons ce que
devient l'erreur.

\begin{exemple}[Trois pas, trois destins]
Nous partons toujours de $a = 0$, $b = 0$, avec une erreur de $31$ et des
pentes de $-24$ et $-10$. Une seule étape suffit à tout révéler.

\begin{center}
\begin{tabular}{cccc}
\toprule
\textbf{Pas} $\pas$ & \textbf{nouveau } $a$ & \textbf{nouveau } $b$ &
\textbf{erreur après une étape} \\
\midrule
$0{,}001$ & $0{,}024$ & $0{,}01$ & $30{,}33$ \\
$0{,}1$   & $2{,}4$   & $1$      & $\bm{0{,}88}$ \\
$1$       & $24$      & $10$     & $3\,103$ \\
\bottomrule
\end{tabular}
\end{center}

Avec un pas de $0{,}001$, l'erreur passe de $31$ à $30{,}33$ : elle a bougé
d'à peine deux pour cent. À ce rythme, il faudrait des milliers d'étapes.
Avec un pas de $0{,}1$, elle tombe à $0{,}88$ : c'est le bon réglage. Avec
un pas de $1$, elle grimpe à $3\,103$, soit cent fois pire qu'au départ ! Et
si l'on s'obstine, l'étape suivante donne une erreur de $316\,063$. La
machine ne se contente pas d'échouer : elle s'envole.
\end{exemple}

\begin{piege}
Un taux d'apprentissage trop grand est la panne numéro un des débutants, et
elle est facile à reconnaître : au lieu de descendre doucement, l'erreur
grossit d'étape en étape, jusqu'à afficher des nombres absurdes, voire le
mot \co{python}{nan}, qui veut dire « ce n'est plus un nombre du tout ».
Le réflexe correct, dans ce cas, est presque toujours le même : diviser le
taux d'apprentissage par dix et relancer.
\end{piege}

\begin{surlesdonnees}
V-JEPA~2.1 ne garde pas le même pas du début à la fin. Pendant les douze
mille premières étapes, il l'augmente doucement de $10^{-4}$ à
$5{,}25 \times 10^{-4}$ : on appelle cela une montée en douceur, et cela
évite de tout casser au tout début, quand les boutons sont encore réglés
n'importe comment. Ensuite il le laisse constant très longtemps. Enfin, sur les douze
mille dernières étapes, il le fait redescendre jusqu'à $10^{-6}$, c'est-à-dire
presque zéro : les pas deviennent minuscules et le modèle se pose tout
doucement au fond de la vallée, au lieu de continuer à sautiller autour.
Cette dernière phase porte un nom que nous retrouverons souvent : la phase
de refroidissement.
\end{surlesdonnees}

Nous savons faire un pas, et le faire de la bonne longueur. Il reste à
savoir sur combien de données on le calcule, et combien de fois on
recommence.

\section{Le rythme de l'entraînement : le lot et l'époque}

Jusqu'ici, nous avions trois enfants, et nous avons calculé l'erreur sur les
trois à chaque étape. Avec trois exemples, c'est facile. Avec cent
quarante-deux millions d'images, c'est impossible : regarder toute la
collection avant de bouger un seul bouton prendrait des jours pour une seule
étape.

\begin{enclair}
La solution est simple et un peu paresseuse : au lieu de regarder toutes les
données, on en prend une petite poignée au hasard, on calcule l'erreur sur
cette poignée-là, et on fait quand même un pas. La direction est un peu
approximative, mais on avance beaucoup plus souvent — et au bout du compte,
on va bien plus vite.
\end{enclair}

\begin{definition}[Le lot][lot]
Un \emph{lot} est la petite poignée d'exemples que l'on regarde ensemble
pour calculer une erreur et faire un pas. En anglais on dit \emph{batch}, et
ce mot revient partout dans le code. Le nombre d'exemples dans la poignée
s'appelle la \emph{taille de lot}.
\end{definition}

\begin{definition}[Le pas, ou itération][iteration]
Un \emph{pas} — on dit aussi une \emph{itération} — est un tour complet de
la boucle : on prend un lot, on prédit, on mesure l'erreur, on calcule les
pentes, on bouge les paramètres. Un pas, un lot.
\end{definition}

\begin{piege}
Le mot « pas » sert malheureusement à deux choses, et il faut les séparer
dans votre tête une bonne fois pour toutes. \emph{Faire un pas}, c'est faire
un tour de boucle : c'est ce que nous venons de définir, et cela se compte
($135\,000$ pas pour V-JEPA~2.1). \emph{La longueur du pas}, c'est le taux
d'apprentissage $\pas$ de la section précédente : c'est de combien on bouge
les boutons pendant ce tour, et cela se règle. On peut donc faire beaucoup
de pas très courts, ou peu de pas très longs. Quand un texte parle de
« pas », demandez-vous toujours s'il compte des tours ou s'il mesure une
longueur.
\end{piege}

\begin{definition}[L'époque][epoque]
Une \emph{époque} est le temps qu'il faut pour que le modèle ait vu une fois
chacune des données disponibles. Si la collection contient mille images et
que les lots en contiennent cent, alors une époque vaut dix pas.
\end{definition}

\begin{schema}{Toutes les données de la collection, découpées en lots. Le
modèle en prend un, fait un pas, en prend un autre, refait un pas. Quand il
a épuisé tous les lots, il a vu chaque donnée une fois : cela s'appelle une
époque, et il peut alors recommencer un nouveau tour.}[lot-epoque]
  \begin{scope}[x=1cm, y=1cm]
    \foreach \mbg in {0,1,2}{
      \begin{scope}[xshift=\mbg*3.55cm]
        \foreach \mbk in {0,1,2,3}{
          \draw[draw=mbBlue!70, fill=mbBlue!14, line width=0.7pt,
                rounded corners=1.5pt]
                (\mbk*0.78,0) rectangle ++(0.62,0.62);
        }
        \draw[accolade bas] (0,-0.12) -- (3.02,-0.12);
        \node[font=\sffamily\scriptsize, text=mbSlate] at (1.51,-0.66)
             {lot n\textsuperscript{o}\,\the\numexpr\mbg+1\relax};
        \draw[fleche accent] (1.51,0.95) -- (1.51,0.72);
        \node[font=\sffamily\scriptsize\itshape, text=mbAmber!75!black]
              at (1.51,1.22) {un pas};
      \end{scope}
    }
    \draw[accolade bas] (0,-1.05) -- (9.62,-1.05);
    \node[font=\sffamily\small\bfseries, text=mbGreen!55!black]
          at (4.81,-1.7) {une époque : chaque donnée a été vue une fois};
  \end{scope}
\end{schema}

\begin{exemple}[Compter des époques sur un petit jeu]
Vous disposez de $1\,000$ photos et vous choisissez des lots de $100$
photos. Chaque pas consomme $100$ photos, il faut donc
\[
  \frac{1\,000}{100} = 10 \text{ pas pour faire une époque.}
\]
Si vous entraînez pendant $50$ époques, votre programme aura effectué
$50 \times 10 = 500$ pas au total, et chaque photo aura été vue $50$ fois.
\end{exemple}

\begin{surlesdonnees}
Faisons le même calcul pour V-JEPA~2.1, et le résultat va sans doute vous
surprendre. Le modèle fait $135\,000$ pas dans sa phase principale, et
chaque lot contient $2\,304$ images. Il a donc regardé
\[
  135\,000 \times 2\,304 = 311\,040\,000 \text{ images en tout,}
\]
soit trois cent onze millions. Or sa collection d'images en compte $142$
millions. Le nombre d'époques vaut donc
\[
  \frac{311\,040\,000}{142\,000\,000} \approx 2{,}2 .
\]
Chaque image n'est vue qu'un peu plus de deux fois dans toute la vie du
modèle. Du côté des vidéos, c'est encore plus frappant : $135\,000$ pas
avec $128$ clips par lot font $17\,280\,000$ clips, pour une collection de
$19$ millions de vidéos. Le modèle ne voit donc même pas une fois chaque
vidéo !

C'est pour cela que les articles de recherche modernes ne parlent presque
plus d'époques, mais comptent en pas. Quand les données sont quasi
inépuisables, repasser deux fois sur la même image n'a plus grand intérêt :
mieux vaut en voir une nouvelle.
\end{surlesdonnees}

\begin{plusloin}
Prendre un lot au hasard plutôt que toutes les données rend le gradient un
peu bruité : la direction indiquée n'est pas tout à fait la bonne, elle
tremble. On appelle cela la descente de gradient \emph{stochastique}, un mot
savant qui veut simplement dire « au hasard ». Ce tremblement, contre toute
attente, aide souvent : il empêche la bille de rester coincée dans un petit
creux sans intérêt et lui fait trouver de meilleures vallées.
\end{plusloin}

Nous savons entraîner. Il reste un dernier danger, et c'est le plus sournois
de tous : une machine peut très bien faire baisser son erreur tout en
n'ayant strictement rien compris.

\section{Apprendre par cœur n'est pas apprendre}

\begin{enclair}
Un élève peut apprendre par cœur les réponses des exercices du livre sans
comprendre la leçon. Tant qu'on l'interroge sur ces exercices-là, il a tout
juste. Le jour où on lui pose une question nouvelle, il s'effondre. Une
machine peut faire exactement pareil.
\end{enclair}

Revenons à nos trois enfants. Notre droite n'a que deux boutons : elle est
trop simple pour tricher, elle est bien obligée de trouver une règle
générale. Mais donnons à un modèle beaucoup, beaucoup de boutons, et il
pourra faire passer une courbe tordue exactement par chacun des trois
points. Son erreur tombera à zéro. Et pourtant, si l'on glisse quatre pièces
dans la machine, ce modèle-là annoncera n'importe quoi.

\begin{definition}[Le sur-apprentissage][surapprentissage]
Le \emph{sur-apprentissage} est ce qui arrive quand un modèle apprend par
cœur les exemples qu'on lui a montrés, y compris leurs petits défauts, au
lieu d'apprendre la règle générale qui se cache derrière. Il est alors
excellent sur ce qu'il a déjà vu, et mauvais sur tout le reste.
\end{definition}

\begin{schema}{Les mêmes points, expliqués de trois façons. À gauche, le
modèle est trop simple : il ne suit même pas la tendance, il passe à côté de
tout le monde. Au milieu, il a saisi la règle générale : il ne touche pas
tous les points, et c'est très bien ainsi. À droite, il passe exactement par
chaque point, son erreur est nulle — mais entre les points il fait n'importe
quoi, et sur un point nouveau il se trompera lourdement.}[surapprentissage]
  \begin{scope}[x=0.76cm, y=0.66cm]
    % ============ panneau 1 : trop simple ============
    \begin{scope}
      \draw[draw=mbSlate!35, line width=0.6pt] (0,0) rectangle (4.4,4.6);
      \draw[draw=mbRose, line width=1.2pt] (0.25,2.45) -- (4.15,2.45);
      \foreach \mbx/\mby in {0.3/0.9, 0.9/1.9, 1.5/1.7, 2.1/2.9,
                             2.7/2.6, 3.3/3.7, 3.9/3.5}{
        \fill[mbBlue] (\mbx,\mby) circle (0.075cm);
      }
      \node[font=\sffamily\scriptsize\bfseries, text=mbRose]
            at (2.2,-0.55) {trop simple};
      \node[font=\sffamily\scriptsize, text=mbSlate!85, align=center]
            at (2.2,-1.4) {il n'a rien\\ appris du tout};
    \end{scope}
    % ============ panneau 2 : juste ============
    \begin{scope}[shift={(5.1,0)}]
      \draw[draw=mbSlate!35, line width=0.6pt] (0,0) rectangle (4.4,4.6);
      \draw[draw=mbGreen, line width=1.2pt] (0.25,0.9) -- (4.15,4.0);
      \foreach \mbx/\mby in {0.3/0.9, 0.9/1.9, 1.5/1.7, 2.1/2.9,
                             2.7/2.6, 3.3/3.7, 3.9/3.5}{
        \fill[mbBlue] (\mbx,\mby) circle (0.075cm);
      }
      \node[font=\sffamily\scriptsize\bfseries, text=mbGreen!55!black]
            at (2.2,-0.55) {bien réglé};
      \node[font=\sffamily\scriptsize, text=mbSlate!85, align=center]
            at (2.2,-1.4) {il a compris\\ la tendance};
    \end{scope}
    % ============ panneau 3 : par cœur ============
    \begin{scope}[shift={(10.2,0)}]
      \draw[draw=mbSlate!35, line width=0.6pt] (0,0) rectangle (4.4,4.6);
      \draw[draw=mbAmber!85!black, line width=1.2pt, smooth]
            plot coordinates {(0.3,0.9) (0.9,1.9) (1.5,1.7) (2.1,2.9)
                              (2.7,2.6) (3.3,3.7) (3.9,3.5)};
      \foreach \mbx/\mby in {0.3/0.9, 0.9/1.9, 1.5/1.7, 2.1/2.9,
                             2.7/2.6, 3.3/3.7, 3.9/3.5}{
        \fill[mbBlue] (\mbx,\mby) circle (0.075cm);
      }
      \node[font=\sffamily\scriptsize\bfseries, text=mbAmber!80!black]
            at (2.2,-0.55) {appris par cœur};
      \node[font=\sffamily\scriptsize, text=mbSlate!85, align=center]
            at (2.2,-1.4) {erreur nulle,\\ mais rien compris};
    \end{scope}
  \end{scope}
\end{schema}

Comment se protéger de ce piège ? En trichant le moins possible : on cache
au modèle une partie des données, et on ne les lui montre jamais pendant
l'entraînement. Elles serviront uniquement à le juger.

\begin{definition}[Jeu d'entraînement et jeu de validation][jeux]
Le \emph{jeu d'entraînement} rassemble les données que le modèle a le droit
de voir pour régler ses boutons. Le \emph{jeu de validation}, parfois appelé
jeu de test, rassemble des données qu'il ne voit jamais pendant
l'entraînement, et qui servent uniquement à mesurer honnêtement ce qu'il
vaut. C'est la différence entre les exercices du cahier et l'examen de fin
d'année.
\end{definition}

\begin{piege}
La faute la plus grave, dans tout l'apprentissage automatique, est de régler
quoi que ce soit en regardant le jeu de validation. Si vous essayez trente
taux d'apprentissage et que vous gardez celui qui donne le meilleur score de
validation, alors ce score n'est plus honnête : vous avez, sans le vouloir,
appris par cœur votre examen. Les vrais résultats se mesurent toujours sur
des données que rien ni personne n'a servi à choisir.
\end{piege}

\begin{enia}
L'article de V-JEPA~2.1 prend cette précaution très au sérieux, et cela se
voit dans ses annexes. Pour le suivi d'objets dans les vidéos, par exemple,
les auteurs cherchent les meilleurs réglages sur la partie entraînement du
jeu DAVIS, puis appliquent tels quels ces réglages à tous les jeux
d'évaluation, sans plus jamais y toucher. C'est un détail d'annexe, et
pourtant c'est ce détail qui rend leurs chiffres croyables.
\end{enia}

Nous avons maintenant tout le vocabulaire. Le moment est venu de vérifier
que nos calculs de papier tiennent devant une vraie machine.

\section{Passons au code}

Nous allons refaire, en Python, les deux étapes que nous venons de calculer
à la main. Nous utiliserons \co{python}{torch}, la bibliothèque PyTorch, qui
sait faire deux choses très utiles : manipuler des tenseurs, et surtout
calculer les pentes toute seule. Nous n'aurons donc plus jamais à écrire les
formules du gradient nous-mêmes.

Commençons par le plus important : vérifier que la petite poussée du début
du chapitre donne bien ce que nous avions trouvé.

\begin{code}{python}[Mesurer une pente en poussant un bouton]
import torch

# Les trois enfants : pieces glissees et bonbons recus.
x = torch.tensor([1.0, 2.0, 3.0])
y = torch.tensor([2.0, 5.0, 8.0])

def erreur(a, b):
    prediction = a * x + b
    return ((prediction - y) ** 2).mean().item()

avant = erreur(0.00, 0.0)   # les boutons a zero
apres = erreur(0.01, 0.0)   # on pousse le bouton a de 0,01

print(f"erreur avant la poussee : {avant:.4f}")   # -> 31.0000
print(f"erreur apres la poussee : {apres:.4f}")   # -> 30.7605
print(f"pente mesuree           : {(apres - avant) / 0.01:.2f}")  # -> -23.95
\end{code}

Les trois nombres affichés sont exactement ceux de notre exemple :
$31$, $30{,}7605$ et $-23{,}95$. Le calcul à la main et la machine sont
d'accord.

Voici maintenant la boucle d'entraînement complète. Elle tient en une
quinzaine de lignes, et c'est littéralement la même boucle que celle qui
entraîne V-JEPA~2.1 : seules les tailles changent.

\begin{code}{python}[Deux étapes de descente de gradient]
import torch

x = torch.tensor([1.0, 2.0, 3.0])
y = torch.tensor([2.0, 5.0, 8.0])

# Nos deux boutons. « requires_grad » dit a PyTorch :
# surveille ce nombre, je te demanderai sa pente.
a = torch.tensor(0.0, requires_grad=True)
b = torch.tensor(0.0, requires_grad=True)
pas = 0.1

# 1. je prends mes donnees : ici, les trois enfants d'un seul coup.
for etape in range(2):
    prediction = a * x + b                       # 2. je calcule ma prediction
    erreur = ((prediction - y) ** 2).mean()      # 3. je mesure mon erreur
    erreur.backward()                            # 4. je calcule les pentes

    print(f"etape {etape} : a = {a.item():.2f}, b = {b.item():.2f}, "
          f"erreur = {erreur.item():.4f}, "
          f"pente_a = {a.grad.item():.2f}, pente_b = {b.grad.item():.2f}")

    with torch.no_grad():                        # 5. je bouge les boutons
        a -= pas * a.grad
        b -= pas * b.grad

    a.grad.zero_()          # puis j'efface les pentes pour le tour suivant
    b.grad.zero_()

print(f"apres 2 etapes : a = {a.item():.2f}, b = {b.item():.2f}")
\end{code}

Ce programme affiche exactement ceci :

\begin{sortie}
etape 0 : a = 0.00, b = 0.00, erreur = 31.0000, pente_a = -24.00, pente_b = -10.00
etape 1 : a = 2.40, b = 1.00, erreur = 0.8800, pente_a = 2.40, pente_b = 1.60
apres 2 etapes : a = 2.16, b = 0.84
\end{sortie}

Comparez ligne à ligne avec le tableau de la section précédente : le $31$,
le $-24$, le $-10$, le $2{,}4$, le $0{,}88$, le $2{,}16$ et le $0{,}84$ y
sont tous. Rien n'a été caché, et la machine n'a rien fait de plus que ce
que nous avons fait au crayon.

\begin{piege}
La ligne \co{python}{a.grad.zero_()} paraît anodine, et pourtant l'oublier
est l'une des erreurs les plus courantes chez les débutants. PyTorch
\emph{additionne} les pentes qu'il calcule au lieu de les remplacer. Si vous
n'effacez pas, l'étape suivante utilisera la somme de deux pentes, puis de
trois, et vos pas grandiront tout seuls jusqu'à faire exploser
l'entraînement. Le symptôme est le même que celui d'un taux d'apprentissage
trop grand, ce qui rend la panne difficile à diagnostiquer.
\end{piege}

Laissons enfin tourner la même boucle bien plus longtemps, pour voir la
machine retrouver la règle des bonbons.

\begin{code}{python}[Deux cents étapes, et la règle apparaît]
import torch

x = torch.tensor([1.0, 2.0, 3.0])
y = torch.tensor([2.0, 5.0, 8.0])
a = torch.tensor(0.0, requires_grad=True)
b = torch.tensor(0.0, requires_grad=True)

for _ in range(200):
    erreur = ((a * x + b - y) ** 2).mean()
    erreur.backward()
    with torch.no_grad():
        a -= 0.1 * a.grad
        b -= 0.1 * b.grad
    a.grad.zero_()
    b.grad.zero_()

print(f"a = {a.item():.3f}   b = {b.item():.3f}")   # -> a = 2.993   b = -0.985
print(f"erreur = {erreur.item():.6f}")              # -> erreur = 0.000034
\end{code}

Le bouton $a$ vaut $2{,}993$ et le bouton $b$ vaut $-0{,}985$. La machine a
retrouvé, toute seule, la règle que nous avions devinée de tête au début du
chapitre : trois bonbons par pièce, moins un pour la machine. Personne ne
lui a jamais soufflé la réponse. Elle a seulement descendu la vallée, deux
cents fois de suite.

\begin{aretenir}
Retenez cette boucle en cinq gestes, car c'est la seule qu'il y ait dans
tout l'apprentissage automatique, du plus petit exemple au plus gros modèle
du monde :

\begin{enumerate}
  \item je prends un lot de données ;
  \item je calcule ma prédiction ;
  \item je mesure mon erreur ;
  \item je calcule les pentes ;
  \item je bouge chaque bouton d'un petit pas dans le sens opposé à sa pente.
\end{enumerate}

Puis je recommence. Cent trente-cinq mille fois, dans le cas de V-JEPA~2.1.
\end{aretenir}

Il manque pourtant quelque chose d'énorme à tout ce que nous venons de
faire, et vous l'avez peut-être déjà remarqué. À chaque étape, nous avions
besoin de connaître la bonne réponse : les deux bonbons de Léa, les cinq de
Sami, les huit de Nina. Sans elles, impossible de calculer le moindre écart,
donc impossible de calculer la moindre erreur, donc impossible d'apprendre
quoi que ce soit.

Mais qui va écrire la bonne réponse à côté d'un milliard de vidéos ?
Personne. Aucune équipe au monde, aucun budget, aucune quantité de patience
ne suffirait. C'est précisément le mur contre lequel s'arrête tout ce que
nous venons d'apprendre — et c'est le mur que le chapitre suivant va
traverser.

\phantomsection
\section*{Exercices}
\addcontentsline{toc}{section}{Exercices}

\begin{exercice}[Calculer une erreur]
Un modèle prédit $\pred{y} = 4$, $\pred{y} = 7$ et $\pred{y} = 9$, alors que
les vraies valeurs sont $3$, $9$ et $9$. Quelle est l'erreur, au sens de la
moyenne des carrés des écarts ?
\end{exercice}

\begin{solution}
Les écarts valent $4 - 3 = 1$, puis $7 - 9 = -2$, puis $9 - 9 = 0$. Leurs
carrés valent $1$, $4$ et $0$, dont la somme fait $5$. Comme il y a trois
exemples,
\[
  \mathcal{L} = \frac{5}{3} \approx 1{,}67 .
\]
Remarquez que le deuxième écart, bien que négatif, contribue positivement à
l'erreur : c'est tout l'intérêt du carré.
\end{solution}

\begin{exercice}[Faire un pas]
Un paramètre vaut actuellement $p = 5$. On calcule sa pente et l'on trouve
$+3$. Le taux d'apprentissage vaut $0{,}2$. Que vaut le paramètre après la
mise à jour ? A-t-il augmenté ou diminué, et pourquoi était-ce le bon sens ?
\end{exercice}

\begin{solution}
On applique la règle $p_{\text{nouveau}} = p_{\text{ancien}} - \pas \times
\text{pente}$, soit
\[
  p_{\text{nouveau}} = 5 - 0{,}2 \times 3 = 5 - 0{,}6 = 4{,}4 .
\]
Le paramètre a diminué. C'était le bon sens, car une pente positive signifie
« si j'augmente ce bouton, l'erreur monte ». Pour faire descendre l'erreur,
il fallait donc baisser le bouton.
\end{solution}

\begin{exercice}[Compter des époques]
Un jeu de données contient $50\,000$ images et l'on choisit des lots de
$250$ images. Combien de pas faut-il pour une époque ? Et si l'on entraîne
pendant $30\,000$ pas, combien d'époques cela représente-t-il ?
\end{exercice}

\begin{solution}
Une époque demande $50\,000 / 250 = 200$ pas. Avec $30\,000$ pas au total,
on effectue
\[
  \frac{30\,000}{200} = 150 \text{ époques,}
\]
et chaque image aura donc été vue $150$ fois. C'est un régime très différent
de celui de V-JEPA~2.1, qui ne dépasse pas deux époques et demie : quand les
données sont peu nombreuses, on repasse dessus souvent, et le risque de
sur-apprentissage augmente d'autant.
\end{solution}

\begin{exercice}[Diagnostiquer une panne]
Vous lancez un entraînement et vous voyez défiler ces erreurs :
$18{,}4$ puis $91{,}7$ puis $2\,340$ puis \co{python}{nan}. Que se passe-t-il
très probablement, et quel est le premier réglage à corriger ?
\end{exercice}

\begin{solution}
L'erreur ne descend pas, elle explose : le modèle saute par-dessus le fond
de la vallée et remonte de plus en plus haut à chaque étape, jusqu'à sortir
des nombres représentables par la machine, ce qui produit le
\co{python}{nan}. C'est la signature d'un taux d'apprentissage trop grand.
Le premier geste est de le diviser par dix et de relancer ; si le problème
persiste, on divise encore par dix. Attention toutefois : oublier
d'effacer les pentes avec \co{python}{zero_()} produit exactement le même
symptôme, et vaut la peine d'être vérifié dans la foulée.
\end{solution}

\begin{exercice}[Le piège de l'examen]
Vous disposez d'un jeu d'entraînement et d'un jeu de validation. Vous
essayez cinquante taux d'apprentissage différents et vous annoncez, comme
résultat final de votre travail, le meilleur score obtenu sur le jeu de
validation. Pourquoi ce chiffre est-il trompeur, et que fallait-il faire ?
\end{exercice}

\begin{solution}
En choisissant le réglage qui donne le meilleur score de validation, vous
avez utilisé le jeu de validation pour prendre une décision : il a donc servi
à l'entraînement, au sens large. Parmi cinquante essais, le meilleur score
est en partie dû à la chance, et il sera plus flatteur que la réalité. La
bonne pratique consiste à garder un troisième jeu de données, mis de côté
dès le départ, que l'on ne regarde qu'une seule et unique fois, tout à la
fin, pour annoncer le résultat.
\end{solution}
                   % poids, perte, descente de gradient
\chapter{Apprendre sans étiquettes : l'auto-supervision}
\label{chap:auto-supervision}

\begin{quote}\itshape
Au chapitre précédent, la machine avait toujours la bonne réponse écrite à
côté de la question. Nous allons voir maintenant que, pour les vidéos, cette
bonne réponse n'existe pas et n'existera jamais. Nous allons donc apprendre
à la machine un nouveau jeu : cacher elle-même un morceau de ce qu'elle
regarde, et essayer de le deviner. Personne ne lui souffle rien. Et
pourtant, elle apprend.
\end{quote}

\begin{prerequis}
Les chapitres~\ref{chap:image-video} et~\ref{chap:apprendre}. Il faut savoir
qu'une image est une grille de nombres et connaître sa forme, et il faut
avoir compris la boucle en cinq gestes : je prédis, je mesure mon erreur, je
calcule les pentes, je bouge les boutons, je recommence. Nous n'aurons besoin
de rien d'autre. Aucun mot nouveau ne sera employé sans être expliqué juste
avant.
\end{prerequis}

Reprenons la machine à bonbons du chapitre précédent. Pour qu'elle apprenne,
il nous fallait trois choses : les questions (le nombre de pièces), les
réponses (le nombre de bonbons), et la boucle qui corrige les boutons. Sans
les réponses, la deuxième étape de la boucle s'arrête net : impossible de
mesurer un écart, donc impossible de savoir dans quel sens tourner quoi que
ce soit.

Tout ce chapitre part de cette question, et d'elle seule : \emph{et si
personne n'écrit les réponses ?}

\section{La bonne réponse écrite à côté : l'étiquette}

\begin{enclair}
Une étiquette, c'est le petit mot qu'un humain colle sur une donnée pour
dire ce qu'elle contient. « Ceci est un chat. » Sans ce petit mot, la
méthode du chapitre précédent ne fonctionne pas.
\end{enclair}

Quand nous avons entraîné la machine à bonbons, nous connaissions les trois
réponses : deux bonbons pour Léa, cinq pour Sami, huit pour Nina. Ces trois
nombres n'étaient pas tombés du ciel. Trois enfants les avaient comptés, et
quelqu'un les avait notés sur un papier. Ce travail de notation a un nom.

\begin{definition}[L'étiquette][etiquette]
Une \emph{étiquette} est la bonne réponse que l'on écrit à côté d'une donnée,
pour dire à la machine ce qu'elle aurait dû répondre. Sur une photo, ce peut
être le mot « chat ». Sur un enregistrement de voix, ce peut être la phrase
prononcée. Sur une image médicale, ce peut être « sain » ou « malade ». Une
étiquette est toujours écrite par quelqu'un : elle ne se fabrique pas toute
seule.
\end{definition}

\begin{definition}[L'apprentissage supervisé][supervise]
On parle d'\emph{apprentissage supervisé} quand on entraîne une machine avec
des données accompagnées de leurs étiquettes. Le mot « supervisé » veut dire
« surveillé » : à chaque essai, quelqu'un — au travers de l'étiquette qu'il a
écrite à l'avance — corrige la machine. C'est exactement ce que nous avons
fait au chapitre~\ref{chap:apprendre}.
\end{definition}

\begin{exemple}[Reconnaître un chat d'un chien]
Vous voulez une machine qui, devant une photo, dise « chat » ou « chien ».
Vous rassemblez dix mille photos. Puis, une par une, vous les ouvrez, vous
regardez, et vous tapez le bon mot à côté. Dix mille fois. Quand c'est fini,
vous avez dix mille étiquettes, et vous pouvez lancer la boucle du chapitre
précédent : la machine prédit, on compare sa prédiction à l'étiquette, on
mesure l'écart, on corrige les boutons.

Retenez bien le prix payé ici : \emph{dix mille regards humains}. C'est ce
prix-là que nous allons regarder de plus près.
\end{exemple}

Tant qu'on parle de dix mille photos, tout va bien. Une petite équipe y passe
quelques semaines et le travail est fait. Voyons ce qui se passe quand on
change d'échelle.

\section{Le mur : ce que coûte vraiment une étiquette}

\begin{enclair}
Étiqueter est un travail humain. Un travail humain prend du temps. Et quand
le nombre de données devient très grand, ce temps devient si long qu'aucune
entreprise, aucun pays, aucune génération ne pourrait le fournir.
\end{enclair}

Faisons le calcul ensemble, avec des nombres simples et vérifiables.

Imaginons la tâche la plus facile qui soit : regarder une photo et taper un
seul mot. Une personne rapide met environ cinq secondes par photo, en
comptant le temps de regarder et le temps de taper. Maintenant, comptons le
temps de travail d'une personne dans une année : disons sept heures par jour,
deux cent vingt jours par an — les week-ends et les congés en moins. Cela
fait $7 \times 220 = 1\,540$ heures de travail par an.

\begin{exemple}[Étiqueter cent quarante-deux millions d'images]
V-JEPA~2.1 s'entraîne, entre autres, sur un ensemble de $142\,000\,000$
d'images. Étiquetons-les à cinq secondes pièce :

\[
  142\,000\,000 \times 5 = 710\,000\,000 \text{ secondes.}
\]

Convertissons en heures en divisant par $3\,600$ :

\[
  \frac{710\,000\,000}{3\,600} \approx 197\,000 \text{ heures.}
\]

Puis en années de travail d'une personne, en divisant par $1\,540$ :

\[
  \frac{197\,000}{1\,540} \approx 128 \text{ ans.}
\]

Cent vingt-huit ans de travail, pour une seule personne, sans jamais
s'arrêter. Et c'était la tâche la plus \emph{facile} : un mot par image.
\end{exemple}

Et les vidéos ? La plus grosse source de V-JEPA~2.1 représente à elle seule
$1\,600\,000$ heures de film. Notez bien : il ne s'agit plus du temps qu'il
faut pour étiqueter, mais du temps qu'il faudrait rien que pour
\emph{regarder} ces vidéos une fois, à vitesse normale.

\[
  \frac{1\,600\,000}{1\,540} \approx 1\,039 \text{ ans.}
\]

Plus de mille ans de visionnage. Il faudrait avoir commencé à regarder à
l'époque des premiers châteaux forts pour finir aujourd'hui — et on n'aurait
encore rien écrit.

\begin{surlesdonnees}
Le mur devient encore bien plus haut dès qu'on demande une étiquette plus
fine qu'un simple mot. Pour apprendre à une machine à dire, \emph{pour chaque
pixel}, s'il appartient à la route, au trottoir ou à un piéton, il faut
colorier l'image à la main, contour par contour. Sur des photos de rue, ce
travail prend de l'ordre d'une heure et demie par image.

Sur nos cent quarante-deux millions d'images, cela ferait
$142\,000\,000 \times 1{,}5 = 213\,000\,000$ heures, soit environ
$213\,000\,000 / 1\,540 \approx 138\,000$ ans. Cent trente-huit mille ans de
coloriage. Autant dire : jamais.
\end{surlesdonnees}

La figure~\ref{fig:montagne-heures} met ces trois durées côte à côte. Le
troisième bâton ne tient pas sur la page, et c'est bien le message du dessin.

\begin{schema}{Le temps de travail humain qu'il faudrait pour étiqueter les
données de V-JEPA~2.1, comparé à une carrière entière. Le premier bâton est
tout ce qu'une personne peut fournir dans sa vie. Le deuxième est le prix
d'un simple mot par image. Le troisième, celui du coloriage pixel par pixel,
sort de la page : il faudrait mille fois plus de temps que ce que la feuille
peut montrer.}[montagne-heures]
  \begin{scope}[x=1cm, y=1cm]
    \draw[draw=mbSlate!45, line width=0.7pt] (0,-0.25) -- (0,3.45);
    % Les trois légendes sont posées de la même façon : juste au-dessus de
    % leur bâton, alignées à gauche sur l'axe.
    % --- une carrière ---
    \node[font=\sffamily\scriptsize, text=mbSlate, anchor=west]
         at (0.05,3.15) {une carrière entière : $40$ ans};
    \draw[draw=mbTeal, fill=mbTeal!25, line width=0.8pt]
          (0,2.4) rectangle (0.9,2.9);
    % --- un mot par image ---
    \node[font=\sffamily\scriptsize, text=mbSlate, anchor=west]
         at (0.05,2.05) {étiqueter les images, un mot chacune : $128$ ans};
    \draw[draw=mbAmber!85!black, fill=mbAmber!25, line width=0.8pt]
          (0,1.3) rectangle (2.88,1.8);
    % --- le coloriage : le bâton sort de la page ---
    \node[font=\sffamily\scriptsize, text=mbSlate, anchor=west]
         at (0.05,0.95) {les colorier pixel par pixel : $138\,000$ ans};
    \draw[draw=mbRose, fill=mbRose!25, line width=0.8pt]
          (0,0.2) rectangle (7.9,0.7);
    \draw[draw=white, line width=1.6pt] (7.45,0.1) -- (7.6,0.8);
    \draw[draw=white, line width=1.6pt] (7.7,0.1) -- (7.85,0.8);
    \draw[fleche accent] (8.0,0.45) -- (8.7,0.45);
    \node[annotation, text width=2.4cm, anchor=west] at (8.85,0.45)
         {et cela continue\\ bien au-delà};
    \node[font=\sffamily\scriptsize\itshape, text=mbSlate, anchor=west]
         at (0.05,-0.65) {longueur des bâtons : temps de travail d'une seule
                          personne};
  \end{scope}
\end{schema}

Vous voyez le problème. Nous savons faire apprendre une machine, mais notre
méthode réclame un carburant — les étiquettes — que nous ne pouvons pas
produire en quantité suffisante. Or les données brutes, elles, sont partout :
des milliards de photos et de vidéos existent déjà, sans le moindre petit mot
écrit à côté.

\begin{aretenir}
Les données brutes sont abondantes et presque gratuites. Les étiquettes sont
rares et très chères. Toute la suite du livre découle de ce déséquilibre :
il faut trouver un moyen d'apprendre à partir de ce qui est abondant.
\end{aretenir}

La question devient donc : peut-on fabriquer la bonne réponse \emph{sans
personne} ? La réponse est oui, et elle est étonnamment simple.

\section{L'idée qui change tout : cacher un morceau}

\begin{intuition}
Si je cache un bout de quelque chose que j'ai sous les yeux, je connais déjà
la bonne réponse : c'est le bout que je viens de cacher. Je peux donc me
poser une question à moi-même, et me corriger tout seul. Aucun humain n'est
nécessaire.
\end{intuition}

Faisons l'expérience sur vous, tout de suite. Lisez cette phrase :

\begin{center}
\itshape « Le chat boit du \underline{\hphantom{xxxxxxx}}. »
\end{center}

Vous avez pensé « lait ». Personne ne vous a donné la réponse. Vous l'avez
trouvée parce que vous connaissez les chats, les boissons, et la façon dont
les phrases sont construites. Autrement dit : \emph{pour remplir le trou, il
a fallu comprendre le reste}. C'est toute l'idée.

Maintenant, regardez ce que cela donne si l'on cache un morceau d'une image.
Si je masque le museau d'un chien sur une photo, vous pouvez encore le
dessiner de tête, à peu près : vous savez où il va, de quelle couleur il est,
quelle forme il a. Pour cela, vous vous êtes servi des oreilles, des yeux, de
la fourrure autour. Vous avez compris la photo.

\begin{definition}[L'apprentissage auto-supervisé][ssl]
On parle d'\emph{apprentissage auto-supervisé} quand la machine fabrique
elle-même ses questions et ses réponses à partir des données brutes, sans
aucune étiquette écrite par un humain. Le préfixe « auto » veut dire
« par elle-même » : c'est elle qui se surveille. On l'abrège souvent en
\emph{SSL}, initiales de l'anglais \emph{self-supervised learning}.
\end{definition}

La figure~\ref{fig:etiquette-vs-ssl} met les deux façons de faire côte à côte.
Regardez surtout la différence de rythme entre la colonne de gauche et celle
de droite.

\begin{schema}{Deux façons d'obtenir la bonne réponse. À gauche,
l'apprentissage supervisé : un humain regarde chaque image et écrit une
étiquette, une par une. À droite, l'apprentissage auto-supervisé : la machine
recouvre elle-même un morceau de l'image, et le morceau caché devient la
bonne réponse. À gauche, il faut un humain à chaque image ; à droite,
personne.}[etiquette-vs-ssl]
  \begin{scope}[x=0.94cm, y=0.94cm]
    %================ colonne de gauche : supervisé ================
    \node[font=\sffamily\bfseries\small, text=mbSlate] at (2.2,4.95)
         {Apprentissage supervisé};
    % le personnage : une tête et un buste très simples
    \begin{scope}[shift={(0.2,2.6)}]
      \draw[draw=mbSlate, fill=mbMist, line width=0.8pt] (0,0.55) circle (0.32);
      \draw[draw=mbSlate, fill=mbMist, line width=0.8pt]
            (-0.42,-0.5) -- (0.42,-0.5) -- (0.3,0.18) -- (-0.3,0.18) -- cycle;
      \node[font=\sffamily\scriptsize, text=mbSlate] at (0,-0.95) {un humain};
    \end{scope}
    % trois images à étiqueter
    \foreach \mbk/\mby in {0/3.6, 1/2.6, 2/1.6}{
      \draw[draw=mbSlate!55, fill=white, line width=0.7pt]
            (1.6,\mby-0.35) rectangle (2.5,\mby+0.35);
      \draw[fleche] (0.75,\mby) -- (1.5,\mby);
      \node[font=\sffamily\scriptsize, text=mbInk] at (3.5,\mby)
           {« chat »};
      \draw[fleche] (2.6,\mby) -- (2.95,\mby);
    }
    \node[annotation, text width=4cm] at (2.2,0.7)
         {une image, un regard, un mot\\ puis on recommence
          142 millions de fois};
    % séparateur
    \draw[draw=mbSlate!30, dashed, line width=0.7pt] (5.2,0.3) -- (5.2,5.25);
    %================ colonne de droite : auto-supervisé ============
    \node[font=\sffamily\bfseries\small, text=mbSlate] at (8.6,4.95)
         {Apprentissage auto-supervisé};
    % l'image complète
    \draw[draw=mbSlate!55, fill=white, line width=0.7pt]
          (6.3,2.85) rectangle (7.5,3.75);
    \node[font=\sffamily\scriptsize, text=mbSlate] at (6.9,2.55)
         {l'image};
    \draw[fleche accent] (7.65,3.3) -- (8.35,3.3);
    % l'image masquée
    \draw[draw=mbSlate!55, fill=white, line width=0.7pt]
          (8.5,2.85) rectangle (9.7,3.75);
    \fill[mbSlate!35] (8.8,3.05) rectangle (9.4,3.55);
    \node[font=\sffamily\scriptsize, text=mbSlate, align=center,
          text width=3.6cm] at (9.1,4.15)
         {la machine recouvre un carré};
    % la réponse
    \draw[fleche] (9.85,3.3) -- (10.6,3.3);
    \draw[draw=mbAmber!85!black, fill=mbAmber!18, line width=0.9pt]
          (10.7,3.05) rectangle (11.3,3.55);
    \node[annotation, text width=3.2cm] at (11.1,2.25)
         {le morceau caché\\ \emph{est} la bonne réponse};
    \node[annotation, text width=4.6cm] at (8.9,1.0)
         {aucun humain n'intervient :\\ on peut le faire des milliards
          de fois, jour et nuit};
  \end{scope}
\end{schema}

Ce petit changement de point de vue a une conséquence énorme, et il faut
prendre le temps de la mesurer. Dans la colonne de gauche, le nombre
d'exemples d'entraînement est limité par le nombre d'heures que des humains
acceptent d'y passer. Dans la colonne de droite, il est limité par le nombre
de vidéos qui existent et par la puissance des machines. Ce ne sont pas du
tout les mêmes limites.

\begin{exemple}[Le même carré, deux rôles]
Prenons une image de $8 \times 8$ cases. Je recouvre les seize cases du coin
en haut à droite. Ce que je donne à la machine, c'est l'image avec son trou.
Ce que je lui demande, c'est de reconstituer les seize cases manquantes. Et
comme je viens de les recouvrir moi-même, je les connais parfaitement : je
peux donc comparer sa réponse à la vérité et calculer une erreur, exactement
comme au chapitre~\ref{chap:apprendre}.

Le même carré joue donc deux rôles à la fois : il est \emph{la question}
quand on le cache, et \emph{la réponse} quand on le compare. Voilà pourquoi
on n'a plus besoin de personne.
\end{exemple}

Ce jeu de cache-cache porte un nom, et il faut le connaître, car nous allons
l'employer jusqu'à la fin du livre.

\section{La tâche prétexte : un jeu que l'on jette après}

\begin{definition}[La tâche prétexte][pretexte]
Une \emph{tâche prétexte} est un exercice que l'on invente pour faire
travailler la machine, alors que la réponse à cet exercice ne nous intéresse
pas du tout. Le mot « prétexte » a ici son sens courant : une bonne excuse. Ce
qui nous intéresse, ce n'est pas que la machine sache remplir des trous ; c'est
tout ce qu'elle a été \emph{obligée} de comprendre pour y arriver.
\end{definition}

Cette idée surprend souvent au début, alors prenons une comparaison de la vie
de tous les jours.

\begin{intuition}
Un professeur de musique demande à son élève de jouer des gammes pendant des
mois. Personne n'a jamais acheté de billet pour aller écouter des gammes. La
gamme est un prétexte : ce que le professeur veut vraiment, ce sont des
doigts agiles et une oreille juste. Le jour du concert, l'élève ne jouera
aucune gamme, mais tout ce qu'il jouera en profitera.

Pour notre machine, c'est pareil. Remplir des trous est la gamme. Reconnaître
une action, mesurer une distance, suivre un objet : voilà le concert.
\end{intuition}

Les chercheurs ont inventé beaucoup de gammes différentes. La
figure~\ref{fig:pretextes} en montre quatre parmi les plus connues. Elles ont
toutes le même point commun : la bonne réponse se fabrique toute seule, en
partant de l'image entière.

\begin{schema}{Quatre tâches prétextes. Dans les quatre cas, on abîme
volontairement l'image d'une façon dont on garde le secret, et l'on demande à
la machine de retrouver ce qui a été fait. Comme c'est nous qui avons abîmé
l'image, nous connaissons toujours la bonne réponse, sans avoir eu besoin de
la moindre étiquette.}[pretextes]
  \begin{scope}[x=0.87cm, y=0.87cm]
    \foreach \mbx/\mbtitre/\mbquestion in {%
      0/{le puzzle}/{remets les morceaux\\ dans l'ordre},
      3.6/{la rotation}/{de quel angle\\ ai-je tourné ?},
      7.2/{la couleur}/{remets les couleurs\\ sur le gris},
      10.8/{le trou}/{devine ce qui\\ manque}}{
      \draw[draw=mbSlate!55, fill=white, line width=0.7pt]
            (\mbx,0) rectangle (\mbx+2.4,1.8);
      \node[font=\sffamily\small\bfseries, text=mbBlue] at (\mbx+1.2,2.15)
           {\mbtitre};
      \node[font=\sffamily\scriptsize, text=mbSlate, align=center,
            text width=2.6cm] at (\mbx+1.2,-0.75) {\mbquestion};
    }
    % 1. puzzle : quatre cases mélangées, numérotées dans le désordre
    \foreach \mbi/\mbj/\mbn in {0/0/3, 1/0/1, 0/1/4, 1/1/2}{
      \draw[draw=mbSlate!45, fill=mbBlue!8, line width=0.5pt]
            (0.25+\mbi*1.05, 0.2+\mbj*0.75)
            rectangle (1.2+\mbi*1.05, 0.85+\mbj*0.75);
      \node[font=\sffamily\scriptsize, text=mbSlate]
           at (0.72+\mbi*1.05, 0.52+\mbj*0.75) {\mbn};
    }
    % 2. rotation : une petite maison penchée, et la flèche qui tourne
    %    reste à l'intérieur du cadre
    \begin{scope}[shift={(4.8,0.72)}, rotate=-35, scale=0.82]
      \draw[draw=mbTeal, fill=mbTeal!12, line width=0.8pt]
            (-0.45,-0.4) rectangle (0.45,0.25);
      \draw[draw=mbTeal, fill=mbTeal!22, line width=0.8pt]
            (-0.58,0.25) -- (0,0.7) -- (0.58,0.25) -- cycle;
    \end{scope}
    \draw[fleche accent] (4.15,1.42) to[bend left=28] (5.45,1.42);
    % 3. couleur : la même petite image, grise à gauche, colorée à droite
    \fill[mbSlate!25] (7.45,0.2) rectangle (8.35,1.6);
    \fill[mbAmber!45] (8.35,0.2) rectangle (9.25,1.6);
    \foreach \mbj in {0.55,0.9,1.25}{
      \draw[draw=white, line width=0.5pt] (7.45,\mbj) -- (9.25,\mbj);
    }
    \draw[draw=white, line width=0.5pt] (7.9,0.2) -- (7.9,1.6);
    \draw[draw=white, line width=0.5pt] (8.8,0.2) -- (8.8,1.6);
    \draw[draw=mbSlate!70, line width=0.8pt] (8.35,0.2) -- (8.35,1.6);
    \draw[draw=mbSlate!45, line width=0.5pt] (7.45,0.2) rectangle (9.25,1.6);
    % 4. trou : un carré bien franc posé au milieu de l'image
    \fill[mbBlue!10] (11.05,0.2) rectangle (12.95,1.6);
    \draw[draw=mbSlate!45, line width=0.5pt] (11.05,0.2) rectangle (12.95,1.6);
    \fill[mbSlate!60] (11.7,0.55) rectangle (12.4,1.15);
    \node[font=\sffamily\small\bfseries, text=white] at (12.05,0.85) {?};
  \end{scope}
\end{schema}

V-JEPA~2.1 a choisi la quatrième : le trou. On recouvre une partie de la
vidéo et l'on demande de deviner ce qui se trouvait dessous. Nous verrons
dans la partie~II que ce choix n'est pas un hasard : c'est celui qui oblige le
plus la machine à comprendre en même temps ce qu'il y a dans la scène et
comment cela bouge.

\begin{piege}
Ne cherchez jamais à juger un modèle auto-supervisé sur sa tâche prétexte.
Une machine qui remplit très bien les trous n'est pas forcément une machine
utile ; et une machine qui les remplit imparfaitement peut être excellente
pour tout le reste. C'est même parfois le contraire de ce qu'on croit : à la
fin de ce chapitre, notre petit programme atteindra un plancher d'erreur qu'il
est \emph{impossible} de franchir, et ce sera une très bonne nouvelle.

La seule question qui compte est celle-ci : \emph{ce que la machine a appris
en jouant sert-il à autre chose ?}
\end{piege}

Ce « quelque chose » que la machine garde après avoir joué, il faut
maintenant lui donner un nom. C'est le vrai produit de tout le travail, et
c'est le sujet de la section suivante.

\section{Ce que l'on garde vraiment : la représentation}

\begin{enclair}
Quand la machine regarde une image, elle en fabrique un petit résumé en
nombres. Ce résumé s'appelle une représentation. C'est lui, et lui seul, que
l'on garde à la fin. Tout le reste — le jeu des trous, les réponses données —
part à la poubelle.
\end{enclair}

Repensez à ce que vous faites quand un ami vous raconte un film. Il ne vous
récite pas les deux heures image par image. Il vous dit : « c'est l'histoire
d'un pêcheur qui cherche son fils ». En quelques mots, il a gardé ce qui
compte et jeté le reste. Ce petit paquet d'informations essentielles, c'est
exactement ce que la machine fabrique, sauf qu'elle l'écrit avec des nombres
au lieu de mots.

\begin{definition}[L'encodeur][encodeur]
Un \emph{encodeur} est la partie du modèle qui prend une image ou une vidéo
en entrée, et qui en ressort un résumé beaucoup plus petit. On peut le voir
comme un entonnoir : beaucoup de nombres entrent, peu de nombres sortent.
C'est la pièce maîtresse de tout le livre : à la fin, c'est l'encodeur, et
presque rien d'autre, que l'on met en production.
\end{definition}

\begin{definition}[La représentation][representation]
La \emph{représentation} d'une image est la liste de nombres que l'encodeur
produit à partir d'elle. On dit aussi les \emph{caractéristiques}, ou en
anglais les \emph{features}. Deux images qui se ressemblent doivent donner
deux représentations proches ; deux images très différentes, deux
représentations éloignées.
\end{definition}

\begin{definition}[L'espace latent][latent]
L'ensemble de tous les résumés possibles s'appelle l'\emph{espace latent}.
« Latent » veut dire « caché », « qui ne se voit pas directement » : ces
nombres ne sont ni des pixels ni des mots, ce sont les notes internes de la
machine. Dire qu'une représentation « vit dans l'espace latent », c'est
simplement dire que c'est une liste de nombres inventée par le modèle, et non
quelque chose que l'on peut afficher tel quel à l'écran.
\end{definition}

La figure~\ref{fig:entonnoir} montre l'encodeur au travail.

\begin{schema}{L'encodeur, dessiné comme un entonnoir. Une image entre par le
haut, avec ses centaines de milliers de nombres. Un petit paquet de nombres
sort par le bas : c'est la représentation. Ce petit paquet vit dans ce qu'on
appelle l'espace latent. Tout l'enjeu du livre est de le rendre aussi riche
que possible.}[entonnoir]
  \begin{scope}[x=1cm, y=1cm]
    % l'image d'entrée, dessinée en petite grille
    \foreach \mbi in {0,...,5}{
      \foreach \mbj in {0,...,3}{
        \pgfmathsetmacro{\mbton}{8 + 12*mod(\mbi+2*\mbj,4)}
        \draw[draw=mbSlate!35, fill=mbBlue!\mbton, line width=0.3pt]
              (\mbi*0.42, 3.9+\mbj*0.42) rectangle
              (\mbi*0.42+0.42, 3.9+\mbj*0.42+0.42);
      }
    }
    \node[font=\sffamily\scriptsize, text=mbSlate, align=center]
         at (-1.9,4.7) {une image\\ $196\,608$ nombres};
    \draw[fleche] (-0.75,4.7) -- (-0.15,4.7);
    % l'entonnoir
    \draw[draw=mbBlue, fill=mbBlue!8, line width=1pt]
          (0,3.6) -- (2.52,3.6) -- (1.6,1.7) -- (0.92,1.7) -- cycle;
    \node[font=\sffamily\scriptsize\bfseries, text=mbBlue] at (1.26,3.25)
         {l'encodeur};
    % la sortie : une petite barre de nombres
    \foreach \mbk/\mbv in {0/0.31, 1/-0.72, 2/1.04, 3/0.15, 4/-0.48}{
      \draw[draw=mbAmber!85!black, fill=mbAmber!14, line width=0.7pt]
            (-0.35+\mbk*0.72, 0.85) rectangle (0.32+\mbk*0.72, 1.35);
      \node[font=\sffamily\scriptsize, text=mbInk]
           at (-0.015+\mbk*0.72, 1.1) {\mbv};
    }
    \draw[fleche accent] (1.26,1.62) -- (1.26,1.42);
    \node[font=\sffamily\scriptsize\bfseries, text=mbAmber!75!black,
          align=center] at (1.26,0.45) {la représentation};
    \node[annotation, text width=3.6cm] at (5.7,1.1)
         {beaucoup moins de nombres,\\ mais ceux qui comptent};
    \draw[trait annotation] (3.9,1.1) -- (3.15,1.1);
    \node[annotation, text width=3.4cm] at (5.7,4.4)
         {rien n'est étiqueté ici :\\ c'est une image brute};
    \draw[trait annotation] (4.0,4.4) -- (2.75,4.4);
  \end{scope}
\end{schema}

Mettons des nombres sur ce dessin, car c'est plus parlant.

\begin{exemple}[De deux cent mille nombres à mille]
Une image en couleur de $256 \times 256$ pixels contient
$256 \times 256 \times 3 = 196\,608$ nombres : nous l'avions calculé au
chapitre~\ref{chap:image-video}. Supposons que l'encodeur la résume en
$1\,024$ nombres. Le rapport vaut

\[
  \frac{196\,608}{1\,024} = 192 .
\]

Le résumé est donc cent quatre-vingt-douze fois plus petit que l'image. Toute
la difficulté du métier tient dans cette division : il faut jeter cent
quatre-vingt-onze nombres sur cent quatre-vingt-douze, et jeter les bons.
\end{exemple}

\begin{remarque}
Vous vous demandez peut-être : pourquoi jeter ? Pourquoi ne pas tout garder ?
Parce qu'un résumé qui garde tout n'est pas un résumé : c'est une copie, et
une copie n'a rien compris. La contrainte de place est précisément ce qui
force la machine à distinguer l'important de l'accessoire. Dans le résumé
d'un film, « un pêcheur cherche son fils » tient en six mots parce qu'on a
renoncé à la couleur des rideaux.
\end{remarque}

Nous avons donc un jeu (la tâche prétexte), et un produit (la représentation).
Il nous manque encore une chose : un moyen de savoir si ce produit est bon.

\section{Comment savoir si une représentation est bonne}

\begin{enclair}
On bloque l'encodeur pour qu'il ne bouge plus, on branche derrière lui une
toute petite machine à apprendre, et on regarde jusqu'où elle arrive. Si
même une machine minuscule s'en sort bien, c'est que le résumé était déjà
bien rangé.
\end{enclair}

Voici l'idée, en deux gestes. Premier geste : on décide que l'encodeur ne
changera plus. Ses boutons sont fixés une fois pour toutes ; on dit qu'on le
\emph{gèle}.

\begin{definition}[L'encodeur gelé][gele]
Geler un encodeur, c'est décider que ses paramètres ne seront plus modifiés.
On s'en sert alors comme d'une machine toute faite : on lui donne une image,
il rend une représentation, et rien à l'intérieur ne bouge. C'est ainsi que
V-JEPA~2.1 est évalué du début à la fin du papier — et c'est aussi, nous le
verrons dans la partie~VII, ce qui rend son déploiement, c'est-à-dire sa mise
au travail pour de vrai, si peu coûteux.
\end{definition}

\begin{exemple}[Ce que le gel permet]
Imaginons que vous ayez dix mille photos et trois questions différentes à
leur poser : à quelle distance est chaque objet, où sont les voitures, et de
quelle sorte de scène il s'agit. Comme l'encodeur ne bouge plus, vous lui
faites lire vos dix mille photos une seule fois, vous rangez les dix mille
résumés dans un fichier, et vous n'y touchez plus. Vos trois questions se
règlent ensuite sur ces résumés déjà calculés, sans jamais rouvrir
l'encodeur. Chacune vous coûte quelques minutes, au lieu de recommencer un
entraînement complet à chaque fois.
\end{exemple}

Deuxième geste : on branche derrière lui une toute petite machine à apprendre,
qui, elle, a le droit d'apprendre. C'est elle qui aura besoin d'étiquettes —
mais très peu, puisqu'elle est minuscule.

\begin{definition}[La sonde][sonde]
Une \emph{sonde} est une petite machine à apprendre que l'on branche derrière
un encodeur gelé pour mesurer la qualité de ses représentations. Quand elle
se réduit à une seule \emph{couche} — c'est-à-dire une seule rangée de boutons
qui multiplient les nombres reçus et les additionnent — on parle de
\emph{sonde linéaire}. C'est un instrument de mesure, comme un thermomètre :
elle est volontairement faible, parce qu'une sonde faible ne peut pas
rattraper un mauvais résumé.
\end{definition}

\begin{schema}{Le principe de l'évaluation par sonde. L'encodeur porte un
cadenas : il est gelé, ses boutons ne bougent plus. Derrière lui, une toute
petite sonde apprend, elle, à répondre à une question précise. Si elle y
arrive avec très peu d'exemples, c'est que l'information était déjà bien
rangée dans la représentation.}[sonde]
  \begin{scope}[x=1cm, y=1cm]
    \draw[draw=mbSlate!55, fill=white, line width=0.7pt]
          (0,0.55) rectangle (1.1,1.55);
    \node[font=\sffamily\scriptsize, text=mbSlate] at (0.55,0.2) {une image};
    \draw[fleche] (1.25,1.05) -- (1.9,1.05);
    % encodeur gelé
    \node[bloc, minimum width=24mm, minimum height=13mm] (enc) at (3.2,1.05)
         {l'encodeur};
    % le cadenas, dessiné à la main : un corps et une anse
    \begin{scope}[shift={(3.2,2.15)}]
      \draw[draw=mbSlate, line width=0.9pt] (-0.16,0.1) arc (180:0:0.16);
      \draw[draw=mbSlate, fill=mbMist, line width=0.9pt]
            (-0.24,-0.28) rectangle (0.24,0.12);
      \fill[mbSlate] (0,-0.08) circle (0.05);
    \end{scope}
    \node[font=\sffamily\scriptsize\itshape, text=mbSlate] at (3.2,2.7)
         {gelé : plus aucun bouton ne bouge};
    \draw[fleche accent] (4.45,1.05) -- (5.15,1.05);
    \node[font=\sffamily\scriptsize, text=mbAmber!75!black]
         at (4.8,0.2) {la représentation};
    % la sonde, volontairement petite
    \node[bloc accent, minimum width=13mm, minimum height=8mm] (sonde)
         at (5.9,1.05) {\scriptsize la sonde};
    \draw[fleche] (6.6,1.05) -- (7.25,1.05);
    \node[font=\sffamily\small, text=mbInk, align=center] at (8.15,1.05)
         {la réponse\\ \scriptsize « chat »};
    \node[annotation, text width=3.1cm] at (8.0,2.35)
         {toute petite, et c'est\\ la seule qui apprend};
    \draw[trait annotation] (7.0,2.1) -- (6.15,1.55);
  \end{scope}
\end{schema}

\begin{intuition}
Pourquoi prendre une sonde volontairement faible ? Parce qu'on cherche à
mesurer la qualité du \emph{rangement}, pas la puissance de l'instrument de
mesure.

Imaginez deux ateliers. Dans le premier, les outils sont rangés par taille,
étiquetés, à portée de main. Dans le second, tout est jeté en vrac dans une
caisse. Envoyez un enfant chercher une clé de douze : dans le premier
atelier, il la trouve en dix secondes ; dans le second, jamais. Si vous
envoyez un mécanicien professionnel, il finira par la trouver dans les deux
cas — et vous n'aurez rien appris sur le rangement. La sonde, c'est l'enfant.
\end{intuition}

Nous reviendrons longuement sur ces mesures dans la partie~III, où nous
verrons comment on note vraiment une représentation. Pour l'instant, retenez
le principe : \emph{encodeur gelé, petite sonde, peu d'étiquettes}.

Ces deux gestes découpent le travail en deux temps. Le premier porte un nom
qu'il faut absolument connaître, car il revient dans tous les articles et
dans tout le reste du livre.

\begin{definition}[Le pré-entraînement][preentrainement]
Le \emph{pré-entraînement} est le long entraînement de départ, celui où la
machine joue à la tâche prétexte sur des montagnes de données brutes, sans
aucune étiquette. Le préfixe « pré » veut dire « avant » : c'est ce qui se
passe \emph{avant} qu'on lui demande quoi que ce soit d'utile. C'est de loin
la partie la plus coûteuse, et elle ne se fait qu'une seule fois.
\end{definition}

\begin{exemple}[Les deux temps, en chiffres]
Pour V-JEPA~2.1, le pré-entraînement, c'est cent trente-cinq mille pas sur
cent soixante-trois millions d'images et de vidéos, sans une seule étiquette.
C'est un travail que l'on ne fait qu'une fois, et les auteurs en ont publié
le résultat pour que personne n'ait à le recommencer.

L'utilisation, c'est tout autre chose : une petite sonde, quelques milliers
d'exemples étiquetés, et l'affaire est réglée. C'est ce deuxième temps, et
lui seul, que vous ferez tourner sur votre machine — y compris à la fin de ce
chapitre.
\end{exemple}

\begin{aretenir}
L'auto-supervision se déroule toujours en deux temps, et il ne faut jamais
les confondre :

\begin{enumerate}
  \item \textbf{Le pré-entraînement.} La machine joue au jeu des trous sur des
        montagnes de données brutes, sans la moindre étiquette. C'est long,
        c'est cher en calcul, et cela ne se fait qu'une fois. C'est ici que
        naît la représentation.
  \item \textbf{L'utilisation}, que l'on appelle aussi l'évaluation quand on
        cherche à noter le modèle. On gèle l'encodeur et l'on branche derrière
        une petite sonde, entraînée avec un tout petit nombre d'étiquettes,
        pour la tâche qui nous intéresse vraiment.
\end{enumerate}

Les étiquettes n'ont donc pas complètement disparu : elles ont été repoussées
à la toute fin, en quantité minuscule.
\end{aretenir}

Voyons maintenant à quelle échelle tout cela se joue dans le papier que nous
expliquons.

\section{À quelle échelle : les chiffres de V-JEPA 2.1}

Jusqu'ici, nous avons parlé de dix mille photos, puis de cent quarante-deux
millions. Mettons les vrais chiffres du papier sur la table, pour que vous
sachiez de quoi il retourne.

\begin{surlesdonnees}
Le jeu de données de V-JEPA~2.1 s'appelle \emph{VisionMix163M}. Son nom vient
tout simplement du nombre d'\emph{échantillons} qu'il contient — un
échantillon, c'est un élément de la collection, ici une image ou une vidéo :

\begin{center}
\footnotesize
\begin{tabular}{lrl}
\toprule
\textbf{Source} & \textbf{Échant.} & \textbf{Ce que c'est} \\
\midrule
LVD-142M    & $142{,}0$ M & des images choisies avec soin \\
YT-1B       & $19{,}0$ M  & des vidéos prises sur internet \\
HowTo100M   & $1{,}1$ M   & des vidéos qui expliquent \\
Kinetics    & $0{,}733$ M & des actions filmées de loin \\
SSv2        & $0{,}168$ M & des gestes filmés de près \\
\midrule
\textbf{Total} & $\approx 163$ M & d'où le nom \\
\bottomrule
\end{tabular}
\end{center}

Faites l'addition vous-même :
$142 + 19 + 1{,}1 + 0{,}733 + 0{,}168 = 163{,}0$ millions. Mises bout à bout,
ces vidéos durent environ $1{,}7$ million d'heures, dont $1\,600\,000$ pour la
seule source YT-1B.

Et voici le point qui doit rester : \textbf{aucun de ces cent soixante-trois
millions d'échantillons n'a été étiqueté à la main pour cet entraînement}.
Le calcul de la section~2 nous a montré qu'il aurait fallu, au minimum, cent
vingt-huit ans pour les images seules. Ce livre entier raconte comment on
s'en passe.
\end{surlesdonnees}

\begin{remarque}
Vous avez peut-être remarqué que certains de ces jeux de données possèdent
bel et bien des étiquettes : Kinetics et SSv2 sont même nés pour cela. Elles
existent, mais V-JEPA~2.1 ne les regarde pas pendant le pré-entraînement. Il
ne prend que les pixels. Les étiquettes ne réapparaissent qu'à la toute fin,
pour entraîner les sondes et pour noter les résultats.
\end{remarque}

Assez parlé. Il est temps de faire tourner tout cela pour de vrai, sur votre
machine, avec un programme que vous pouvez lire en entier.

\section{Passons au code}

Nous allons construire le plus petit exemple d'auto-supervision possible, et
vérifier de nos yeux qu'il fonctionne. Le programme tient en une page, et il
suit exactement les deux temps de l'encadré précédent : d'abord jouer au jeu
des trous sans aucune étiquette, ensuite geler et brancher une sonde.

Voici notre monde miniature. Chaque image fait $8 \times 8$ cases, et contient
une seule barre : soit une barre \emph{couchée}, c'est-à-dire horizontale,
soit une barre \emph{debout}, c'est-à-dire verticale. Elle est posée
n'importe où. Nous dirons toujours « couchée » et « debout », y compris dans
le programme.

\begin{schema}{Notre monde miniature et le jeu que nous allons y jouer.
Chaque image contient une barre, couchée ou debout. On coupe l'image en deux :
la moitié gauche est donnée à la machine, la moitié droite est cachée et
devient la bonne réponse. La machine ne saura jamais qu'il existe des mots
comme « couchée » et « debout ».}[monde-barres]
  \begin{scope}[x=1cm, y=1cm]
    % --- image 1 : barre horizontale ---
    \begin{scope}[shift={(0,0)}]
      \foreach \mbi in {0,...,7}{
        \foreach \mbj in {0,...,7}{
          \draw[draw=mbSlate!30, fill=white, line width=0.3pt]
                (\mbi*0.3,\mbj*0.3) rectangle (\mbi*0.3+0.3,\mbj*0.3+0.3);
        }
      }
      \foreach \mbi in {0,...,7}{
        \draw[draw=mbSlate!30, fill=mbBlue!45, line width=0.3pt]
              (\mbi*0.3,1.5) rectangle (\mbi*0.3+0.3,1.8);
      }
      \node[font=\sffamily\scriptsize, text=mbSlate] at (1.2,-0.4)
           {barre couchée};
    \end{scope}
    % --- image 2 : barre verticale ---
    \begin{scope}[shift={(3.1,0)}]
      \foreach \mbi in {0,...,7}{
        \foreach \mbj in {0,...,7}{
          \draw[draw=mbSlate!30, fill=white, line width=0.3pt]
                (\mbi*0.3,\mbj*0.3) rectangle (\mbi*0.3+0.3,\mbj*0.3+0.3);
        }
      }
      \foreach \mbj in {0,...,7}{
        \draw[draw=mbSlate!30, fill=mbTeal!45, line width=0.3pt]
              (0.6,\mbj*0.3) rectangle (0.9,\mbj*0.3+0.3);
      }
      \node[font=\sffamily\scriptsize, text=mbSlate] at (1.2,-0.4)
           {barre debout};
    \end{scope}
    % --- la coupe en deux ---
    \begin{scope}[shift={(6.6,0)}]
      \foreach \mbi in {0,...,7}{
        \foreach \mbj in {0,...,7}{
          \draw[draw=mbSlate!30, fill=white, line width=0.3pt]
                (\mbi*0.3,\mbj*0.3) rectangle (\mbi*0.3+0.3,\mbj*0.3+0.3);
        }
      }
      \foreach \mbi in {0,...,7}{
        \draw[draw=mbSlate!30, fill=mbBlue!45, line width=0.3pt]
              (\mbi*0.3,0.9) rectangle (\mbi*0.3+0.3,1.2);
      }
      \fill[mbSlate!45, opacity=0.85] (1.2,0) rectangle (2.4,2.4);
      \draw[draw=mbSlate, line width=1pt, dashed] (1.2,-0.1) -- (1.2,2.5);
      \node[font=\sffamily\scriptsize, text=mbSlate] at (0.6,-0.4)
           {donné};
      \node[font=\sffamily\scriptsize, text=white] at (1.8,1.2) {caché};
      \node[font=\sffamily\scriptsize, text=mbSlate] at (1.8,-0.4)
           {à deviner};
    \end{scope}
    \draw[accolade] (6.6,2.6) -- (9.0,2.6);
    \node[annotation, text width=3.4cm] at (7.8,3.4)
         {la question et la réponse\\ sortent de la même image :\\
          aucune étiquette};
  \end{scope}
\end{schema}

Commençons par fabriquer ces images. Le programme ci-dessous ne contient rien
que nous n'ayons expliqué : des tenseurs, comme au
chapitre~\ref{chap:image-video}, et des tirages au hasard.

\begin{code}{python}[Fabriquer le monde miniature]
import torch

torch.manual_seed(0)          # pour que vous obteniez les memes nombres

def fabrique_images(n):
    """n images de 8x8 : soit une barre couchee, soit une barre debout."""
    images = torch.zeros(n, 8, 8)
    genre = torch.randint(0, 2, (n,))    # 0 = couchee, 1 = debout
    place = torch.randint(0, 8, (n,))    # a quelle ligne, ou a quelle colonne
    for i in range(n):
        if genre[i] == 0:
            images[i, place[i], :] = 1.0     # une ligne entiere
        else:
            images[i, :, place[i]] = 1.0     # une colonne entiere
    return images, genre.float()

images, genre = fabrique_images(2000)
gauche = images[:, :, :4].reshape(-1, 32)   # 8 lignes x 4 colonnes = 32
droite = images[:, :, 4:].reshape(-1, 32)   # l'autre moitie
\end{code}

Deux mots sur les deux dernières lignes, car elles contiennent un geste que
nous referons souvent. L'écriture \co{python}{images[:, :, :4]} veut dire :
« garde toutes les images, toutes les lignes, mais seulement les quatre
premières colonnes » — c'est notre moitié gauche. Quant à
\co{python}{reshape(-1, 32)}, il met chaque petite grille de $8 \times 4$ à
plat, en une seule rangée de $8 \times 4 = 32$ nombres, comme si l'on
découpait les lignes d'un tableau pour les recoller bout à bout. Le $-1$ dit
simplement à la machine : « débrouille-toi pour le nombre d'images, tu sais
compter ».

Notez bien que \co{python}{genre} existe. C'est une \emph{variable},
c'est-à-dire un nom que nous donnons à une valeur pour pouvoir la retrouver
plus loin dans le programme ; celle-ci retient, pour chaque image, si la
barre était couchée ou debout. Nous allons pourtant l'ignorer complètement
pendant tout le pré-entraînement : elle attendra sagement la toute fin, pour
la sonde. C'est notre façon d'imiter la réalité : les étiquettes existent
quelque part dans le monde, mais on ne s'en sert presque pas.

Construisons maintenant les deux pièces du modèle. La première, nous la
connaissons : c'est l'encodeur, qui résume les trente-deux nombres de la
moitié gauche en seize nombres. La seconde est nouvelle, et il faut lui
donner son nom.

\begin{definition}[Le prédicteur][predicteur]
Le \emph{prédicteur} est la pièce qui part du résumé fabriqué par l'encodeur
et qui tente de deviner ce qui a été caché. C'est lui qui joue, à proprement
parler, au jeu des trous. À la fin de l'entraînement, c'est pourtant
l'encodeur que l'on garde précieusement ; le prédicteur, lui, n'a servi qu'à
faire travailler son voisin. Nous verrons dans la partie~V qu'il sait tout de
même rendre un dernier service, quand on veut deviner la suite d'une vidéo.
\end{definition}

Chez nous, ce prédicteur part des seize nombres du résumé et tente de
reconstituer les trente-deux nombres cachés.

Trois outils de PyTorch apparaissent pour la première fois. Prenons-les un par
un, car nous les retrouverons dans tout le livre.

\begin{enclair}
\co{python}{torch.nn.Linear(32, 16)} fabrique une \emph{couche} : une rangée
de boutons qui prend trente-deux nombres en entrée, les multiplie chacun par
un poids, additionne le tout, et rend seize nombres. C'est exactement la
machine à bonbons du chapitre~\ref{chap:apprendre}, en beaucoup plus large :
au lieu des deux boutons $a$ et $b$, celle-ci en possède
$32 \times 16 + 16 = 528$.

\co{python}{torch.nn.ReLU()} est une règle minuscule appliquée après la
couche : tout nombre négatif est remplacé par zéro, les autres passent tels
quels. Sans elle, empiler deux couches ne servirait à rien, car deux
multiplications à la suite se ramènent toujours à une seule. C'est cette
petite règle qui permet à un modèle d'apprendre autre chose que des droites.
Quant à \mbox{\co{python}{Sequential}}, il sert simplement à mettre ces deux
pièces bout à bout.

\co{python}{torch.optim.Adam} remplace la ligne
\co{python}{a -= pas * a.grad} que nous écrivions à la main. C'est la même
descente de gradient, avec une amélioration : Adam règle tout seul la
longueur du pas de chaque bouton, en se souvenant des pentes passées. Le
réglage \co{python}{lr=0.01} est le taux d'apprentissage de départ, celui du
chapitre~\ref{chap:apprendre}. Restent deux lignes de service.
\co{python}{optim.step()} bouge tous les boutons d'un coup, et
\co{python}{optim.zero_grad()} efface les pentes avant le tour suivant. Ce
sont les deux gestes que nous écrivions bouton par bouton au chapitre
précédent.
\end{enclair}

Ces trois outils paraissent abstraits tant qu'on ne les a pas vus tourner sur
de vrais nombres. Faisons-le en tout petit.

\begin{exemple}[Une couche minuscule, calculée à la main]
Prenons une couche qui reçoit deux nombres et n'en rend qu'un seul. Elle
possède donc deux poids — disons $0{,}5$ et $-2$ — plus un décalage de $4$,
qui est exactement le bouton $b$ de la machine à bonbons. Donnons-lui les
nombres $3$ et $1$ :
\[
  0{,}5 \times 3 \;+\; (-2) \times 1 \;+\; 4
  \;=\; 1{,}5 - 2 + 4 \;=\; 3{,}5 .
\]
La couche rend donc $3{,}5$, et la règle ReLU le laisse passer tel quel
puisqu'il est positif.

Refaisons le calcul avec les nombres $1$ et $4$ :
\[
  0{,}5 \times 1 \;+\; (-2) \times 4 \;+\; 4
  \;=\; 0{,}5 - 8 + 4 \;=\; -3{,}5 .
\]
Cette fois le résultat est négatif : ReLU le remplace par $0$, et c'est ce
zéro qui poursuivra le voyage. Voilà tout le mystère. Notre vrai encodeur
fait la même chose, mais avec trente-deux nombres à l'entrée, seize à la
sortie, et $528$ boutons au lieu de trois.
\end{exemple}

\begin{code}{python}[Le pré-entraînement, sans la moindre étiquette]
encodeur = torch.nn.Sequential(torch.nn.Linear(32, 16), torch.nn.ReLU())
predicteur = torch.nn.Linear(16, 32)

params = list(encodeur.parameters()) + list(predicteur.parameters())
optim = torch.optim.Adam(params, lr=0.01)

for etape in range(1, 601):
    erreur = ((predicteur(encodeur(gauche)) - droite) ** 2).mean()
    optim.zero_grad()
    erreur.backward()
    optim.step()
    if etape in (1, 100, 300, 600):
        print(f"etape {etape:3d} : erreur = {erreur.item():.4f}")
\end{code}

Vous reconnaissez la boucle en cinq gestes du chapitre~\ref{chap:apprendre},
à un détail près : la bonne réponse, \co{python}{droite}, n'a été écrite par
personne. Elle a été découpée dans l'image elle-même.

Voici ce que le programme affiche :

\begin{sortie}
etape   1 : erreur = 0.1648
etape 100 : erreur = 0.0487
etape 300 : erreur = 0.0483
etape 600 : erreur = 0.0483
\end{sortie}

L'erreur tombe de $0{,}1648$ à $0{,}0483$, puis ne bouge plus. Cette
stagnation n'est pas une panne, et c'est même le passage le plus instructif
du chapitre.

\begin{piege}
Quand une barre est debout dans la moitié \emph{droite} de l'image, la moitié
gauche est entièrement vide. On demande alors à la machine de deviner une
barre dont absolument rien, dans ce qu'elle voit, n'indique la position.
Vous-même, devant une moitié gauche toute blanche, ne feriez pas mieux.

Dans notre tirage, cela arrive pour $25{,}8\,\%$ des images. Il existe donc
un plancher sous lequel l'erreur ne peut pas descendre, et le modèle l'a
simplement atteint. Ne confondez jamais « le modèle a cessé de progresser »
avec « le modèle est mauvais » : il faut d'abord se demander si ce qu'on lui
demande est seulement possible.
\end{piege}

Passons au deuxième temps : geler, brancher une sonde, et compter les points.
Nous n'accorderons à la sonde que \emph{vingt} exemples étiquetés — c'est
volontairement très peu.

Une dernière nouveauté, et ce sera la seule. Jusqu'ici, notre erreur était la
moyenne des carrés des écarts, parce que nous prédisions des nombres. Ici, la
sonde doit répondre à une question par oui ou par non : cette barre est-elle
debout ? Pour ce genre de question, on emploie une autre mesure d'erreur,
\co{python}{binary_cross_entropy_with_logits}. Son nom fait peur, mais son
idée tient en une phrase : elle punit d'autant plus la machine qu'elle s'est
trompée \emph{avec assurance}. Répondre « oui » à moitié convaincu quand la
réponse était non coûte peu ; répondre « oui, absolument certain » quand la
réponse était non coûte très cher.

La sonde rend un seul nombre par image. On lit ce nombre ainsi : s'il est
plus grand que zéro, la machine dit « debout » ; sinon, elle dit « couchée ».
C'est le rôle du test \co{python}{> 0} à l'avant-dernière ligne. Au passage,
le mot \emph{logits}, dans le nom de la mesure d'erreur, désigne simplement
ce nombre brut, celui qu'on lit avant de le traduire en « oui » ou en
« non ». Quant au petit \co{python}{squeeze(1)}, il ne fait que retirer un
axe devenu inutile, pour que les formes des deux tenseurs se correspondent :
rien de plus qu'un rangement, au sens du chapitre~\ref{chap:image-video}.

\begin{code}{python}[Geler, puis brancher une sonde sur vingt étiquettes]
def sonde(entrees, cibles, entrees_test, cibles_test):
    """Une seule couche, qui apprend a repondre : couchee ou debout ?"""
    tete = torch.nn.Linear(entrees.shape[1], 1)
    opt = torch.optim.Adam(tete.parameters(), lr=0.05)
    for _ in range(300):
        p = tete(entrees).squeeze(1)
        perte = torch.nn.functional.binary_cross_entropy_with_logits(p, cibles)
        opt.zero_grad()
        perte.backward()
        opt.step()
    with torch.no_grad():
        prediction = (tete(entrees_test).squeeze(1) > 0).float()
        return (prediction == cibles_test).float().mean().item()

test_images, test_genre = fabrique_images(1000)
test_gauche = test_images[:, :, :4].reshape(-1, 32)

with torch.no_grad():                  # l'encodeur est gele : aucun gradient
    h = encodeur(gauche)
    h_test = encodeur(test_gauche)

torch.manual_seed(1)
print("sonde sur la representation :",
      f"{sonde(h[:20], genre[:20], h_test, test_genre)*100:.1f} %")
torch.manual_seed(1)
print("sonde sur les pixels bruts  :",
      f"{sonde(gauche[:20], genre[:20], test_gauche, test_genre)*100:.1f} %")
\end{code}

Et voici le résultat, mesuré sur mille images que la machine n'avait jamais
vues :

\begin{sortie}[Le verdict]
sonde sur la representation : 80.4 %
sonde sur les pixels bruts  : 61.8 %
\end{sortie}

Arrêtons-nous, car ces deux lignes résument tout le chapitre.

La même sonde, avec les mêmes vingt étiquettes et le même nombre d'étapes,
obtient $80{,}4\,\%$ de bonnes réponses quand on la branche sur la
représentation, contre $61{,}8\,\%$ quand on la branche directement sur les
pixels. La différence ne vient pas de la sonde, qui est identique dans les
deux cas. Elle vient de ce que l'encodeur a appris \emph{en jouant à un jeu
qui n'avait rien à voir}. Personne ne lui a jamais montré le mot « debout »
ni le mot « couchée ».

\begin{remarque}
Et $80{,}4\,\%$, est-ce beaucoup ? Refaisons le calcul du piège précédent. Dans
$25{,}8\,\%$ des cas, la moitié gauche est vide et personne ne peut répondre
autrement qu'en tirant à pile ou face, ce qui rapporte la moitié des points,
soit environ $12{,}9\,\%$. Le meilleur score atteignable est donc de l'ordre de
$74{,}2 + 12{,}9 \approx 87\,\%$. Avec $80{,}4\,\%$, notre encodeur n'est pas
loin du maximum possible, et il l'a obtenu avec vingt étiquettes.
\end{remarque}

Vous n'avez pas à recopier ces morceaux à la main : le programme entier, tel
que nous venons de le construire, est rangé dans le fichier
\term{scripts/03-auto-supervision.py} qui accompagne ce livre. Lancez-le par
la commande \co{text}{python 03-auto-supervision.py} et vous devez retrouver,
à la décimale près, les nombres imprimés ci-dessus.

Si vous obtenez autre chose, regardez d'abord la toute première ligne,
\co{python}{torch.manual_seed(0)}. Un ordinateur ne sait pas tirer au sort
comme on lance un dé : il fabrique des nombres qui \emph{ressemblent} à du
hasard, en partant d'un nombre de départ qu'on appelle la \emph{graine}.
Avec la même graine, il refait exactement les mêmes tirages, dans le même
ordre. C'est ce qui vous permet de retrouver nos barres, nos vingt étiquettes
et nos résultats à l'identique ; changez le $0$ en $1$, et tout le monde
miniature change avec.

\begin{enia}
Changez maintenant les échelles, et vous avez V-JEPA~2.1. Les images de
$8 \times 8$ deviennent des vidéos de soixante-quatre images en
$384 \times 384$. Les seize nombres de la représentation deviennent des
milliers. Les six cents étapes deviennent cent trente-cinq mille. Les deux
mille images deviennent cent soixante-trois millions d'images et de vidéos.
Le prédicteur, lui, devient une machine à calculer bien plus grosse, que nous
démonterons pièce par pièce le moment venu.

Mais la structure du programme, elle, ne change pas d'une ligne : cacher,
deviner, mesurer, corriger — puis geler et brancher une sonde.
\end{enia}

\begin{aretenir}
Les trois idées du chapitre, en trois phrases.

\begin{enumerate}
  \item Les étiquettes coûtent trop cher pour être produites à l'échelle des
        vidéos : il fallait une autre source de bonnes réponses.
  \item Cette source, c'est la donnée elle-même : on cache un morceau, et le
        morceau caché fait office de bonne réponse.
  \item Le jeu ne compte pas ; ce qui compte, c'est la représentation que la
        machine fabrique en y jouant, et que l'on mesure avec une petite
        sonde derrière un encodeur gelé.
\end{enumerate}
\end{aretenir}

Il reste pourtant une question que nous avons soigneusement contournée, et
elle est bien plus profonde qu'elle n'en a l'air. Dans notre petit programme,
nous avons demandé à la machine de deviner les \emph{pixels} cachés : les
trente-deux nombres de la moitié droite, un par un. Est-ce vraiment ce qu'il
faut demander ?

Réfléchissez un instant à ce que cela implique. Pour redessiner exactement
les pixels d'une zone cachée dans une vraie vidéo, il faudrait deviner chaque
brin d'herbe, chaque grain de la peau, chaque reflet dans une flaque. Or
personne au monde ne peut prédire cela, et surtout, personne n'en a besoin.
La machine passerait le plus clair de son temps sur un travail inutile.

Il existe une autre possibilité, beaucoup plus fine : lui demander de deviner
non pas les pixels cachés, mais \emph{le résumé} de la zone cachée. Non pas
l'herbe brin par brin, mais l'idée « ici, il y a de l'herbe ». C'est-à-dire
prédire dans l'espace latent que nous venons de définir.

Cette question — \emph{dans quel espace faut-il deviner : celui des pixels,
ou celui des idées ?} — est le cœur de tout le papier que ce livre
explique. Sa réponse porte un nom, JEPA, et c'est le sujet de la partie~II.

Avant d'y arriver, il nous manque encore une pièce. Nous savons ce que la
machine doit apprendre, mais nous n'avons pas encore vu \emph{avec quoi} elle
regarde. Quel est l'œil de la machine, celui qui découpe une vidéo en
morceaux et les compare entre eux ? C'est le sujet du chapitre suivant.

\phantomsection
\section*{Exercices}
\addcontentsline{toc}{section}{Exercices}

\begin{exercice}[Le prix d'une étiquette]
Vous devez faire étiqueter $500\,000$ images. Chaque étiquette prend huit
secondes. Vous employez cinq personnes, qui travaillent chacune sept heures
par jour. Combien de jours faudra-t-il ?
\end{exercice}

\begin{solution}
Le travail total vaut $500\,000 \times 8 = 4\,000\,000$ secondes, soit
$4\,000\,000 / 3\,600 \approx 1\,111$ heures. Chaque jour, les cinq personnes
fournissent ensemble $5 \times 7 = 35$ heures. Il faut donc

\[
  \frac{1\,111}{35} \approx 32 \text{ jours.}
\]

Un mois de travail à cinq, pour un demi-million d'images seulement. Multipliez
par deux cent quatre-vingt-quatre pour atteindre les cent quarante-deux
millions d'images de V-JEPA~2.1, et vous retrouvez l'ordre de grandeur de la
section~2.
\end{solution}

\begin{exercice}[Fabriquer une tâche prétexte]
On vous donne un million de vidéos sans aucune étiquette. Proposez deux jeux
que la machine pourrait jouer toute seule, et dites à chaque fois d'où sort
la bonne réponse.
\end{exercice}

\begin{solution}
Premier jeu : remettre les images dans l'ordre. On mélange l'ordre des images
d'un clip et l'on demande de le rétablir ; la bonne réponse est l'ordre
d'origine, que l'on connaît puisque c'est nous qui l'avons mélangé.

Deuxième jeu : deviner la suite. On coupe le clip en deux et l'on demande de
prédire la seconde moitié à partir de la première ; la bonne réponse est la
seconde moitié, que l'on possède déjà.

Dans les deux cas, la règle est la même : on abîme volontairement la donnée
d'une façon dont on garde la trace, et cette trace est la bonne réponse.
\end{solution}

\begin{exercice}[Le taux de compression]
Un clip de $16$ images en $256 \times 256$ pixels couleur est résumé par un
encodeur en $512$ nombres. Par combien l'information a-t-elle été divisée ?
\end{exercice}

\begin{solution}
Le clip contient $16 \times 256 \times 256 \times 3 = 3\,145\,728$ nombres,
comme nous l'avions calculé au chapitre~\ref{chap:image-video}. Le rapport
vaut

\[
  \frac{3\,145\,728}{512} = 6\,144 .
\]

Le résumé est six mille cent quarante-quatre fois plus petit que le clip. Toute la
question est de savoir si l'on a jeté les bons nombres — et c'est précisément
ce que mesure une sonde.
\end{solution}

\begin{exercice}[Lire un résultat de sonde]
Deux encodeurs sont évalués avec exactement la même sonde et les mêmes
étiquettes. Le premier obtient $78\,\%$, le second $62\,\%$. Que peut-on en
conclure, et que ne peut-on \emph{pas} en conclure ?
\end{exercice}

\begin{solution}
On peut conclure que le premier encodeur range mieux, dans ses
représentations, l'information dont cette tâche précise a besoin : une simple
sonde y accède plus facilement.

On ne peut pas conclure que le premier encodeur est meilleur en général. Il
pourrait être moins bon sur une autre tâche, par exemple mesurer des
distances plutôt que reconnaître des objets. Nous verrons dans la partie~V
que c'est exactement ce qui arrive à V-JEPA~2 : excellent pour reconnaître une
action, faible pour décrire chaque pixel. Un chiffre de sonde ne vaut que
pour la question posée.
\end{solution}

\begin{exercice}[Le plancher d'erreur]
Dans le programme de ce chapitre, l'erreur cesse de descendre à $0{,}0483$.
Un collègue vous dit : « ton modèle est bloqué, augmente le taux
d'apprentissage ». Que lui répondez-vous ?
\end{exercice}

\begin{solution}
Qu'il faut d'abord vérifier si la tâche est possible. Ici, dans $25{,}8\,\%$
des cas, la moitié gauche de l'image est entièrement vide alors qu'une barre
se cache à droite : aucune information ne permet de la placer. Une partie de
l'erreur est donc irréductible, et le modèle a simplement atteint ce plancher.

Augmenter le taux d'apprentissage ne servirait à rien, et risquerait même de
faire exploser l'entraînement, comme nous l'avons vu au
chapitre~\ref{chap:apprendre}. La bonne réaction est de se demander ce que la
machine \emph{pourrait} savoir, avant de lui reprocher de ne pas le savoir.
\end{solution}
            % apprendre sans étiquettes
\chapter{Le transformeur de vision, l'œil de la machine}
\label{chap:vit}

\begin{quote}\itshape
Nous savons ce qu'est une image pour un ordinateur, nous savons ce que veut
dire apprendre, et nous savons jouer au jeu des trous. Il nous manque une
chose, et une seule : l'appareil qui regarde. Dans ce chapitre, nous le
construisons pièce par pièce, sans en sauter aucune. À la fin, vous saurez
exactement ce qui se passe entre le moment où une vidéo entre dans la
machine et le moment où elle en ressort sous forme de résumé.
\end{quote}

\begin{prerequis}
Les chapitres~\ref{chap:image-video}, \ref{chap:apprendre}
et~\ref{chap:auto-supervision}. Il faut savoir qu'une image est une grille
de nombres et connaître sa forme ; il faut savoir qu'un modèle possède des
boutons que l'on tourne pour réduire une erreur ; et il faut se souvenir
qu'un encodeur transforme une image en un petit paquet de nombres appelé
représentation. Nous n'aurons besoin de rien d'autre. Les seules opérations
employées seront l'addition, la multiplication et la division.
\end{prerequis}

Au chapitre précédent, nous avons dit qu'une machine cache un morceau de ce
qu'elle regarde et essaie de le deviner. Nous avons même écrit un petit
programme qui le faisait. Mais nous avons triché sur un point : notre monde
tenait dans une grille de huit cases sur huit, et notre modèle était une
simple couche de calcul. Une vraie vidéo, elle, contient des millions de
nombres. Il faut donc une machine à regarder beaucoup plus solide.

Cette machine porte un nom : le \emph{transformeur de vision}. Tout le
monde l'appelle par son surnom anglais, \emph{ViT}, qui se prononce
«~vi-té~» et veut dire \emph{Vision Transformer}. C'est l'œil de V-JEPA~2.1,
et c'est aussi l'œil de presque tous les grands modèles d'images
d'aujourd'hui. Nous allons le monter devant vous, morceau par morceau.

\section{Pourquoi on ne peut pas regarder les pixels un par un}

\begin{enclair}
L'idée du transformeur, c'est de laisser chaque morceau de l'image regarder
tous les autres morceaux. Le problème, c'est que si les morceaux sont des
pixels, il y en a beaucoup trop : le nombre de regards à échanger devient
si grand qu'aucune machine au monde ne pourrait le calculer.
\end{enclair}

Commençons par la question qui commande tout le reste. Pour comprendre une
image, il ne suffit pas de regarder chaque endroit tout seul dans son coin.
Il faut les mettre en rapport les uns avec les autres.

Prenons un exemple. Sur une photo, vous voyez une oreille pointue en haut à
gauche, des moustaches au milieu, et une queue touffue en bas à droite.
Aucun de ces trois morceaux ne dit «~chat~» à lui tout seul : une oreille
pointue, cela peut être un renard ; des moustaches, cela peut être un phoque.
C'est le \emph{rapprochement} des trois qui donne la réponse. Une bonne
machine à regarder doit donc permettre à chaque endroit de l'image de
consulter tous les autres endroits.

Comptons maintenant le prix de cette consultation. Imaginons que chaque
pixel consulte chaque autre pixel. Dans une image de $256 \times 256$, il y
a $65\,536$ pixels. Le nombre de consultations vaut alors

\[
  65\,536 \times 65\,536 = 4\,294\,967\,296 ,
\]

soit plus de quatre milliards de consultations, et cela pour \emph{une
seule} image, à \emph{un seul} endroit du modèle. Un clip de seize images
contient seize fois plus de pixels, donc $16 \times 16 = 256$ fois plus de
consultations encore. C'est tout simplement impossible.

\begin{intuition}
Imaginez une salle où chaque personne doit regarder toutes les autres, une
par une. À dix personnes, cela fait $10 \times 10 = 100$ regards : on s'en
sort en une minute. À mille personnes, cela fait $1\,000 \times 1\,000 =
1\,000\,000$ de regards : la réunion durerait des semaines. Le nombre de
personnes a été multiplié par cent, mais le travail a été multiplié par dix
mille. C'est ce qu'on appelle un coût \emph{qui grimpe au carré}, et c'est
l'ennemi numéro un de tout ce chapitre.
\end{intuition}

\begin{definition}[Un coût qui grimpe au carré][cout-carre]
On dit qu'un coût \emph{grimpe au carré} quand il suit le nombre d'éléments
multiplié par lui-même. Appelons $n$ le nombre d'éléments — écrire une
lettre à la place d'un nombre permet simplement de parler de tous les cas à
la fois. Si chacun de ces $n$ éléments doit consulter tous les autres, le
nombre de consultations vaut $n \times n$, ce qui s'écrit aussi $n^2$ et se
lit «~$n$ au carré~».
\end{definition}

\begin{exemple}[Voir le coût s'emballer]
Mettons quelques nombres côte à côte, et regardons la colonne de droite
grandir :

\begin{center}
\footnotesize
\begin{tabular}{rr}
\toprule
\textbf{Nombre d'éléments} & \textbf{Travail à faire} \\
\midrule
$10$    & $100$ \\
$100$   & $10\,000$ \\
$1\,000$ & $1\,000\,000$ \\
\bottomrule
\end{tabular}
\end{center}

À chaque fois que la colonne de gauche est multipliée par dix, celle de
droite est multipliée par cent. Voilà pourquoi on dit que le coût
\emph{s'emballe} : il grandit bien plus vite que ce qu'on lui donne.
\end{exemple}

La solution est celle que tout le monde emploie quand une réunion est trop
nombreuse : au lieu de faire parler chaque personne, on fait parler chaque
\emph{table}. On regroupe. C'est exactement ce que fait le transformeur de
vision, et c'est le sujet de la section suivante.

\section{Le patch : découper l'image en petits carrés}

\begin{enclair}
Plutôt que de laisser dialoguer les pixels un par un, on découpe l'image en
petits carrés de taille fixe. Chaque carré devient un porte-parole, et ce
sont ces porte-parole qui dialoguent entre eux. Il y en a beaucoup moins, et
tout devient calculable.
\end{enclair}

\begin{definition}[Le patch][patch]
Un \emph{patch} est un petit carré de pixels découpé dans l'image, toujours
de la même taille. Le mot vient de l'anglais et veut dire «~pièce~» ou
«~morceau~», comme une pièce de tissu cousue sur un pantalon. Dans
V-JEPA~2.1, un patch mesure $16 \times 16$ pixels. On dit aussi que la
\emph{taille de patch} vaut $16$.
\end{definition}

Le découpage se fait sans laisser de trou et sans que deux morceaux se
chevauchent : on pose une grille régulière sur l'image, exactement comme on
découpe une plaque de chocolat en carrés le long des rainures. Chaque carré
est ensuite mis à la queue leu leu, dans un ordre convenu une fois pour
toutes : la première ligne de gauche à droite, puis la deuxième ligne, et
ainsi de suite. On obtient une longue file de morceaux.

\begin{schema}{Le découpage en patchs, puis la mise en file. À gauche,
l'image est recouverte d'une grille régulière ; chaque case est un patch et
porte son numéro. À droite, les mêmes patchs sont rangés les uns derrière
les autres, dans l'ordre de lecture d'un texte : ligne par ligne, de gauche
à droite. Cette file est ce que la machine manipulera désormais.}[decoupage-patchs]
  \begin{scope}[x=1cm, y=1cm]
    \node[font=\sffamily\small\bfseries, text=mbSlate] at (2.0,4.5)
         {L'image et sa grille};
    \foreach \ri in {0,...,3}{
      \foreach \ci in {0,...,3}{
        \pgfmathtruncatemacro{\num}{\ri*4+\ci+1}
        \pgfmathtruncatemacro{\tone}{14+4*\num}
        \fill[mbAzure!\tone] (\ci*1.0,3.9-\ri*1.0) rectangle ++(1.0,-1.0);
        \draw[draw=mbSlate!55, line width=0.6pt]
              (\ci*1.0,3.9-\ri*1.0) rectangle ++(1.0,-1.0);
        % au-delà de la moitié de la grille, le fond devient trop foncé
        % pour du texte sombre : on bascule en blanc.
        \ifnum\num>10\relax
          \node[font=\sffamily\scriptsize\bfseries, text=white]
                at (\ci*1.0+0.5,3.9-\ri*1.0-0.5) {\num};
        \else
          \node[font=\sffamily\scriptsize\bfseries, text=mbInk]
                at (\ci*1.0+0.5,3.9-\ri*1.0-0.5) {\num};
        \fi
      }
    }
    \node[font=\sffamily\scriptsize, text=mbSlate] at (2.0,-0.45)
         {$4 \times 4 = 16$ patchs};

    \draw[fleche accent] (4.35,1.9) -- (5.55,1.9);
    \node[etiquette, above] at (4.95,1.98) {on met en file};

    \begin{scope}[xshift=5.9cm, yshift=1.55cm]
      \node[font=\sffamily\small\bfseries, text=mbSlate] at (2.56,2.95)
           {La file de patchs};
      \foreach \k in {0,...,4}{
        \pgfmathtruncatemacro{\num}{\k+1}
        \pgfmathtruncatemacro{\tone}{14+4*\num}
        \fill[mbAzure!\tone] (\k*0.72,0) rectangle ++(0.72,0.72);
        \draw[draw=mbSlate!55, line width=0.6pt]
              (\k*0.72,0) rectangle ++(0.72,0.72);
        \node[font=\sffamily\scriptsize\bfseries, text=mbInk]
              at (\k*0.72+0.36,0.36) {\num};
      }
      \node[font=\sffamily\small, text=mbSlate] at (4.05,0.36) {$\cdots$};
      \fill[mbAzure!78] (4.4,0) rectangle ++(0.72,0.72);
      \draw[draw=mbSlate!55, line width=0.6pt] (4.4,0) rectangle ++(0.72,0.72);
      \node[font=\sffamily\scriptsize\bfseries, text=white] at (4.76,0.36) {16};
      \draw[accolade bas] (0,-0.12) -- (5.12,-0.12);
      \node[font=\sffamily\scriptsize, text=mbSlate] at (2.56,-0.75)
           {une file de $16$ morceaux};
    \end{scope}
  \end{scope}
\end{schema}

Combien de morceaux obtient-on ? Le calcul est le plus simple du monde : on
divise la largeur par la taille du patch, on divise la hauteur par la taille
du patch, et on multiplie les deux résultats.

\begin{exemple}[Compter les patchs d'une image de V-JEPA 2.1]
Pendant la phase principale de l'entraînement, les images font
$256 \times 256$ pixels et les patchs mesurent $16 \times 16$. Sur la
largeur, on peut donc poser

\[
  \frac{256}{16} = 16 \text{ patchs,}
\]

et autant sur la hauteur. Le nombre total de patchs vaut

\[
  16 \times 16 = 256 .
\]

Retenez bien ce nombre, $256$ : nous allons le retrouver sans cesse.
Souvenez-vous aussi que cette image contenait $65\,536$ pixels. Nous venons
donc de ramener le nombre d'éléments à consulter de $65\,536$ à $256$,
c'est-à-dire de le diviser par $256$. Et comme le coût grimpe au carré, le
travail, lui, est divisé par $256 \times 256$, soit par plus de soixante-cinq
mille.
\end{exemple}

\begin{piege}
La division doit tomber juste. Si vous prenez une image de $250$ pixels de
côté avec des patchs de $16$, la division $250/16$ ne donne pas un nombre
entier, et il vous restera une bande de pixels sur le bord que personne ne
sait où ranger. C'est pour cela que les tailles d'images choisies dans
l'article sont toujours des multiples de $16$ : $256$, $384$, $512$. Quand
vous changerez de résolution dans vos propres essais, vérifiez toujours ce
point avant de lancer quoi que ce soit.
\end{piege}

\begin{enia}
Ce découpage est la toute première opération de V-JEPA~2.1, celle qui a lieu
avant tout le reste. L'article lui donne un nom que nous retrouverons tout
au long de la partie~IV : le \emph{tokeniseur}. C'est simplement la pièce
qui fabrique les morceaux.
\end{enia}

Nous avons donc une file de petits carrés. Mais un petit carré, c'est encore
une grille de nombres. Or un transformeur ne sait pas manipuler des grilles :
il manipule des listes de nombres. Il faut donc transformer chaque carré en
une liste. C'est ce que nous faisons tout de suite.

\section{Du carré au jeton : le plongement}

\begin{enclair}
Chaque petit carré est d'abord mis à plat, c'est-à-dire écrit comme une
simple liste de nombres. Puis on lui applique un calcul qui le remplace par
une autre liste, plus courte et mieux rangée. Cette nouvelle liste
s'appelle un jeton, et c'est elle que le modèle manipulera du début à la
fin.
\end{enclair}

Reprenons un patch de $16 \times 16$ pixels en couleur. Combien de nombres
contient-il ? Nous savons le calculer depuis le chapitre~\ref{chap:image-video} :
la hauteur, fois la largeur, fois les trois canaux de couleur.

\[
  16 \times 16 \times 3 = 768 \text{ nombres.}
\]

On écrit ces $768$ nombres les uns derrière les autres, en une seule ligne.
Cette opération s'appelle \emph{mettre à plat}, ou \emph{aplatir}. Rien n'est
perdu, rien n'est ajouté : on change seulement la façon de ranger.

Ensuite vient le vrai travail. On ne garde pas ces $768$ nombres bruts : on
les remplace par une autre liste, obtenue grâce à des \emph{coefficients}.
Un coefficient, c'est tout simplement un nombre par lequel on en multiplie
un autre. Ceux dont nous parlons ici sont \emph{appris} : ce sont exactement
les boutons du chapitre~\ref{chap:apprendre}. Au début ils valent n'importe
quoi, et l'entraînement les règle petit à petit.

\begin{definition}[Le plongement, ou embedding][embedding]
Un \emph{plongement} est une liste de nombres qui représente un morceau de
donnée. On dit aussi \emph{embedding}, qui est le mot anglais employé
partout. On l'obtient en multipliant les nombres du morceau par des
coefficients appris, puis en additionnant les résultats. La longueur de cette
liste s'appelle la \emph{dimension} du plongement : c'est un choix que fait
celui qui construit le modèle, et elle est la même pour tous les morceaux.
\end{definition}

\begin{schema}{Du carré au jeton, en deux gestes. On part d'un patch de
$16 \times 16$ pixels en couleur. On le met d'abord à plat, c'est-à-dire
qu'on écrit ses nombres les uns derrière les autres : cela en fait $768$, et
aucun n'a été touché. On leur applique ensuite un calcul dont les
coefficients sont appris, et qui rend une liste bien plus courte. Cette
liste est le jeton.}[carre-vers-jeton]
  \begin{scope}[x=1cm, y=1cm]
    %---- 1. le patch, centré sur la ligne y = 0 ---------------------
    \foreach \ri in {0,...,3}{
      \foreach \ci in {0,...,3}{
        \pgfmathtruncatemacro{\tone}{18+8*(\ri+\ci)}
        \fill[mbAzure!\tone] (\ci*0.45,0.9-\ri*0.45) rectangle ++(0.45,-0.45);
        \draw[draw=mbSlate!45, line width=0.4pt]
              (\ci*0.45,0.9-\ri*0.45) rectangle ++(0.45,-0.45);
      }
    }
    \node[font=\sffamily\scriptsize, text=mbSlate, align=center] at (0.9,-1.3)
         {un patch\\ $16 \times 16 \times 3$};

    %---- 2. on met à plat -------------------------------------------
    \draw[fleche accent] (2.1,0) -- (3.3,0);
    \node[etiquette, above] at (2.7,0.08) {on met à plat};

    \foreach \k in {0,...,8}{
      \pgfmathtruncatemacro{\tone}{18+9*\k}
      \fill[mbAzure!\tone] (3.6+\k*0.32,-0.22) rectangle ++(0.32,0.44);
      \draw[draw=mbSlate!45, line width=0.4pt]
            (3.6+\k*0.32,-0.22) rectangle ++(0.32,0.44);
    }
    \node[font=\sffamily\small, text=mbSlate] at (6.72,0) {$\cdots$};
    \node[font=\sffamily\scriptsize, text=mbSlate, align=center] at (5.4,-1.3)
         {$768$ nombres,\\ aucun n'a changé};

    %---- 3. le calcul appris ----------------------------------------
    \draw[fleche accent] (7.1,0) -- (8.3,0);
    \node[etiquette, above] at (7.7,0.08) {calcul appris};

    \foreach \k in {0,...,3}{
      \pgfmathtruncatemacro{\tone}{30+16*\k}
      \fill[mbTeal!\tone] (8.6+\k*0.32,-0.22) rectangle ++(0.32,0.44);
      \draw[draw=mbTeal!80!black, line width=0.5pt]
            (8.6+\k*0.32,-0.22) rectangle ++(0.32,0.44);
    }
    \node[font=\sffamily\scriptsize, text=mbTeal!55!black, align=center]
         at (9.24,-1.3) {le jeton :\\ des mesures,\\ pas des pixels};

    %---- l'accolade qui montre que la liste a raccourci --------------
    \draw[accolade] (3.6,0.32) -- (7.04,0.32);
    \draw[accolade] (8.6,0.32) -- (9.88,0.32);
    \node[annotation] at (9.24,0.95) {plus court};
  \end{scope}
\end{schema}

\begin{definition}[Le jeton, ou token][jeton]
Un \emph{jeton} est le plongement d'un patch, c'est-à-dire la liste de
nombres qui remplace ce patch dans toute la suite des calculs. On dit aussi
\emph{token}, en anglais. Le mot est bien choisi : comme un jeton de casino
remplace une somme d'argent, un jeton remplace ici un morceau d'image. Le
découpage suivi du plongement s'appelle la \emph{tokenisation}.
\end{definition}

\begin{exemple}[Un plongement en tout petit]
Pour voir le calcul en entier, prenons un patch minuscule de $2 \times 2$
pixels en niveaux de gris, donc quatre nombres, et une dimension de
plongement de $2$ seulement. Le patch vaut

\[
  \begin{pmatrix} 10 & 20 \\ 30 & 40 \end{pmatrix}
  \quad \longrightarrow \quad
  \text{mis à plat : } (10,\ 20,\ 30,\ 40).
\]

Le modèle possède, pour cette opération, huit coefficients appris : quatre
pour fabriquer le premier nombre du jeton, quatre pour le second. Disons
qu'ils valent

\[
  (0{,}1 \; ; \; 0{,}0 \; ; \; -0{,}1 \; ; \; 0{,}2)
  \quad \text{et} \quad
  (0{,}0 \; ; \; 0{,}5 \; ; \; 0{,}0 \; ; \; -0{,}5).
\]

Le premier nombre du jeton s'obtient en multipliant terme à terme puis en
additionnant :

\[
  10 \times 0{,}1 + 20 \times 0{,}0 + 30 \times (-0{,}1) + 40 \times 0{,}2
  = 1 + 0 - 3 + 8 = 6 .
\]

Le second nombre, de la même façon :

\[
  10 \times 0{,}0 + 20 \times 0{,}5 + 30 \times 0{,}0 + 40 \times (-0{,}5)
  = 0 + 10 + 0 - 20 = -10 .
\]

Le jeton vaut donc $(6\ ;\ -10)$. Quatre nombres sont devenus deux nombres,
et ces deux nombres ne sont plus des intensités de gris : ce sont des
mesures fabriquées par le modèle. Regardez la seconde liste de
coefficients, $(0 \; ; \; 0{,}5 \; ; \; 0 \; ; \; -0{,}5)$ : elle ajoute
le pixel en haut à droite et retranche celui d'en bas à droite. Elle mesure
donc une différence entre le haut et le bas. Personne ne le lui a demandé :
c'est l'entraînement qui a réglé ces quatre nombres ainsi.
\end{exemple}

\begin{intuition}
Le plongement est un résumé. Là où les $768$ nombres bruts disaient «~voici
l'éclat de chacun de mes pixels~», le jeton dit plutôt «~je contiens un bord
oblique, plutôt sombre, plutôt vert~». Personne n'a écrit ces mots nulle
part : ce sont les coefficients appris qui décident, tout seuls, de ce
qu'il vaut la peine de mesurer.
\end{intuition}

\begin{surlesdonnees}
Dans V-JEPA~2.1, ce calcul n'est pas fait patch par patch dans une boucle,
mais d'un seul coup par une opération appelée \emph{convolution}. Une
convolution, c'est une petite fenêtre que l'on fait glisser sur l'image :
à chaque arrêt, elle regarde ce qui se trouve dessous et en tire quelques
nombres. Ici, la fenêtre mesure $16 \times 16$ pixels et elle avance de $16$
pixels à chaque pas. Avancer d'exactement la largeur de la fenêtre, c'est
précisément découper sans chevauchement. Le résultat est le même que ce que
nous venons de calculer à la main, mais il est obtenu bien plus vite. Nous
reverrons la convolution de près au chapitre~13 ; pour l'instant, retenez
seulement l'image de la fenêtre qui glisse.
\end{surlesdonnees}

Nous savons transformer une image en une file de jetons. Reste une question
que nous avons soigneusement mise de côté : et le temps ?

\section{Le tubelet : découper aussi dans le temps}

\begin{enclair}
Une vidéo, c'est une pile d'images. Si l'on découpait chaque image
séparément, on obtiendrait beaucoup trop de jetons. On découpe donc aussi
dans l'épaisseur de la pile : chaque morceau prend deux images d'un coup.
Ce petit pavé s'appelle un tubelet.
\end{enclair}

\begin{definition}[Le tubelet][tubelet]
Un \emph{tubelet} est un patch qui possède en plus une épaisseur dans le
temps. Au lieu d'être un carré plat découpé dans une seule image, c'est un
petit pavé qui traverse plusieurs images de suite. Dans V-JEPA~2.1, ce pavé
mesure $16 \times 16$ pixels et $2$ images d'épaisseur : on dit que la
\emph{taille de tubelet} vaut $2$.
\end{definition}

\begin{schema}{Le patch et le tubelet. À gauche, un patch : un carré plat,
découpé dans une seule image. À droite, un tubelet : le même carré, pris sur
deux images qui se suivent. Les traits pointillés montrent qu'il s'agit bien
du même endroit. Le tubelet contient donc deux fois plus de pixels, et il
porte en lui un petit bout de mouvement, puisqu'il voit cet endroit à deux
instants.}[patch-tubelet]
  \begin{scope}[x=1cm, y=1cm]
    %================ à gauche : le patch, plat =======================
    \node[font=\sffamily\small\bfseries, text=mbSlate] at (1.2,3.65)
         {Un patch : plat};
    \foreach \ri in {0,...,3}{
      \foreach \ci in {0,...,3}{
        \draw[draw=mbSlate!40, fill=white, line width=0.4pt]
              (\ci*0.6,2.4-\ri*0.6) rectangle ++(0.6,-0.6);
      }
    }
    \fill[mbBlue!45] (0.6,1.8) rectangle ++(0.6,-0.6);
    \draw[draw=mbBlue, line width=1pt] (0.6,1.8) rectangle ++(0.6,-0.6);
    \node[font=\sffamily\scriptsize, text=mbSlate] at (1.2,-0.45)
         {une seule image};

    % le trait qui sépare les deux moitiés du dessin
    \draw[draw=mbSlate!25, line width=0.5pt] (3.6,-0.9) -- (3.6,3.35);

    %================ à droite : le tubelet ===========================
    % Les deux images sont posées côte à côte, sans se recouvrir : ainsi
    % la case retenue reste visible sur l'une comme sur l'autre.
    \node[font=\sffamily\small\bfseries, text=mbSlate] at (7.5,3.65)
         {Un tubelet : le même carré, sur deux images};

    \foreach \img/\dx/\teinte in {1/4.6/55, 2/8.0/35}{
      \begin{scope}[shift={(\dx,0)}]
        \foreach \ri in {0,...,3}{
          \foreach \ci in {0,...,3}{
            \draw[draw=mbSlate!40, fill=white, line width=0.4pt]
                  (\ci*0.6,2.4-\ri*0.6) rectangle ++(0.6,-0.6);
          }
        }
        \fill[mbTeal!\teinte] (0.6,1.8) rectangle ++(0.6,-0.6);
        \draw[draw=mbTeal!85!black, line width=1pt]
              (0.6,1.8) rectangle ++(0.6,-0.6);
        \node[font=\sffamily\scriptsize, text=mbSlate] at (1.2,-0.45)
             {image \img};
      \end{scope}
    }

    % les pointillés : c'est bien le même endroit des deux côtés
    \draw[draw=mbTeal!85!black, line width=0.7pt, dashed]
          (5.8,1.8) -- (8.6,1.8);
    \draw[draw=mbTeal!85!black, line width=0.7pt, dashed]
          (5.8,1.2) -- (8.6,1.2);

    \draw[accolade bas] (4.6,-0.75) -- (10.4,-0.75);
    \node[font=\sffamily\scriptsize, text=mbTeal!55!black] at (7.5,-1.4)
         {le tubelet, c'est les deux carrés ensemble};

    \draw[-{Stealth[round,length=2.2mm,width=1.7mm]},
          draw=mbAmber!85!black, line width=1pt] (4.6,2.72) -- (10.4,2.72);
    \node[font=\sffamily\scriptsize, text=mbAmber!80!black, above]
         at (7.5,2.75) {le temps passe};
  \end{scope}
\end{schema}

Pourquoi coller deux images ensemble plutôt que de les traiter séparément ?
Pour deux raisons, et les deux sont bonnes.

La première est une simple affaire de nombre. Un clip de seize images
découpé image par image donnerait $16 \times 256 = 4\,096$ jetons. Avec des
tubelets de deux images, on en obtient deux fois moins. Comme le coût grimpe
au carré, le travail est divisé par quatre.

La seconde raison est plus intéressante. Un tubelet voit le même endroit à
deux instants différents. Il contient donc, dès sa fabrication, une petite
trace de ce qui a bougé à cet endroit. Le mouvement n'est pas quelque chose
que le modèle devra reconstituer plus tard : il est déjà là, dans le
morceau.

\begin{exemple}[Compter les jetons d'un clip de la phase principale]
Pendant la phase principale, un clip contient $16$ images de
$256 \times 256$ pixels. Comptons dans les trois directions.

Dans le temps, avec des tubelets de $2$ images :
\[ \frac{16}{2} = 8 \text{ tranches.} \]

En hauteur, avec des patchs de $16$ pixels :
\[ \frac{256}{16} = 16 \text{ rangées.} \]

En largeur, de même : $16$ colonnes. Le nombre total de jetons vaut donc le
produit des trois :

\[
  8 \times 16 \times 16 = 2\,048 \text{ jetons.}
\]

Voilà le deuxième nombre à retenir par cœur : \textbf{$2\,048$ jetons pour
un clip de la phase principale}.
\end{exemple}

\begin{exemple}[Compter les jetons d'un clip de refroidissement]
Vous connaissez déjà cette phase. Le chapitre~\ref{chap:image-video} vous a
donné la taille de ses clips, $64$ images de $384 \times 384$ pixels, et le
chapitre~\ref{chap:apprendre} vous a dit d'où lui vient son nom : sur les
douze mille dernières étapes, le taux d'apprentissage redescend jusqu'à
presque zéro, et le modèle se pose en douceur. Refaisons pour elle le calcul
des jetons :

\[
  \frac{64}{2} = 32, \qquad \frac{384}{16} = 24, \qquad \frac{384}{16} = 24,
\]
\[
  32 \times 24 \times 24 = 18\,432 \text{ jetons.}
\]

C'est le troisième nombre à retenir. Comparons-le au précédent :

\[
  \frac{18\,432}{2\,048} = 9 .
\]

Neuf fois plus de jetons : c'est exactement le rapport que nous avions
trouvé au chapitre~\ref{chap:image-video} en comptant les pixels bruts, et
ce n'est pas un hasard, puisque le découpage garde les proportions.

Mais voici ce que le chapitre~\ref{chap:image-video} ne pouvait pas encore
nous dire. Là-bas, neuf fois plus de nombres voulait dire neuf fois plus de
mémoire, et nous étions passés de $1{,}6$ à $14{,}5$ gigaoctets par lot. Ici,
c'est le \emph{calcul} qui est en jeu, et le calcul grimpe au carré :
\[ 9 \times 9 = 81 . \]
Un clip de refroidissement coûte donc \emph{quatre-vingt-une fois} plus de
travail qu'un clip de la phase principale. La mémoire était multipliée par
neuf, le calcul l'est par quatre-vingt-un. Voilà pourquoi l'article ne fait
durer cette phase que $12\,000$ pas, alors que la phase principale en compte
$135\,000$.
\end{exemple}

\begin{aretenir}
Quatre lignes à garder en tête pour tout le livre :

\begin{center}
\footnotesize
\begin{tabular}{llr}
\toprule
\textbf{Entrée} & \textbf{Calcul} & \textbf{Jetons} \\
\midrule
Image $256 \times 256$        & $16 \times 16$               & $256$ \\
Image $512 \times 512$        & $32 \times 32$               & $1\,024$ \\
Clip $16 \times 256 \times 256$ & $8 \times 16 \times 16$    & $2\,048$ \\
Clip $64 \times 384 \times 384$ & $32 \times 24 \times 24$   & $18\,432$ \\
\bottomrule
\end{tabular}
\end{center}

La règle est toujours la même : on divise chaque côté par la taille du
découpage, puis on multiplie les résultats entre eux.
\end{aretenir}

Nos morceaux sont prêts et rangés en file. Mais en les mettant en file, nous
avons perdu quelque chose d'important, et il faut le récupérer sans tarder.

\section{Dire où l'on se trouve : l'encodage de position}

\begin{enclair}
Une fois les jetons mis en file, plus rien ne dit lequel venait du coin en
haut à gauche et lequel venait du milieu. On doit donc écrire l'adresse de
chaque morceau à l'intérieur même de son jeton.
\end{enclair}

Voici le problème, et il est sérieux. Le calcul de la section suivante,
l'attention, traite la file de jetons comme un \emph{sac} : il ne regarde
jamais l'ordre. Mélangez la file, il rendra exactement le même résultat. Or
une image dont on mélange les morceaux n'est plus la même image du tout.

\begin{intuition}
Prenez un puzzle terminé, photographiez-le, puis jetez toutes les pièces
dans une boîte. Les pièces sont toujours là, aucune n'a été perdue, et
pourtant le dessin a disparu. Ce qui manque, ce n'est pas la matière, c'est
la \emph{place} de chaque pièce. Voilà exactement ce que nous devons rendre
au modèle.
\end{intuition}

\begin{definition}[L'encodage de position][encodage-position]
L'\emph{encodage de position} est la façon dont on inscrit, dans le jeton
lui-même, l'endroit d'où venait son morceau. Sans lui, le modèle ne saurait
pas distinguer le haut du bas, la gauche de la droite, ni le début de la fin
d'une vidéo.
\end{definition}

V-JEPA~2.1 emploie une méthode particulière, appelée \emph{RoPE}. Ce nom est
une abréviation anglaise de \emph{Rotary Position Embedding}, ce qui veut
dire «~encodage de position par rotation~». Le mot rotation dit tout : au
lieu d'ajouter un numéro au jeton, on \emph{tourne} le jeton d'un angle qui
dépend de sa position.

\begin{schema}{L'idée de RoPE, expliquée avec une aiguille d'horloge. On
prend les nombres du jeton deux par deux et on les voit comme une flèche.
Cette flèche est ensuite tournée d'un angle proportionnel à la position du
morceau : rien à la position $0$, un cran à la position $1$, deux crans à la
position $2$, et ainsi de suite. La longueur de la flèche ne change jamais :
seule sa direction porte l'information de place.}[rope-horloge]
  \begin{scope}[x=1cm, y=1cm]
    \foreach \k/\ang in {0/0, 1/30, 2/60, 3/90}{
      \begin{scope}[shift={(\k*3.4,0)}]
        \draw[draw=mbSlate!30, line width=0.5pt] (0,0) circle (1.05);
        \draw[draw=mbSlate!25, line width=0.4pt] (-1.25,0) -- (1.25,0);
        \draw[draw=mbSlate!25, line width=0.4pt] (0,-1.25) -- (0,1.25);
        % la position de départ, rappelée en pointillé sur chaque cadran
        \draw[draw=mbSlate!45, line width=0.6pt, dashed] (0,0) -- (1.05,0);
        % l'angle parcouru depuis le départ
        \ifnum\ang>0\relax
          \draw[draw=mbAmber!70, line width=0.7pt] (0.45,0) arc (0:\ang:0.45);
        \fi
        \draw[vecteur accent] (0,0) -- (\ang:1.05);
        \fill[mbInk] (0,0) circle (1.4pt);
        \node[font=\sffamily\scriptsize\bfseries, text=mbSlate] at (0,1.55)
             {position \k};
        \node[font=\sffamily\scriptsize, text=mbAmber!80!black] at (0,-1.62)
             {tourné de $\ang^\circ$};
      \end{scope}
    }
  \end{scope}
\end{schema}

\begin{exemple}[Tourner une flèche, avec les nombres]
Prenons deux nombres d'un jeton, disons $(1\ ;\ 0)$, et voyons-les comme une
flèche qui pointe vers la droite. Convenons que chaque position fait tourner
la flèche de trente degrés. Mesurons à la règle où arrive la pointe de la
flèche après chaque rotation. Voici ce que l'on relève :

\begin{center}
\footnotesize
\begin{tabular}{crr}
\toprule
\textbf{Position} & \textbf{Angle} & \textbf{Les deux nombres} \\
\midrule
$0$ & $0^\circ$  & $(1{,}000 \ ; \ 0{,}000)$ \\
$1$ & $30^\circ$ & $(0{,}866 \ ; \ 0{,}500)$ \\
$2$ & $60^\circ$ & $(0{,}500 \ ; \ 0{,}866)$ \\
$3$ & $90^\circ$ & $(0{,}000 \ ; \ 1{,}000)$ \\
\bottomrule
\end{tabular}
\end{center}

Les deux nombres changent, alors que le contenu du morceau, lui, n'a pas
bougé. C'est la position, et elle seule, qui a fait tourner la flèche.
\end{exemple}

Pourquoi tourner plutôt qu'ajouter un numéro ? Parce que la rotation a un
avantage qui tombe à pic. Quand deux jetons se comparent, ce qui compte
n'est pas leur angle à chacun, mais l'\emph{écart} entre leurs deux angles.
Or cet écart ne dépend que de la \emph{distance} entre leurs positions.

\begin{exemple}[Ce qui compte, c'est l'écart]
Reprenons nos trente degrés par position. Deux jetons placés aux positions
$1$ et $3$ sont séparés de deux crans, donc de $60$ degrés. Deux autres
jetons placés aux positions $4$ et $6$ sont eux aussi séparés de deux crans,
donc des mêmes $60$ degrés. Pour le modèle, les deux paires sont dans la
même situation.

C'est très exactement ce que l'on veut. Un bord vertical reste un bord
vertical, qu'il se trouve à gauche ou à droite de la photo. Ce qui importe,
c'est la place des choses \emph{les unes par rapport aux autres}, et non
l'endroit exact où elles se trouvent dans l'image.
\end{exemple}

Dans une vidéo, il y a trois directions à repérer : la ligne, la colonne et
l'instant. On fait donc tourner une partie des nombres du jeton selon la
hauteur, une autre selon la largeur, et une troisième selon le temps. Le
principe ne change pas d'un pouce : il est simplement appliqué trois fois,
une fois par direction.

Les auteurs appellent cela un \emph{RoPE en trois dimensions}. Attention à
ne pas confondre ce mot avec la «~dimension du plongement~» vue plus haut :
ici, «~dimension~» veut dire «~direction dans laquelle on se repère~», et
rien d'autre.

\begin{enia}
Voici maintenant le point qui vous servira le jour où vous mettrez ce modèle
en service, et l'article y insiste. Comme la position se traduit par un
angle, on peut décider de \emph{resserrer} ou d'\emph{écarter} les crans.

Supposons que le modèle ait été entraîné sur des images de $16$ patchs de
côté, avec trente degrés par cran. On veut maintenant lui donner des images
de $32$ patchs de côté. Si l'on garde trente degrés par cran, la flèche fera
plusieurs tours complets et le modèle sera perdu. Mais si l'on passe à
quinze degrés par cran, la traversée complète de l'image correspondra
toujours au même angle total : le modèle retrouve ses repères.

Cette opération porte un nom, l'\emph{interpolation des fréquences}, et
l'article s'en sert vraiment, pour deux travaux que nous découvrirons au
chapitre~7 : la \emph{segmentation}, qui consiste à dire à quel objet
appartient chaque pixel, et l'\emph{estimation de profondeur}, qui consiste
à dire à quelle distance se trouve chaque pixel.

Ce que cela permet est énorme : \textbf{on peut donner au modèle des images
d'une autre taille que celles de son entraînement, sans le réentraîner}.
Nous y reviendrons au chapitre~26, quand nous choisirons nos réglages de
production.
\end{enia}

\begin{piege}
«~Sans le réentraîner~» ne veut pas dire «~sans rien vérifier~». Changer de
résolution modifie toujours un peu la qualité des résultats. La règle de
métier est simple : vous avez le droit de changer la résolution, mais vous
avez le devoir de mesurer ce que cela change sur vos propres données. Nous
verrons dans la partie~VII que l'article lui-même a mesuré cet effet, et
qu'il n'est pas nul.
\end{piege}

Nos jetons contiennent maintenant ce qu'ils voient et l'endroit d'où ils
viennent. Il est temps de les faire dialoguer.

\section{L'attention : chaque jeton écoute les autres}

\begin{enclair}
L'attention est le cœur du transformeur. Chaque jeton pose une question, tous
les autres présentent une étiquette, et celui qui pose la question écoute
surtout ceux dont l'étiquette répond le mieux. Il repart alors avec un
contenu enrichi de tout ce qu'il vient d'entendre.
\end{enclair}

Nous arrivons à la pièce maîtresse. Prenons le temps qu'il faut, car tout le
reste du livre s'appuie dessus.

Reprenons notre photo de chat. Le jeton qui contient les moustaches aimerait
bien savoir s'il y a des oreilles pointues quelque part dans l'image. Il ne
peut pas aller les chercher lui-même : il ne voit que son propre morceau. Il
va donc lancer un appel à toute la file, et voir qui répond.

Pour organiser cet appel, chaque jeton fabrique trois listes de nombres à
partir de son contenu, exactement comme il fabriquait son plongement : en
multipliant ses nombres par des coefficients appris. Ces trois listes ont
chacune un rôle bien précis.

\begin{definition}[La question, la clé et le contenu][qkv]
À partir de son propre jeton, chacun fabrique trois listes de nombres :

\begin{itemize}
  \item la \emph{question}, que l'on appelle \emph{query} en anglais : elle
        dit «~voici ce que je cherche~» ;
  \item la \emph{clé}, ou \emph{key} : elle dit «~voici de quoi je
        parle~» ;
  \item le \emph{contenu}, ou \emph{value} : c'est ce que le jeton donnera
        réellement à ceux qui l'auront écouté.
\end{itemize}

Les coefficients qui fabriquent ces trois listes sont appris pendant
l'entraînement. Personne ne les écrit à la main.
\end{definition}

\begin{intuition}
Imaginez une salle de classe. Un élève lève la main et demande : «~qui sait
quelque chose sur les oreilles pointues ?~» — c'est sa question. Chaque
autre élève porte une pancarte autour du cou où il a écrit son sujet : «~je
parle de fourrure~», «~je parle d'oreilles~», «~je parle de flaques d'eau~» —
ce sont les clés. L'élève qui a posé la question regarde toutes les
pancartes, repère celles qui collent le mieux à sa demande, et va écouter
ces camarades-là en priorité. Ce qu'ils lui racontent alors, c'est leur
contenu.
\end{intuition}

\begin{schema}{L'attention vue comme une salle de classe. Le jeton du bas
pose sa question. Chaque autre jeton présente sa clé. Plus la clé répond à
la question, plus la flèche est épaisse et plus le jeton sera écouté. Le
résultat est un mélange des contenus, dosé par l'épaisseur des flèches. Ici,
le jeton «~oreille~» est écouté fort, tandis que le jeton «~flaque~» est
presque ignoré.}[attention-classe]
  \begin{scope}[x=1cm, y=1cm]
    \node[bloc calcul, minimum width=20mm] (j1) at (0,3.2)   {fourrure};
    \node[bloc calcul, minimum width=20mm] (j2) at (3.2,3.2) {oreille};
    \node[bloc calcul, minimum width=20mm] (j3) at (6.4,3.2) {flaque};
    \node[font=\sffamily\scriptsize, text=mbSlate, above] at (0,3.72)
         {clé + contenu};
    \node[font=\sffamily\scriptsize, text=mbSlate, above] at (3.2,3.72)
         {clé + contenu};
    \node[font=\sffamily\scriptsize, text=mbSlate, above] at (6.4,3.72)
         {clé + contenu};

    \node[bloc accent, minimum width=30mm] (q) at (3.2,0) {moustaches};
    \node[font=\sffamily\scriptsize, text=mbAmber!80!black, below] at (3.2,-0.55)
         {pose la question : «~qui a des oreilles ?~»};

    \draw[draw=mbBlue!55, line width=1.1pt, -{Stealth[round,length=2.2mm]}]
          (j1.south) -- ($(q.north)+(-0.9,0)$);
    \draw[draw=mbBlue!85, line width=3.2pt, -{Stealth[round,length=3.2mm]}]
          (j2.south) -- (q.north);
    \draw[draw=mbBlue!30, line width=0.5pt, -{Stealth[round,length=1.8mm]}]
          (j3.south) -- ($(q.north)+(0.9,0)$);

    \node[etiquette] at (1.05,1.75) {$0{,}117$};
    \node[etiquette] at (3.62,1.75) {$0{,}867$};
    \node[etiquette] at (5.6,1.75) {$0{,}016$};
    \node[annotation] at (8.6,1.75) {les trois poids\\ font $1$ en tout};
  \end{scope}
\end{schema}

Comment mesure-t-on qu'une clé «~répond bien~» à une question ? On multiplie
les deux listes terme à terme et on additionne le tout. Cette opération porte
un nom en mathématiques : le \emph{produit scalaire}. Elle donne un grand
nombre quand les deux listes se ressemblent, et un petit nombre quand elles
n'ont rien à voir.

Le nombre obtenu s'appelle un \emph{score}. On en obtient un par jeton
écouté. Il faut ensuite transformer ces scores en \emph{poids},
c'est-à-dire en parts qui s'additionnent pour faire un, comme des parts d'un
gâteau. L'opération qui fait cela s'appelle la \emph{fonction softmax}.

\begin{definition}[La fonction softmax][softmax]
La \emph{softmax} transforme une liste de nombres quelconques en une liste de
parts, toutes positives, dont la somme fait exactement $1$. Elle procède en
deux temps : on calcule d'abord l'\emph{exponentielle} de chaque nombre, puis
on divise chacun par la somme de tous.

L'exponentielle est une machine à agrandir. On l'écrit $e^{x}$, ce qui se lit
«~e puissance $x$~». Elle prend un nombre et en rend un autre, toujours
positif, et qui grossit très vite : $e^{0} = 1$, $e^{1} \approx 2{,}7$,
$e^{2} \approx 7{,}4$, $e^{4} \approx 54{,}6$. Vous n'avez pas besoin de
savoir la calculer — toute calculatrice le fait. Retenez ses deux effets, qui
sont exactement ceux dont la softmax a besoin : elle rend tous les nombres
positifs, et elle écarte fortement les grands des petits.

Le nom «~softmax~» veut dire «~maximum en douceur~» : elle donne la plus
grosse part au plus grand score, sans pour autant écraser complètement les
autres.
\end{definition}

\begin{exemple}[Une attention calculée du début à la fin]
Faisons le calcul en entier, avec de tout petits nombres. Reprenons les
trois jetons de la figure~\ref{fig:attention-classe} : «~fourrure~»,
«~oreille~» et «~flaque~». Toutes les listes auront quatre nombres, et nous
conviendrons que le premier de ces quatre nombres veut dire «~je parle
d'oreilles~».

\textbf{Premier temps : la question.} Le jeton «~moustaches~» cherche des
oreilles. Il met donc du poids sur le premier nombre :
\[ q = (2 \ ; \ 0 \ ; \ 0 \ ; \ 0). \]

\textbf{Deuxième temps : les clés.} Les trois autres présentent les leurs.
«~Oreille~» parle beaucoup d'oreilles, «~fourrure~» un peu, «~flaque~» pas
du tout :
\[
  \begin{aligned}
    k_{\text{fourrure}} &= (2 \ ; \ 0 \ ; \ 0 \ ; \ 0), \\
    k_{\text{oreille}}  &= (4 \ ; \ 0 \ ; \ 0 \ ; \ 0), \\
    k_{\text{flaque}}   &= (0 \ ; \ 3 \ ; \ 0 \ ; \ 0).
  \end{aligned}
\]

\textbf{Troisième temps : les scores.} On fait le produit scalaire de la
question avec chaque clé :
\[
  2 \times 2 = 4, \qquad
  2 \times 4 = 8, \qquad
  0 .
\]

\textbf{Quatrième temps : on remet à l'échelle.} On divise chaque score par
la racine carrée de la longueur des listes, ici $\sqrt{4} = 2$. La racine
carrée
d'un nombre, c'est celui qui, multiplié par lui-même, redonne ce nombre :
$\sqrt{4} = 2$ parce que $2 \times 2 = 4$, et $\sqrt{9} = 3$ parce que
$3 \times 3 = 9$. Cette division empêche les scores de devenir énormes quand
les listes sont longues :
\[
  \frac{4}{2} = 2, \qquad \frac{8}{2} = 4, \qquad \frac{0}{2} = 0 .
\]

\textbf{Cinquième temps : la softmax.} On prend l'exponentielle de chacun :
\[
  e^{2} \approx 7{,}389, \qquad e^{4} \approx 54{,}598, \qquad e^{0} = 1 .
\]
On les additionne : $7{,}389 + 54{,}598 + 1 = 62{,}987$. Puis on divise
chacun par cette somme :
\[
  \frac{7{,}389}{62{,}987} \approx 0{,}117, \quad
  \frac{54{,}598}{62{,}987} \approx 0{,}867, \quad
  \frac{1}{62{,}987} \approx 0{,}016 .
\]
Vérifions : $0{,}117 + 0{,}867 + 0{,}016 = 1{,}000$. Ce sont bien des parts
de gâteau, et ce sont exactement les trois nombres écrits sur les flèches de
la figure~\ref{fig:attention-classe}.

\textbf{Sixième temps : le mélange des contenus.} Disons que les trois
contenus valent
\[
  \begin{aligned}
    v_{\text{fourrure}} &= (10 \ ; \ 0 \ ; \ 0 \ ; \ 0), \\
    v_{\text{oreille}}  &= (0 \ ; \ 10 \ ; \ 0 \ ; \ 0), \\
    v_{\text{flaque}}   &= (0 \ ; \ 0 \ ; \ 10 \ ; \ 0).
  \end{aligned}
\]
Le résultat est la somme des contenus, chacun multiplié par sa part :
\[
  \begin{aligned}
    &0{,}117 \times v_{\text{fourrure}}
     + 0{,}867 \times v_{\text{oreille}}
     + 0{,}016 \times v_{\text{flaque}} \\
    &\qquad = (1{,}17 \ ; \ 8{,}67 \ ; \ 0{,}16 \ ; \ 0).
  \end{aligned}
\]

Voilà le nouveau jeton «~moustaches~». Il ne contient plus seulement ce
qu'il voyait tout seul : il contient surtout ce que lui a dit «~oreille~»,
un peu de ce que lui a dit «~fourrure~», et presque rien de «~flaque~». Il a
posé une question sur les oreilles, et c'est bien le jeton «~oreille~» qu'il
a le plus écouté.
\end{exemple}

\begin{piege}
N'oubliez jamais la division par la racine carrée. Sans elle, avec des
listes de mille nombres, les scores atteignent des valeurs si grandes que la
softmax donne une part de $1$ à un seul jeton et $0$ à tous les autres. Le
modèle n'écoute plus qu'un seul camarade, l'entraînement se bloque, et rien
dans les messages d'erreur ne vous dira pourquoi.
\end{piege}

\begin{aretenir}
L'attention se résume à six gestes, toujours les mêmes : je fabrique ma
question, chacun fabrique sa clé et son contenu, je compare ma question à
chaque clé, je remets à l'échelle, je transforme les scores en parts, je
mélange les contenus selon ces parts. Tout ce qui suit dans ce chapitre
n'est que de l'organisation autour de ces six gestes.
\end{aretenir}

Cette façon de faire est puissante, mais il lui manque quelque chose. Voyons
quoi, et comment on y remédie.

\section{Plusieurs têtes valent mieux qu'une}

\begin{enclair}
Avec une seule attention, chaque jeton ne peut poser qu'une seule question à
la fois. On en fait donc tourner plusieurs en même temps, chacune avec
ses propres réglages, et on met leurs réponses bout à bout. Chacune de ces
attentions qui travaillent côte à côte s'appelle une tête.
\end{enclair}

Notre jeton «~moustaches~» cherchait des oreilles. Mais il aurait aussi aimé
savoir quelle est la couleur dominante autour de lui, et si quelque chose
bouge à sa gauche. Avec une seule question, il doit choisir. C'est dommage.

\begin{definition}[La tête d'attention][tete-attention]
Une \emph{tête d'attention} est un exemplaire complet de tout ce que nous
venons de décrire, avec ses propres coefficients pour fabriquer les
questions, les clés et les contenus. Plusieurs têtes travaillent en même
temps sur les mêmes jetons, mais chacune pose sa propre question. On parle
alors d'\emph{attention multi-têtes}.
\end{definition}

Le partage se fait ainsi : on ne donne pas la totalité des nombres du jeton
à chaque tête, on les répartit entre elles. Chaque tête travaille sur sa
part, produit sa réponse, et l'on recolle toutes les réponses bout à bout à
la fin. La liste finale a donc la même longueur qu'au départ.

\begin{schema}{Le partage entre les têtes. Les nombres du jeton sont répartis
en parts égales, une par tête. Chaque tête pose sa propre question et ne
travaille que sur sa part. On recolle ensuite les réponses bout à bout, et
l'on retrouve une liste de la même longueur qu'au départ. Quatre têtes sont
dessinées ici ; le modèle en emploie seize.}[tetes-attention]
  \begin{scope}[x=1cm, y=1cm]
    \node[font=\sffamily\scriptsize, text=mbSlate] at (-1.5,3.3) {le jeton};
    \foreach \k in {0,...,3}{
      \pgfmathtruncatemacro{\tone}{22+16*\k}
      \fill[mbAzure!\tone] (\k*1.9,3.05) rectangle ++(1.9,0.5);
      \draw[draw=white, line width=1.6pt] (\k*1.9,3.05) rectangle ++(1.9,0.5);
      \draw[draw=mbSlate!50, line width=0.5pt] (\k*1.9,3.05) rectangle ++(1.9,0.5);
    }

    \foreach \k/\quoi in {0/{les formes}, 1/{les couleurs},
                          2/{les mouvements}, 3/{les bords}}{
      \draw[fleche] (\k*1.9+0.95,3.0) -- (\k*1.9+0.95,2.05);
      \node[bloc calcul, minimum width=17mm, minimum height=8mm,
            font=\sffamily\scriptsize] (t\k) at (\k*1.9+0.95,1.6) {tête \the\numexpr\k+1\relax};
      \node[font=\sffamily\scriptsize, text=mbAmber!80!black]
            at (\k*1.9+0.95,0.85) {\quoi};
      \draw[fleche] (\k*1.9+0.95,0.55) -- (\k*1.9+0.95,-0.2);
    }

    \foreach \k in {0,...,3}{
      \pgfmathtruncatemacro{\tone}{22+16*\k}
      \fill[mbTeal!\tone] (\k*1.9,-0.75) rectangle ++(1.9,0.5);
      \draw[draw=white, line width=1.6pt] (\k*1.9,-0.75) rectangle ++(1.9,0.5);
      \draw[draw=mbSlate!50, line width=0.5pt] (\k*1.9,-0.75) rectangle ++(1.9,0.5);
    }
    \node[font=\sffamily\scriptsize, text=mbSlate] at (-1.6,-0.5)
         {on recolle};
    \draw[accolade bas] (0,-0.9) -- (7.6,-0.9);
    \node[font=\sffamily\scriptsize, text=mbSlate] at (3.8,-1.55)
         {la même longueur qu'au départ};
  \end{scope}
\end{schema}

\begin{exemple}[Répartir mille vingt-quatre nombres entre seize têtes]
Prenons un jeton de $1\,024$ nombres et $16$ têtes d'attention. Chaque tête
reçoit
\[
  \frac{1\,024}{16} = 64 \text{ nombres.}
\]
Chacune produit une réponse de $64$ nombres. En recollant les seize
réponses, on retrouve
\[
  16 \times 64 = 1\,024 \text{ nombres.}
\]
La taille du jeton n'a pas changé, mais seize questions différentes ont été
posées au lieu d'une seule, et pour à peu près le même prix.
\end{exemple}

\begin{intuition}
Retournons dans la salle de classe. Avec une seule tête, l'élève ne peut
poser qu'une question et doit attendre la réponse avant d'en poser une
autre. Avec seize têtes, c'est comme s'il avait seize paires d'oreilles :
l'une écoute les formes, l'autre les couleurs, une troisième les
mouvements. Il rassemble ensuite tout ce qu'il a entendu en une seule fiche.
\end{intuition}

Nous tenons la pièce centrale. Il faut maintenant l'emballer proprement, et
c'est le rôle du bloc de transformeur.

\section{Le bloc de transformeur, ouvert en coupe}

\begin{enclair}
Un bloc de transformeur, c'est l'attention plus trois pièces qui la rendent
utilisable : une remise au propre des nombres avant chaque étage, un petit
réseau qui réfléchit sur chaque jeton séparément, et deux raccourcis qui
laissent passer l'information sans la déformer.
\end{enclair}

L'attention toute seule ne suffit pas. Elle mélange les jetons entre eux,
mais elle ne fait aucun raisonnement à l'intérieur d'un jeton. Et surtout,
si l'on empile beaucoup d'attentions à la suite, les nombres partent dans
tous les sens et l'entraînement échoue. Trois pièces règlent ces problèmes.

\subsection*{Première pièce : la normalisation de couche}

\begin{definition}[La normalisation de couche, ou LayerNorm][layernorm]
La \emph{normalisation de couche} remet les nombres d'un jeton à une échelle
raisonnable. Pour chaque jeton pris séparément, on calcule la moyenne de ses
nombres, on la retire de chacun, puis on divise le tout par un nombre qui
mesure à quel point ils s'écartaient de cette moyenne. Après cette
opération, les nombres du jeton sont centrés autour de zéro et ne sont ni
minuscules ni gigantesques. On l'appelle presque toujours par son nom
anglais, \emph{LayerNorm}.
\end{definition}

\begin{exemple}[Normaliser quatre nombres]
Un jeton contient les nombres $(2 \ ; \ 4 \ ; \ 6 \ ; \ 8)$.

Sa moyenne vaut $(2+4+6+8)/4 = 20/4 = 5$. On la retire de chacun :
\[ (-3 \ ; \ -1 \ ; \ 1 \ ; \ 3). \]
On mesure ensuite de combien les nombres s'écartent en moyenne. On multiplie
chaque écart par lui-même — c'est ce qu'on appelle l'élever au carré — ce
qui donne $9$, $1$, $1$, $9$. Multiplier par soi-même sert ici à effacer les
signes moins, car un nombre négatif multiplié par lui-même devient positif.
La moyenne de ces carrés vaut $(9+1+1+9)/4 = 5$, et sa racine carrée vaut
$\sqrt{5} \approx 2{,}236$. On divise :
\[
  \left( \frac{-3}{2{,}236} \ ; \ \frac{-1}{2{,}236} \ ; \
         \frac{1}{2{,}236} \ ; \ \frac{3}{2{,}236} \right)
  \approx (-1{,}342 \ ; \ -0{,}447 \ ; \ 0{,}447 \ ; \ 1{,}342).
\]

Les nombres tiennent maintenant à peu près entre $-1{,}5$ et $1{,}5$. Le
jeton dit toujours la même chose — les rapports entre ses nombres n'ont pas
changé — mais il le dit à voix normale au lieu de crier ou de chuchoter.
\end{exemple}

Voilà nos nombres remis d'aplomb. Passons à la pièce qui va s'en servir.

\subsection*{Deuxième pièce : le petit réseau de réflexion}

Après avoir écouté les autres, chaque jeton a besoin d'un moment pour
digérer ce qu'il vient d'entendre. C'est le rôle du second étage.

\begin{definition}[Le MLP][mlp]
Un \emph{MLP} est un petit réseau appliqué à chaque jeton séparément, sans
aucun regard sur les autres. Ces trois lettres viennent de l'anglais
\emph{Multi-Layer Perceptron}. Ce nom ne vous servira jamais : retenez
seulement «~petit réseau de réflexion~». Il procède en
trois temps : il agrandit d'abord la liste de nombres, en général en la
multipliant par quatre ; il applique ensuite à chaque nombre une petite
opération de tri appelée \emph{fonction d'activation} ; enfin il ramène la
liste à sa taille de départ.
\end{definition}

\begin{definition}[La fonction d'activation][activation]
Vous connaissez déjà cette pièce, sans connaître son nom. Dans le programme
du chapitre~\ref{chap:auto-supervision}, nous avions glissé une petite règle
appelée \co{python}{ReLU}, qui remplaçait par zéro tout nombre négatif et
laissait passer les autres. Nous avions dit alors que, sans elle, empiler
deux couches ne servirait à rien.

Cette règle porte un nom de famille : on l'appelle une \emph{fonction
d'activation}. C'est une opération appliquée à chaque nombre séparément, et
qui ne se contente pas de le multiplier ou de lui ajouter quelque chose. La
\co{python}{ReLU} est la plus simple de la famille ; celle qu'emploient les
transformeurs, la \co{python}{GELU}, arrondit un peu les angles mais rend le
même service. Sans l'une d'elles, toute la pile de blocs se ramènerait à une
seule multiplication suivie d'une seule addition.
\end{definition}

\begin{exemple}[Un MLP en tout petit]
Prenons un jeton de deux nombres, $(3 \ ; \ -2)$, et un MLP qui agrandit la
liste en la multipliant par deux, donc qui passe de deux à quatre nombres.

\textbf{On agrandit.} Après le premier calcul, disons que l'on obtient
\[ (4 \ ; \ -1 \ ; \ 6 \ ; \ -5). \]

\textbf{On trie.} La fonction d'activation remplace les négatifs par zéro :
\[ (4 \ ; \ 0 \ ; \ 6 \ ; \ 0). \]

Regardez ce qui vient de se passer : deux des quatre mesures ont été
éteintes, deux sont restées allumées. Le jeton ne réagit donc qu'à certaines
situations et reste muet devant les autres. C'est exactement ce qu'une suite
de multiplications ne saurait pas faire.

\textbf{On revient à la taille de départ.} Un dernier calcul ramène ces
quatre nombres à deux, et le résultat est ajouté au jeton d'origine — nous
allons voir tout de suite pourquoi.
\end{exemple}

\begin{intuition}
L'attention, c'est le moment où l'on discute avec ses voisins. Le MLP, c'est
le moment où l'on retourne à sa place pour réfléchir tout seul à ce qu'on
vient d'apprendre. Un bloc de transformeur fait donc alterner, encore et
encore, un temps de discussion et un temps de réflexion personnelle.
\end{intuition}

Il nous reste une pièce, la plus discrète des trois, et pourtant celle sans
laquelle rien ne tiendrait debout.

\subsection*{Troisième pièce : la connexion résiduelle}

\begin{definition}[La connexion résiduelle][residuelle]
Une \emph{connexion résiduelle} est un raccourci : au lieu de remplacer le
jeton par le résultat d'un étage, on \emph{ajoute} ce résultat au jeton
d'origine. L'ancienne valeur n'est donc jamais perdue ; l'étage ne fait
qu'apporter une correction. On dit aussi \emph{connexion saute-mouton},
parce que l'information peut sauter par-dessus l'étage.
\end{definition}

\begin{exemple}[Ce que le raccourci change, en chiffres]
Un jeton vaut $(3 \ ; \ -2)$ et l'étage qu'il traverse produit
$(0{,}4 \ ; \ 0{,}1)$.

\textbf{Sans raccourci}, le jeton devient $(0{,}4 \ ; \ 0{,}1)$ : tout ce
qu'il contenait a disparu, remplacé par le travail de l'étage.

\textbf{Avec raccourci}, on additionne :
\[ (3 \ ; \ -2) + (0{,}4 \ ; \ 0{,}1) = (3{,}4 \ ; \ -1{,}9). \]

Le jeton est toujours reconnaissable, simplement corrigé un peu. Imaginez
maintenant quarante-huit étages à la suite : dans le premier cas, il ne reste
plus rien du jeton de départ après quelques étages ; dans le second, il est
encore là, augmenté de quarante-huit petites corrections.
\end{exemple}

\begin{intuition}
Vous relisez le devoir d'un élève. Vous pouvez soit tout réécrire de votre
main — et si vous vous trompez, le travail de l'élève est perdu — soit
écrire vos remarques en marge, à côté de son texte. La connexion résiduelle,
c'est la deuxième façon de faire. Elle explique pourquoi l'on peut empiler
quarante-huit blocs les uns sur les autres sans que le message du départ se
dissolve en chemin.
\end{intuition}

Assemblons maintenant les trois pièces autour de l'attention.

\begin{schema}{Un bloc de transformeur ouvert en coupe. Le jeton entre par
le bas. Il passe dans deux étages successifs. Chaque étage commence par une
normalisation, puis fait son travail — écouter les autres pour le premier,
réfléchir tout seul pour le second — et son résultat est \emph{ajouté} à ce
qui entrait dans l'étage, grâce au raccourci dessiné sur le côté. Le jeton
ressort par le haut, de la même taille qu'à l'entrée.}[bloc-transformeur]
  \begin{scope}[x=1cm, y=1cm]
    \node[font=\sffamily\small\bfseries, text=mbSlate] at (0,-0.85)
         {jeton qui entre};
    \draw[fleche] (0,-0.55) -- (0,0.15);

    %---- premier étage : le jeton écoute les autres ------------------
    \node[bloc, minimum width=48mm, text width=42mm] (n1) at (0,0.7)
         {normalisation};
    \node[bloc calcul, minimum width=48mm, text width=42mm] (a1) at (0,2.3)
         {attention multi-têtes};
    \node[circle, draw=mbAmber!85!black, fill=mbAmber!12, line width=0.9pt,
          minimum size=8mm, inner sep=0pt, font=\small] (s1) at (0,3.9) {$+$};

    \draw[fleche] (n1) -- (a1);
    \draw[fleche] (a1) -- (s1);
    \draw[fleche accent] (-3.5,0.15) -- (-3.5,3.9) -- (s1);
    \draw[draw=mbSlate!45, line width=0.6pt] (0,0.15) -- (-3.5,0.15);
    \node[annotation, rotate=90] at (-3.82,2.05) {raccourci};

    %---- second étage : le jeton réfléchit tout seul -----------------
    \node[bloc, minimum width=48mm, text width=42mm] (n2) at (0,5.3)
         {normalisation};
    \node[bloc calcul, minimum width=48mm, text width=42mm] (a2) at (0,6.9)
         {petit réseau de réflexion};
    \node[circle, draw=mbAmber!85!black, fill=mbAmber!12, line width=0.9pt,
          minimum size=8mm, inner sep=0pt, font=\small] (s2) at (0,8.5) {$+$};

    \draw[fleche] (s1) -- (n2);
    \draw[fleche] (n2) -- (a2);
    \draw[fleche] (a2) -- (s2);
    \draw[fleche accent] (-3.5,4.65) -- (-3.5,8.5) -- (s2);
    \draw[draw=mbSlate!45, line width=0.6pt] (0,4.65) -- (-3.5,4.65);
    \node[annotation, rotate=90] at (-3.82,6.6) {raccourci};

    \draw[fleche] (s2) -- (0,9.3);
    \node[font=\sffamily\small\bfseries, text=mbSlate] at (0,9.6)
         {jeton qui sort};

    %---- les bulles, posées à distance des cadres --------------------
    \node[annotation bulle, text width=2.9cm, anchor=west] at (3.1,2.3)
         {ici, le jeton\\ discute avec\\ les autres};
    \node[annotation bulle, text width=2.9cm, anchor=west] at (3.1,6.9)
         {ici, il réfléchit\\ tout seul};
  \end{scope}
\end{schema}

\begin{definition}[Le bloc de transformeur][bloc-transformeur]
Un \emph{bloc de transformeur} est l'ensemble que nous venons d'assembler :
une normalisation puis une attention multi-têtes, le tout ajouté à l'entrée ;
puis une normalisation puis un petit réseau de réflexion, le tout ajouté de
nouveau. Un bloc reçoit une file de jetons et rend une file de jetons de la
même taille. C'est pour cela qu'on peut en empiler autant qu'on veut.
\end{definition}

\begin{exemple}[Un bloc, en chiffres]
Donnons à un bloc les $2\,048$ jetons d'un clip de la phase principale,
chacun fait de $1\,024$ nombres. À la sortie, nous récupérons encore
$2\,048$ jetons de $1\,024$ nombres.

Rien n'a changé en taille, tout a changé en contenu : chaque jeton a écouté
les $2\,047$ autres, puis a réfléchi de son côté. Comme l'entrée et la
sortie ont exactement la même forme, on peut brancher la sortie de ce bloc
sur l'entrée du suivant, sans le moindre ajustement. C'est ce détail, et lui
seul, qui rend les empilements possibles.
\end{exemple}

Et c'est justement ce que l'on fait dans la section suivante.

\section{Empiler les blocs : la profondeur et les quatre tailles}

\begin{enclair}
On met les blocs les uns au-dessus des autres, et la file de jetons les
traverse tous. Le nombre de blocs empilés s'appelle la profondeur. Plus on
monte, plus ce que voient les jetons devient général.
\end{enclair}

\begin{definition}[La profondeur][profondeur]
La \emph{profondeur} d'un transformeur est le nombre de blocs empilés les uns
sur les autres. Les jetons traversent le premier bloc, puis le deuxième, et
ainsi de suite jusqu'au dernier. Chaque bloc reçoit le travail du précédent.
\end{definition}

Ce qui se passe le long de cette pile vaut la peine qu'on s'y arrête, et
nous nous en servirons énormément dans la partie~IV. Les premiers blocs
repèrent des choses très simples : un bord, un changement de couleur, un
léger déplacement. Les blocs du milieu assemblent ces petites choses en
morceaux reconnaissables : un coin, une roue, une main. Les derniers blocs
assemblent ces morceaux en idées complètes : une voiture, quelqu'un qui
verse de l'eau.

\begin{schema}{Ce que voient les blocs, du bas vers le haut. En bas, on
repère des traits. Un peu plus haut, des coins. Plus haut encore, des
morceaux d'objets. Tout en haut, l'objet entier et l'action. C'est cette
montée en généralité qui explique pourquoi, plus tard, nous aurons parfois
envie d'aller chercher des informations à mi-hauteur plutôt qu'au
sommet.}[montee-blocs]
  \begin{scope}[x=1cm, y=1cm]
    \foreach \k/\nom/\quoi in {
      0/{blocs du bas}/{des traits},
      1/{blocs du milieu}/{des coins},
      2/{blocs du haut}/{des roues, des mains},
      3/{dernier bloc}/{une voiture, un geste}}{
      \pgfmathsetmacro{\yy}{\k*1.35}
      \pgfmathtruncatemacro{\tone}{12+16*\k}
      \node[rectangle, rounded corners=2.5pt, draw=mbBlue, fill=mbBlue!\tone,
            line width=0.8pt, minimum width=38mm, minimum height=9mm,
            font=\sffamily\small, text=mbInk] (b\k) at (0,\yy) {\nom};
      \node[font=\sffamily\small, text=mbSlate, anchor=west] at (2.4,\yy)
           {il y voit \quoi};
    }
    \draw[fleche] (b0) -- (b1);
    \draw[fleche] (b1) -- (b2);
    \draw[fleche] (b2) -- (b3);
    \draw[fleche accent] (-2.5,-0.55) -- (-2.5,4.6);
    \node[annotation, rotate=90] at (-2.85,2.0) {de plus en plus général};
  \end{scope}
\end{schema}

\begin{piege}
Cette montée en généralité a un revers, et c'est exactement le problème que
V-JEPA~2.1 vient résoudre. Si le sommet ne parle plus que d'idées générales,
il a pu oublier \emph{où} les choses se trouvaient. Or pour dire à quelle
distance est chaque pixel, on a justement besoin du «~où~». Gardez cette
tension en tête : elle est le sujet de toute la partie~III.
\end{piege}

\begin{piege}
Il faut maintenant corriger une image que nous vous avions donnée au
chapitre~\ref{chap:auto-supervision}. Nous y avions dessiné l'encodeur comme
un \emph{entonnoir} : beaucoup de nombres entrent, un petit paquet sort. Cette
image était commode pour comprendre l'idée de résumé, mais elle est trop
simple.

La vérité, vous venez de la voir : l'encodeur ne rend pas \emph{un} paquet
pour toute l'image. Il rend \emph{un paquet par jeton}. Un clip de la phase
principale entre avec $2\,048$ jetons, et il ressort avec $2\,048$ jetons,
un pour chaque endroit du clip. L'entonnoir n'a pas rétréci la file : il a
rétréci chaque morceau, chacun de son côté.

Retenez bien cette différence, car tout le livre repose dessus. S'il vous
faut une seule réponse pour toute la vidéo — « que se passe-t-il ici ? » —
il faudra mélanger ces $2\,048$ paquets en un seul. Mais s'il vous faut une
réponse \emph{par endroit} — « à quelle distance est ce pixel-là ? » — il
faudra que chaque paquet, pris tout seul, décrive bien son propre morceau.
Ces deux besoins tirent le modèle dans des directions différentes, et c'est
précisément le conflit que raconte la partie~III.
\end{piege}

Combien de blocs empile-t-on, et de quelle taille ? La communauté s'est mise
d'accord sur quelques tailles standards, désignées par une lettre. V-JEPA~2.1
en emploie quatre.

\begin{surlesdonnees}
Les quatre tailles employées dans l'article, de la plus petite à la plus
grande :

\begin{center}
\footnotesize
\begin{tabular}{llrl}
\toprule
\textbf{Nom} & \textbf{Se lit} & \textbf{Boutons} & \textbf{Rôle dans l'article} \\
\midrule
ViT-B & «~bé~», pour \emph{Base}     & $80$ M            & le petit modèle de service \\
ViT-L & «~èl~», pour \emph{Large}    & $300$ M           & le modèle des essais \\
ViT-g & «~gé~» minuscule, \emph{giant} & $1$ milliard    & la taille de V-JEPA~2 \\
ViT-G & «~gé~» majuscule, \emph{Giant} & $2$ milliards   & le grand modèle final \\
\bottomrule
\end{tabular}
\end{center}

Les deux plus petits ne sont pas entraînés depuis le début : ils apprennent
en copiant le grand, par un procédé appelé \emph{distillation} que nous
verrons au chapitre~15. Retenez seulement, pour l'instant, qu'un petit modèle
peut hériter de presque toute la qualité d'un grand.

Attention à la lettre \textbf{g}~: la minuscule et la majuscule désignent
deux modèles différents, et l'article emploie les deux à quelques lignes
d'écart. Le grand modèle final, celui qui obtient tous les meilleurs
résultats, est le ViT-G majuscule et compte $48$ blocs empilés. Nous le
savons parce que l'article indique qu'il va chercher des informations aux
blocs numéro $12$, $24$, $36$ et $48$ — nous verrons au
chapitre~11 pourquoi il fait cela.
\end{surlesdonnees}

\begin{exemple}[Ce que «~deux milliards de boutons~» veut dire]
Reprenez le chapitre~\ref{chap:apprendre} : un bouton, c'est un nombre que
l'entraînement règle. Le ViT-G en possède deux milliards. Si vous les
imprimiez à raison de cinquante nombres par ligne et cinquante lignes par
page, il vous faudrait huit cent mille pages, soit une pile de papier de
quatre-vingts mètres de haut. Et l'entraînement les règle \emph{tous},
$147\,000$ fois de suite.

C'est aussi pour cela que nous consacrerons tout le chapitre~15 à rétrécir
ce modèle : personne ne veut faire tourner quatre-vingts mètres de papier
dans un serveur, et l'article montre comment ramener presque toute cette
qualité dans un modèle six fois plus petit.
\end{exemple}

Il nous reste une dernière chose à regarder en face avant de coder : le prix
de tout cela.

\section{Ce que l'attention coûte vraiment}

\begin{enclair}
Dans un transformeur, le nombre de jetons est le bouton le plus important
pour la vitesse. Le doubler ne double pas le travail : il le multiplie par
quatre. C'est la première chose à regarder quand on choisit une résolution.
\end{enclair}

Nous l'avons dit au début du chapitre, mais nous pouvons maintenant le dire
avec de vrais chiffres. Dans l'attention, chaque jeton compare sa question à
la clé de tous les autres. Avec $n$ jetons, cela fait $n \times n$
comparaisons.

\begin{exemple}[Le prix des deux configurations de V-JEPA 2.1]
Phase principale, $2\,048$ jetons :
\[ 2\,048 \times 2\,048 = 4\,194\,304 \text{ comparaisons.} \]

Phase de refroidissement, $18\,432$ jetons :
\[ 18\,432 \times 18\,432 = 339\,738\,624 \text{ comparaisons.} \]

Le rapport vaut
\[ \frac{339\,738\,624}{4\,194\,304} = 81 . \]

Quatre-vingt-une fois plus de travail, et cela dans \emph{chacun} des
quarante-huit blocs. Voilà le chiffre qui explique tout le calendrier de
l'entraînement décrit dans l'article.
\end{exemple}

\begin{aretenir}
Retenez la règle du praticien : \textbf{ce que vous paierez suit à peu près
le carré du nombre de jetons}. Quand votre modèle sera trop lent
en service, le premier bouton à tourner ne sera ni le processeur ni la
bibliothèque : ce sera la résolution d'entrée, parce que c'est elle qui
décide du nombre de jetons. Nous referons ce calcul, pour de vrai, au
chapitre~29.
\end{aretenir}

Nous avons monté l'œil de la machine en entier. Vérifions maintenant, avec
un programme, que tout ce que nous avons calculé à la main est bien ce qui se
passe pour de vrai.

\section{Passons au code}

Nous allons écrire quatre petits programmes. Les deux premiers découpent une
image et un clip en jetons, et retrouvent les quatre nombres du tableau. Le
troisième refait, ligne par ligne, l'exemple d'attention calculé à la main.
Le quatrième assemble un bloc de transformeur complet. Vous pouvez les taper
tels quels : ils tournent sur n'importe quel ordinateur, même sans le
processeur spécialisé dont les gros entraînements ont besoin.

Commençons par le découpage. Comme annoncé, on l'obtient avec une
convolution dont la fenêtre et le pas ont la même taille.

\begin{code}{python}[Découper une image et un clip en jetons]
import torch
import torch.nn as nn

# --- La voie image : une convolution 2D de 16 x 16, qui avance de 16 ---
tokeniseur_image = nn.Conv2d(
    in_channels=3,       # les trois couleurs
    out_channels=1024,   # la taille d'un jeton
    kernel_size=16,      # la fenetre : un patch de 16 x 16
    stride=16,           # elle avance de 16 : aucun chevauchement
)

image = torch.zeros(1, 3, 256, 256)          # (lot, couleurs, hauteur, largeur)
sortie = tokeniseur_image(image)             # -> (1, 1024, 16, 16)
jetons_image = sortie.flatten(2).transpose(1, 2)

print(jetons_image.shape)   # -> torch.Size([1, 256, 1024])
# 256 jetons : c'est bien le nombre calcule a la main.

# --- La voie video : une convolution 3D, avec 2 images d'epaisseur ---
tokeniseur_video = nn.Conv3d(
    in_channels=3,
    out_channels=1024,
    kernel_size=(2, 16, 16),   # (tubelet, hauteur, largeur)
    stride=(2, 16, 16),
)

clip = torch.zeros(1, 3, 16, 256, 256)       # (lot, couleurs, temps, H, W)
sortie = tokeniseur_video(clip)              # -> (1, 1024, 8, 16, 16)
jetons_video = sortie.flatten(2).transpose(1, 2)

print(jetons_video.shape)   # -> torch.Size([1, 2048, 1024])
# 2048 jetons : 8 x 16 x 16, exactement notre calcul.
\end{code}

\begin{piege}
Regardez bien l'ordre des axes : PyTorch attend les couleurs en deuxième
position, pas en dernière. C'est le changement d'ordre que nous avions déjà
rencontré au chapitre~\ref{chap:image-video}. Si vous lui donnez un tenseur
rangé dans l'ordre de l'affichage, il ne se plaindra pas toujours : il
calculera quelque chose de faux, en silence. Prenez l'habitude d'afficher la
forme avant et après chaque étape.
\end{piege}

Vérifions maintenant les quatre nombres du tableau, tous d'un coup.

\begin{code}{python}[Vérifier les quatre nombres de jetons]
def compte_jetons(images, hauteur, largeur, patch=16, tubelet=2):
    """Rend le nombre de jetons. Mettre images=None pour une image seule."""
    n = (hauteur // patch) * (largeur // patch)
    if images is not None:
        n = n * (images // tubelet)
    return n

print(compte_jetons(None, 256, 256))   # -> 256
print(compte_jetons(None, 512, 512))   # -> 1024
print(compte_jetons(16,  256, 256))    # -> 2048
print(compte_jetons(64,  384, 384))    # -> 18432

# Et le prix de l'attention, qui grimpe au carre :
print(18432**2 / 2048**2)              # -> 81.0
\end{code}

Passons à l'attention. Nous reprenons exactement les nombres de l'exemple
calculé à la main : vous devez retrouver les mêmes parts, à la troisième
décimale.

\begin{code}{python}[L'attention, refaite avec nos nombres]
import torch

# Notre question, puis les cles et les contenus de fourrure, oreille, flaque.
q = torch.tensor([2.0, 0.0, 0.0, 0.0])
k = torch.tensor([[2.0, 0.0, 0.0, 0.0],    # fourrure
                  [4.0, 0.0, 0.0, 0.0],    # oreille
                  [0.0, 3.0, 0.0, 0.0]])   # flaque
v = torch.tensor([[10.0,  0.0,  0.0, 0.0],
                  [ 0.0, 10.0,  0.0, 0.0],
                  [ 0.0,  0.0, 10.0, 0.0]])

scores = k @ q                      # les produits scalaires
print(scores)                       # -> tensor([4., 8., 0.])

scores = scores / (q.shape[0] ** 0.5)   # on divise par racine de 4, donc par 2
print(scores)                       # -> tensor([2., 4., 0.])

poids = torch.softmax(scores, dim=0)
print(poids)                        # -> tensor([0.1173, 0.8668, 0.0159])
print(poids.sum())                  # -> tensor(1.)

resultat = poids @ v
print(resultat)                     # -> tensor([1.1731, 8.6681, 0.1588, 0.])
\end{code}
Les trois parts sont bien $0{,}117$, $0{,}867$ et $0{,}016$, dans cet ordre,
et le résultat est bien celui que nous avions trouvé à la main. Le programme
ne fait rien d'autre que ce que nous avons écrit sur le papier.

Terminons par un bloc de transformeur complet. Comptez les lignes : il n'y
en a pas vingt. Tout ce chapitre tient là-dedans.

\begin{code}{python}[Un bloc de transformeur, en entier]
import torch
import torch.nn as nn

class BlocTransformeur(nn.Module):
    def __init__(self, dimension=1024, nb_tetes=16):
        super().__init__()
        self.norme1 = nn.LayerNorm(dimension)
        self.attention = nn.MultiheadAttention(
            dimension, nb_tetes, batch_first=True)
        self.norme2 = nn.LayerNorm(dimension)
        self.reflexion = nn.Sequential(
            nn.Linear(dimension, 4 * dimension),   # on agrandit
            nn.GELU(),                             # la fonction d'activation
            nn.Linear(4 * dimension, dimension),   # on revient a la taille de depart
        )

    def forward(self, jetons):
        # Premier etage : on ecoute les autres, et on AJOUTE le resultat.
        propre = self.norme1(jetons)
        ecoute, _ = self.attention(propre, propre, propre)
        jetons = jetons + ecoute

        # Second etage : on reflechit tout seul, et on AJOUTE encore.
        jetons = jetons + self.reflexion(self.norme2(jetons))
        return jetons

bloc = BlocTransformeur()
jetons = torch.zeros(1, 2048, 1024)     # un clip de la phase principale
print(bloc(jetons).shape)               # -> torch.Size([1, 2048, 1024])

# La taille ne change pas : c'est ce qui permet d'empiler 48 blocs.
encodeur = nn.Sequential(*[BlocTransformeur() for _ in range(48)])
\end{code}

\begin{piege}
Repérez les deux \co{python}{jetons = jetons + ...} : ce sont les connexions
résiduelles. Si vous écrivez \co{python}{jetons = ecoute} au lieu de
\co{python}{jetons = jetons + ecoute}, le programme tournera sans erreur,
mais l'entraînement d'un modèle profond ne convergera jamais. C'est l'une
des erreurs les plus fréquentes, et l'une des plus difficiles à repérer,
parce que rien ne casse : les résultats sont simplement mauvais.
\end{piege}

\begin{aretenir}
Un transformeur de vision, c'est quatre gestes et rien de plus. On découpe
l'image ou la vidéo en petits morceaux. On remplace chaque morceau par une
liste de nombres, dans laquelle on inscrit aussi sa position. On empile des
blocs, où les jetons discutent entre eux puis réfléchissent chacun de leur
côté. Et l'on récupère à la sortie une file de jetons de la même taille
qu'à l'entrée, mais dont chacun a été enrichi par tous les autres.
\end{aretenir}

\vspace{0.5em}

Nous avons maintenant tout ce qu'il faut. Nous savons ce qu'est une image
pour une machine, ce que veut dire apprendre, comment jouer au jeu des trous
sans étiquettes, et de quoi est fait l'œil qui regarde. Ces quatre chapitres
formaient les fondations, et elles sont posées.

Une seule question reste en suspens, celle que nous avions laissée à la fin
du chapitre~\ref{chap:auto-supervision} : quand la machine devine ce qui est
caché, doit-elle deviner les pixels, ou l'idée ? La réponse à cette question
porte un nom, JEPA, et c'est par elle que s'ouvre la partie~II.

\phantomsection
\section*{Exercices}
\addcontentsline{toc}{section}{Exercices}

\begin{exercice}[Compter des jetons]
Vous donnez au modèle une image de $768 \times 768$ pixels, avec des patchs
de $16$. Combien de jetons obtenez-vous ? Et un clip de $32$ images en
$256 \times 256$, avec des tubelets de $2$ ?
\end{exercice}

\begin{solution}
Pour l'image : $768 / 16 = 48$ sur chaque côté, donc
$48 \times 48 = 2\,304$ jetons.

Pour le clip : $32 / 2 = 16$ tranches de temps, et $256/16 = 16$ patchs sur
chaque côté, donc $16 \times 16 \times 16 = 4\,096$ jetons. Remarquez que ce
clip coûte $(4\,096 / 2\,048)^2 = 4$ fois plus cher qu'un clip de la phase
principale, alors qu'il ne contient que deux fois plus d'images.
\end{solution}

\begin{exercice}[Choisir une résolution sous contrainte]
Votre serveur ne peut pas dépasser $5\,000$ jetons par clip. Vous voulez des
clips de $16$ images. Quelle est la plus grande résolution carrée, multiple
de $16$, que vous pouvez vous permettre ?
\end{exercice}

\begin{solution}
Avec $16$ images et des tubelets de $2$, il y a $8$ tranches de temps. Le
budget par tranche vaut donc $5\,000 / 8 = 625$ patchs. Comme l'image est
carrée, il faut au plus $\sqrt{625} = 25$ patchs de côté, ce qui donne
$25 \times 16 = 400$ pixels. Et $400$ est bien un multiple de $16$ :
la réponse est donc $400 \times 400$, qui consomme
$8 \times 25 \times 25 = 5\,000$ jetons, exactement le budget.

Un mot de prudence, quand même. Tomber pile sur la limite ne laisse aucune
marge, et l'article, lui, a retenu $384 \times 384$ pour ses vidéos, ce qui
donne $8 \times 24 \times 24 = 4\,608$ jetons. Choisir la résolution juste
en dessous de son plafond est une habitude qui évite bien des mauvaises
surprises le jour où le serveur est chargé.
\end{solution}

\begin{exercice}[Répartir entre les têtes]
Un modèle a des jetons de $1\,408$ nombres et $22$ têtes d'attention.
Combien de nombres chaque tête reçoit-elle ? Que se passerait-il avec $20$
têtes ?
\end{exercice}

\begin{solution}
Avec $22$ têtes : $1\,408 / 22 = 64$ nombres chacune. Le compte tombe juste.

Avec $20$ têtes : $1\,408 / 20 = 70{,}4$, ce qui n'est pas un nombre entier.
On ne peut pas donner $70{,}4$ nombres à une tête. La bibliothèque refusera
de construire le modèle. C'est pour cela que le nombre de têtes doit toujours
diviser exactement la taille du jeton.
\end{solution}

\begin{exercice}[Une softmax à la main]
Trois scores valent $1$, $1$ et $3$. Calculez les trois parts. Que
remarquez-vous sur les deux premières ?
\end{exercice}

\begin{solution}
Les exponentielles valent $e^1 \approx 2{,}718$, $e^1 \approx 2{,}718$ et
$e^3 \approx 20{,}086$. Leur somme vaut $25{,}522$. Les parts valent donc

\[
  \frac{2{,}718}{25{,}522} \approx 0{,}106, \qquad
  \frac{2{,}718}{25{,}522} \approx 0{,}106, \qquad
  \frac{20{,}086}{25{,}522} \approx 0{,}787 .
\]

Les deux premières sont égales, ce qui est normal : elles avaient le même
score. Et l'on voit l'effet de l'exponentielle : le score $3$ n'est que trois
fois plus grand que le score $1$, mais il récolte sept fois et demie plus de
part, puisque $0{,}787 / 0{,}106 \approx 7{,}4$. La softmax écarte bel et
bien les grands des petits.
\end{solution}

\begin{exercice}[Le mélange des positions]
Un collègue vous dit : «~j'ai mélangé l'ordre de mes jetons avant de les
donner au modèle, mais comme ils contiennent tous la même information
qu'avant, cela ne change rien~». A-t-il raison ?
\end{exercice}

\begin{solution}
Il a raison sur un point et tort sur l'essentiel. L'attention traite bien la
file comme un sac : à elle seule, elle rendrait le même résultat.

Mais les jetons ne contiennent pas que leur contenu : l'encodage de position
y a inscrit l'endroit d'où ils viennent. En mélangeant la file, il a donc
associé chaque morceau à une mauvaise adresse. Pour le modèle, c'est comme
si l'on avait recollé les pièces du puzzle au hasard : les pièces sont
justes, le dessin ne l'est plus. Le résultat sera faux, et sans le moindre
message d'erreur.
\end{solution}

\begin{exercice}[Le prix de la finesse]
Un ingénieur propose de passer les patchs de $16$ à $8$ pixels, «~pour voir
plus de détails~», en gardant des images de $256 \times 256$. Par combien le
travail de l'attention est-il multiplié ?
\end{exercice}

\begin{solution}
Avec des patchs de $16$, on avait $16 \times 16 = 256$ jetons. Avec des
patchs de $8$, on obtient $256/8 = 32$ patchs de côté, donc
$32 \times 32 = 1\,024$ jetons : quatre fois plus.

Comme le coût grimpe au carré, le travail est multiplié par
$4 \times 4 = 16$. Sa proposition améliorera sans doute la finesse, mais elle
coûtera seize fois plus cher. C'est un choix parfaitement défendable — à
condition de le faire les yeux ouverts, et de vérifier que le gain de qualité
vaut la dépense.
\end{solution}
      % patchs, jetons, attention, RoPE

\part{Comprendre JEPA et V-JEPA 2, l'ancêtre direct}
\chapter{L'idée JEPA : deviner l'idée, pas les pixels}
\label{chap:idee-jepa}

\begin{quote}\itshape
Nous avons laissé une question en suspens à la fin de la partie~I. Quand la
machine cache un morceau de ce qu'elle regarde et essaie de le deviner,
que doit-elle deviner au juste~? Les pixels qui manquent, un par un~? Ou
bien autre chose~? Ce chapitre répond à cette question, et la réponse est
le cœur de tout le reste du livre. Nous allons monter, pièce par pièce, la
machine qui porte cette réponse. Elle s'appelle JEPA.
\end{quote}

\begin{prerequis}
Les quatre chapitres de la partie~I, et rien d'autre. Il faut savoir qu'une
image est une grille de nombres (chapitre~\ref{chap:image-video})~; qu'un
modèle possède des boutons que l'on tourne pour faire baisser une erreur
(chapitre~\ref{chap:apprendre})~; qu'on peut apprendre sans étiquettes en
cachant un morceau de l'image, et qu'un encodeur transforme une image en un
petit paquet de nombres appelé représentation
(chapitre~\ref{chap:auto-supervision})~; et enfin comment une image devient
une file de jetons qui discutent entre eux
(chapitre~\ref{chap:vit}). Les seules opérations employées ici seront
l'addition, la soustraction, la multiplication et la division.
\end{prerequis}

Au chapitre~\ref{chap:auto-supervision}, nous avons trouvé un moyen
d'apprendre sans que personne n'écrive la bonne réponse~: on cache un
morceau de l'image, et on demande à la machine de deviner ce qui manque.
C'était une belle idée, et elle marchait déjà sur notre petit monde de huit
cases sur huit.

Puis, à la fin de ce même chapitre, nous avons buté sur une question que
nous avons laissée ouverte~: fallait-il deviner les \emph{pixels} cachés, ou
seulement \emph{l'idée} de ce qui manquait~? Nous avions même donné un avant-goût
de la réponse, et lâché un nom~: JEPA. Le moment est venu de tenir cette
promesse. Nous allons reprendre cette question, la trancher pour de bon, et
surtout monter la machine qui en découle.

\section{Deviner les pixels, ou deviner l'idée}

\begin{enclair}
Quand on cache un morceau d'image, on peut demander deux choses très
différentes à la machine. Soit on lui demande de \emph{redessiner} le
morceau manquant, pixel par pixel, comme un peintre qui refait un bout de
tableau. Soit on lui demande seulement de \emph{décrire} ce qu'il y avait,
comme quelqu'un qui dirait «~là, il y avait de l'herbe~». La première
demande est terriblement difficile. La seconde est raisonnable. V-JEPA
choisit la seconde.
\end{enclair}

Reprenons la pelouse du chapitre~\ref{chap:auto-supervision}. On a une photo
de jardin, on cache un carré au milieu, et ce carré tombe sur l'herbe.
Posons cette fois les deux questions côte à côte, bien en face l'une de
l'autre.

\textbf{Première question, celle du peintre.} «~Redessine-moi ce carré,
exactement comme il était.~» Nous savons déjà où cela mène~: il faudrait
retrouver chaque brin d'herbe, celui qui penche à gauche, celui qui est un
peu jauni, la petite ombre entre les deux. Ces détails-là ne se déduisent de
rien, ils sont le fruit du hasard, et la machine s'épuiserait à deviner ce
qui ne peut pas se deviner.

\textbf{Deuxième question, celle du témoin.} «~Dis-moi seulement ce qu'il y
avait dans ce carré.~» Cette fois, la machine peut répondre~: «~de l'herbe, verte, en
pleine lumière, et c'est la suite de la pelouse d'à côté~». Cette réponse-là,
elle peut la trouver. Elle est même obligée de comprendre la scène pour la
trouver. Et surtout, elle a le droit de \emph{ne pas savoir} où était
exactement chaque brin d'herbe, parce qu'on ne le lui demande pas.

\begin{intuition}
Imaginez qu'un ami vous téléphone depuis une pièce que vous ne voyez pas, et
qu'il vous demande de deviner ce qu'il y a derrière lui. Si vous devez dire
«~il y a une bibliothèque, avec des livres~», vous avez une chance de tomber
juste. Si vous devez dire le titre exact de chaque livre, dans l'ordre, vous
n'avez aucune chance. Le premier jeu vous apprend à réfléchir. Le second
ne vous apprend rien~: il vous fait seulement perdre à tous les coups.
\end{intuition}

\begin{definition}[Méthode générative][generative]
Une méthode est dite \emph{générative} quand elle demande au modèle de
reconstruire l'entrée elle-même, c'est-à-dire de ressortir des pixels. Le
mot vient de «~générer~», qui veut dire «~fabriquer~». On compare alors les
pixels fabriqués aux vrais pixels.
\end{definition}

\begin{definition}[Méthode prédictive dans l'espace latent][predictif-latent]
Une méthode est dite \emph{prédictive dans l'espace latent} quand elle
demande au modèle de prédire, non pas les pixels, mais la
\emph{représentation} de la zone manquante — c'est-à-dire le petit paquet de
nombres qui la résume. On rappelle que l'\emph{espace latent} est simplement
le nom savant du monde des représentations~: l'endroit où vivent ces paquets
de nombres, une fois que l'encodeur a fait son travail
(chapitre~\ref{chap:auto-supervision}).
\end{definition}

Voilà les deux chemins nommés. Le chapitre~\ref{chap:auto-supervision} nous
avait mis sur la piste du second~; nous savons maintenant comment il
s'appelle, et la figure~\ref{fig:jepa-deux-voies} les met côte à côte pour
qu'on ne les confonde plus jamais.

\begin{schema}{Les deux façons de remplir un trou. En haut, la méthode
générative redessine les pixels, et doit donc inventer chaque brin d'herbe.
En bas, la méthode JEPA se contente de décrire ce qui manquait, avec un
petit paquet de nombres.}[jepa-deux-voies]

% ---------- Ligne du haut : la voie générative ----------
\node[font=\sffamily\bfseries\scriptsize, text=mbInk, anchor=west]
      at (0,2.05) {La façon du peintre : redessiner les pixels};

\fill[mbAzure!18] (0,0.8) rectangle (2.2,1.6);
\fill[mbGreen!30] (0,0)   rectangle (2.2,0.8);
\draw[mbSlate, line width=0.7pt] (0,0) rectangle (2.2,1.6);
\fill[mbSlate!25] (0.75,0.4) rectangle (1.45,1.1);
\draw[mbSlate, dashed, line width=0.7pt] (0.75,0.4) rectangle (1.45,1.1);
\node[font=\sffamily\scriptsize, text=mbSlate] at (1.1,-0.35) {la photo trouée};

\draw[fleche] (2.45,0.8) -- (3.65,0.8);
\node[etiquette] at (3.05,1.1) {on redessine};

\fill[mbAzure!18] (3.9,0.8) rectangle (6.1,1.6);
\fill[mbGreen!30] (3.9,0)   rectangle (6.1,0.8);
\draw[mbSlate, line width=0.7pt] (3.9,0) rectangle (6.1,1.6);
\fill[mbGreen!30] (4.65,0.4) rectangle (5.35,1.1);
\draw[mbGreen!60!black, line width=0.6pt] (4.65,0.4) rectangle (5.35,1.1);
\foreach \k in {0,1,...,6}{
  \draw[draw=mbGreen!70!black, line width=0.5pt]
        (4.72+\k*0.09, 0.45) -- (4.67+\k*0.09, 0.95);
}
\node[font=\sffamily\scriptsize, text=mbSlate] at (5,-0.35) {chaque brin inventé};

\node[annotation bulle, align=center, anchor=west] (bulhaut) at (7.15,0.8)
     {il faudrait deviner\\ce qui ne se devine pas};
\draw[trait annotation, line width=0.7pt, draw=mbAmber!80]
      (6.25,0.8) -- (bulhaut.west);

% ---------- Ligne du bas : la voie JEPA ----------
\begin{scope}[yshift=-3.3cm]
\node[font=\sffamily\bfseries\scriptsize, text=mbInk, anchor=west]
      at (0,2.05) {La façon de JEPA : décrire l'idée};

\fill[mbAzure!18] (0,0.8) rectangle (2.2,1.6);
\fill[mbGreen!30] (0,0)   rectangle (2.2,0.8);
\draw[mbSlate, line width=0.7pt] (0,0) rectangle (2.2,1.6);
\fill[mbSlate!25] (0.75,0.4) rectangle (1.45,1.1);
\draw[mbSlate, dashed, line width=0.7pt] (0.75,0.4) rectangle (1.45,1.1);
\node[font=\sffamily\scriptsize, text=mbSlate] at (1.1,-0.35) {la photo trouée};

% Le nuage est posé AVANT les flèches : celles-ci s'accrochent ensuite à
% ses ancrages (nuage.west, nuage.east) au lieu de viser des coordonnées
% devinées à la main. Sans cela, leurs pointes se retrouvent cachées sous
% le nuage, qui déborde largement de sa largeur annoncée.
\node[cloud, cloud puffs=13, cloud puff arc=110, aspect=1.9,
      draw=mbAmber!80!black, fill=mbAmber!12, line width=0.7pt,
      minimum width=2.7cm, minimum height=1.6cm,
      font=\sffamily\scriptsize, text=mbInk] (nuage) at (5.05,0.8)
     {$0{,}8$ \ $-0{,}3$ \ $1{,}2$};

\draw[fleche accent, shorten >=2pt] (2.45,0.8) -- (nuage.west);
\node[etiquette] at (3.05,1.1) {on décrit};

\node[font=\sffamily\scriptsize, text=mbSlate] at (5.05,-0.35)
     {un petit paquet de nombres};

\node[annotation bulle, align=center, anchor=west] (bulbas) at (7.15,0.8)
     {« de l'herbe verte,\\en pleine lumière » :\\c'est faisable};
\draw[trait annotation, line width=0.7pt, draw=mbAmber!80, shorten <=2pt]
      (nuage.east) -- (bulbas.west);
\end{scope}
\end{schema}

\begin{piege}
Une erreur revient tout le temps, alors autant l'écarter tout de suite~: on
croit souvent que la voie JEPA est meilleure parce qu'elle aurait
\emph{moins de nombres} à deviner. C'est faux, et il faut le savoir. Un
carré de $64 \times 64$ pixels en couleur, cela fait
$64 \times 64 \times 3 = 12\,288$ nombres à redessiner. La représentation de
ce même carré peut très bien en contenir davantage. Ce n'est donc pas une
question de \emph{quantité}. C'est une question de \emph{nature}~: les
pixels sont imposés de l'extérieur, avec leur part de hasard, alors que la
représentation est fabriquée par la machine elle-même, qui a donc le droit
de laisser de côté ce qui n'est pas devinable.
\end{piege}

\begin{aretenir}
Prédire des pixels oblige à deviner du hasard. Prédire une représentation
permet de l'oublier. C'est là toute la différence, et c'est elle qui donne
son nom à toute la famille JEPA.
\end{aretenir}

Nous savons maintenant \emph{quoi} deviner. Voyons à partir de quoi, et
mettons enfin des noms sur les deux images qui entrent dans la machine.

\section{Les deux vues : la trouée et la complète}

\begin{enclair}
La machine regarde toujours la même vidéo deux fois. La première fois, on
lui en cache des morceaux~: c'est la vue trouée. La deuxième fois, on la lui
montre en entier~: c'est la vue complète. La première sert de question, la
seconde sert de corrigé.
\end{enclair}

Reprenons le jeu du chapitre~\ref{chap:auto-supervision}, mais en donnant
cette fois un nom précis à chaque chose. On part d'une vidéo. On en fabrique
deux versions.

\begin{definition}[La vue trouée, notée $x$][vue-trouee]
La \emph{vue trouée} est la vidéo à laquelle on a retiré des morceaux. On la
note $x$. Les morceaux retirés s'appellent les zones \emph{masquées}, et les
morceaux que l'on a laissés s'appellent le \emph{contexte}. On dit aussi que
$x$ est la vue \emph{corrompue}, ce qui veut simplement dire «~abîmée
exprès~».
\end{definition}

\begin{definition}[La vue complète, notée $y$][vue-complete]
La \emph{vue complète} est la même vidéo, sans rien retirer du tout. On la
note $y$. On l'appelle aussi la vue \emph{propre}, au sens de «~intacte~».
\end{definition}

Attention à ne pas se tromper sur un point~: $x$ et $y$ ne sont pas deux
vidéos différentes. C'est la même vidéo, la même seconde, la même scène. La
seule différence, c'est qu'on a caché des morceaux dans l'une et pas dans
l'autre.

\begin{exemple}[Compter les jetons cachés]
Reprenons un clip de la phase principale d'entraînement, celui du
chapitre~\ref{chap:vit}~: seize images en $256 \times 256$, avec des patchs
de $16$ et des tubelets de $2$. Nous avions compté
$8 \times 16 \times 16 = 2\,048$ jetons.

Disons que le masque en recouvre $1\,200$. Alors~:

\begin{center}
\footnotesize
\begin{tabular}{lr}
\toprule
\textbf{Ce que l'on compte} & \textbf{Combien} \\
\midrule
Jetons en tout dans le clip        & $2\,048$ \\
Jetons masqués (le trou)           & $1\,200$ \\
Jetons de contexte (ce qui reste)  & $848$ \\
\bottomrule
\end{tabular}
\end{center}

La machine ne verra donc que $848$ jetons, et devra parler des $1\,200$
autres. Retenez ces nombres~: ils vont nous resservir tout au long du
chapitre.
\end{exemple}

\begin{schema}{La même vidéo, vue deux fois. À gauche la vue trouée $x$,
où seuls les jetons bleus sont visibles. À droite la vue complète $y$, où
tout est visible.}[jepa-vues]

% --- vue trouée ---
\foreach \i in {0,1,2,3,4,5}{
  \foreach \j in {0,1,2,3,4}{
    \fill[mbAzure!25, draw=mbAzure!60, line width=0.4pt]
          (\i*0.5, \j*0.5) rectangle ++(0.46,0.46);
  }
}
\foreach \i in {1,2,3}{
  \foreach \j in {1,2,3}{
    \fill[mbSlate!30, draw=mbSlate!70, line width=0.4pt, dashed]
          (\i*0.5, \j*0.5) rectangle ++(0.46,0.46);
  }
}
\node[font=\sffamily\small, text=mbInk] at (1.5,-0.55)
     {la vue trouée $x$};
\node[etiquette, text=mbAzure!80!black] at (1.5,-1.05) {bleu : le contexte};
\node[etiquette, text=mbSlate] at (1.5,-1.45) {gris : les zones masquées};

% --- vue complète ---
\begin{scope}[xshift=5.5cm]
\foreach \i in {0,1,2,3,4,5}{
  \foreach \j in {0,1,2,3,4}{
    \fill[mbTeal!22, draw=mbTeal!60, line width=0.4pt]
          (\i*0.5, \j*0.5) rectangle ++(0.46,0.46);
  }
}
\node[font=\sffamily\small, text=mbInk] at (1.5,-0.55)
     {la vue complète $y$};
\node[etiquette, text=mbTeal!80!black] at (1.5,-1.05) {tout est visible};
\end{scope}
\end{schema}

Nous avons nos deux vues. Il faut maintenant deux appareils pour les
regarder. Et là, une petite surprise nous attend~: ces deux appareils sont
bâtis de la même façon, mais leurs boutons ne sont pas réglés pareil.

\section{L'élève et le professeur : deux encodeurs, presque jumeaux}

\begin{enclair}
La vue trouée est regardée par un premier encodeur, que nous appellerons
l'\emph{élève}. La vue complète est regardée par un second encodeur, que
nous appellerons le \emph{professeur}. L'élève est celui qui apprend~: on
lui tourne ses boutons à chaque pas. Le professeur, lui, ne travaille pas~:
il se contente de recopier l'élève, très lentement.
\end{enclair}

Avant d'aller plus loin, il faut une convention d'écriture. Nous avons vu au
chapitre~\ref{chap:apprendre} qu'un modèle est une machine à boutons, et
qu'apprendre revient à tourner ces boutons. Un gros modèle en a des
millions. Plutôt que de les nommer un par un, on leur donne un seul nom
collectif, avec une lettre grecque.

\begin{notation}[Les lettres grecques des boutons][lettres-grecques]
\begin{itemize}
  \item $\theta$ se lit «~thêta~». C'est le nom collectif de \emph{tous} les
        boutons de l'encodeur élève. Écrire $E_{\theta}$ veut donc dire
        «~l'encodeur, réglé avec les boutons $\theta$~».
  \item $\bar{\theta}$ se lit «~thêta barre~». C'est le nom collectif de
        tous les boutons de l'encodeur professeur. La petite barre au-dessus
        rappelle qu'il s'agit d'une \emph{moyenne} de l'élève~: nous verrons
        cela en détail dans la section~\ref{sec:ema}.
  \item $\varphi$ se lit «~phi~». C'est le nom collectif des boutons du
        prédicteur, la pièce que nous montons juste après.
\end{itemize}
Aucune de ces lettres ne cache une difficulté. Ce sont des étiquettes, rien
de plus. On aurait pu les appeler $A$, $B$ et $C$.
\end{notation}

Voici ce que fait chacun des deux encodeurs.

\textbf{L'encodeur élève, $E_{\theta}$.} Il reçoit la vue trouée $x$,
c'est-à-dire seulement les $848$ jetons de contexte de notre exemple. Il ne
voit rien d'autre. Il les fait discuter entre eux, comme nous l'avons vu au
chapitre~\ref{chap:vit}, et il ressort $848$ jetons enrichis. C'est lui, et
lui seul, qui apprend. À chaque pas d'entraînement, ses boutons bougent.

\textbf{L'encodeur professeur, $E_{\bar{\theta}}$.} Il reçoit la vue
complète $y$, donc les $2\,048$ jetons, sans aucun trou. Il ressort $2\,048$
jetons enrichis. Ces jetons-là serviront de \emph{corrigé}~: ce sont les
réponses que l'élève doit essayer de retrouver. Le professeur, lui,
n'apprend pas au sens habituel~: aucun calcul d'erreur ne vient tourner ses
boutons. Il se contente de suivre l'élève de très loin.

\begin{intuition}
Pensez à un élève qui fait un devoir avec la moitié des pages arrachées, et
à un professeur qui a le sujet complet sous les yeux. Le professeur ne fait
pas d'effort~: il lit. Mais c'est sa copie qui sert de référence. Et le
professeur n'est pas quelqu'un d'autre~: c'est le même élève, quelques
milliers de pas plus tôt, une version un peu plus posée de lui-même.
\end{intuition}

\begin{notation}[Les autres noms de ces deux encodeurs][x-encodeur]
Dans l'article, et dans le code publié sur internet, ces deux encodeurs
portent des noms tirés des lettres des deux vues. L'encodeur élève, celui
qui regarde la vue trouée $x$, s'appelle le \emph{$x$-encodeur}. L'encodeur
professeur, celui qui regarde la vue complète $y$, s'appelle le
\emph{$y$-encodeur}. Ce sont exactement les mêmes pièces que notre élève et
notre professeur~: seul le nom change. Nous emploierons les deux façons de
dire, pour que vous soyez à l'aise le jour où vous ouvrirez le code.
\end{notation}

\begin{schema}{La paire d'encodeurs. Les deux appareils sont bâtis de la
même façon, mais un seul des deux apprend~: le professeur se contente de
recopier l'élève, très lentement.}[jepa-paire]

\node[bloc, font=\sffamily\scriptsize, minimum width=28mm] (el) at (0,0)
     {encodeur élève\\$E_{\theta}$\\[1pt]\scriptsize il apprend};
\node[bloc calcul, font=\sffamily\scriptsize, minimum width=28mm]
     (pr) at (0,-2.6)
     {encodeur professeur\\$E_{\bar{\theta}}$\\[1pt]\scriptsize il recopie};

\node[donnee, font=\sffamily\scriptsize, minimum width=15mm]
     (vx) at (-3.7,0) {vue trouée\\$x$};
\node[donnee, font=\sffamily\scriptsize, minimum width=15mm]
     (vy) at (-3.7,-2.6) {vue complète\\$y$};
\draw[fleche] (vx) -- (el);
\draw[fleche] (vy) -- (pr);

\draw[fleche] (el.east) -- ++(1.0,0)
      node[right, font=\sffamily\scriptsize, text=mbInk, align=left]
      {848 jetons\\enrichis};
\draw[fleche] (pr.east) -- ++(1.0,0)
      node[right, font=\sffamily\scriptsize, text=mbInk, align=left]
      {2 048 cibles};

\draw[fleche pointillee] (el) -- (pr);
\node[etiquette, align=center] at (-1.35,-1.3)
     {recopie\\très lentement};
\end{schema}

\begin{definition}[Cible][cible]
On appelle \emph{cible} la sortie du professeur~: c'est le paquet de nombres
que la machine essaie de retrouver. Le mot est celui du tir à l'arc~: c'est
ce que l'on vise. Attention à ne pas la confondre avec les \emph{étiquettes}
du chapitre~\ref{chap:auto-supervision}, qui étaient écrites par des
humains~: ici, la cible est fabriquée par la machine elle-même, et personne
ne l'a jamais écrite nulle part.
\end{definition}

Une question doit vous venir à l'esprit, et elle est excellente~: pourquoi
faire deux encodeurs, alors qu'un seul suffirait~? Après tout, on pourrait
donner $x$ et $y$ au même appareil. Gardez cette question bien au chaud~: la
réponse arrive à la section~\ref{sec:effondrement}, et c'est une des plus
jolies idées du livre.

Pour l'instant, il nous manque encore une pièce, et c'est celle qui fait le
lien entre les deux côtés.

\section{Le prédicteur : la pièce qui fait le pont}

\begin{enclair}
L'élève ne voit que les morceaux visibles. Il ne peut donc rien dire des
morceaux cachés. Il faut une pièce supplémentaire, qui prenne ce que l'élève
a compris et qui en tire une réponse pour les endroits qu'il n'a jamais vus.
Cette pièce s'appelle le \emph{prédicteur}.
\end{enclair}

Cette pièce, nous l'avons déjà croisée au
chapitre~\ref{chap:auto-supervision}~: c'est elle qui, dans notre petit
programme, partait du résumé et tentait de deviner la moitié cachée de
l'image. Nous la retrouvons ici, en beaucoup plus grand. Il est donc temps
de lui donner son écriture officielle.

\begin{notation}[Le prédicteur s'écrit $P_{\varphi}$][notation-predicteur]
On note le prédicteur $P_{\varphi}$, ce qui se lit «~$P$ de phi~» et veut
dire~: le prédicteur, réglé avec les boutons $\varphi$. Le
chapitre~\ref{chap:auto-supervision} nous avait déjà prévenus d'une chose
qu'il ne faut pas oublier~: à la fin de l'entraînement, c'est l'encodeur que
l'on garde précieusement, pas le prédicteur. Une chose change tout de même
ici~: notre prédicteur ne rend plus une seule réponse, mais une réponse par
jeton reçu.
\end{notation}

Le prédicteur reçoit deux choses, et il faut bien les distinguer.

\textbf{Un~: les jetons de contexte}, c'est-à-dire les $848$ jetons que
l'élève vient d'enrichir. Ce sont les indices dont on dispose.

\textbf{Deux~: les jetons de masque.} Voilà une nouveauté qu'il faut
expliquer, car elle surprend toujours. Pour les $1\,200$ endroits cachés, on
ne peut évidemment pas donner le contenu, puisque c'est justement ce que
l'on cache. On donne alors à la place un jeton spécial, qui ne dit qu'une
seule chose~: «~ici, à cet endroit précis, il manquait quelque chose~».

\begin{definition}[Jeton de masque, noté $\Delta_y$][jeton-masque]
Un \emph{jeton de masque} est un jeton qui ne contient aucune information
sur l'image, mais qui porte une adresse~: la position spatiale et temporelle
du morceau caché. Son contenu est un paquet de nombres appris pendant
l'entraînement, le même pour tous les trous~; seule l'adresse change. On
note l'ensemble de ces jetons $\Delta_y$, et $\Delta$ se lit «~delta~».
\end{definition}

\begin{intuition}
Un jeton de masque, c'est une enveloppe vide sur laquelle on a écrit
l'adresse. On la tend au prédicteur en disant~: «~remplis-la~». Sans
l'adresse, le prédicteur ne saurait pas de quel trou on parle, et il
répondrait la même chose pour tous.
\end{intuition}

Le prédicteur met bout à bout ces deux paquets de jetons, les fait discuter
entre eux — c'est encore un transformeur, exactement comme au
chapitre~\ref{chap:vit} — et ressort un jeton pour chaque jeton d'entrée.

\begin{exemple}[Le compte exact des jetons dans le prédicteur]
Avec nos chiffres, le prédicteur reçoit~:

\[
  \underbrace{848}_{\text{jetons de contexte}}
  \;+\;
  \underbrace{1\,200}_{\text{jetons de masque}}
  \;=\;
  2\,048 \ \text{jetons en entrée},
\]

et il ressort $2\,048$ jetons en sortie. Notez bien ce résultat, car il va
compter énormément dans quelques chapitres~: \textbf{le prédicteur produit
une réponse pour tout le monde, y compris pour les endroits qui étaient déjà
visibles.} Nous verrons au chapitre suivant que V-JEPA~2 jette une bonne
partie de ces réponses, et au chapitre~10 que c'est précisément là que se
cache le défaut que V-JEPA~2.1 vient corriger.
\end{exemple}

Nous avons enfin toutes les pièces. Assemblons-les.

\section{Le schéma complet, monté sous vos yeux}

Voici la machine entière, figure~\ref{fig:jepa-central}. C'est le dessin le
plus important du livre~: tous les chapitres suivants ne feront que le
compléter ou le corriger. Prenez le temps de le suivre du doigt, de gauche à
droite.

\begin{schema}{L'architecture JEPA au complet. En haut, le chemin de
l'élève, qui part de la vue trouée. En bas, celui du professeur, qui part de
la vue complète. À droite, on compare les deux~: cet écart est la perte. La
barre rouge marque l'endroit où la correction n'a pas le droit de
passer.}[jepa-central]

% ---- chemin de l'élève (en haut) ----
\node[donnee, font=\sffamily\scriptsize, minimum width=16mm] (x) at (0,0)
     {vue trouée\\$x$};
\node[bloc, font=\sffamily\scriptsize, minimum width=19mm] (enc) at (2.9,0)
     {encodeur élève\\$E_{\theta}$};
\node[bloc accent, font=\sffamily\scriptsize, minimum width=19mm]
     (pred) at (6.0,0) {prédicteur\\$P_{\varphi}$};
\node[bloc, font=\sffamily\scriptsize, minimum width=16mm, draw=mbAmber,
      fill=mbAmber!8] (sortie) at (8.9,0) {prédiction};

\draw[fleche] (x) -- (enc);
\draw[fleche] (enc) -- (pred);
\node[etiquette] at (4.45,0.42) {848 jetons};
\draw[fleche] (pred) -- (sortie);

\node[bloc, font=\sffamily\scriptsize, minimum width=19mm, minimum height=6mm,
      draw=mbSlate!70, fill=mbMist] (masque) at (6.0,1.7)
     {jetons de masque $\Delta_y$};
\draw[fleche] (masque) -- (pred);

% ---- chemin du professeur (en bas) ----
\node[donnee, font=\sffamily\scriptsize, minimum width=16mm] (y) at (0,-3.3)
     {vue complète\\$y$};
\node[bloc calcul, font=\sffamily\scriptsize, minimum width=19mm]
     (encp) at (2.9,-3.3) {encodeur professeur\\$E_{\bar{\theta}}$};
\node[bloc calcul, font=\sffamily\scriptsize, minimum width=16mm]
     (target) at (8.9,-3.3) {cible};

\draw[fleche] (y) -- (encp);
\draw[fleche] (encp) -- (target);
\node[etiquette] at (5.4,-2.88) {2 048 jetons};

% ---- la comparaison ----
\draw[<->, draw=mbRose!85!black, line width=1pt] (sortie) -- (target);
\node[annotation bulle, align=center, anchor=east] at (8.35,-1.65)
     {on compare :\\cet écart,\\c'est la perte};

% ---- la copie lente ----
\draw[fleche pointillee] (enc) -- (encp);
\node[etiquette, align=center, anchor=east] at (2.72,-1.65)
     {copie très lente\\(moyenne mobile)};

% ---- le gradient bloqué ----
\draw[line width=2pt, draw=mbRose!85!black] (7.3,-3.75) -- (7.3,-2.85);
\node[annotation, text=mbRose!85!black, align=center] at (7.3,-4.35)
     {gradient bloqué :\\rien ne revient par ici};
\end{schema}

Récapitulons le trajet en cinq phrases, dans l'ordre du dessin.

\begin{enumerate}
  \item La vue trouée $x$ entre dans l'encodeur élève, qui en ressort les
        jetons de contexte enrichis.
  \item Ces jetons sont mis bout à bout avec les jetons de masque, qui
        portent les adresses des trous.
  \item Le prédicteur fait discuter tout ce beau monde et produit une
        prédiction pour chaque position.
  \item Pendant ce temps, la vue complète $y$ est passée dans l'encodeur
        professeur, qui produit les cibles.
  \item On compare la prédiction à la cible. L'écart est la \emph{perte}, et
        c'est lui que l'entraînement fait baisser.
\end{enumerate}

Il reste, sur ce dessin, deux choses que nous n'avons pas encore
expliquées~: la flèche en pointillé de la copie lente, et la barre rouge du
gradient bloqué. Elles ont l'air de petits détails de dessin. Ce sont en
réalité deux \emph{garde-fous}, et sans eux tout s'écroule.

Le mot mérite une explication, car nous allons l'employer sans arrêt. Un
garde-fou, dans la vie de tous les jours, c'est la barrière que l'on pose au
bord d'un pont pour empêcher les gens de tomber. Elle ne sert à rien tant
que tout va bien~; elle sert le jour où quelqu'un glisse. Nos deux
garde-fous sont exactement cela~: deux barrières qui ne gênent personne tant
que la machine travaille honnêtement, et qui l'empêchent de tomber le jour
où elle essaie de tricher. Nous allons voir pourquoi, mais donnons d'abord
son vrai nom à cette machine.

\section{JEPA : le nom décodé mot à mot}

\begin{enclair}
Cette architecture porte un nom très court, fait de quatre lettres~: JEPA.
Chacune de ces lettres est la première lettre d'un mot anglais. Nous allons
les traduire un par un, et vous verrez qu'il n'y a rien d'autre, derrière
ces quatre lettres, que la machine que nous venons de monter.
\end{enclair}

Ce nom s'écrit JEPA et se prononce «~jé-pa~». Il vient de \emph{Joint
Embedding Predictive Architecture}, une idée proposée
par~\textcite{lecun2022path}. Voici le décodage.

\begin{center}
\footnotesize
\begin{tabular}{llp{6.6cm}}
\toprule
\textbf{Lettre} & \textbf{Mot anglais} & \textbf{Ce que cela veut dire ici} \\
\midrule
J & \emph{Joint} &
«~conjoint~», c'est-à-dire à deux. Il y a deux chemins, celui de l'élève et
celui du professeur, et ils travaillent ensemble. \\[2pt]
E & \emph{Embedding} &
le \emph{plongement}, c'est-à-dire le fait de transformer une image en un
paquet de nombres. Nous l'avons vu au chapitre~\ref{chap:vit}. \\[2pt]
P & \emph{Predictive} &
«~prédictive~»~: il y a une pièce qui devine, le prédicteur. \\[2pt]
A & \emph{Architecture} &
la façon dont les pièces sont assemblées. Ni plus, ni moins. \\
\bottomrule
\end{tabular}
\end{center}

\begin{definition}[JEPA][jepa]
Une \emph{architecture prédictive à plongements conjoints}, ou JEPA, est un
assemblage où deux encodeurs transforment deux vues d'une même donnée en
paquets de nombres, et où un prédicteur essaie de retrouver le paquet du
second à partir du paquet du premier. La comparaison se fait toujours entre
paquets de nombres, jamais entre pixels.
\end{definition}

\begin{enia}
Cette idée n'est pas propre à la vidéo. On l'a appliquée au son, aux images
fixes avec I-JEPA~\parencite{assran2023ijepa}, puis à la vidéo avec la
famille V-JEPA, dont le «~V~» veut dire \emph{video}. V-JEPA~2.1, le modèle
de ce livre, est le dernier membre de cette famille. Toute la partie~IV
consistera à ajouter des pièces à ce squelette, sans jamais le changer.
\end{enia}

Le nom est posé, la machine est montée. Nous pouvons maintenant aborder le
danger qui menace tout ce montage — et comprendre enfin pourquoi il faut
deux encodeurs plutôt qu'un.

\section{Le grand danger : l'effondrement}
\label{sec:effondrement}

\begin{enclair}
Notre machine a une faille énorme. Si on la laisse faire, elle va trouver
une solution qui donne une erreur parfaitement nulle, sans rien apprendre du
tout. Il lui suffit de répondre toujours la même chose. C'est ce qu'on
appelle l'\emph{effondrement}, et c'est le cauchemar de tous ceux qui
entraînent ce genre de modèles.
\end{enclair}

Reprenons ce que nous demandons à la machine~: «~fais en sorte que ta
prédiction ressemble à la cible~». Et rappelons-nous du
chapitre~\ref{chap:apprendre}~: une machine qui apprend ne cherche pas à
bien faire. Elle cherche à faire baisser un nombre. Si elle trouve un moyen
malhonnête de faire baisser ce nombre, elle le prendra, sans le moindre
remords.

Or, il existe ici un moyen malhonnête, et il est même très facile à trouver.

\begin{definition}[Effondrement des représentations][effondrement]
Il y a \emph{effondrement} quand l'encodeur se met à sortir le même paquet
de nombres quelle que soit l'image qu'on lui donne. Les représentations ne
distinguent alors plus rien, et le modèle, bien que très content de lui,
est devenu complètement inutile.
\end{definition}

\begin{exemple}[Pourquoi la triche donne un zéro parfait]
Imaginons que les deux encodeurs se mettent d'accord pour toujours sortir le
paquet $(0 \,;\, 0 \,;\, 0)$, quoi qu'on leur montre. Et que le prédicteur,
de son côté, réponde toujours $(0 \,;\, 0 \,;\, 0)$ lui aussi.

Calculons l'erreur. On mesure l'écart en additionnant les différences prises
en valeur absolue, c'est-à-dire sans se soucier du signe~:

\[
  |0 - 0| \;+\; |0 - 0| \;+\; |0 - 0| \;=\; 0 .
\]

L'erreur vaut zéro. C'est le meilleur score possible, celui qu'aucun modèle
honnête n'atteindra jamais. Et pourtant ce modèle-là ne sait strictement
rien~: il donne la même réponse pour un chat, une rue et une plage.
\end{exemple}

\begin{schema}{L'effondrement. Trois vidéos qui n'ont rien à voir entre
elles ressortent de l'encodeur avec exactement le même paquet de nombres.
L'erreur est nulle, et le modèle est bon à jeter.}[jepa-effondrement]

\node[donnee, font=\sffamily\scriptsize, minimum width=15mm] (a) at (0,1.5)
     {un chat};
\node[donnee, font=\sffamily\scriptsize, minimum width=15mm] (b) at (0,0)
     {une rue};
\node[donnee, font=\sffamily\scriptsize, minimum width=15mm] (c) at (0,-1.5)
     {une plage};

\node[bloc, font=\sffamily\scriptsize, minimum width=21mm, minimum height=24mm]
     (e) at (3.5,0) {encodeur\\effondré};

\draw[fleche] (a) -- (e);
\draw[fleche] (b) -- (e);
\draw[fleche] (c) -- (e);

\node[bloc, font=\sffamily\scriptsize, draw=mbRose, fill=mbRose!8,
      minimum width=19mm] (s1) at (7.1,1.5) {$(0\,;0\,;0)$};
\node[bloc, font=\sffamily\scriptsize, draw=mbRose, fill=mbRose!8,
      minimum width=19mm] (s2) at (7.1,0)   {$(0\,;0\,;0)$};
\node[bloc, font=\sffamily\scriptsize, draw=mbRose, fill=mbRose!8,
      minimum width=19mm] (s3) at (7.1,-1.5){$(0\,;0\,;0)$};

\draw[fleche] (e) -- (s1);
\draw[fleche] (e) -- (s2);
\draw[fleche] (e) -- (s3);

\draw[accolade] (8.2,-1.9) -- (8.2,1.9);
\node[annotation, align=center, anchor=west] at (8.6,0)
     {toujours\\la même\\réponse};

\node[annotation, align=center] at (3.5,-2.4)
     {erreur nulle, savoir nul};
\end{schema}

\begin{intuition}
C'est l'élève qui a compris que le professeur note la ressemblance entre les
deux copies, et pas leur contenu. Il rend une feuille blanche. Le
professeur, qui est en réalité le même élève, rend lui aussi une feuille
blanche. Les deux copies se ressemblent parfaitement. Note maximale, et
aucun apprentissage.
\end{intuition}

Voilà pourquoi la question posée tout à l'heure était si bonne. Si l'on
utilisait \emph{un seul} encodeur des deux côtés, et qu'on le laissait
corriger ses boutons librement à partir des deux chemins, il découvrirait
cette triche en quelques centaines de pas. Nous le verrons se produire pour
de vrai, en direct, dans la section de code de ce chapitre.

Il faut donc empêcher la machine de tricher. Les auteurs emploient pour
cela deux garde-fous, et nous les prenons l'un après l'autre.

\section{Premier garde-fou : le gradient bloqué}

\begin{enclair}
Quand la machine se trompe, une correction remonte à l'envers dans tout le
montage pour tourner les boutons dans le bon sens. Le premier garde-fou
consiste à poser une barrière sur le chemin du professeur~: la correction a
le droit de remonter du côté de l'élève, mais pas du côté du corrigé. On
interdit ainsi à la machine de bouger la cible qu'elle est censée viser.
\end{enclair}

Rappelons d'abord d'où vient ce mot de gradient, vu au
chapitre~\ref{chap:apprendre}. Après avoir mesuré l'erreur, on calcule pour
chaque bouton dans quel sens et de combien il faudrait le tourner pour que
l'erreur diminue. Cette information voyage en sens inverse du calcul, de la
sortie vers l'entrée. C'est le gradient.

Maintenant, réfléchissons à ce que la machine ferait si on la laissait
libre. Elle doit rapprocher deux choses~: la prédiction et la cible. Il y a
deux façons de rapprocher deux points~:

\begin{enumerate}
  \item déplacer la prédiction vers la cible — c'est ce que l'on veut,
        puisque cela oblige l'élève à comprendre la scène~;
  \item déplacer la cible vers la prédiction — c'est la triche, puisqu'il
        suffit alors d'écraser tout le monde au même endroit.
\end{enumerate}

\begin{exemple}[Les deux façons de refermer un écart]
Prenons un seul nombre, pour voir la chose en face. La prédiction vaut
$0{,}5$ et la cible vaut $0{,}9$. L'écart entre les deux vaut donc
$0{,}9 - 0{,}5 = 0{,}4$.

\textbf{La façon honnête}~: on fait monter la prédiction de $0{,}5$ jusqu'à
$0{,}9$. L'écart tombe à $0$. Pour y arriver, l'élève a bien été obligé de
mieux regarder l'image.

\textbf{La façon tricheuse}~: on fait descendre la cible de $0{,}9$ jusqu'à
$0{,}5$. L'écart tombe à $0$ lui aussi, et pourtant personne n'a rien
appris. Pire~: si on laisse faire, la prédiction et la cible vont continuer
à se rapprocher l'une de l'autre jusqu'à valoir toutes les deux zéro.

Les deux façons donnent le même score. Voilà pourquoi il faut interdire la
seconde~: la machine, elle, ne voit aucune différence entre les deux.
\end{exemple}

La deuxième façon est de très loin la plus facile. La machine la
choisirait tout de suite. On la lui interdit donc, tout simplement.

\begin{definition}[Gradient bloqué, noté $\mathrm{sg}(\cdot)$][stop-gradient]
Le \emph{gradient bloqué} est une instruction qui dit~: «~calcule cette
valeur normalement, mais ne fais remonter aucune correction par ce
chemin~». On l'écrit $\mathrm{sg}(\cdot)$, d'après l'anglais
\emph{stop-gradient}, et l'on met entre les parenthèses ce que l'on veut
protéger. Dans notre montage, on protège la cible.
\end{definition}

\begin{schema}{Le trajet de la correction. Les flèches grises du haut sont
le calcul, qui va de gauche à droite. Les flèches orange du dessous sont la
correction, qui revient en arrière. La barre rouge lui interdit d'entrer
dans le chemin du professeur.}[jepa-stopgrad]

\node[bloc, font=\sffamily\scriptsize, minimum width=21mm] (e1) at (0,0)
     {encodeur élève};
\node[bloc accent, font=\sffamily\scriptsize, minimum width=21mm]
     (p1) at (3.7,0) {prédicteur};
\node[bloc, font=\sffamily\scriptsize, minimum width=16mm, draw=mbAmber,
      fill=mbAmber!8] (l1) at (7.4,0) {la perte};
\node[bloc calcul, font=\sffamily\scriptsize, minimum width=21mm]
     (e2) at (3.7,-3.0) {encodeur professeur};

% --- le calcul, vers la droite ---
\draw[fleche] (e1) -- (p1);
\draw[fleche] (p1) -- (l1);
\draw[fleche] (e2) -| (l1);
\node[etiquette] at (1.85,0.42) {le calcul avance};

% --- la correction, vers la gauche, en dessous ---
\draw[fleche accent, dashed] (l1.south) to[bend left=22] (p1.south);
\draw[fleche accent, dashed] (p1.south) to[bend left=22] (e1.south);
\node[annotation, align=center] at (3.7,-1.5)
     {la correction revient par ici};

% --- la barrière ---
\draw[line width=2pt, draw=mbRose!85!black] (6.6,-3.28) -- (6.6,-2.72);
\node[annotation, text=mbRose!85!black, align=center] at (6.6,-3.85)
     {barrière : la correction\\n'a pas le droit de passer};
\end{schema}

\begin{piege}
C'est l'erreur la plus fréquente, et surtout la plus silencieuse, de tout ce
domaine. Si vous oubliez le gradient bloqué, votre programme ne plantera
pas. Il n'affichera aucun message. La courbe d'erreur descendra même
magnifiquement vers zéro, ce qui vous rendra très fier. Et vous aurez
entraîné, pendant des jours, un modèle qui répond la même chose à tout.
Nous verrons au chapitre~25 comment repérer ce désastre à temps.
\end{piege}

Ce premier garde-fou est indispensable, mais il ne suffit pas tout à fait. Il
empêche la cible de reculer, mais il ne dit rien de la manière dont le
professeur, lui, doit évoluer. C'est le rôle du second garde-fou.

\section{Second garde-fou : le professeur qui suit de loin}
\label{sec:ema}

\begin{enclair}
Le professeur n'apprend pas tout seul~: il recopie l'élève. Mais il le
recopie très lentement, en n'attrapant à chaque pas qu'une toute petite
part du chemin. Résultat~: il est toujours en retard sur l'élève, et
c'est exactement ce retard qui empêche les deux de se mettre d'accord sur
une réponse vide, c'est-à-dire une réponse qui ne dit rien de ce que l'on
regarde.
\end{enclair}

La recette porte un nom un peu impressionnant, mais l'opération tient en une
seule ligne de calcul.

\begin{definition}[Moyenne mobile exponentielle (EMA)][ema]
La \emph{moyenne mobile exponentielle} est une façon de mettre à jour une
valeur en gardant beaucoup de l'ancienne et très peu de la nouvelle. On
choisit un nombre $m$ compris entre $0$ et $1$, appelé
\emph{coefficient}, et à chaque pas on remplace l'ancienne valeur par~:

\[
  \bar{\theta} \;\leftarrow\; m \times \bar{\theta} \;+\; (1 - m) \times \theta .
\]

La flèche $\leftarrow$ se lit «~devient~». À gauche du signe~$+$, ce que
l'on garde du professeur. À droite, ce que l'on prend de l'élève. Les trois
lettres EMA sont les premières lettres des trois mots anglais
\emph{Exponential Moving Average}, qui veulent dire exactement cela~:
moyenne mobile exponentielle.
\end{definition}

Ce nom fait peur, alors regardons-le de près~: il contient trois mots, et
chacun dit une chose simple. \emph{Moyenne}, parce que la nouvelle valeur
est un mélange de l'ancienne et de celle de l'élève. \emph{Mobile}, parce
que l'on refait ce mélange à chaque pas, sans jamais s'arrêter. Et
\emph{exponentielle}, parce qu'un vieux réglage ne disparaît pas d'un coup~:
il est multiplié par $m$ à chaque pas, donc il compte de moins en moins,
mais toujours un tout petit peu. Au bout de dix pas, ce qu'il en restait a
été multiplié dix fois de suite par $m$. Voilà tout ce que ce grand mot veut
dire.

Dans V-JEPA~2.1, ce coefficient vaut $m = 0{,}99925$. Traduisons ce nombre
en français courant~: à chaque pas, le professeur garde
$99{,}925\,\%$ de ce qu'il était, et n'attrape que $0{,}075\,\%$ de l'élève.
C'est très, très peu.

\begin{exemple}[Trois pas de moyenne mobile, calculés à la main]
Prenons un seul bouton, pour que tout tienne sur une ligne. Le professeur
est réglé sur $1{,}0$ et l'élève sur $2{,}0$. L'élève ne bouge pas, pour
simplifier. Faisons trois pas.

\textbf{Premier pas}~:
\[
  0{,}99925 \times 1{,}0 \;+\; 0{,}00075 \times 2{,}0
  \;=\; 0{,}99925 + 0{,}0015 \;=\; 1{,}0007500 .
\]

\textbf{Deuxième pas}~: on recommence, mais en partant de $1{,}00075$~:
\[
  0{,}99925 \times 1{,}0007500 \;+\; 0{,}0015 \;=\; 1{,}0014994 .
\]

\textbf{Troisième pas}~:
\[
  0{,}99925 \times 1{,}0014994 \;+\; 0{,}0015 \;=\; 1{,}0022483 .
\]

Regardons maintenant l'écart qui sépare le professeur de l'élève~:

\begin{center}
\footnotesize
\begin{tabular}{lrr}
\toprule
& \textbf{Professeur} & \textbf{Écart avec l'élève} \\
\midrule
Au départ      & $1{,}0000000$ & $1{,}0000000$ \\
Après 1 pas    & $1{,}0007500$ & $0{,}9992500$ \\
Après 2 pas    & $1{,}0014994$ & $0{,}9985006$ \\
Après 3 pas    & $1{,}0022483$ & $0{,}9977517$ \\
\bottomrule
\end{tabular}
\end{center}

En trois pas, l'écart n'a diminué que de deux millièmes. Autrement dit~: le
professeur n'a presque pas bougé. Et ce n'est pas un défaut, c'est
précisément ce que l'on cherche.
\end{exemple}

Pour se faire une idée du rythme, deux repères sont bien utiles.

\textbf{Repère un~: le temps de demi-parcours.} Il faut environ $924$ pas
pour que l'écart entre le professeur et l'élève soit divisé par deux. Dans
notre exemple, le professeur passe alors de $1{,}0$ à $1{,}5$, soit la
moitié du chemin vers $2{,}0$.

\textbf{Repère deux~: la mémoire.} On calcule $1 \div 0{,}00075 \approx
1\,333$. On dit que le professeur «~se souvient~» d'environ $1\,333$ pas
d'entraînement~: sa valeur est une sorte de moyenne des mille trois
cents derniers réglages de l'élève.

\begin{graphe}{Le professeur rattrape l'élève, mais très lentement.
L'élève est fixé à $2$ (trait orange). Le professeur part de $1$~: il lui
faut $924$ pas pour faire la moitié du chemin, et près de cinq mille pas
pour arriver tout près.}[jepa-ema]
\begin{axis}[mbaxis,
    xlabel={pas d'entraînement},
    ylabel={valeur du bouton},
    xmin=0, xmax=5000, ymin=0.9, ymax=2.15,
    legend pos=north west,
]
% repères du demi-parcours, tracés avant les courbes
\addplot[mbSlate!60, dashed, thin, forget plot]
        coordinates {(924,0.9) (924,1.5)};
\addplot[mbSlate!60, dashed, thin, forget plot]
        coordinates {(0,1.5) (924,1.5)};

\addplot[mbBlue, thick, domain=0:5000, samples=200, no marks]
        {2 - exp(-0.00075*x)};
\addlegendentry{le professeur}
\addplot[mbAmber, thick, dashed, domain=0:5000, samples=2, no marks] {2};
\addlegendentry{l'élève}

\node[circle, fill=mbRose, inner sep=1.6pt] at (axis cs:924,1.5) {};
\node[anchor=west, font=\sffamily\scriptsize, text=mbRose!85!black]
     at (axis cs:1100,1.42) {924 pas : la moitié du chemin};
\end{axis}
\end{graphe}

Reste à comprendre \emph{pourquoi} cette lenteur protège de l'effondrement.
L'explication tient en une phrase~: pour que la triche fonctionne, il
faudrait que l'élève et le professeur s'écrasent \emph{ensemble}, au même
moment, sur la même réponse vide. Mais le professeur ne suit qu'avec un
millier de pas de retard. Dès que l'élève se met à donner des réponses toutes pareilles, le
professeur, lui, continue de fournir des cibles variées, héritées d'un passé
où l'élève ne trichait pas encore. L'élève reçoit donc immédiatement une
erreur, et la triche ne lui rapporte plus rien.

\begin{piege}
Le coefficient $m$ se règle avec prudence, car les deux extrêmes sont
mauvais, et de façons opposées.

Si $m$ est trop petit — disons $0{,}5$ — le professeur rattrape l'élève
presque instantanément. Les deux deviennent le même appareil, et l'on
retombe droit dans l'effondrement.

Si $m$ est trop grand — disons $0{,}99999999$ — le professeur ne bouge
pratiquement plus. Il continue à donner les cibles de son premier jour, où
ses boutons étaient encore tirés au hasard. L'élève passe alors son temps à
imiter du bruit.
\end{piege}

\begin{aretenir}
Les deux garde-fous ne font pas la même chose, et il faut les deux. Le
gradient bloqué empêche la machine de \emph{tirer la cible vers elle}. La
moyenne mobile fait avancer la cible \emph{avec du retard}, pour que
personne ne puisse s'entendre sur une réponse vide.
\end{aretenir}

Nos deux garde-fous sont en place. Voyons maintenant à quoi ressemble un
pas d'entraînement complet, du début à la fin.

\section{Un pas d'entraînement, du début à la fin}

Voici l'enchaînement exact, dans l'ordre où vous l'écrirez dans votre
programme. Ces neuf étapes se répètent plus de cent mille fois. Nous les
retrouverons, presque mot pour mot, au chapitre~24, quand nous coderons pour
de bon.

\begin{enumerate}
  \item On prend un lot de vidéos.
  \item On tire un masque au hasard, et l'on fabrique la vue trouée $x$. La
        vue complète $y$, elle, est la vidéo telle quelle.
  \item L'encodeur élève regarde $x$ et produit les jetons de contexte.
  \item On ajoute les jetons de masque, qui portent les adresses des trous.
  \item Le prédicteur produit une réponse pour chaque position.
  \item L'encodeur professeur regarde $y$ et produit les cibles. On pose ici
        le gradient bloqué.
  \item On mesure l'écart entre les réponses et les cibles~: c'est la perte.
  \item La correction remonte, et l'on tourne les boutons de l'élève et du
        prédicteur.
  \item On met à jour le professeur avec la moyenne mobile. Puis on
        recommence au point~1.
\end{enumerate}

Ces neuf étapes forment une boucle~: la neuvième ramène à la première, et
l'on repart pour un tour. La figure~\ref{fig:jepa-boucle} la montre d'un
seul coup d'œil, en regroupant les neuf étapes en quatre moments.

\begin{schema}{La boucle d'entraînement. Les neuf étapes se rangent en
quatre moments, et le quatrième ramène au premier. C'est ce tour complet
que la machine répète plus de cent mille fois.}[jepa-boucle]

\node[bloc, font=\sffamily\scriptsize, minimum width=30mm] (a) at (0,0)
     {\textbf{1--2. Préparer}\\on cache des morceaux\\et l'on obtient $x$ et $y$};
\node[bloc accent, font=\sffamily\scriptsize, minimum width=30mm]
     (b) at (5.4,0)
     {\textbf{3--5. Deviner}\\l'élève regarde, le\\prédicteur répond};
\node[bloc, font=\sffamily\scriptsize, minimum width=30mm, draw=mbAmber,
      fill=mbAmber!8] (c) at (5.4,-2.6)
     {\textbf{6--7. Comparer}\\on mesure l'écart\\avec les cibles};
\node[bloc calcul, font=\sffamily\scriptsize, minimum width=30mm]
     (d) at (0,-2.6)
     {\textbf{8--9. Corriger}\\on tourne les boutons,\\puis le professeur suit};

\draw[fleche] (a) -- (b);
\draw[fleche] (b) -- (c);
\draw[fleche] (c) -- (d);
\draw[fleche accent] (d) -- (a);

\node[annotation, align=center] at (2.7,-1.3)
     {et l'on recommence\\encore et encore};
\end{schema}

\begin{surlesdonnees}
Pour donner l'échelle réelle~: dans V-JEPA~2.1, cette boucle tourne
$135\,000$ fois pendant la phase principale, puis $12\,000$ fois de plus
pendant une phase finale à plus haute résolution. Cela fait $147\,000$
répétitions au total, sur des centaines de machines à la fois. Avec un
coefficient de $0{,}99925$ et un temps de demi-parcours de $924$ pas, l'écart
entre le professeur et l'élève aura eu le temps d'être divisé par deux
environ cent soixante fois, puisque $147\,000 \div 924 \approx 159$. Le
professeur est donc très lent à l'échelle d'un pas, mais pas du tout à
l'échelle de l'entraînement complet.
\end{surlesdonnees}

Assez de théorie. Faisons tourner tout cela.

\section{Passons au code}

Nous allons écrire trois petits programmes. Le premier refait le calcul de
la moyenne mobile, pour vérifier les trois nombres que nous venons de poser
à la main. Le deuxième
monte un entraînement \emph{sans} garde-fou, et nous regardons le modèle
s'effondrer sous nos yeux. Le troisième remet les deux garde-fous, et nous
vérifions que tout tient. Tout cela tourne en quelques secondes sur
n'importe quel ordinateur.

Commençons par la moyenne mobile.

\begin{code}{python}[La moyenne mobile du professeur, pas à pas]
m = 0.99925
professeur, eleve = 1.0, 2.0

for pas in range(1, 4):
    professeur = m * professeur + (1 - m) * eleve
    print(f"pas {pas} : professeur = {professeur:.7f}"
          f"   ecart = {eleve - professeur:.7f}")
\end{code}

À l'exécution, on obtient exactement les nombres du tableau de tout à
l'heure~:

\begin{sortie}
pas 1 : professeur = 1.0007500   ecart = 0.9992500
pas 2 : professeur = 1.0014994   ecart = 0.9985006
pas 3 : professeur = 1.0022483   ecart = 0.9977517
\end{sortie}

Si l'on laisse tourner la boucle plus longtemps, on retrouve aussi nos deux
repères~: au pas $924$, le professeur vaut $1{,}5001$, soit presque exactement
la moitié du chemin~; au pas $5\,000$, il vaut $1{,}9765$.

Passons maintenant à l'effondrement. Pour le voir arriver, nous avons besoin
d'un instrument de mesure, et c'est important~: la perte, à elle seule, ne
nous préviendra de rien, puisqu'elle descendra joliment. Nous allons donc
mesurer autre chose, que nous appellerons la \emph{diversité}~: à quel point
les sorties de l'encodeur diffèrent les unes des autres. Si ce nombre tombe
vers zéro, c'est que toutes les réponses se ressemblent, et donc que le
modèle s'est effondré.

Cette mesure repose sur une petite opération que nous n'avons pas encore
rencontrée, et qu'il faut donc expliquer avant de s'en servir.

\begin{definition}[L'écart-type][ecart-type]
L'\emph{écart-type} d'une liste de nombres dit à quel point ces nombres sont
étalés. On regarde d'abord leur moyenne, puis de combien chacun s'en écarte.
Si tous les nombres de la liste sont égaux, l'écart-type vaut exactement
zéro. Plus ils sont éparpillés, plus il est grand. En Python, on l'obtient
avec \co{python}{.std()}.
\end{definition}

\begin{exemple}[Deux listes, deux écarts-types]
La liste $(5 \,;\, 5 \,;\, 5)$ a pour moyenne $5$, et aucun de ses nombres
ne s'écarte de cette moyenne. Son écart-type vaut donc $0$. C'est exactement
la situation de l'effondrement~: tout le monde répond la même chose.

La liste $(3 \,;\, 5 \,;\, 7)$ a elle aussi pour moyenne $5$, mais ses
nombres s'en écartent de $2$, de $0$ et de $2$. Son écart-type vaut $2$, et
c'est bien ce que rend Python si vous le lui demandez. C'est la situation
d'un modèle en bonne santé~: les réponses sont variées.
\end{exemple}

Notre diversité n'est rien d'autre que cela~: l'écart-type des sorties de
l'encodeur, mesuré d'un exemple à l'autre, puis moyenné pour n'avoir qu'un
seul nombre à surveiller.

\begin{code}{python}[Un entraînement sans le moindre garde-fou]
import torch
import torch.nn as nn

torch.manual_seed(0)

encodeur = nn.Sequential(nn.Linear(8, 16), nn.GELU(), nn.Linear(16, 16))
predicteur = nn.Sequential(nn.Linear(16, 16), nn.GELU(), nn.Linear(16, 16))
optim = torch.optim.Adam(
    list(encodeur.parameters()) + list(predicteur.parameters()), lr=1e-3)

x = torch.randn(64, 8)             # a la place de la vue trouee : du hasard
y = x + 0.1 * torch.randn(64, 8)   # a la place de la vue complete

for pas in range(1, 301):
    cible = encodeur(y)                # AUCUN gradient bloque ici
    prediction = predicteur(encodeur(x))
    perte = (prediction - cible).abs().mean()

    optim.zero_grad()
    perte.backward()
    optim.step()

    if pas % 100 == 0 or pas == 1:
        diversite = cible.std(dim=0).mean().item()
        print(f"pas {pas:4d} : perte = {perte.item():.6f}"
              f"   diversite = {diversite:.6f}")
\end{code}

Voici ce que cela donne~:

\begin{sortie}[Le modèle s'effondre]
pas    1 : perte = 0.273260   diversite = 0.219396
pas  100 : perte = 0.039338   diversite = 0.066842
pas  200 : perte = 0.007948   diversite = 0.017721
pas  300 : perte = 0.002699   diversite = 0.007204
\end{sortie}

Lisons ces quatre lignes ensemble, car elles racontent toute l'histoire de
ce chapitre. La colonne de la perte descend vers zéro~: du point de vue de
la machine, tout va très bien, elle fait exactement ce qu'on lui a demandé.
Mais celle de la diversité descend elle aussi, et c'est la catastrophe~:
partie de
$0{,}22$, la diversité est tombée à $0{,}007$, soit trente fois moins. Les
sorties se sont écrasées les unes sur les autres. En trois cents pas à
peine, sur un ordinateur ordinaire, nous venons de reproduire le désastre
que nous décrivions plus haut.

Remettons maintenant les deux garde-fous. Trois choses changent, et
seulement trois~: un professeur séparé, dont les boutons ne sont jamais
tournés par la correction~; un gradient bloqué autour de la cible~; et la mise à
jour par moyenne mobile à la fin de chaque pas.

\begin{code}{python}[Le même entraînement, avec les deux garde-fous]
eleve = nn.Sequential(nn.Linear(8, 16), nn.GELU(), nn.Linear(16, 16))
predicteur = nn.Sequential(nn.Linear(16, 16), nn.GELU(), nn.Linear(16, 16))

professeur = nn.Sequential(nn.Linear(8, 16), nn.GELU(), nn.Linear(16, 16))
professeur.load_state_dict(eleve.state_dict())   # il part de l'eleve
for p in professeur.parameters():
    p.requires_grad_(False)                      # il n'apprendra jamais seul

optim = torch.optim.Adam(
    list(eleve.parameters()) + list(predicteur.parameters()), lr=1e-3)
m = 0.996                                        # coefficient de la moyenne

x = torch.randn(64, 8)                           # de nouvelles donnees
y = x + 0.1 * torch.randn(64, 8)

for pas in range(1, 301):
    with torch.no_grad():                        # <-- le gradient bloque
        cible = professeur(y)

    prediction = predicteur(eleve(x))
    perte = (prediction - cible).abs().mean()

    optim.zero_grad()
    perte.backward()
    optim.step()

    with torch.no_grad():                        # <-- la moyenne mobile
        for pp, pe in zip(professeur.parameters(), eleve.parameters()):
            pp.mul_(m).add_(pe, alpha=1 - m)

    if pas % 100 == 0 or pas == 1:
        diversite = cible.std(dim=0).mean().item()
        print(f"pas {pas:4d} : perte = {perte.item():.6f}"
              f"   diversite = {diversite:.6f}")
\end{code}

Et le résultat~:

\begin{sortie}[Le modèle tient debout]
pas    1 : perte = 0.247933   diversite = 0.171066
pas  100 : perte = 0.069328   diversite = 0.188252
pas  200 : perte = 0.048201   diversite = 0.232273
pas  300 : perte = 0.038985   diversite = 0.275197
\end{sortie}

Comparons les deux tableaux, colonne par colonne. La perte descend dans les
deux cas — c'est bien la preuve qu'elle ne nous apprend rien à elle seule.
Mais la diversité, elle, fait le contraire~: elle tombait à $0{,}007$ sans
garde-fou, elle \emph{monte} à $0{,}275$ avec. Les représentations, au lieu
de s'écraser, se sont écartées les unes des autres. Le modèle a commencé à
distinguer ses entrées, et c'est exactement le comportement que nous
voulions.

\begin{plusloin}
Vous avez peut-être remarqué que nous avons pris $m = 0{,}996$ dans ce petit
programme, et non $0{,}99925$ comme dans l'article. C'est volontaire, et
c'est une règle générale~: le bon coefficient dépend du nombre de pas que
dure l'entraînement. Ici nous n'en faisons que trois cents, alors que
l'article en fait cent quarante-sept mille. Avec $0{,}99925$, notre
professeur n'aurait pas eu le temps de bouger. Essayez, vous verrez~: la
diversité ne monte alors que de $0{,}171$ à $0{,}189$, au lieu de $0{,}275$.
Elle ne s'effondre pas, mais elle n'avance presque plus, faute de temps.
\end{plusloin}

\begin{aretenir}
Une architecture JEPA tient en cinq pièces et deux garde-fous. Les pièces~:
la vue trouée $x$, la vue complète $y$, l'encodeur élève $E_{\theta}$,
l'encodeur professeur $E_{\bar{\theta}}$ et le prédicteur $P_{\varphi}$. Les
garde-fous~: le gradient bloqué, qui empêche de tirer la cible vers soi, et
la moyenne mobile, qui fait avancer le professeur avec du retard. Et une
seule idée au-dessus de tout cela~: on compare des représentations, jamais
des pixels.
\end{aretenir}

\vspace{0.5em}

Nous savons désormais comment la machine est bâtie, et pourquoi elle ne
s'écroule pas. Il nous reste à écrire, noir sur blanc, ce que
l'entraînement cherche à faire baisser. Cela tient en une seule ligne de
symboles~: c'est la formule officielle de V-JEPA~2. Nous la découperons
symbole par symbole, sans en laisser un seul de côté, au chapitre suivant. Nous y verrons non seulement
\emph{quoi} la machine doit prédire, mais aussi, et c'est le point capital,
\emph{à quels endroits} on lui demande des comptes. Car un détail se cache
là — un tout petit détail, presque invisible — et c'est de ce détail que
naîtra V-JEPA~2.1.

\phantomsection
\section*{Exercices}
\addcontentsline{toc}{section}{Exercices}

\begin{exercice}[Deux pas de moyenne mobile]
Un professeur est réglé sur $3{,}0$, un élève sur $5{,}0$, et le coefficient
vaut $m = 0{,}99$. Calculez la valeur du professeur après un pas, puis après
deux pas. De combien l'écart a-t-il diminué~?
\end{exercice}

\begin{solution}
Premier pas~: $0{,}99 \times 3{,}0 + 0{,}01 \times 5{,}0 = 2{,}97 + 0{,}05
= 3{,}02$.

Deuxième pas~: $0{,}99 \times 3{,}02 + 0{,}05 = 2{,}9898 + 0{,}05 =
3{,}0398$.

L'écart valait $5{,}0 - 3{,}0 = 2{,}0$ au départ. Il vaut maintenant
$5{,}0 - 3{,}0398 = 1{,}9602$. Il a donc diminué de $0{,}0398$, soit environ
deux pour cent en deux pas. Remarquez au passage que $m = 0{,}99$ est
beaucoup plus rapide que le $0{,}99925$ de l'article~: le professeur y
attrape un centième de l'élève à chaque pas, contre moins d'un millième dans
V-JEPA~2.1.
\end{solution}

\begin{exercice}[Compter les jetons du prédicteur]
Un clip contient $4\,096$ jetons. Le masque en recouvre $2\,500$. Combien de
jetons l'encodeur élève reçoit-il~? Combien le prédicteur reçoit-il en
entrée~? Combien en ressort-il~?
\end{exercice}

\begin{solution}
L'encodeur élève ne voit que le contexte, soit $4\,096 - 2\,500 = 1\,596$
jetons.

Le prédicteur reçoit ces $1\,596$ jetons de contexte, plus $2\,500$ jetons
de masque, soit $1\,596 + 2\,500 = 4\,096$ jetons en entrée.

Et il en ressort $4\,096$, puisqu'il produit une réponse par jeton reçu.
Gardez bien en tête que $1\,596$ de ces réponses concernent des endroits qui
étaient déjà visibles~: nous verrons bientôt que V-JEPA~2 les jette, et que
c'est là son point faible.
\end{solution}

\begin{exercice}[Trouver l'erreur]
Un collègue vous montre sa boucle d'entraînement. Elle contient la ligne
suivante~:

\begin{center}
\term{cible = professeur(y)}
\end{center}

Sa courbe d'erreur descend magnifiquement vers zéro en deux heures, et il
est très content. Que lui dites-vous~?
\end{exercice}

\begin{solution}
Qu'il a très probablement oublié le gradient bloqué, et qu'il faut écrire
cette ligne à l'intérieur d'un \term{with torch.no\_grad():}.

Il faut aussi lui expliquer pourquoi sa jolie courbe est un mauvais signe et
non un bon~: une erreur qui descend vers zéro en deux heures, sur un
problème difficile, ne veut pas dire que le modèle a compris quelque chose.
Cela veut dire, presque toujours, qu'il a trouvé une réponse vide.
Conseillez-lui de tracer la diversité de ses représentations à côté de sa
perte~: c'est ce deuxième nombre qui dit la vérité.
\end{solution}

\begin{exercice}[Le compte des nombres]
Un lecteur affirme~: «~JEPA est meilleur parce qu'il a moins de nombres à
deviner qu'une méthode qui redessine les pixels~». A-t-il raison~?
\end{exercice}

\begin{solution}
Non, et c'est le piège classique de ce chapitre. Rien ne garantit qu'une
représentation contienne moins de nombres qu'une zone de pixels~; elle peut
même en contenir beaucoup plus.

Ce qui change, ce n'est pas la quantité, c'est la nature de ce qu'on
demande. Les pixels sont imposés du dehors, avec tout leur hasard~: la
position exacte de chaque brin d'herbe doit être devinée, alors qu'elle est
indevinable. La représentation, elle, est fabriquée par le modèle lui-même,
qui a donc le droit d'y laisser tomber ce qui n'est pas prévisible et de n'y
garder que ce qui compte.
\end{solution}

\begin{exercice}[Régler le coefficient]
Vous choisissez $m = 0{,}5$ pour la moyenne mobile de votre professeur. Que
va-t-il se passer, et pourquoi~?
\end{exercice}

\begin{solution}
Avec $m = 0{,}5$, le professeur attrape la moitié de l'élève à chaque pas.
En cinq ou six pas, il lui est devenu presque identique.

Or, c'est justement le retard du professeur qui empêchait l'effondrement.
Sans ce retard, l'élève et le professeur peuvent s'écraser ensemble sur une
réponse vide, et plus rien ne les en empêche. Vous verrez donc la perte
plonger vers zéro et la diversité des représentations s'effondrer, exactement
comme dans notre deuxième programme.
\end{solution}

\begin{exercice}[À quoi servent les adresses]
Que se passerait-il si les jetons de masque ne portaient pas l'adresse du
trou, mais étaient tous rigoureusement identiques~?
\end{exercice}

\begin{solution}
Le prédicteur recevrait plusieurs enveloppes vides toutes
pareilles, et il n'aurait aucun moyen de savoir laquelle correspond
au coin en haut à gauche et laquelle correspond au bas de l'image.

N'ayant pas de raison de les traiter différemment, il donnerait à peu près
la même réponse pour tous les trous — sans doute une sorte de moyenne de ce
qu'il y a d'habitude dans une vidéo. La perte ne descendrait pas très bas, et
surtout le modèle n'apprendrait rien de la façon dont la scène est rangée
dans l'espace.
L'adresse n'est donc pas un détail de rangement~: c'est elle qui rend la
question possible à poser.
\end{solution}

% \chapter{La perte de V-JEPA~2, expliquée symbole par symbole}
\label{chap:perte-vjepa2}

\begin{quote}\itshape
Nous avons monté la machine au chapitre précédent. Nous savons qui regarde
quoi, qui devine, et qui donne le corrigé. Il reste à écrire, noir sur
blanc, le nombre que l'entraînement fait baisser. Ce nombre s'écrit en une
seule ligne de symboles, et cette ligne fait peur à beaucoup de monde. Nous
allons la découper en morceaux, et donner à chaque symbole sa ligne à lui.
Aucun ne sera laissé de côté. Et tout à la fin, nous découvrirons dans cette
ligne un tout petit détail, presque invisible, qui est le sujet de tout le
reste du livre.
\end{quote}

\begin{prerequis}
Le chapitre~\ref{chap:idee-jepa}, et lui seul. Voici ce qu'il faut en avoir
retenu. Il faut savoir qu'une vidéo est regardée deux fois, une
fois trouée et une fois complète ; que la version trouée s'appelle $x$ et la
version complète $y$ ; que l'encodeur élève $E_{\theta}$ regarde $x$, que
l'encodeur professeur $E_{\bar{\theta}}$ regarde $y$, et que le prédicteur
$P_{\varphi}$ fait le pont entre les deux. Il faut aussi se souvenir du
chapitre~\ref{chap:apprendre} : une perte est un nombre qui dit de combien
la machine se trompe. Les seules opérations employées ici seront
l'addition, la soustraction et la division.
\end{prerequis}

Au chapitre précédent, nous avons dit une phrase un peu vague, et nous
allons enfin la remplacer par quelque chose de précis. Cette phrase, la
voici : « on compare la prédiction à la cible, et l'écart entre les deux est
la perte ».

C'est une bonne phrase, mais elle ne suffit pas pour écrire un programme. Comparer,
comment ? Sur quoi, exactement ? Et pour quels morceaux de la vidéo ? Tant
que ces trois questions restent sans réponse, personne ne peut coder quoi
que ce soit.

La réponse à ces trois questions tient en une seule ligne de symboles. Cette
ligne s'appelle, dans l'article, l'équation numéro un. Nous allons la
construire ensemble, morceau par morceau, en commençant par le début.

\section{Pourquoi on écrit une formule plutôt qu'une phrase}

\begin{enclair}
Une formule n'est rien d'autre qu'une phrase, écrite en abrégé. Chaque
symbole remplace un mot ou un groupe de mots. On l'écrit ainsi pour deux
raisons : c'est plus court, et surtout c'est plus précis. Une phrase peut
souvent se comprendre de deux façons différentes. Une formule, non.
\end{enclair}

Prenons un exemple loin de notre sujet, pour bien voir l'idée. Imaginez que
je vous dise : « pour connaître le prix, on prend le nombre de pommes et on
compte deux euros chacune ». C'est clair, mais c'est long. On peut écrire la
même chose ainsi :

\[
  \text{prix} = 2 \times \text{nombre de pommes} .
\]

Rien n'a changé. On a seulement remplacé des mots par des signes. Personne
ne trouve cette ligne effrayante, et pourtant c'est déjà une formule.

Maintenant, essayons une phrase un peu plus longue : « on regarde de combien
la machine s'est trompée à chaque endroit caché, on additionne toutes ces
erreurs, et on divise par le nombre d'endroits cachés ». Cette phrase-là
tient, elle aussi, en une ligne de symboles. C'est exactement la ligne que
nous allons écrire.

\begin{intuition}
Une formule, c'est une recette de cuisine écrite en abrégé. La recette
dit : « prendre trois œufs, les battre, ajouter cent grammes de sucre ». La
formule dit la même chose avec des signes. Si vous savez lire les signes,
vous savez lire la recette. Et personne n'est né en sachant lire les signes :
cela s'apprend, un signe à la fois. C'est ce que nous allons faire.
\end{intuition}

\begin{piege}
Une formule ne cache jamais d'idée nouvelle. Elle ne fait que résumer des
idées que l'on a déjà expliquées avec des mots. Si une formule vous
paraît incompréhensible, ce n'est presque jamais parce que l'idée est
difficile : c'est parce qu'un des symboles ne vous a pas été présenté. Le
bon réflexe n'est donc pas de relire la formule dix fois, mais de repérer le
symbole que l'on ne connaît pas, et d'aller chercher ce qu'il veut dire.
\end{piege}

Voilà pourquoi ce chapitre est bâti d'une drôle de façon. Nous allons
d'abord raconter la formule en français, tranquillement. Ensuite nous
l'écrirons. Et ensuite seulement nous prendrons chaque symbole, un par un,
pour lui donner une ligne à lui tout seul.

Commençons par la version en français.

\section{La formule, d'abord racontée avec des mots}

\begin{enclair}
La machine a donné une réponse pour chaque endroit caché. Pour chacun de ces
endroits, on regarde de combien elle s'est trompée. On additionne toutes ces
erreurs. Puis on divise par le nombre d'endroits cachés, pour obtenir une
moyenne. Ce nombre-là est la perte.
\end{enclair}

Un mot, avant de commencer, sur ce que veut dire « faire une moyenne ». Vous
le faites déjà sans y penser. Trois enfants ont des billes : le premier en a
$2$, le deuxième $4$, le troisième $6$. On les met toutes en tas, ce qui en
fait $12$, puis on partage ce tas en trois parts égales. Chacun repart avec
$4$ billes. Voilà une moyenne : on met tout ensemble, puis on redistribue en
parts égales.

Pourquoi diviser, ici, plutôt que de garder la simple addition ? Parce que
sans cette division, une vidéo où l'on cache beaucoup de morceaux donnerait
toujours une erreur plus grosse qu'une vidéo où l'on en cache peu — même si
la machine travaillait aussi bien dans les deux cas. La division remet tout
le monde sur un pied d'égalité.

Ce point étant réglé, reprenons la machine du chapitre~\ref{chap:idee-jepa}
et suivons-la du doigt. Vous vous souvenez peut-être de la boucle en neuf
étapes que nous avions dressée à la fin de ce chapitre-là. Tout ce que nous
allons faire ici, c'est ouvrir la septième de ces étapes — celle qui disait
simplement « on mesure l'écart entre les réponses et les cibles » — et
regarder ce qu'elle contient vraiment. Voici donc ce qui se passe, dans
l'ordre.

D'abord, la vue trouée $x$ entre dans l'encodeur élève. Celui-ci ne voit que
les morceaux visibles, et il en ressort une version enrichie.

Ensuite, on ajoute les jetons de masque, ces enveloppes vides qui portent
l'adresse de chaque trou. Le prédicteur reçoit le tout et produit une
réponse pour chaque position.

Pendant ce temps, la vue complète $y$ est passée dans l'encodeur professeur,
qui produit les cibles, c'est-à-dire les bonnes réponses.

Enfin — et c'est le moment qui nous occupe dans ce chapitre — on prend
chaque endroit caché, on met côte à côte la réponse de la machine et la
bonne réponse, et on mesure l'écart. Puis on fait la moyenne de tous ces
écarts.

\begin{exemple}[La perte racontée sur trois trous]
Supposons, pour l'instant, qu'il n'y ait que trois endroits cachés dans
notre vidéo, et que l'on sache déjà mesurer l'écart à chacun d'eux. Disons
que ces trois écarts valent $1{,}5$, $0{,}9$ et $0{,}6$.

On les additionne : $1{,}5 + 0{,}9 + 0{,}6 = 3{,}0$. Puis on divise par
trois, puisqu'il y a trois endroits :

\[
  \frac{3{,}0}{3} = 1{,}0 .
\]

La perte vaut donc $1{,}0$. C'est tout. Le reste du chapitre ne fera
qu'expliquer comment on obtient ces trois écarts, et comment on écrit ce
calcul en abrégé.
\end{exemple}

\begin{aretenir}
Retenez la phrase, elle vous accompagnera jusqu'à la fin du livre :
\textbf{la perte est la moyenne des écarts, mesurés aux endroits cachés}.
Tout le reste n'est que de l'écriture.
\end{aretenir}

Nous avons la phrase. Écrivons-la maintenant en abrégé.

\section{La formule écrite, et son tableau de décodage}

Voici la ligne de symboles. Ne cherchez pas à la comprendre tout de suite :
regardez-la comme on regarde une photo de groupe avant de reconnaître les
visages un par un.

\begin{equation}
  \mathcal{L}_{\text{predict}}
  \;=\;
  \frac{1}{|M|}
  \sum_{i \in M}
  \Bigl\|
    \; P_{\varphi}\bigl(E_{\theta}(x),\, \Delta_y\bigr)_i
    \;-\;
    \mathrm{sg}\bigl(E_{\bar{\theta}}(y)_i\bigr) \;
  \Bigr\|_{1}
  \label{eq:perte-vjepa2}
\end{equation}

Voilà. C'est la formule complète de V-JEPA~2, celle que l'article appelle
son équation numéro un, recopiée sans rien enlever.

Elle contient exactement treize symboles différents. Pas un de plus. Et vous
en connaissez déjà la plupart, puisque nous les avons rencontrés dans les
chapitres précédents. Voici le tableau complet, sans en laisser un seul de
côté : gardez-le sous les yeux, nous allons le parcourir ligne par ligne.

\begin{center}
\footnotesize
\begin{tabular}{llp{6.1cm}}
\toprule
\textbf{Symbole} & \textbf{Se lit} & \textbf{Veut dire} \\
\midrule
$\mathcal{L}_{\text{predict}}$ & « L de prédiction » &
le nom du nombre que l'on calcule : c'est la perte. \\[3pt]

$\dfrac{1}{|M|}$ & « un sur taille de M » &
on divise par le nombre d'endroits cachés : c'est ce qui fait la
moyenne. \\[3pt]

$\displaystyle\sum_{i \in M}$ & « somme, pour $i$ dans M » &
« refais le calcul pour chaque endroit caché, et additionne tout ». \\[3pt]

$x$ & « iks » &
la vue trouée : la vidéo dont on a retiré des morceaux. \\[3pt]

$y$ & « igrec » &
la vue complète : la même vidéo, entière. \\[3pt]

$E_{\theta}$ & « E de thêta » &
l'encodeur élève, celui qui apprend. Ses boutons s'appellent $\theta$. \\[3pt]

$E_{\bar{\theta}}$ & « E de thêta barre » &
l'encodeur professeur, la copie lente de l'élève. \\[3pt]

$P_{\varphi}$ & « P de phi » &
le prédicteur, la pièce qui devine. Ses boutons s'appellent $\varphi$. \\[3pt]

$\Delta_y$ & « delta y » &
les jetons de masque : des enveloppes vides qui portent l'adresse des
trous. \\[3pt]

$\mathrm{sg}$ & « esse-gé » &
le gradient bloqué : « ne fais remonter aucune correction par ce
chemin ». \\[3pt]

$M$ & « M » &
la liste des endroits cachés. \\[3pt]

$|M|$ & « taille de M » &
combien il y a d'endroits cachés dans cette liste. \\[3pt]

$\|\,\cdot\,\|_{1}$ & « norme un » &
la façon de mesurer un écart : on additionne les différences sans se
soucier du signe. \\
\bottomrule
\end{tabular}
\end{center}

Il reste un tout petit signe qui ne figure pas dans ce tableau, parce qu'il
n'a pas de nom à lui : le minuscule $i$ posé en bas à droite de
$P_{\varphi}(\ldots)_i$ et de $E_{\bar{\theta}}(y)_i$. Nous lui donnerons
aussi son explication, un peu plus loin. Rien, vraiment rien, ne sera passé
sous silence.

Avant d'aller plus loin, arrêtons-nous une seconde sur le nom écrit à
gauche du signe égal, car il annonce déjà la suite du livre.

\begin{notation}[Le nom de la perte, $\mathcal{L}_{\text{predict}}$][nom-perte]
Nous avions vu au chapitre~\ref{chap:apprendre} que l'on note souvent une
perte par la lettre $\mathcal{L}$, d'après le mot anglais \emph{loss} qui
veut dire « perte ». Ici, cette lettre porte en plus une petite étiquette
écrite en bas : \emph{predict}, c'est-à-dire « prédiction ».

Pourquoi préciser ? Parce que cette perte-ci mesure une chose bien
particulière : ce que la machine a répondu \emph{aux endroits qu'elle devait
prédire}, c'est-à-dire aux endroits cachés. Écrire l'étiquette, c'est déjà
reconnaître qu'il pourrait en exister d'autres, mesurant autre chose.

Dans V-JEPA~2, il n'y en a pas d'autre : cette perte est la seule. Mais
gardez cette étiquette en mémoire, car au chapitre~10 nous en ajouterons une
seconde, et il faudra bien pouvoir les distinguer.
\end{notation}

\begin{exemple}[Une étiquette, c'est juste un prénom]
Dans une famille où tout le monde s'appelle Martin, on ajoute le prénom pour
savoir de qui l'on parle : Martin Paul, Martin Sophie. Le nom de famille ne
change pas, seul le prénom distingue.

C'est exactement ce qui se passe ici. Le nom de famille est $\mathcal{L}$,
et il veut dire « perte ». L'étiquette $\text{predict}$ est le prénom, et il
dit \emph{quelle} perte. Écrire
\[
  \mathcal{L}_{\text{predict}} = 0{,}5
\]
se lit donc : « la perte des endroits à prédire vaut $0{,}5$ ». Nous
calculerons ce nombre-là dans quelques pages.
\end{exemple}

Les trois dernières lignes du tableau sont les seules vraiment nouvelles :
la norme L1, l'ensemble $M$ et sa taille $|M|$. Nous leur consacrerons les
deux sections qui suivent. Mais avant, faisons un premier tour de la formule
à grande vitesse, pour voir sa forme d'ensemble.

\begin{enclair}
Lisez la formule de la droite vers la gauche, comme on démonte une machine
en commençant par l'extérieur.

Tout à droite, entre les deux grandes barres, il y a une soustraction : la
réponse de la machine moins la bonne réponse. C'est l'écart.

Les deux grandes barres autour disent comment on transforme cet écart en un
seul nombre positif.

Le grand signe $\sum$, à gauche, dit qu'on additionne ce nombre pour tous
les endroits cachés.

Et la fraction $\frac{1}{|M|}$, tout à gauche, dit qu'on divise le total par
le nombre d'endroits, pour obtenir une moyenne.

Voilà. La formule dit exactement la même chose que notre phrase en français.
\end{enclair}

La figure~\ref{fig:formule-decoupee} montre ce découpage en quatre morceaux.

\begin{schema}{La formule découpée en quatre morceaux. Les pastilles
numérotées donnent l'ordre dans lequel il faut la lire : on commence par la
droite, avec la soustraction, et l'on remonte vers la gauche jusqu'à la
moyenne. La légende du dessous reprend les quatre numéros et dit, pour
chacun, ce que ce morceau fabrique.}[formule-decoupee]

% ================== la formule, morceau par morceau ==================
% On place d'abord l'écart, qui est le plus large, puis on accroche les
% autres morceaux à sa gauche et à sa droite.
\node[font=\small, inner sep=2pt] (ecart) at (2.6,0)
     {$P_{\varphi}\bigl(E_{\theta}(x), \Delta_y\bigr)_i
       \;-\; \mathrm{sg}\bigl(E_{\bar{\theta}}(y)_i\bigr)$};
\node[font=\normalsize, inner sep=2pt, anchor=east]
      (bg) at ($(ecart.west)+(-0.18,0)$) {$\Bigl\|$};
\node[font=\normalsize, inner sep=2pt, anchor=west]
      (bd) at ($(ecart.east)+(0.18,0)$) {$\Bigr\|_{1}$};
\node[font=\normalsize, inner sep=2pt, anchor=east]
      (som) at ($(bg.west)+(-0.22,0)$) {$\displaystyle\sum_{i \in M}$};
\node[font=\normalsize, inner sep=2pt, anchor=east]
      (moy) at ($(som.west)+(-0.22,0)$) {$\dfrac{1}{|M|}$};

% ---- les fonds colorés, posés derrière ----
\begin{scope}[on background layer]
  \node[fit=(moy), inner sep=3pt, rounded corners=2pt,
        fill=mbBlue!12, draw=mbBlue!45, line width=0.5pt] (fmoy) {};
  \node[fit=(som), inner sep=3pt, rounded corners=2pt,
        fill=mbTeal!12, draw=mbTeal!45, line width=0.5pt] (fsom) {};
  \node[fit=(bg)(bd), inner sep=3pt, rounded corners=2pt,
        fill=mbAmber!10, draw=mbAmber!55, line width=0.5pt] (fbar) {};
  \node[fit=(ecart), inner sep=2pt, rounded corners=2pt,
        fill=mbViolet!10, draw=mbViolet!45, line width=0.5pt, dashed] (fec) {};
\end{scope}

% ---- les pastilles numérotées, au-dessus de chaque morceau ----
\newcommand{\mbPastille}[3]{%
  \node[circle, fill=#2, text=white, font=\sffamily\bfseries\scriptsize,
        inner sep=1.6pt] at (#1) {#3};}
\mbPastille{moy.north |- 0,0.72}{mbBlue}{4}
\mbPastille{som.north |- 0,0.72}{mbTeal}{3}
\mbPastille{bd.north |- 0,0.72}{mbAmber!85!black}{2}
\mbPastille{ecart.north |- 0,0.72}{mbViolet}{1}

% ---- le sens de lecture ----
\draw[-{Stealth[round,length=2.4mm,width=1.8mm]}, draw=mbSlate!70,
      line width=0.7pt] (5.35,1.15) -- (-1.55,1.15);
\node[font=\sffamily\scriptsize\itshape, text=mbSlate] at (1.9,1.42)
     {on lit la formule dans ce sens};

% ================== la légende ==================
\node[circle, fill=mbViolet, text=white, font=\sffamily\bfseries\scriptsize,
      inner sep=1.6pt] at (-1.55,-1.30) {1};
\node[font=\sffamily\scriptsize, text=mbInk, anchor=west] at (-1.25,-1.30)
     {\textbf{la soustraction} — ce que la machine a répondu, moins ce
      qu'il fallait répondre};

\node[circle, fill=mbAmber!85!black, text=white,
      font=\sffamily\bfseries\scriptsize, inner sep=1.6pt] at (-1.55,-1.85) {2};
\node[font=\sffamily\scriptsize, text=mbInk, anchor=west] at (-1.25,-1.85)
     {\textbf{les deux barres} — on efface les signes et on additionne :
      l'écart devient un seul nombre positif};

\node[circle, fill=mbTeal, text=white, font=\sffamily\bfseries\scriptsize,
      inner sep=1.6pt] at (-1.55,-2.40) {3};
\node[font=\sffamily\scriptsize, text=mbInk, anchor=west] at (-1.25,-2.40)
     {\textbf{le signe somme} — on recommence pour chacun des endroits
      cachés, et on additionne tout};

\node[circle, fill=mbBlue, text=white, font=\sffamily\bfseries\scriptsize,
      inner sep=1.6pt] at (-1.55,-2.95) {4};
\node[font=\sffamily\scriptsize, text=mbInk, anchor=west] at (-1.25,-2.95)
     {\textbf{la fraction} — on divise par le nombre d'endroits cachés :
      c'est la moyenne};
\end{schema}

Nous avons vu la forme d'ensemble. Approchons-nous maintenant, et regardons
de tout près la pièce qui se trouve au milieu de la formule : les deux
grandes barres.

\section{Zoom sur les deux grandes barres : la norme L1}

\begin{enclair}
Les deux grandes barres sont une machine à transformer un écart en un seul
nombre positif. Elle travaille en deux temps : on efface les signes moins, et
on additionne. C'est tout.
\end{enclair}

Avant tout le reste, il y a un petit problème à régler. Ce que le
prédicteur produit n'est pas un nombre :
c'est un \emph{jeton}, c'est-à-dire une liste de nombres, comme nous l'avons
vu au chapitre~\ref{chap:vit}. La cible n'est pas un nombre non plus :
c'est aussi une liste de nombres.

Or une perte doit être \emph{un seul nombre}. On ne peut pas descendre une
vallée dans plusieurs directions à la fois. Il faut donc une machine qui
prenne deux listes et rende un seul nombre : c'est exactement le rôle des
deux grandes barres.

Cette machine fait deux choses, l'une après l'autre. Voyons la première.

\subsection*{Premier temps : effacer les signes moins}

Quand on soustrait deux listes, on obtient une liste d'écarts. Certains sont
positifs, d'autres négatifs.

\begin{exemple}[Soustraire deux listes]
Le prédicteur a répondu $(0{,}5 \ ; \ -0{,}2 \ ; \ 1{,}0)$ et la bonne
réponse était $(0{,}4 \ ; \ 0{,}1 \ ; \ 0{,}8)$. On soustrait terme à
terme, c'est-à-dire le premier avec le premier, le deuxième avec le
deuxième, et ainsi de suite :

\[
  0{,}5 - 0{,}4 = 0{,}1 , \qquad
  -0{,}2 - 0{,}1 = -0{,}3 , \qquad
  1{,}0 - 0{,}8 = 0{,}2 .
\]

Les trois écarts valent donc $(0{,}1 \ ; \ -0{,}3 \ ; \ 0{,}2)$. Remarquez
le signe moins du deuxième : sur ce nombre-là, la machine a répondu trop
\emph{bas}.
\end{exemple}

Nous avons déjà rencontré ce problème au chapitre~\ref{chap:apprendre}, avec
la machine à bonbons : si l'on additionne les écarts tels quels, un $+5$ et
un $-5$ s'annulent, et l'on croit à tort que la machine est parfaite. Il
faut donc effacer les signes moins avant d'additionner.

Là-bas, nous avions choisi de mettre chaque écart au carré, c'est-à-dire de
le multiplier par lui-même. Ici, V-JEPA~2 fait autrement, et plus simplement
encore : il prend la \emph{valeur absolue}.

\begin{definition}[La valeur absolue][valeur-absolue]
La \emph{valeur absolue} d'un nombre, c'est ce même nombre, débarrassé de
son signe moins s'il en avait un. On l'écrit avec deux petites barres
verticales autour du nombre. Ainsi $|3| = 3$, et $|-3| = 3$ aussi. On peut
la voir comme une distance : elle dit à quelle distance de zéro se trouve le
nombre, sans dire de quel côté.
\end{definition}

\begin{exemple}[Quelques valeurs absolues]
Rien de plus simple : on regarde le nombre, et s'il porte un signe moins, on
le raye.

\begin{center}
\footnotesize
\begin{tabular}{cc}
\toprule
\textbf{Le nombre} & \textbf{Sa valeur absolue} \\
\midrule
$7$        & $7$ \\
$-7$       & $7$ \\
$0{,}3$    & $0{,}3$ \\
$-0{,}3$   & $0{,}3$ \\
$0$        & $0$ \\
\bottomrule
\end{tabular}
\end{center}

Remarquez les deux premières lignes : $7$ et $-7$ donnent la même réponse.
C'est exactement l'effet que nous cherchions, puisque nous voulions que se
tromper de sept en trop coûte autant que se tromper de sept en moins.
\end{exemple}

\begin{intuition}
Vous habitez au kilomètre zéro d'une route. Un ami habite trois kilomètres
au nord, un autre trois kilomètres au sud. Si vous demandez « à quelle
distance habites-tu ? », les deux répondent « trois kilomètres ». Aucun ne
répond « moins trois ». La valeur absolue, c'est exactement cette
réponse-là : la distance, sans la direction.
\end{intuition}

Les signes moins ont disparu. Il ne reste plus qu'un geste, et c'est le plus
simple des deux.

\subsection*{Second temps : additionner}

Une fois les signes effacés, il ne reste plus qu'à additionner tous les
écarts de la liste. On obtient enfin un seul nombre, et il est forcément
positif.

Cette machine complète — effacer les signes, puis additionner — porte un
nom.

\begin{definition}[La norme L1][norme-l1]
La \emph{norme L1} d'une liste de nombres est la somme des valeurs absolues
de ces nombres. On l'écrit en entourant la liste de deux grandes barres, et
en écrivant un petit $1$ en bas à droite : $\|\,\cdot\,\|_{1}$. On la
prononce « norme un ». Le mot \emph{norme} veut simplement dire « façon de
mesurer la taille de quelque chose ».
\end{definition}

\begin{exemple}[Calculer une norme L1, en entier]
Reprenons nos trois écarts, $(0{,}1 \ ; \ -0{,}3 \ ; \ 0{,}2)$.

On efface les signes : $(0{,}1 \ ; \ 0{,}3 \ ; \ 0{,}2)$.

On additionne : $0{,}1 + 0{,}3 + 0{,}2 = 0{,}6$.

On écrit donc :
\[
  \bigl\| (0{,}1 \ ; \ -0{,}3 \ ; \ 0{,}2) \bigr\|_{1} = 0{,}6 .
\]

Trois nombres sont devenus un seul, et ce seul nombre dit de combien la
machine s'est trompée sur ce jeton-là. Retenez ce $0{,}6$ : il nous servira
tout à l'heure.
\end{exemple}

La figure~\ref{fig:norme-l1} montre les deux temps de l'opération.

\begin{schema}{La norme L1, en deux temps. À gauche, les trois écarts tels
qu'ils sortent de la soustraction : la longueur de chaque barre donne sa
taille, et celle du milieu est rouge parce qu'elle est négative. Au milieu,
les signes ont été effacés : les trois barres sont devenues des longueurs,
toutes de la même couleur. À droite, on les met bout à bout, et la longueur
obtenue est la norme L1.}[norme-l1]

% Échelle : 1 unité de valeur = 4 cm. Ainsi 0,1 -> 0,4 et 0,3 -> 1,2.
% Les trois barres sont empilées, ce qui permet de comparer leurs longueurs
% d'un coup d'œil.

% ================== 1. les écarts, avec leurs signes ==================
\node[font=\sffamily\scriptsize\bfseries, text=mbSlate] at (0.75,1.55)
     {les écarts};
\node[font=\small, text=mbInk, anchor=east] at (-0.15,1.0)  {$0{,}1$};
\node[font=\small, text=mbRose!85!black, anchor=east] at (-0.15,0.5) {$-0{,}3$};
\node[font=\small, text=mbInk, anchor=east] at (-0.15,0.0)  {$0{,}2$};

\fill[mbAzure!30, draw=mbAzure!75, line width=0.6pt]
      (0,0.85) rectangle ++(0.4,0.3);
\fill[mbRose!30, draw=mbRose!80, line width=0.6pt]
      (0,0.35) rectangle ++(1.2,0.3);
\fill[mbAzure!30, draw=mbAzure!75, line width=0.6pt]
      (0,-0.15) rectangle ++(0.8,0.3);

\node[font=\sffamily\scriptsize, text=mbRose!85!black, anchor=north west]
      at (0,-0.42) {celui-ci est négatif};

% ================== la première flèche ==================
\draw[fleche accent] (2.15,0.5) -- (3.55,0.5);
\node[font=\sffamily\scriptsize, text=mbAmber!80!black, align=center]
      at (2.85,0.95) {on efface\\les signes};

% ================== 2. les longueurs, toutes positives ==================
\begin{scope}[shift={(4.35cm,0cm)}]
  \node[font=\sffamily\scriptsize\bfseries, text=mbSlate] at (0.75,1.55)
       {des longueurs};
  \node[font=\small, text=mbInk, anchor=east] at (-0.15,1.0) {$0{,}1$};
  \node[font=\small, text=mbInk, anchor=east] at (-0.15,0.5) {$0{,}3$};
  \node[font=\small, text=mbInk, anchor=east] at (-0.15,0.0) {$0{,}2$};

  \fill[mbAzure!30, draw=mbAzure!75, line width=0.6pt]
        (0,0.85) rectangle ++(0.4,0.3);
  \fill[mbAzure!30, draw=mbAzure!75, line width=0.6pt]
        (0,0.35) rectangle ++(1.2,0.3);
  \fill[mbAzure!30, draw=mbAzure!75, line width=0.6pt]
        (0,-0.15) rectangle ++(0.8,0.3);

  \node[font=\sffamily\scriptsize, text=mbSlate, anchor=north west]
        at (0,-0.42) {plus aucun signe moins};
\end{scope}

% ================== la seconde flèche ==================
\draw[fleche accent] (6.50,0.5) -- (7.90,0.5);
\node[font=\sffamily\scriptsize, text=mbAmber!80!black, align=center]
      at (7.20,0.95) {on les met\\bout à bout};

% ================== 3. la somme ==================
\begin{scope}[shift={(8.70cm,0cm)}]
  \node[font=\sffamily\scriptsize\bfseries, text=mbSlate] at (1.2,1.55)
       {une seule longueur};

  \fill[mbAzure!30, draw=mbAzure!75, line width=0.6pt]
        (0,0.35) rectangle ++(0.4,0.3);
  \fill[mbAzure!45, draw=mbAzure!75, line width=0.6pt]
        (0.4,0.35) rectangle ++(1.2,0.3);
  \fill[mbAzure!30, draw=mbAzure!75, line width=0.6pt]
        (1.6,0.35) rectangle ++(0.8,0.3);

  \draw[accolade bas] (0,0.28) -- (2.4,0.28);
  \node[font=\sffamily\small\bfseries, text=mbAzure!70!black] at (1.2,-0.28)
       {en tout : $0{,}6$};
  \node[font=\sffamily\scriptsize, text=mbSlate] at (1.2,-0.75)
       {c'est la norme L1};
\end{scope}
\end{schema}

\begin{plusloin}
Pourquoi la valeur absolue plutôt que le carré du
chapitre~\ref{chap:apprendre} ? Parce que les deux ne punissent pas de la
même façon.

Prenons deux écarts, un tout petit et un gros : $0{,}1$ et $1{,}0$.

Avec la norme L1, le total vaut $0{,}1 + 1{,}0 = 1{,}1$. Le gros écart
compte pour $1{,}0 / 1{,}1 \approx 91\,\%$ du total.

Avec les carrés, le total vaut $0{,}01 + 1{,}00 = 1{,}01$. Le gros écart
compte cette fois pour $1{,}00 / 1{,}01 \approx 99\,\%$ du total, et le petit
pour à peine $1\,\%$.

Le carré écrase donc les petits écarts et ne s'occupe presque plus que des
gros. Or, dans une vidéo, les gros écarts arrivent souvent aux endroits
imprévisibles : une porte qui claque, quelqu'un qui entre dans le champ.
Avec le carré, la machine passerait son temps sur ces surprises qu'elle ne
pouvait pas deviner. Avec la valeur absolue, elle continue d'écouter les
petits écarts, ceux qui viennent des endroits qu'elle aurait vraiment pu
mieux comprendre.
\end{plusloin}

Voilà les deux grandes barres expliquées. Il nous reste la lettre $M$ et sa
taille $|M|$ : c'est encore plus simple, et c'est le sujet de la section qui
vient.

\section{Zoom sur la lettre M : la liste des endroits cachés}

\begin{enclair}
La lettre $M$ est la liste des adresses de tous les trous. Ce n'est pas un
nombre, c'est une liste. Et $|M|$, avec les deux barres, est la longueur de
cette liste, c'est-à-dire combien il y a de trous.
\end{enclair}

Reprenons notre clip de la phase principale, celui qui nous accompagne
depuis le chapitre~\ref{chap:vit}. Il contient $2\,048$ jetons. On leur a
donné un numéro, de $1$ jusqu'à $2\,048$, pour pouvoir en parler.

Maintenant, on tire un masque au hasard. Certains jetons sont recouverts,
d'autres restent visibles. Il faut bien noter quelque part lesquels sont
recouverts. C'est ce que fait la lettre $M$.

\begin{definition}[Un ensemble][ensemble]
Un \emph{ensemble} est simplement une liste d'objets. Chaque objet n'y
figure qu'une seule fois, et l'ordre dans lequel on les écrit n'a aucune
importance. On note un ensemble entre accolades. Par exemple,
$\{3 \ ; \ 17 \ ; \ 42\}$ est l'ensemble qui contient les trois nombres $3$,
$17$ et $42$. Dire qu'un objet « appartient » à l'ensemble, c'est dire qu'il
est dans la liste, et on l'écrit avec le petit signe $\in$, qui se lit
« appartient à ».
\end{definition}

\begin{definition}[L'ensemble des endroits cachés, noté $M$][ensemble-m]
On note $M$ l'ensemble des numéros des jetons masqués. La lettre $M$ vient
de l'anglais \emph{masked}, qui veut dire « masqué ». Écrire $i \in M$ se lit
« le jeton numéro $i$ fait partie des jetons masqués ».
\end{definition}

\begin{definition}[La taille d'un ensemble, notée $|M|$][taille-ensemble]
Les deux barres autour d'un ensemble donnent le nombre d'objets qu'il
contient. Ainsi, si $M = \{3 \ ; \ 17 \ ; \ 42\}$, alors $|M| = 3$.
Attention : ce sont les mêmes barres que celles de la valeur absolue, mais
elles ne font pas la même chose. Autour d'un nombre, elles effacent le
signe ; autour d'un ensemble, elles comptent combien il y a d'objets dedans.
C'est le contexte
qui dit laquelle des deux on emploie.
\end{definition}

\begin{exemple}[Un masque minuscule, écrit en entier]
Prenons une toute petite vidéo qui ne contiendrait que huit jetons,
numérotés de $1$ à $8$. On tire un masque, et il recouvre les jetons
numéro $3$, $4$, $5$ et $7$. On écrit alors :

\[
  M = \{3 \ ; \ 4 \ ; \ 5 \ ; \ 7\} , \qquad |M| = 4 .
\]

Les quatre jetons restants, ceux numéro $1$, $2$, $6$ et $8$, forment le
contexte. Nous leur donnerons leur propre lettre au chapitre~10 ; pour
l'instant, ils n'ont pas de nom dans la formule, et c'est justement là que
tout se joue.
\end{exemple}

\begin{exemple}[Le masque de notre clip habituel]
Reprenons les chiffres du chapitre~\ref{chap:idee-jepa}. Notre clip contient
$2\,048$ jetons, et le masque en recouvre $1\,200$. On a donc :

\[
  |M| = 1\,200 .
\]

L'ensemble $M$ contient donc mille deux cents numéros. On ne les écrit
évidemment jamais à la main : c'est le programme qui les garde en mémoire.
Les $848$ jetons restants sont le contexte.
\end{exemple}

La figure~\ref{fig:ensemble-m} met cette petite vidéo sous vos yeux, pour
que vous voyiez d'un coup ce que la lettre $M$ contient et ce qu'elle laisse
de côté.

\begin{schema}{Ce que contient la lettre $M$. Notre petite vidéo n'a que
huit jetons, numérotés de $1$ à $8$. Quatre d'entre eux sont cachés : ce sont
les gris, et leurs quatre numéros forment l'ensemble $M$. Les quatre autres,
les bleus, restent visibles : ce sont les jetons de contexte, et aucun de
leurs numéros n'entre dans $M$.}[ensemble-m]

% ---- la file des huit jetons ----
\foreach \k in {1,...,8}{
  \node[font=\sffamily\scriptsize, text=mbSlate] at ({(\k-1)*1.15+0.4},1.05)
       {\k};
}
\foreach \k in {1,2,6,8}{
  \fill[mbAzure!25, draw=mbAzure!70, line width=0.7pt]
        ({(\k-1)*1.15},0) rectangle ++(0.8,0.8);
}
\foreach \k in {3,4,5,7}{
  \fill[mbSlate!35, draw=mbSlate!80, line width=0.7pt, dashed]
        ({(\k-1)*1.15},0) rectangle ++(0.8,0.8);
}

\node[font=\sffamily\scriptsize, text=mbSlate, anchor=east] at (-0.25,0.4)
     {les huit jetons};

% ---- on descend les numeros gris dans l'ensemble ----
\foreach \k in {3,4,5,7}{
  \draw[trait annotation] ({(\k-1)*1.15+0.4},-0.12) -- (4.0,-0.95);
}
\node[font=\small, text=mbInk] at (4.0,-1.35)
     {$M = \{\,3 \ ; \ 4 \ ; \ 5 \ ; \ 7\,\}$};
\node[font=\small, text=mbInk] at (4.0,-2.0)
     {$|M| = 4$};
\node[annotation, anchor=west, align=left] at (6.4,-1.35)
     {les numéros\\des cases cachées};
\node[annotation, anchor=west, align=left] at (6.4,-2.0)
     {combien il y en a};

\node[font=\sffamily\scriptsize, text=mbAzure!80!black, anchor=west]
     at (-0.25,-2.75)
     {les jetons $1$, $2$, $6$ et $8$ restent dehors : ils sont le contexte};
\end{schema}

Il nous reste à comprendre ce que fait le grand signe $\sum$, celui qui est
posé juste devant $M$ dans la formule.

\begin{notation}[Le signe somme, revu][somme-revue]
Nous avions déjà croisé ce signe au chapitre~\ref{chap:apprendre}. Écrire
$\sum_{i \in M}$ veut dire : « fais prendre à $i$ chacune des valeurs de
l'ensemble $M$, calcule à chaque fois ce qui suit, et additionne tous les
résultats ».

Avec notre masque minuscule $M = \{3 \ ; \ 4 \ ; \ 5 \ ; \ 7\}$, cela veut
donc dire : calcule pour le jeton~$3$, puis pour le jeton~$4$, puis pour le
jeton~$5$, puis pour le jeton~$7$, et additionne les quatre résultats.
Et surtout : \textbf{ne calcule rien du tout pour les jetons $1$, $2$, $6$
et $8$}. Ils ne sont pas dans la liste, donc on ne les regarde pas.
\end{notation}

Cette dernière phrase, en gras, a l'air d'un détail d'écriture. Retenez-la :
c'est elle qui va tout faire basculer dans quelques pages.

\begin{notation}[Le petit $i$ en bas à droite][indice-i]
Vous avez sans doute remarqué le petit $i$ posé en bas à droite, dans
$P_{\varphi}(\ldots)_i$ et dans $E_{\bar{\theta}}(y)_i$. Il veut dire
« prends seulement le jeton numéro $i$ dans tout ce qui sort ».

Souvenez-vous du chapitre~\ref{chap:vit} : l'encodeur ne rend pas un seul
paquet de nombres, mais un paquet \emph{par jeton}. Le prédicteur non plus.
Le petit $i$ est donc l'instruction qui dit : « dans ces deux mille
quarante-huit paquets, va me chercher celui qui porte le numéro $i$ ».
\end{notation}

\begin{exemple}[Le petit $i$ au travail]
C'est le même geste que lorsque vous cherchez un appartement dans un
immeuble. Le prédicteur, c'est l'immeuble entier : il contient $2\,048$
appartements, un par jeton. L'écriture
$P_{\varphi}\bigl(E_{\theta}(x), \Delta_y\bigr)_{17}$ veut simplement dire :
« va sonner à la porte numéro $17$ ».

Si l'on remplace $17$ par $500$, on obtient un autre paquet de nombres,
celui d'un tout autre endroit de la vidéo. Le calcul qui a fabriqué
l'immeuble est le même ; seul le numéro de la porte a changé.
\end{exemple}

Nous avons maintenant tous les symboles. Mettons-les au travail sur de vrais
nombres.

\section{Le calcul complet, du début à la fin}
\label{sec:calcul-complet}

\begin{enclair}
Nous allons calculer une perte entière, sans rien sauter. Pour que tout
tienne sur une ligne, nos jetons n'auront que trois nombres au lieu de mille,
et notre masque ne recouvrira que deux endroits au lieu de mille deux cents.
L'idée est exactement la même ; seules les tailles changent.
\end{enclair}

Voici la situation. Notre masque recouvre deux jetons, disons ceux qui
portent les numéros $3$ et $7$. On écrit donc :

\[
  M = \{3 \ ; \ 7\} , \qquad |M| = 2 .
\]

Nous allons traiter ces deux endroits l'un après l'autre, puis rassembler
les deux résultats. Commençons par le premier.

\subsection*{Étape 1 : le jeton numéro 3}

Le prédicteur a répondu, pour cet endroit :
\[
  P_{\varphi}\bigl(E_{\theta}(x), \Delta_y\bigr)_3
  = (0{,}5 \ ; \ -0{,}2 \ ; \ 1{,}0) .
\]

Et le professeur, qui voyait la vidéo complète, avait produit comme cible :
\[
  \mathrm{sg}\bigl(E_{\bar{\theta}}(y)_3\bigr)
  = (0{,}4 \ ; \ 0{,}1 \ ; \ 0{,}8) .
\]

On soustrait terme à terme :
\[
  (0{,}5 - 0{,}4 \ ; \ -0{,}2 - 0{,}1 \ ; \ 1{,}0 - 0{,}8)
  = (0{,}1 \ ; \ -0{,}3 \ ; \ 0{,}2) .
\]

On efface les signes et on additionne :
\[
  0{,}1 + 0{,}3 + 0{,}2 = 0{,}6 .
\]

L'erreur sur le jeton numéro $3$ vaut donc $\bm{0{,}6}$. Un endroit caché
est réglé ; passons au second, en refaisant exactement les mêmes gestes.

\subsection*{Étape 2 : le jeton numéro 7}

Cette fois, le prédicteur a répondu, et le professeur avait produit sa
cible :
\[
  P_{\varphi}\bigl(E_{\theta}(x), \Delta_y\bigr)_7 = (0{,}0 \ ; \ 0{,}0 \ ; \ 0{,}0) ,
  \qquad
  \mathrm{sg}\bigl(E_{\bar{\theta}}(y)_7\bigr) = (0{,}1 \ ; \ -0{,}2 \ ; \ 0{,}1) .
\]

On soustrait :
\[
  (0{,}0 - 0{,}1 \ ; \ 0{,}0 + 0{,}2 \ ; \ 0{,}0 - 0{,}1)
  = (-0{,}1 \ ; \ 0{,}2 \ ; \ -0{,}1) .
\]

On efface les signes et on additionne :
\[
  0{,}1 + 0{,}2 + 0{,}1 = 0{,}4 .
\]

L'erreur sur le jeton numéro $7$ vaut donc $\bm{0{,}4}$. Nous avons nos deux
erreurs : il ne reste plus qu'à les rassembler en un seul nombre.

\subsection*{Étape 3 : la moyenne}

Le signe $\sum$ nous dit d'additionner les deux résultats :
\[
  0{,}6 + 0{,}4 = 1{,}0 .
\]

Et la fraction $\frac{1}{|M|}$ nous dit de diviser par le nombre d'endroits
cachés, c'est-à-dire par deux :
\[
  \mathcal{L}_{\text{predict}} = \frac{1{,}0}{2} = \bm{0{,}5} .
\]

\begin{exemple}[Le calcul rassemblé dans un tableau]
Voici les mêmes calculs, rangés d'un coup d'œil :

\begin{center}
\footnotesize
\begin{tabular}{cccc}
\toprule
\textbf{Jeton} & \textbf{Réponse} & \textbf{Cible} & \textbf{Écart, norme L1} \\
\midrule
$3$ & $(0{,}5 \ ; \ -0{,}2 \ ; \ 1{,}0)$ & $(0{,}4 \ ; \ 0{,}1 \ ; \ 0{,}8)$ & $0{,}6$ \\
$7$ & $(0{,}0 \ ; \ 0{,}0 \ ; \ 0{,}0)$ & $(0{,}1 \ ; \ -0{,}2 \ ; \ 0{,}1)$ & $0{,}4$ \\
\midrule
\multicolumn{3}{r}{\textbf{somme}} & $1{,}0$ \\
\multicolumn{3}{r}{\textbf{divisée par } $|M| = 2$} & $\bm{0{,}5}$ \\
\bottomrule
\end{tabular}
\end{center}
\end{exemple}

\begin{aretenir}
Vous venez de calculer, à la main et en entier, la perte de V-JEPA~2. Dans
l'article, les jetons contiennent mille nombres au lieu de trois, et le
masque couvre mille deux cents endroits au lieu de deux. Mais le calcul ne
fait strictement rien d'autre que ce que vous venez de faire. Si vous savez
refaire ce tableau, vous savez lire l'équation~1 de l'article.
\end{aretenir}

Nous avons fini de lire la formule. Il est temps de regarder ce qu'elle ne
dit pas.

\section{Le détail qui change tout : où la perte regarde-t-elle ?}

\begin{enclair}
Le prédicteur produit une réponse pour \emph{tous} les jetons, y compris
ceux qui étaient déjà visibles. Mais la perte, elle, ne regarde que les
jetons masqués. Toutes les autres réponses sont jetées à la poubelle, sans
être notées. Ce détail est le sujet de tout le reste de ce livre.
\end{enclair}

Revenons une dernière fois sur le signe $\sum_{i \in M}$. Nous avons dit
qu'il fait prendre à $i$ toutes les valeurs de l'ensemble $M$. Posons-nous
donc la question qui fâche : \emph{que se passe-t-il pour les jetons qui ne
sont pas dans $M$ ?}

La réponse est brutale : rien. Absolument rien. Ils ne sont pas dans la
liste, donc on ne les calcule pas, donc ils ne comptent pas dans la perte.

Or, rappelez-vous ce que nous avions compté au chapitre~\ref{chap:idee-jepa}.
Le prédicteur reçoit les jetons de contexte \emph{et} les jetons de masque,
et il ressort une réponse pour chaque jeton reçu. Il travaille donc pour
tout le monde. Mais on ne le note que sur une partie de sa copie.

\begin{exemple}[Combien de réponses partent à la poubelle]
Reprenons nos chiffres habituels.

\begin{center}
\footnotesize
\begin{tabular}{lr}
\toprule
\textbf{Ce que l'on compte} & \textbf{Combien} \\
\midrule
Réponses produites par le prédicteur & $2\,048$ \\
Réponses regardées par la perte (les masquées) & $1\,200$ \\
Réponses jetées sans être regardées (le contexte) & $848$ \\
\bottomrule
\end{tabular}
\end{center}

La part jetée vaut
\[
  \frac{848}{2\,048} \approx 0{,}414 ,
\]
c'est-à-dire un peu plus de \textbf{quarante et un pour cent}. Plus de
quatre réponses sur dix sont donc fabriquées, puis jetées sans que personne
ne les regarde.
\end{exemple}

La figure~\ref{fig:perte-regarde} est le dessin le plus important de cette
partie du livre. Regardez-le longuement, et surtout retenez ses couleurs :
nous le retrouverons au chapitre~10, presque identique, avec un seul
changement — et ce changement sera toute la nouveauté de V-JEPA~2.1.

\begin{schema}{Où la perte de V-JEPA~2 regarde-t-elle. Chaque petit carré
est un jeton, et le prédicteur a produit une réponse pour les quarante. Mais
l'accolade de droite ne ramasse que les vingt-quatre carrés gris, ceux qui
étaient cachés : eux seuls entrent dans la perte. Les seize carrés bleus
étaient visibles, et leurs réponses partent à la poubelle sans avoir été
regardées.}[perte-regarde]

% ================== la grille : 8 colonnes x 5 lignes ==================
% cellule de 0,70 ; la grille occupe donc 5,60 de large et 3,50 de haut.
\foreach \c in {0,...,7}{
  \foreach \r in {0,...,4}{
    \fill[mbAzure!28, draw=mbAzure!70, line width=0.5pt]
          (\c*0.70, -\r*0.70) rectangle ++(0.66,-0.66);
  }
}
% ---- le bloc masqué : colonnes 1 à 6, lignes 1 à 4 (24 carrés) ----
\foreach \c in {1,...,6}{
  \foreach \r in {1,...,4}{
    \fill[mbSlate!38, draw=mbSlate!75, line width=0.5pt]
          (\c*0.70, -\r*0.70) rectangle ++(0.66,-0.66);
  }
}

% ================== l'accolade : elle ne prend QUE le gris ==================
% le bloc gris va de y = -0,70 à y = -3,50 ; la grille s'arrête à x = 5,26.
\draw[accolade, line width=0.9pt, draw=mbRose!85!black]
      (5.95,-3.50) -- (5.95,-0.70);
\node[font=\sffamily\small\bfseries, text=mbRose!85!black, anchor=west,
      align=left, text width=4.6cm] at (6.50,-1.60)
     {les 24 réponses\\que la perte regarde};
\node[font=\sffamily\scriptsize, text=mbSlate, anchor=west,
      align=left, text width=4.6cm] at (6.50,-2.55)
     {ce sont elles, et elles seules,\\que désigne l'ensemble $M$};

% ================== l'accolade du bas : les 40 réponses ==================
\draw[accolade bas] (0,-3.72) -- (5.56,-3.72);
\node[font=\sffamily\small\bfseries, text=mbSlate] at (2.78,-4.32)
     {40 réponses produites par le prédicteur};

% ================== la légende, sous le dessin ==================
\fill[mbSlate!38, draw=mbSlate!75, line width=0.5pt]
      (0,-5.15) rectangle ++(0.42,-0.42);
\node[font=\sffamily\scriptsize, text=mbInk, anchor=west] at (0.60,-5.36)
     {\textbf{24 jetons masqués} — ils étaient cachés, et leurs réponses
      sont notées};

\fill[mbAzure!28, draw=mbAzure!70, line width=0.5pt]
      (0,-5.85) rectangle ++(0.42,-0.42);
\node[font=\sffamily\scriptsize, text=mbInk, anchor=west] at (0.60,-6.06)
     {\textbf{16 jetons de contexte} — ils étaient visibles, et leurs
      réponses sont jetées};
\end{schema}

\begin{intuition}
Imaginez un professeur qui donne un devoir de dix questions. L'élève répond
aux dix. Le professeur, lui, ne corrige que six questions et arrache la page
des quatre autres sans même la lire.

Que va faire l'élève, au bout de quelques devoirs ? Il va comprendre que
quatre questions ne rapportent rien. Il cessera d'y répondre sérieusement.
Il s'en servira peut-être comme brouillon pour préparer les six autres.

Et si un jour on lui demande ce qu'il avait écrit sur ces pages-là, il n'y
aura rien d'utile. Non pas parce qu'il est mauvais, mais parce que personne
ne le lui a jamais demandé.
\end{intuition}

\begin{piege}
Attention à ne pas se tromper de reproche. Ce que nous venons de décrire
n'est pas un bug, ni un oubli des auteurs de V-JEPA~2. C'est un choix, et un
choix qui se défend très bien.

Souvenez-vous du chapitre~\ref{chap:idee-jepa} : le but du jeu est de forcer
l'élève à comprendre la scène pour deviner ce qu'il ne voit pas. Noter les
endroits qu'il voyait déjà semble idiot, puisqu'il lui suffirait de recopier
ce qu'il a sous les yeux. C'est un raisonnement parfaitement logique.

Et pourtant, ce raisonnement a une conséquence que personne n'avait mesurée
avant l'article de V-JEPA~2.1. Nous la découvrirons au chapitre~8, et nous y
mettrons des chiffres. Pour l'instant, contentez-vous de retenir le fait :
\textbf{les jetons de contexte ne reçoivent aucune consigne}.
\end{piege}

\begin{aretenir}
La formule de V-JEPA~2 dit trois choses, et il faut les retenir toutes les
trois :

\begin{enumerate}
  \item on compare des \emph{représentations}, jamais des pixels ;
  \item on mesure l'écart avec la norme L1, c'est-à-dire en effaçant les
        signes et en additionnant ;
  \item et l'on ne regarde que les endroits \emph{masqués}.
\end{enumerate}

Les deux premières sont l'idée JEPA. La troisième est un détail d'écriture
qui tient dans trois symboles, $i \in M$. C'est pourtant elle, et elle
seule, qui va occuper toute la suite de ce livre.
\end{aretenir}

\begin{enia}
Ce petit $i \in M$ n'est pas un détail de mathématicien : c'est une ligne de
programme, et vous l'écrirez vous-même au chapitre~24.

Dans le code, la perte de V-JEPA~2 se calcule en deux temps. On demande
d'abord au prédicteur ses réponses pour tous les jetons, les $2\,048$. Puis
on met de côté celles des endroits masqués, et l'on ne compare que
celles-là. Ce tri, c'est exactement ce que dit $i \in M$.

Deux conséquences, que vous rencontrerez pour de vrai. La première est un
gaspillage : quatre réponses sur dix sont calculées puis jetées, et ce
travail-là, votre machine le paie quand même. La seconde est bien plus
intéressante : \textbf{le premier changement qui mène de V-JEPA~2 à
V-JEPA~2.1 consiste précisément à toucher à ce tri}. On cessera de jeter les
réponses des endroits visibles, et l'on s'en servira au contraire pour
donner une consigne à chaque jeton. Tout le chapitre~10 tient dans cette
seule idée — et il faudra ensuite quelques autres ingrédients, que les
chapitres~11 à~15 vous feront découvrir un par un.
\end{enia}

Il reste une question que nous avons repoussée depuis le début du chapitre.
Cet ensemble $M$, ces trous, ces zones cachées : comment sont-ils choisis ?

\section{Comment on tire le masque au sort}

\begin{enclair}
À chaque pas d'entraînement, on tire un nouveau masque au hasard. Mais pas
n'importe comment : on tire des rectangles, dont on choisit la surface et la
forme dans des plages fixées à l'avance. Et ces rectangles traversent toute
la durée du clip, de la première image à la dernière.
\end{enclair}

Il serait tentant de cacher les jetons un par un, complètement au hasard,
comme on tire des numéros de loto. Ce serait une mauvaise idée, et il faut
comprendre pourquoi.

\begin{intuition}
Si l'on cache un jeton sur deux, en damier, deviner le contenu d'un trou
devient trop facile : il suffit de regarder les quatre voisins, qui sont
tous visibles, et de faire une moyenne. La machine n'a alors besoin de rien
comprendre à la scène : elle apprend seulement à boucher chaque trou avec la
couleur d'à côté.

Pour que le jeu ait un intérêt, il faut cacher des \emph{morceaux entiers},
assez grands pour qu'aucun voisin immédiat ne donne la réponse. La machine
est alors obligée d'aller chercher des indices loin de là, et donc de
comprendre la scène dans son ensemble.
\end{intuition}

Sur ce point, V-JEPA~2.1 n'invente rien : il reprend telle quelle la façon
de faire des modèles V-JEPA qui l'ont précédé. Cette façon de faire repose
sur trois réglages, et nous allons les prendre un par un.

\subsection*{Premier réglage : l'échelle spatiale}

\begin{definition}[L'échelle spatiale du masque][echelle-spatiale]
L'\emph{échelle spatiale} est la part de la surface de l'image que le trou
doit recouvrir. C'est un nombre entre $0$ et $1$ : une échelle de $0{,}3$
veut dire que le trou couvre trente pour cent de la surface.

Dans l'article, cette échelle est tirée au hasard, à chaque pas, entre
$0{,}15$ et $0{,}7$. Autrement dit, le trou recouvre au minimum quinze pour
cent de l'image, et au maximum soixante-dix pour cent.
\end{definition}

\begin{exemple}[De l'échelle à un nombre de cases]
Prenons une image découpée en $16 \times 16 = 256$ cases, ce qui est le
découpage d'une tranche de notre clip habituel.

Avec une échelle de $0{,}15$, le trou doit recouvrir
$0{,}15 \times 256 = 38{,}4$ cases, donc environ trente-huit.

Avec une échelle de $0{,}7$, il doit en recouvrir
$0{,}7 \times 256 = 179{,}2$, donc environ cent soixante-dix-neuf.
\end{exemple}

Nous savons maintenant quelle \emph{surface} le trou doit couvrir. Mais une
même surface peut se ranger de bien des façons : en carré, en rectangle
couché, en rectangle debout. C'est le deuxième réglage qui décide.

\subsection*{Deuxième réglage : le rapport d'aspect}

\begin{definition}[Le rapport d'aspect][rapport-aspect]
Le \emph{rapport d'aspect} d'un rectangle est sa hauteur divisée par sa
largeur. Un rapport de $1$ donne un carré. Un rapport plus petit que $1$
donne un rectangle couché, plus large que haut. Un rapport plus grand que
$1$ donne un rectangle debout, plus haut que large.

Dans l'article, ce rapport est tiré au hasard entre $0{,}75$ et $1{,}5$. Les
trous ne sont donc jamais de longues bandes très fines : ils restent
toujours des rectangles à peu près aussi hauts que larges.
\end{definition}

La figure~\ref{fig:rapport-aspect} montre les trois formes que ce réglage
peut donner, avec les deux valeurs extrêmes que l'article autorise.

\begin{schema}{Ce que le rapport d'aspect change. À gauche, un rapport plus
petit que $1$ : le rectangle est couché, plus large que haut. Au milieu, un
rapport de $1$ : c'est un carré. À droite, un rapport plus grand que $1$ :
le rectangle est debout, plus haut que large. Les trois trous ont pourtant
la même surface : seule leur forme
change.}[rapport-aspect]

% --- couche : rapport 0,75 ---
\fill[mbSlate!35, draw=mbSlate!80, line width=0.7pt] (0,0) rectangle (2.0,1.5);
\draw[<->, draw=mbSlate!70, line width=0.5pt] (0,-0.25) -- (2.0,-0.25);
\node[font=\sffamily\scriptsize, text=mbSlate] at (1.0,-0.55) {largeur};
\draw[<->, draw=mbSlate!70, line width=0.5pt] (-0.25,0) -- (-0.25,1.5);
\node[font=\sffamily\scriptsize, text=mbSlate, rotate=90] at (-0.55,0.75)
     {hauteur};
\node[font=\sffamily\scriptsize\bfseries, text=mbSlate] at (1.0,2.0)
     {rapport $0{,}75$};
\node[font=\sffamily\scriptsize, text=mbAzure!80!black] at (1.0,-1.0)
     {couché};

% --- carre : rapport 1 ---
\begin{scope}[shift={(4.0cm,0cm)}]
  \fill[mbSlate!35, draw=mbSlate!80, line width=0.7pt] (0,0) rectangle (1.73,1.73);
  \node[font=\sffamily\scriptsize\bfseries, text=mbSlate] at (0.87,2.0)
       {rapport $1$};
  \node[font=\sffamily\scriptsize, text=mbAzure!80!black] at (0.87,-1.0)
       {carré};
\end{scope}

% --- debout : rapport 1,5 ---
\begin{scope}[shift={(7.6cm,0cm)}]
  \fill[mbSlate!35, draw=mbSlate!80, line width=0.7pt] (0,0) rectangle (1.41,2.12);
  \node[font=\sffamily\scriptsize\bfseries, text=mbSlate] at (0.7,2.4)
       {rapport $1{,}5$};
  \node[font=\sffamily\scriptsize, text=mbAzure!80!black] at (0.7,-1.0)
       {debout};
\end{scope}
\end{schema}

\begin{exemple}[Lire un rapport d'aspect]
Un trou fait $12$ cases de haut et $16$ cases de large. Son rapport d'aspect
vaut $12 \div 16 = 0{,}75$ : il est plus petit que $1$, donc le trou est
couché. C'est bien ce que disent les deux nombres, puisque $12$ est plus
petit que $16$.

Un autre trou fait $18$ cases de haut et $12$ de large. Son rapport vaut
$18 \div 12 = 1{,}5$ : il est plus grand que $1$, donc le trou est debout.

Et un trou de $14$ cases sur $14$ a pour rapport $14 \div 14 = 1$ : c'est un
carré.
\end{exemple}

Comment passe-t-on de ces deux réglages aux deux côtés du rectangle ? Par
deux petits calculs, toujours les mêmes, et il faut d'abord se rappeler ce
qu'est une racine carrée. Nous l'avions rencontrée au
chapitre~\ref{chap:vit} : la racine carrée d'un nombre, c'est celui qui,
multiplié par lui-même, redonne ce nombre. Ainsi $\sqrt{121} = 11$, puisque
$11 \times 11 = 121$.

Voici les deux calculs. Pour la hauteur, on multiplie la surface par le
rapport, et on prend la racine carrée du résultat. Pour la largeur, on
divise la surface par le rapport, et on prend là aussi la racine carrée.
Ces deux calculs tombent rarement juste : ils donnent des nombres à virgule,
comme $16{,}4$. Or on ne peut pas cacher $16{,}4$ cases : une case est cachée
ou elle ne l'est pas. Il faut donc \emph{arrondir}, c'est-à-dire remplacer le
nombre à virgule par le nombre sans virgule le plus proche. Ainsi $5{,}2$
devient $5$, parce qu'il est plus près de $5$ que de $6$ ; et $7{,}8$ devient
$8$, parce qu'il est plus près de $8$ que de $7$.

\begin{exemple}[Trouver les côtés du rectangle]
Prenons d'abord une surface de $121$ cases avec un rapport de $1$.

\[
  \begin{aligned}
    \text{hauteur} &= \sqrt{121 \times 1} = \sqrt{121} = 11 , \\
    \text{largeur} &= \sqrt{121 \div 1} = \sqrt{121} = 11 .
  \end{aligned}
\]

Le trou fait donc $11$ cases sur $11$. C'est bien un carré, comme le
promettait un rapport de $1$.

Prenons maintenant une surface de $179$ cases avec un rapport de $1{,}5$.

\[
  \begin{aligned}
    \text{hauteur} &= \sqrt{179 \times 1{,}5} = \sqrt{268{,}5} \approx 16{,}4 , \\
    \text{largeur} &= \sqrt{179 \div 1{,}5} = \sqrt{119{,}3} \approx 10{,}9 .
  \end{aligned}
\]

On arrondit : $16{,}4$ devient $16$, et $10{,}9$ devient $11$.

Le trou fait donc $16$ cases de haut sur $11$ de large. Il est bien plus
haut que large, comme le promettait un rapport supérieur à $1$.
\end{exemple}

\begin{piege}
Vérifions la surface obtenue : $16 \times 11 = 176$ cases, alors que nous en
visions $179$. Les deux nombres ne sont pas tout à fait égaux, et c'est
normal : en arrondissant les deux côtés à des nombres entiers, on s'écarte
forcément un peu de la surface visée.

Ne cherchez donc pas à retomber pile sur l'échelle demandée. Elle indique une
taille souhaitée, pas une taille exacte. C'est aussi pourquoi, dans la
figure et dans le programme qui suivent, vous verrez $35$ cases là où
l'échelle $0{,}15$ en demandait $38$ : rien n'est cassé, ce sont seulement
deux arrondis.
\end{piege}

Nos deux premiers réglages dessinent le trou sur \emph{une} image. Reste à
savoir ce qu'il devient quand la vidéo avance. C'est le troisième réglage,
et c'est le plus surprenant des trois.

\subsection*{Troisième réglage : l'échelle temporelle}

\begin{definition}[L'échelle temporelle du masque][echelle-temporelle]
L'\emph{échelle temporelle} est la part de la durée du clip que le trou
recouvre. Dans l'article, elle vaut toujours exactement $1{,}0$ : le trou
traverse donc \emph{toute} la durée du clip, de la première image à la
dernière.
\end{definition}

Ce qui surprend, ici, c'est que ce réglage ne varie jamais. Ce n'est pas un
oubli : c'est un choix, et il est très malin.

\begin{exemple}[Ce que « toujours $1$ » veut dire en nombre d'images]
Un clip de la phase principale contient $16$ images, ce qui fait $8$
tranches de temps une fois les images collées deux par deux.

Une échelle temporelle de $0{,}5$ voudrait dire que le trou ne couvre que
$0{,}5 \times 8 = 4$ tranches, donc la moitié du clip. Avec l'échelle de $1$
retenue par l'article, il en couvre $1 \times 8 = 8$, c'est-à-dire toutes.

Autrement dit, un trou de $150$ cases sur une tranche ne cache pas $150$
jetons, mais $8 \times 150 = 1\,200$ jetons : le même trou est creusé huit
fois, une fois par tranche.
\end{exemple}

\begin{intuition}
Supposons que le trou ne couvre que la moitié du temps. Alors, pour deviner
ce qui est caché à un endroit, la machine n'aurait qu'à regarder le même
endroit un instant plus tôt, quand il était encore visible. Comme les choses
bougent lentement d'une image à l'autre, elle aurait la réponse presque
gratuitement, et sans rien comprendre.

En faisant traverser le trou d'un bout à l'autre du clip, on lui coupe cette
issue de secours. Cet endroit-là n'a \emph{jamais} été visible, à aucun
moment. La machine est donc obligée de raisonner sur ce qui se passe
\emph{autour}, et sur la façon dont la scène évolue. C'est exactement ce
qu'on veut lui apprendre.
\end{intuition}

La figure~\ref{fig:masque-tube} montre cette idée : le trou est une sorte de
tunnel creusé à travers toute la pile d'images.

\begin{schema}{L'échelle temporelle vaut toujours $1$. Les quatre images du
clip sont posées côte à côte, dans l'ordre du temps. Le trou gris tombe
exactement à la même place sur chacune : les traits pointillés le montrent.
Cet endroit de la scène n'est donc visible à aucun moment, et la machine ne
peut pas se contenter de regarder l'image d'avant pour deviner ce qui s'y
cache.}[masque-tube]

% ================== les quatre images, bien séparées ==================
% Chaque image fait 2,1 de large sur 1,6 de haut, et l'on en pose une
% toutes les 2,6 unités. Le trou occupe toujours le même rectangle.
\foreach \k/\nom in {0/{image 1}, 1/{image 2}, 2/{image 3}, 3/{image 4}}{
  \begin{scope}[shift={(\k*2.6,0)}]
    \fill[mbAzure!20, draw=mbAzure!70, line width=0.7pt]
          (0,0) rectangle (2.1,1.6);
    \fill[mbSlate!42, draw=mbSlate!80, line width=0.7pt]
          (0.6,0.45) rectangle (1.5,1.15);
    \node[font=\sffamily\scriptsize, text=mbSlate] at (1.05,-0.32) {\nom};
  \end{scope}
}

% ---- les pointillés : c'est bien le même endroit d'une image à l'autre ----
\draw[draw=mbRose!75, line width=0.7pt, dashed] (0.6,1.15) -- (9.9,1.15);
\draw[draw=mbRose!75, line width=0.7pt, dashed] (0.6,0.45) -- (9.9,0.45);

% ---- la flèche du temps ----
\draw[-{Stealth[round,length=2.6mm,width=2mm]}, draw=mbAmber!85!black,
      line width=1pt] (0,-0.72) -- (9.9,-0.72);
\node[font=\sffamily\scriptsize, text=mbAmber!80!black, anchor=north]
      at (4.95,-0.82) {le temps passe : les images du clip};

% ---- l'annotation, posée au-dessus, sans croiser le dessin ----
\node[font=\sffamily\small\bfseries, text=mbRose!85!black, anchor=west]
      at (0,2.32) {le même trou, à la même place, sur toutes les images};
\node[font=\sffamily\scriptsize, text=mbSlate, anchor=west]
      at (0,1.95) {cet endroit n'est jamais visible, pas une seule fois};
\end{schema}

\begin{surlesdonnees}
Les trois réglages, tels qu'ils sont écrits dans l'annexe de l'article, et
tels que vous les retrouverez au chapitre~24 quand nous coderons :

\begin{center}
\footnotesize
\begin{tabular}{lll}
\toprule
\textbf{Réglage} & \textbf{Valeur} & \textbf{Ce que cela veut dire} \\
\midrule
Échelle spatiale  & de $0{,}15$ à $0{,}7$ &
  le trou couvre de 15 à 70 \% de l'image \\
Échelle temporelle & $1{,}0$ toujours &
  il traverse tout le clip \\
Rapport d'aspect  & de $0{,}75$ à $1{,}5$ &
  ni trop couché, ni trop debout \\
\bottomrule
\end{tabular}
\end{center}

Ces mêmes valeurs sont employées pendant la phase principale et pendant la
phase de refroidissement, sans le moindre changement.
\end{surlesdonnees}

Nous tenons nos trois réglages. Voyons maintenant, de nos yeux, à quoi
ressemblent les trous qu'ils fabriquent.

\subsection*{Trois masques tirés au sort}

Ces trois réglages restent un peu flous tant qu'on n'a pas vu ce qu'ils
donnent. Tirons-en donc trois au hasard, et dessinons les trous obtenus.

La figure~\ref{fig:trois-masques} montre le résultat sur une tranche de
notre clip habituel, c'est-à-dire sur une grille de $16 \times 16$ cases.
Les trois tirages sont ceux que nous ferons calculer par le programme, tout
à l'heure : vous retrouverez donc exactement les mêmes nombres.

\begin{schema}{Trois masques tirés au sort, dessinés sur une tranche de
$16 \times 16$ cases. À gauche, le plus petit trou que l'article autorise ;
au milieu, un trou de taille moyenne, ici carré ; à droite, le plus gros,
un rectangle debout si haut qu'il touche les deux bords de la grille. Les
trois viennent des mêmes deux réglages, tirés au hasard dans leurs plages :
seuls les nombres tirés changent.}[trois-masques]

% ================= panneau 1 : petit trou, 5 x 7 =================
\begin{scope}[shift={(0cm,0cm)}]
  \begin{scope}[x=0.15cm, y=0.15cm]
    \fill[mbAzure!18] (0,0) rectangle (16,-16);
    \foreach \k in {1,...,15}{
      \draw[draw=mbAzure!40, line width=0.2pt] (\k,0) -- (\k,-16);
      \draw[draw=mbAzure!40, line width=0.2pt] (0,-\k) -- (16,-\k);
    }
    \fill[mbSlate!45] (5,-4) rectangle (12,-9);
    \draw[draw=mbSlate!85, line width=0.6pt, dashed] (5,-4) rectangle (12,-9);
    \draw[draw=mbAzure!70, line width=0.7pt] (0,0) rectangle (16,-16);
  \end{scope}
  \node[font=\sffamily\scriptsize\bfseries, text=mbSlate] at (1.2,-2.85)
       {petit trou};
  \node[font=\sffamily\scriptsize, text=mbSlate, align=center] at (1.2,-3.5)
       {échelle $0{,}15$\\rapport $0{,}75$};
  \node[font=\sffamily\scriptsize, text=mbAzure!80!black] at (1.2,-4.2)
       {$5 \times 7 = 35$ cases};
\end{scope}

% ================= panneau 2 : trou moyen, 11 x 11 =================
\begin{scope}[shift={(3.3cm,0cm)}]
  \begin{scope}[x=0.15cm, y=0.15cm]
    \fill[mbAzure!18] (0,0) rectangle (16,-16);
    \foreach \k in {1,...,15}{
      \draw[draw=mbAzure!40, line width=0.2pt] (\k,0) -- (\k,-16);
      \draw[draw=mbAzure!40, line width=0.2pt] (0,-\k) -- (16,-\k);
    }
    \fill[mbSlate!45] (2,-2) rectangle (13,-13);
    \draw[draw=mbSlate!85, line width=0.6pt, dashed] (2,-2) rectangle (13,-13);
    \draw[draw=mbAzure!70, line width=0.7pt] (0,0) rectangle (16,-16);
  \end{scope}
  \node[font=\sffamily\scriptsize\bfseries, text=mbSlate] at (1.2,-2.85)
       {trou moyen};
  \node[font=\sffamily\scriptsize, text=mbSlate, align=center] at (1.2,-3.5)
       {échelle $0{,}47$\\rapport $1{,}00$};
  \node[font=\sffamily\scriptsize, text=mbAzure!80!black] at (1.2,-4.2)
       {$11 \times 11 = 121$ cases};
\end{scope}

% ================= panneau 3 : gros trou, 16 x 11 =================
\begin{scope}[shift={(6.6cm,0cm)}]
  \begin{scope}[x=0.15cm, y=0.15cm]
    \fill[mbAzure!18] (0,0) rectangle (16,-16);
    \foreach \k in {1,...,15}{
      \draw[draw=mbAzure!40, line width=0.2pt] (\k,0) -- (\k,-16);
      \draw[draw=mbAzure!40, line width=0.2pt] (0,-\k) -- (16,-\k);
    }
    \fill[mbSlate!45] (3,0) rectangle (14,-16);
    \draw[draw=mbSlate!85, line width=0.6pt, dashed] (3,0) rectangle (14,-16);
    \draw[draw=mbAzure!70, line width=0.7pt] (0,0) rectangle (16,-16);
  \end{scope}
  \node[font=\sffamily\scriptsize\bfseries, text=mbSlate] at (1.2,-2.85)
       {gros trou};
  \node[font=\sffamily\scriptsize, text=mbSlate, align=center] at (1.2,-3.5)
       {échelle $0{,}70$\\rapport $1{,}50$};
  \node[font=\sffamily\scriptsize, text=mbAzure!80!black] at (1.2,-4.2)
       {$16 \times 11 = 176$ cases};
\end{scope}
\end{schema}

\begin{exemple}[Les trois masques, en jetons]
Chacun de ces trois trous traverse les $8$ tranches de temps du clip,
puisque l'échelle temporelle vaut toujours $1$. Il suffit donc de
multiplier le nombre de cases par $8$ pour obtenir le nombre de jetons
cachés :

\begin{center}
\footnotesize
\begin{tabular}{lrrr}
\toprule
\textbf{Tirage} & \textbf{Cases} & \textbf{Jetons masqués} &
\textbf{Jetons de contexte} \\
\midrule
petit trou  & $35$  & $280$   & $1\,768$ \\
trou moyen  & $121$ & $968$   & $1\,080$ \\
gros trou   & $176$ & $1\,408$ & $640$ \\
\bottomrule
\end{tabular}
\end{center}

D'un tirage à l'autre, le nombre de jetons cachés passe donc de $280$ à
$1\,408$, soit cinq fois plus. La difficulté du jeu change beaucoup d'un pas
d'entraînement au suivant, et c'est voulu : la machine doit savoir se
débrouiller aussi bien avec quelques indices qu'avec beaucoup.
\end{exemple}

\begin{surlesdonnees}
Regardez bien le panneau de droite. Le calcul demandait une hauteur de
$16{,}4$ cases, qui une fois arrondie donne $16$ : c'est exactement la
hauteur de la grille. Le trou touche donc le bord du haut et le bord du bas.

Il s'en est fallu de peu que le rectangle ne dépasse. C'est pour ce cas-là
que notre programme contiendra un \co{python}{min} : si jamais le calcul
réclame plus de cases que la grille n'en possède, on ramène simplement le
rectangle à la taille de la grille. Ce n'est ni une erreur ni un cas rare,
et le programme continue sans se plaindre.
\end{surlesdonnees}

\begin{remarque}
Dans la pratique, on ne se contente pas d'un seul rectangle. On en tire
plusieurs, et l'on cache tout ce qui est recouvert par au moins l'un d'eux.
Quand deux rectangles se chevauchent, la partie commune n'est évidemment
comptée qu'une seule fois : c'est ce qu'on appelle prendre leur
\emph{réunion}. Le résultat est un trou de forme moins régulière, ce qui
rend le jeu plus intéressant. Nous le ferons tourner dans la section de code,
et vous verrez apparaître les nombres qui nous accompagnent depuis le
chapitre~\ref{chap:idee-jepa}.
\end{remarque}

Nous avons la formule et nous avons le masque. Vérifions tout cela sur une
machine.

\section{Passons au code}

Nous allons écrire quatre petits programmes. Le premier refait, ligne par
ligne, le calcul de perte que nous venons de mener à la main : vous devez
retrouver le $0{,}6$, le $0{,}4$ et le $0{,}5$. Le deuxième compare la norme
L1 aux carrés du chapitre~\ref{chap:apprendre}. Les deux derniers s'occupent
du masque : l'un calcule la taille d'un trou à partir de ses réglages,
l'autre pose deux trous sur une grille et compte les jetons cachés.

Commençons par la perte.

\begin{code}{python}[La perte de V-JEPA~2, sur nos deux jetons]
import torch

# Deux jetons masques, de trois nombres chacun.
# Ligne 0 : le jeton numero 3. Ligne 1 : le jeton numero 7.
prediction = torch.tensor([[0.5, -0.2, 1.0],
                           [0.0,  0.0, 0.0]])
cible      = torch.tensor([[0.4,  0.1, 0.8],
                           [0.1, -0.2, 0.1]])

# 1. l'ecart, terme a terme, puis on efface les signes
ecarts = (prediction - cible).abs()
print(ecarts)

# 2. la norme L1 : on additionne les nombres de chaque jeton
par_jeton = ecarts.sum(dim=1)
print(par_jeton)

# 3. la moyenne sur les jetons masques : c'est la perte
perte = par_jeton.mean()
print(perte)
\end{code}

Voici ce que le programme affiche :

\begin{sortie}
tensor([[0.1000, 0.3000, 0.2000],
        [0.1000, 0.2000, 0.1000]])
tensor([0.6000, 0.4000])
tensor(0.5000)
\end{sortie}

Reprenez le tableau de la section~\ref{sec:calcul-complet} et comparez. Le premier affichage donne
les écarts, une fois les signes effacés, exactement comme nous les avions
calculés. Le deuxième donne les deux normes L1, soit $0{,}6$ puis $0{,}4$.
Le troisième donne la perte, $0{,}5$. La machine n'a rien fait de plus que
ce que nous venons de faire au crayon.

\begin{piege}
Regardez bien la ligne \co{python}{ecarts.sum(dim=1)}. Le réglage
\co{python}{dim=1} dit : « additionne le long du deuxième axe », c'est-à-dire
à l'intérieur de chaque jeton. Si vous écriviez \co{python}{dim=0}, vous
additionneriez d'un jeton à l'autre, ce qui n'a aucun sens ici. Le
programme ne se plaindrait pas : il rendrait trois nombres au lieu de deux,
et la perte serait fausse en silence. Reprenez le réflexe du
chapitre~\ref{chap:image-video} : affichez la forme, et demandez-vous si
elle est celle que vous attendiez.
\end{piege}

Passons à la comparaison entre la norme L1 et les carrés. C'est le calcul de
l'encadré « pour aller plus loin », vérifié sur machine.

\begin{code}{python}[Norme L1 contre carrés : qui punit le plus les gros écarts ?]
import torch

ecarts = torch.tensor([0.1, 1.0])     # un petit ecart, et un gros

total_l1     = ecarts.abs().sum()
total_carres = (ecarts ** 2).sum()

print(f"norme L1 : total = {total_l1:.2f}")
print(f"   le petit ecart compte pour {100 * 0.1 / total_l1:.0f} %")
print(f"   le gros  ecart compte pour {100 * 1.0 / total_l1:.0f} %")

print(f"carres   : total = {total_carres:.2f}")
print(f"   le petit ecart compte pour {100 * 0.01 / total_carres:.0f} %")
print(f"   le gros  ecart compte pour {100 * 1.00 / total_carres:.0f} %")
\end{code}

Et le résultat :

\begin{sortie}
norme L1 : total = 1.10
   le petit ecart compte pour 9 %
   le gros  ecart compte pour 91 %
carres   : total = 1.01
   le petit ecart compte pour 1 %
   le gros  ecart compte pour 99 %
\end{sortie}

Avec les carrés, le petit écart ne pèse plus qu'un pour cent : autant dire
qu'il a disparu. Avec la norme L1, il pèse encore neuf fois plus. C'est
précisément ce que V-JEPA~2 cherche : continuer d'écouter les petits écarts,
au lieu de ne s'occuper que des grosses surprises imprévisibles.

Terminons par le masque. Nous allons d'abord calculer la taille d'un
rectangle à partir de ses deux réglages, puis en poser deux sur une grille
et compter.

\begin{code}{python}[Calculer la taille d'un trou à partir de ses réglages]
def taille_du_trou(echelle, rapport, H=16, W=16):
    """Rend la hauteur et la largeur du rectangle, en cases."""
    aire = echelle * H * W                       # la surface visee
    hauteur = round((aire * rapport) ** 0.5)
    largeur = round((aire / rapport) ** 0.5)
    return min(hauteur, H), min(largeur, W)      # on ne sort pas de la grille

for echelle, rapport in [(0.15, 0.75), (0.473, 1.0), (0.70, 1.5)]:
    h, w = taille_du_trou(echelle, rapport)
    print(f"echelle {echelle:.3f}, rapport {rapport:.2f} -> "
          f"{h:2d} x {w:2d} = {h * w:3d} cases, soit {8 * h * w:4d} jetons")
\end{code}

Le programme affiche :

\begin{sortie}[Trois trous, du plus petit au plus grand]
echelle 0.150, rapport 0.75 ->  5 x  7 =  35 cases, soit  280 jetons
echelle 0.473, rapport 1.00 -> 11 x 11 = 121 cases, soit  968 jetons
echelle 0.700, rapport 1.50 -> 16 x 11 = 176 cases, soit 1408 jetons
\end{sortie}

Le facteur $8$ de la dernière colonne, c'est le nombre de tranches de temps
d'un clip de la phase principale, que nous avions calculé au
chapitre~\ref{chap:vit} : $16$ images divisées par des tubelets de $2$ font
bien $8$ tranches. Comme le trou traverse tout le clip, on multiplie le
nombre de cases par ce $8$.

Remarquez la dernière ligne : la hauteur voulue dépassait la grille, et le
\co{python}{min} l'a ramenée à $16$. C'est normal, et le programme continue
sans se plaindre.

Posons enfin deux rectangles sur la grille, et comptons.

\begin{code}{python}[Fabriquer un masque et compter les jetons cachés]
import torch

T, H, W = 8, 16, 16          # tranches de temps, lignes, colonnes

masque = torch.zeros(H, W, dtype=torch.bool)   # au depart, rien n'est cache
masque[2:13, 2:13] = True                      # premier bloc : 11 x 11
masque[9:16, 10:16] = True                     # second bloc  :  7 x  6

par_tranche = int(masque.sum())      # la reunion des deux blocs
total = par_tranche * T              # le trou traverse tout le clip

print(f"cases cachees sur une tranche : {par_tranche}")
print(f"jetons masques  (|M|)         : {total}")
print(f"jetons visibles (le contexte) : {T * H * W - total}")
print(f"part jetee par la perte       : "
      f"{100 * (T * H * W - total) / (T * H * W):.1f} %")
\end{code}

Le résultat :

\begin{sortie}[Le compte final]
cases cachees sur une tranche : 151
jetons masques  (|M|)         : 1208
jetons visibles (le contexte) : 840
part jetee par la perte       : 41.0 %
\end{sortie}

Arrêtons-nous sur ces quatre lignes, car elles bouclent le chapitre.

Le premier bloc couvre $11 \times 11 = 121$ cases, le second $7 \times 6 =
42$. Cela ferait $163$ si on les additionnait bêtement. Mais ils se
chevauchent sur $12$ cases, et ces cases-là ne doivent être comptées qu'une
fois : le total est donc $121 + 42 - 12 = 151$. C'est bien ce qu'affiche le
programme, et c'est ce que veut dire « prendre la réunion ».

Ces $151$ cases, multipliées par les $8$ tranches de temps, donnent $1\,208$
jetons masqués et $840$ jetons de contexte. Ce sont, à quelques unités près,
les $1\,200$ et $848$ qui nous accompagnent depuis le
chapitre~\ref{chap:idee-jepa} — et nous continuerons d'employer ces
chiffres ronds, plus faciles à retenir.

Et la dernière ligne dit tout : \textbf{quarante et un pour cent des réponses
du prédicteur sont fabriquées, puis jetées}.

\begin{aretenir}
L'équation~1 de V-JEPA~2, en une phrase de français : \emph{on mesure, aux
seuls endroits masqués, l'écart entre ce que le prédicteur a répondu et ce
que le professeur avait produit, on efface les signes, on additionne, et on
divise par le nombre d'endroits masqués}.

Et le mot qui compte dans cette phrase est le mot « seuls ».
\end{aretenir}

\vspace{0.5em}

Nous voici au bout de la partie~II. Nous savons comment la machine est
bâtie, pourquoi elle ne s'écroule pas, et ce qu'elle cherche exactement à
faire baisser. Toutes les pièces de V-JEPA~2 sont maintenant entre vos
mains.

Cette formule marche d'ailleurs très bien pour reconnaître une action dans
une vidéo, et l'article de V-JEPA~2 l'a prouvé sur des dizaines de mesures.
Mais que se passe-t-il si on lui demande autre chose ? Si, au lieu de
demander « que se passe-t-il dans cette vidéo ? », on demande « à quelle
distance se trouve ce pixel-là ? », ou bien « à quel objet appartient ce
pixel-là ? »

Ces questions-là ne réclament pas une seule réponse pour toute l'image :
elles en réclament une \emph{par endroit}. Pour y répondre, il faudrait que
chaque jeton décrive bien sa propre place — y compris les jetons de
contexte, ceux dont personne, nous venons de le voir, ne demande jamais
rien.

Vous devinez déjà où cela nous mène : sur ces questions-là, la formule que
nous venons d'écrire échoue. Et c'est précisément cet échec qui a donné
naissance à V-JEPA~2.1. Avant de le constater nous-mêmes, il faut apprendre
à poser ces nouvelles questions et à noter les réponses : c'est le sujet du
chapitre suivant, qui ouvre la partie~III.

\phantomsection
\section*{Exercices}
\addcontentsline{toc}{section}{Exercices}

\begin{exercice}[Calculer une norme L1]
Un prédicteur répond $(1{,}2 \ ; \ -0{,}5 \ ; \ 0{,}3 \ ; \ 0{,}0)$ alors
que la cible valait $(1{,}0 \ ; \ -0{,}1 \ ; \ 0{,}3 \ ; \ -0{,}2)$. Quelle
est la norme L1 de l'écart ?
\end{exercice}

\begin{solution}
On soustrait terme à terme :
\[
  1{,}2 - 1{,}0 = 0{,}2 , \quad
  -0{,}5 + 0{,}1 = -0{,}4 , \quad
  0{,}3 - 0{,}3 = 0 , \quad
  0{,}0 + 0{,}2 = 0{,}2 .
\]

On efface les signes, ce qui donne $(0{,}2 \ ; \ 0{,}4 \ ; \ 0 \ ; \
0{,}2)$, puis on additionne :
\[
  0{,}2 + 0{,}4 + 0 + 0{,}2 = 0{,}8 .
\]

La norme L1 vaut donc $0{,}8$. Remarquez le troisième nombre : la machine
était tombée parfaitement juste dessus, et il ne coûte donc rien du tout.
\end{solution}

\begin{exercice}[Une perte complète sur trois jetons]
Le masque recouvre trois jetons. Les normes L1 de leurs écarts valent
$1{,}2$, $0{,}9$ et $0{,}3$. Que vaut la perte ? Et si un quatrième jeton
était masqué, avec un écart de $0{,}0$, que deviendrait-elle ?
\end{exercice}

\begin{solution}
Avec trois jetons, $|M| = 3$ et la somme vaut $1{,}2 + 0{,}9 + 0{,}3 =
2{,}4$. La perte vaut donc
\[
  \frac{2{,}4}{3} = 0{,}8 .
\]

Avec un quatrième jeton d'écart nul, la somme ne bouge pas — elle vaut
toujours $2{,}4$ — mais $|M|$ passe à $4$. La perte devient
\[
  \frac{2{,}4}{4} = 0{,}6 .
\]

Elle a baissé, alors que la machine ne s'est pas améliorée d'un pouce sur
les trois premiers jetons. C'est l'effet de la division : la perte est une
\emph{moyenne}, et ajouter une bonne réponse tire la moyenne vers le bas.
C'est aussi pourquoi on ne compare jamais deux pertes calculées sur des
masques de tailles différentes.
\end{solution}

\begin{exercice}[Compter les réponses jetées]
Un clip de la phase de refroidissement contient $18\,432$ jetons. Le masque
en recouvre $11\,000$. Combien de réponses le prédicteur produit-il ?
Combien la perte en regarde-t-elle ? Combien sont jetées, et quelle part
cela représente-t-il ?
\end{exercice}

\begin{solution}
Le prédicteur produit une réponse par jeton reçu, et il reçoit le contexte
plus les masques, c'est-à-dire tout le clip : il produit donc $18\,432$
réponses.

La perte ne regarde que les jetons masqués, soit $11\,000$ réponses.

Il en reste donc $18\,432 - 11\,000 = 7\,432$ qui sont jetées, ce qui
représente
\[
  \frac{7\,432}{18\,432} \approx 0{,}403 ,
\]
soit un peu plus de quarante pour cent. On retrouve, à très peu de chose
près, la même proportion que sur un clip de la phase principale : ce n'est
pas la taille du clip qui décide de cette part, mais la taille du masque.
\end{solution}

\begin{exercice}[Retrouver la taille d'un trou]
Sur une grille de $16 \times 16$ cases, on tire une échelle spatiale de
$0{,}25$ et un rapport d'aspect de $1{,}0$. Combien de cases le trou
couvre-t-il, et quelles sont ses dimensions ? Et comme le clip compte $8$
tranches de temps, combien de jetons cela fait-il en tout ?
\end{exercice}

\begin{solution}
La surface visée vaut $0{,}25 \times 256 = 64$ cases. On applique ensuite
les deux calculs habituels. Hauteur : $\sqrt{64 \times 1} = 8$. Largeur :
$\sqrt{64 \div 1} = 8$. Le trou fait donc $8$ cases sur $8$, et c'est bien
un carré, comme le voulait un rapport de $1$.

Comme l'échelle temporelle vaut toujours $1$, ce même carré est creusé sur
les huit tranches de temps. Cela fait
\[
  8 \times 64 = 512 \text{ jetons masqués,}
\]
et il reste donc $2\,048 - 512 = 1\,536$ jetons de contexte.
\end{solution}

\begin{exercice}[Pourquoi le trou traverse tout le clip]
Un collègue propose de ne masquer que la moitié des images du clip, pour
« laisser un peu d'aide à la machine ». Pourquoi est-ce une mauvaise idée ?
\end{exercice}

\begin{solution}
Parce que cette « aide » est bien trop grande, et qu'elle vide le jeu de son
intérêt.

Si un endroit est visible sur la première moitié du clip et caché sur la
seconde, la machine n'a plus besoin de comprendre quoi que ce soit : il lui
suffit de recopier ce qu'elle a vu au même endroit un instant plus tôt.
Comme les choses bougent peu d'une image à l'autre, elle aura presque juste
à tous les coups, et la perte descendra très bas sans qu'elle ait rien
appris.

C'est exactement pour fermer cette porte que l'article fixe l'échelle
temporelle à $1{,}0$ et ne la fait jamais varier. L'endroit caché ne doit
avoir été visible à aucun moment.
\end{solution}

\begin{exercice}[Le mot qui compte]
Résumez en une phrase ce que la formule de V-JEPA~2 fait des réponses que le
prédicteur produit pour les jetons de contexte. Puis dites, en une phrase,
pourquoi ce choix pouvait sembler raisonnable à ses auteurs.
\end{exercice}

\begin{solution}
Elle les jette : le signe $\sum_{i \in M}$ ne fait prendre à $i$ que les
numéros des jetons masqués, si bien que les réponses produites aux endroits
visibles ne sont jamais comparées à quoi que ce soit, et ne comptent donc
pour rien dans la perte.

Ce choix semblait raisonnable parce que noter la machine là où elle voit
déjà tout paraît sans intérêt : il lui suffirait de recopier son entrée pour
avoir juste, et l'on n'apprendrait rien. Le raisonnement est logique — mais
nous verrons au chapitre~8 qu'il a une conséquence que personne n'avait
mesurée, et au chapitre~10 comment V-JEPA~2.1 la corrige.
\end{solution}


\part{Le problème que V-JEPA 2.1 vient résoudre}
% \input{chapters/07-global-contre-dense}
% \input{chapters/08-diagnostic}

\part{V-JEPA 2.1, ingrédient par ingrédient}
% \input{chapters/09-tableau-de-bord}
% \input{chapters/10-perte-dense}
% \input{chapters/11-supervision-profonde}
% \input{chapters/12-donnees}
% \input{chapters/13-tokeniseur-multimodal}
% \input{chapters/14-taille-et-refroidissement}
% \input{chapters/15-distillation}

\part{Les résultats, lus avec honnêteté}
% \input{chapters/16-prediction-du-futur}
% \input{chapters/17-robot}
% \input{chapters/18-taches-denses}

\part{Construire le modèle et le programme d'entraînement}
% \input{chapters/19-atelier}
% \input{chapters/20-tokeniseur-en-code}
% \input{chapters/21-encodeur}
% \input{chapters/22-masques}
% \input{chapters/23-predicteur}
% \input{chapters/24-pertes-et-boucle}
% \input{chapters/25-surveiller-entrainement}

\part{De l'entraînement à la production}
% \input{chapters/26-choisir-son-modele}
% \input{chapters/27-encodeur-gele}
% \input{chapters/28-entrainer-sa-sonde}
% \input{chapters/29-servir-le-modele}
% \input{chapters/30-surveiller-en-service}
% \input{chapters/31-pieges}

%----------------------------------------------------------------------
% Pages finales
%----------------------------------------------------------------------
\backmatter
\cleardoublepage
\printbibliography[heading=bibintoc]

\end{document}
